% =====================================================================
%  ScholarPath — الدليل العربي الشارح للنظام (كورس من الصفر)
%  الملف الرئيسي — يجمّع كل الفصول
%  Compile:  xelatex ScholarPath-AR-Guide.tex   (مرّتين للفهرس)
% =====================================================================
\documentclass[12pt,a4paper]{article}
% =====================================================================
%  ScholarPath — دليل الشرح العربي/الإنجليزي — الـ preamble المشترك
%  Compile with:  xelatex  (NOT pdflatex — Arabic needs XeTeX)
%  عربي أساسي (RTL) + مصطلحات إنجليزية داخلية + صناديق mdframed
% =====================================================================

% ── اللغة والخطوط ─────────────────────────────────────────────────────
\usepackage{fontspec}
\usepackage{polyglossia}
% numerals=maghrib → أرقام لاتينية (0-9) للأقسام والصفحات، أنسب لمستند تقني ثنائي اللغة
\setmainlanguage[numerals=maghrib]{arabic}
\setotherlanguage{english}
\newfontfamily\arabicfont[Script=Arabic,Scale=1.06]{Amiri}
\newfontfamily\englishfont{TeX Gyre Termes}
\setmainfont{TeX Gyre Termes}
% خط أحادي المسافة للأكواد/المعرّفات
\newfontfamily\monofont{Consolas}[Scale=0.92]

% ── الهندسة والألوان ─────────────────────────────────────────────────
\usepackage[a4paper,margin=2.3cm,top=2.7cm,bottom=2.5cm,headheight=26pt,headsep=12pt]{geometry}
\usepackage[table]{xcolor}
\definecolor{spInk}{HTML}{1F2A44}
\definecolor{spPrimary}{HTML}{2A4D9B}
\definecolor{spAccent}{HTML}{0E7C7B}
\definecolor{spHow}{HTML}{1565C0}
\definecolor{spWhy}{HTML}{2E7D32}
\definecolor{spSec}{HTML}{B23A00}
\definecolor{spQA}{HTML}{5E35B1}
\definecolor{spGrayBg}{HTML}{F4F6FB}
\definecolor{spRule}{HTML}{C9D2E8}
\definecolor{spLightBlue}{HTML}{E3F2FD}
\definecolor{spLightGreen}{HTML}{E8F5E9}
\definecolor{spLightPurple}{HTML}{EDE7F6}
\definecolor{spLightOrange}{HTML}{FBE9E7}
\definecolor{spLightTeal}{HTML}{E0F2F1}
\definecolor{spLightYellow}{HTML}{FFF9C4}
\definecolor{spYellow}{HTML}{F9A825}
\definecolor{spOrange}{HTML}{E65100}

% ── حزم عامة ─────────────────────────────────────────────────────────
\usepackage{booktabs}
\usepackage{array}
\usepackage{tabularx}
\usepackage{longtable}
\usepackage{enumitem}
\usepackage{amsmath}
\usepackage{amssymb}
\usepackage{graphicx}
\usepackage{pdflscape}   % صفحات landscape لإدماج الديجرامات الرسمية العريضة
\usepackage{listings}
\usepackage{tikz}
\usetikzlibrary{arrows.meta,positioning,backgrounds,calc,fit}
% ── أنماط عُقَد الرسومات (تُستخدَم داخل \begin{LTR}) — ألوان متّسقة مع الكتاب ──
\tikzset{
  pbox/.style={rectangle,rounded corners=4pt,text width=2.7cm,minimum height=0.8cm,
    align=center,draw=spPrimary!65,fill=spLightBlue,font=\small\bfseries,text=spPrimary!80},
  sbox/.style={rectangle,rounded corners=3pt,text width=2.7cm,minimum height=0.72cm,
    align=center,draw=spPrimary!40,fill=spGrayBg,font=\small,text=spInk},
  gbox/.style={rectangle,rounded corners=4pt,text width=2.7cm,minimum height=0.8cm,
    align=center,draw=spQA!70,fill=spLightPurple,font=\small\bfseries,text=spQA!80},
  fbox/.style={rectangle,rounded corners=3pt,text width=2.7cm,minimum height=0.72cm,
    align=center,draw=spSec!65,fill=spLightOrange,font=\small,text=spSec!80},
  tbox/.style={rectangle,rounded corners=3pt,text width=2.7cm,minimum height=0.72cm,
    align=center,draw=spAccent!65,fill=spLightTeal,font=\small,text=spAccent!80},
  srcbox/.style={rectangle,rounded corners=4pt,text width=2.6cm,minimum height=0.66cm,
    align=center,draw=spWhy!60,fill=spLightGreen,font=\scriptsize,text=spWhy!80},
  dbbox/.style={rectangle,rounded corners=3pt,text width=2.1cm,minimum height=0.8cm,
    align=center,draw=spAccent!70,fill=spLightTeal,font=\small\bfseries,text=spAccent!80,
    double,double distance=1.4pt},
  dbox/.style={rectangle,rounded corners=2pt,text width=2.6cm,minimum height=0.7cm,
    align=center,draw=spYellow!85!black,fill=spLightYellow,font=\scriptsize\itshape,text=spOrange!90},
  wbox/.style={rectangle,rounded corners=3pt,text width=11.4cm,minimum height=0.68cm,
    align=center,draw=spPrimary!35,fill=spGrayBg,font=\small,text=spInk},
  actor/.style={rectangle,rounded corners=2pt,text width=2.2cm,minimum height=0.7cm,
    align=center,draw=spInk!45,fill=white,font=\small\bfseries,text=spInk},
  arr/.style={->,>=Stealth,semithick,color=spInk!55},
  arrd/.style={->,>=Stealth,semithick,dashed,color=spInk!40},
  arrb/.style={<->,>=Stealth,semithick,color=spInk!55},
  lbl/.style={font=\scriptsize,text=spInk!60},
}
\usepackage[framemethod=default]{mdframed}
\usepackage{titlesec}
\usepackage{fancyhdr}
\usepackage[hidelinks]{hyperref}

% ── أدوات نصية مساعدة ────────────────────────────────────────────────
% مصطلح إنجليزي داخل جملة عربية: اتجاه LTR + خط لاتيني نظيف
\newcommand{\en}[1]{\textenglish{#1}}
% معرّف كود/تقني بخط أحادي المسافة
\newcommand{\code}[1]{{\monofont #1}}

% ── تنسيق العناوين (RTL: الرقم على اليمين) ───────────────────────────
% أرقام الأقسام المركّبة (2.1) تُجبَر على الاتجاه LTR حتى تُقرأ صحيحة لا معكوسة
\renewcommand{\thesection}{\textenglish{\arabic{section}}}
\renewcommand{\thesubsection}{\textenglish{\arabic{section}.\arabic{subsection}}}
\renewcommand{\thesubsubsection}{\textenglish{\arabic{section}.\arabic{subsection}.\arabic{subsubsection}}}
\titleformat{\section}{\Large\bfseries\color{spPrimary}}{\thesection.}{0.5em}{}
\titleformat{\subsection}{\large\bfseries\color{spAccent}}{\thesubsection}{0.5em}{}
\titleformat{\subsubsection}{\normalsize\bfseries\color{spInk}}{\thesubsubsection}{0.5em}{}
\titlespacing*{\section}{0pt}{16pt}{6pt}
\titlespacing*{\subsection}{0pt}{11pt}{4pt}
\setcounter{secnumdepth}{3}
\setcounter{tocdepth}{2}
\setlist{topsep=2pt,itemsep=1pt,parsep=0pt,leftmargin=1.6em,rightmargin=0.2em}
\setlength{\parindent}{0pt}
\setlength{\parskip}{5pt plus 1pt minus 1pt}

% ── الترويسة والتذييل ────────────────────────────────────────────────
\pagestyle{fancy}\fancyhf{}
\renewcommand{\headrulewidth}{0.4pt}
\fancyhead[R]{\small\color{spPrimary}\textbf{ScholarPath}}
% الترويسة اليسرى: «الفصل N» — نتجنّب عرض عنوان القسم في الترويسة لأن العناوين
% المختلطة (عربي + إنجليزي بين أقواس) تتبعثر في وضع الـ mark مع bidi.
\fancyhead[L]{\small\color{spInk}\ifnum\value{section}>0 {\arabicfont الفصل}~\thesection\fi}
\fancyfoot[C]{\small\thepage}

% ── تنسيق الأكواد (listings) — تُستخدَم داخل بيئة english ────────────
\lstset{basicstyle=\monofont\small,backgroundcolor=\color{spGrayBg},
  frame=single,rulecolor=\color{spRule},breaklines=true,breakatwhitespace=true,
  keywordstyle=\color{spPrimary}\bfseries,commentstyle=\color{spWhy}\itshape,
  stringstyle=\color{spSec},columns=fullflexible,showstringspaces=false,
  xleftmargin=4pt,framexleftmargin=2pt,tabsize=2}
% ملاحظة: الأكواد تُكتب دائمًا داخل بيئة english يدويًّا (نمط مثبَّت مع RTL):
%   \begin{english}\begin{lstlisting}[language={[Sharp]C}] ... \end{lstlisting}\end{english}
% لا تستخدم lstnewenvironment مع english (يبتلع باقي المستند في وضع bidi).

% ── الجداول ──────────────────────────────────────────────────────────
\newcolumntype{L}[1]{>{\raggedright\arraybackslash}p{#1}}
\newcolumntype{C}[1]{>{\centering\arraybackslash}p{#1}}
\newcolumntype{Y}{>{\raggedright\arraybackslash}X}
\arrayrulecolor{spRule}
\renewcommand{\arraystretch}{1.28}
\setlength{\tabcolsep}{5pt}

% ── الصناديق الملوّنة (mdframed آمن مع RTL — لا تستخدم tcolorbox) ────
% نمط أساسي مشترك لكل الصناديق المعنونة
\global\mdfdefinestyle{spbase}{%
  linewidth=0.8pt,roundcorner=3pt,
  innerleftmargin=9pt,innerrightmargin=9pt,
  innertopmargin=7pt,innerbottommargin=7pt,
  skipabove=7pt,skipbelow=7pt,
  frametitlefont=\bfseries\color{white},
  frametitlerule=false,
  frametitleaboveskip=4pt,frametitlebelowskip=4pt}

% مولّد بيئة صندوق: #1 مفتاح لوني نصّي، #2 خلفية، #3 خلفية العنوان، #4 عنوان افتراضي
% كل بيئة تقبل عنوانًا اختياريًا لتخصيص السطر (مثل «الفكرة: RAG»)
\newenvironment{idea}[1][الفكرة]%
  {\begin{mdframed}[style=spbase,linecolor=spHow!75!black,backgroundcolor=spLightBlue,
     frametitlebackgroundcolor=spHow,frametitle={\strut #1}]}%
  {\end{mdframed}}
\newenvironment{why}[1][ليه بالطريقة دي؟]%
  {\begin{mdframed}[style=spbase,linecolor=spWhy!70!black,backgroundcolor=spLightGreen,
     frametitlebackgroundcolor=spWhy!80!black,frametitle={\strut #1}]}%
  {\end{mdframed}}
\newenvironment{keyp}[1][نقطة مهمة]%
  {\begin{mdframed}[style=spbase,linecolor=spQA!70!black,backgroundcolor=spLightPurple,
     frametitlebackgroundcolor=spQA!80!black,frametitle={\strut #1}]}%
  {\end{mdframed}}
\newenvironment{secnote}[1][زاوية الأمان]%
  {\begin{mdframed}[style=spbase,linecolor=spSec!75!black,backgroundcolor=spLightOrange,
     frametitlebackgroundcolor=spSec!80!black,frametitle={\strut #1}]}%
  {\end{mdframed}}
\newenvironment{qa}[1][سؤال متوقع في المناقشة]%
  {\begin{mdframed}[style=spbase,linecolor=spQA!70!black,backgroundcolor=spLightPurple,
     frametitlebackgroundcolor=spQA!70!black,frametitle={\strut #1}]}%
  {\end{mdframed}}
\newenvironment{note}[1][ملاحظة]%
  {\begin{mdframed}[style=spbase,linecolor=spAccent!70!black,backgroundcolor=spLightTeal,
     frametitlebackgroundcolor=spAccent!80!black,frametitle={\strut #1}]}%
  {\end{mdframed}}
% صندوق «ملخّص الأفكار بالعربي» — الأبرز؛ حدّ أعرض
\newenvironment{summary}[1][الخلاصة في سطور]%
  {\begin{mdframed}[style=spbase,linewidth=1.1pt,linecolor=spPrimary,
     backgroundcolor=spGrayBg,frametitlebackgroundcolor=spPrimary,frametitle={\strut #1}]}%
  {\end{mdframed}}
% صندوق «ركّز هنا لإبهار اللجنة» — تمييز ذهبيّ بعنوان داكن لتباين واضح
\newenvironment{spotlight}[1][ركّز هنا لإبهار اللجنة]%
  {\begin{mdframed}[style=spbase,linewidth=1.1pt,linecolor=spYellow!85!black,
     backgroundcolor=spLightYellow,frametitlebackgroundcolor=spYellow,
     frametitlefont=\bfseries\color{spInk},frametitle={\strut #1}]}%
  {\end{mdframed}}

% صندوق العنوان (الغلاف) لكل مستند
\newcommand{\coverbox}[2]{%
  \begin{mdframed}[style=spbase,linewidth=0pt,roundcorner=4pt,
     backgroundcolor=spPrimary,innertopmargin=12pt,innerbottommargin=12pt]%
  {\color{white}\Huge\bfseries ScholarPath}\\[3pt]
  {\color{white}\LARGE #1}\\[6pt]
  {\color{white}\normalsize #2}%
  \end{mdframed}}


\begin{document}

% ── الغلاف ───────────────────────────────────────────────────────────
\coverbox{الدليل الشارح للنظام — من الصفر}{%
شرحٌ عربيٌّ متدرّج لِمنصّة \en{ScholarPath} كاملةً: المعمارية، الأمان،
الذكاء الاصطناعي، الجلسات المرئية، والمدفوعات — من غير أيّ خلفية سابقة عن
النظام. كلّ فكرة مشروحة واحدةً واحدة، ومعها قِطَعٌ حقيقيّة من كود التطبيق.}

\vspace{6pt}
{\small\color{spInk}
\textbf{المنصّة باختصار:} \en{.NET 10} خلفيّة \en{(backend)} $\cdot$
\en{React 19} واجهة \en{(frontend)} $\cdot$ \en{SQL Server} قاعدة بيانات $\cdot$
\en{Azure} سحابة $\cdot$ ثنائيّة اللغة \en{EN/AR} باتّجاه \en{RTL} كامل.}

\vspace{4pt}\hrule height 0.4pt\vspace{8pt}

% ── عن هذا الكتاب + مسار الثلاث ساعات ───────────────────────────────
\begin{summary}[كيف تقرأ هذا الكتاب — مسار الـ 3 ساعات]
الكتاب مُرتَّب كَـ«رحلة» من الصفر. لو مُبتدئ تمامًا اقرأه بالترتيب؛ كلّ فصل
يبني على اللي قبله. الجدول ده تقسيمٌ زمنيٌّ مُقترَح:
\begin{center}
\begin{tabular}{C{1.7cm} L{7.4cm} C{2.4cm}}
\toprule
\textbf{الفصل} & \textbf{الموضوع} & \textbf{زمن مُقترَح} \\
\midrule
1 & من الصفر: يعني إيه نظام برمجيّ أصلًا؟ & 20 دقيقة \\
2--5 & المعمارية و\en{CQRS} والتوثيق والأمان & 60 دقيقة \\
6--7 & الذكاء الاصطناعي: المفاهيم ثمّ البناء & 45 دقيقة \\
8--9 & الجلسات المرئية والإشعارات والشات (زمن حقيقيّ) & 25 دقيقة \\
10--11 & المدفوعات والتحليلات والخلاصة & 15 دقيقة \\
12 & دليل المناقشة والعرض (تحضير الـ \en{defense}) & 15 دقيقة \\
\bottomrule
\end{tabular}
\end{center}
كلّ فصل يُغلَق بصندوق \textbf{«الخلاصة في سطور»} — لو مستعجل اقرأ صناديق
الخلاصة الأول، وبعدين ارجع للتفاصيل.
\end{summary}

\vspace{4pt}
{\small\tableofcontents}
\vspace{4pt}\hrule height 0.4pt

\clearpage
% ── الفصول ───────────────────────────────────────────────────────────
% ch00-primer
\section{من الصفر: يعني إيه نظام برمجيّ؟}

أهلًا بك في أوّل فصل من الكتاب. لو أنت لا تعرف أيّ شيء عن البرمجة، فأنت في المكان
الصحيح تمامًا. هدفنا هنا بسيط: أن تُغلق هذا الفصل وأنت فاهم \emph{إيه هو} نظام
\en{ScholarPath}، ومن أيّ قطع كبيرة يتكوّن، وكيف تتكلّم هذه القطع مع بعضها. مش هنغرق
في تفاصيل تقنية عميقة دلوقتي؛ ده مجرّد تمهيد نبني عليه باقي الفصول. خُد نفَسًا، وتعال
نمشي \emph{واحدة واحدة}.

\subsection{إيه هو \en{ScholarPath} أصلًا؟}

تخيّل طالبًا مصريًّا في سنته الأخيرة بالجامعة، عايز يكمّل دراسته بمنحة في الخارج. هو
يعرف أنّ المنح موجودة، لكنّه \emph{تايه}: المنح متناثرة على عشرات المواقع بلغات
مختلفة، ومواعيد التقديم متضاربة، والأوراق المطلوبة غامضة، ولا يوجد من يرشده هل هو
مؤهَّل أصلًا لهذه المنحة أم لا. النتيجة المعتادة: يضيع الميعاد، أو يقدّم بشكل ناقص،
أو ييأس ويستسلم.

\begin{idea}[الفكرة: ما هو \en{ScholarPath}؟]
\en{ScholarPath} هو \textbf{منصّة إلكترونية} (موقع ويب) تجمع رحلة المنحة الدراسية
كلّها في مكان واحد. المنصّة \textbf{تكتشف} المنح المناسبة للطالب، و\textbf{تتابع} خطوات
تقديمه عليها، وتتيح له \textbf{حجز جلسة فيديو} مع استشاريّ متخصّص (جلسة مدفوعة)، وفيها
\textbf{مجتمع} يسأل فيه ويتناقش، و\textbf{مساعد ذكاء اصطناعيّ} يجاوب على أسئلته.
والمنصّة \textbf{ثنائية اللغة}: تعمل بالعربية (من اليمين لليسار) وبالإنجليزية.
\end{idea}

المنصّة تخدم \textbf{أربعة أنواع من المستخدمين}، ونسمّي كلّ نوع \en{«role»} أي «دور».
كلّ دور يرى شاشات مختلفة وله صلاحيّات مختلفة:

\begin{center}
\begin{tabular}{L{3.2cm} L{9cm}}
\toprule
\textbf{\color{spPrimary}الدور} & \textbf{\color{spPrimary}بيعمل إيه على المنصّة} \\
\midrule
الطالب \en{(Student)} & يبحث عن المنح، يتابع تقديماته، يحجز جلسات استشاريّ، يشارك في المجتمع، ويسأل المساعد الذكيّ. \\
الاستشاريّ \en{(Consultant)} & خبير يعرض أوقاته المتاحة، يستقبل الحجوزات، ويقدّم جلسات فيديو مدفوعة للطلاب. \\
الجهة المانحة \en{(Scholarship Provider)} & تنشر منحها على المنصّة، وتراجع طلبات الطلاب الخاصّة بها. \\
الأدمن \en{(Admin)} & المشرف الذي يوافق على المحتوى، يدير المستخدمين، ويراقب سلامة المنصّة. \\
\bottomrule
\end{tabular}
\end{center}

\begin{keyp}[نقطة مهمة: حجم النظام]
\en{ScholarPath} ليس تطبيقًا صغيرًا. هو مقسَّم إلى \textbf{18 وحدة كبرى} (نسمّيها
\en{modules} أو \en{epics}) مرقّمة من \en{PB-001} إلى \en{PB-018}، تغطّي الرحلة كاملة
من إنشاء الحساب حتى تحليلات البيانات. وبلغة الأرقام: 211 \textbf{متطلَّبًا وظيفيًّا}
(\en{functional requirement})، و203 \textbf{حالة استخدام} (\en{use case})،
و159 \textbf{قصّة مستخدم} (\en{user story}). لا تقلق من هذه المصطلحات الآن — هي مجرّد
طرق منظَّمة لكتابة «النظام لازم يعمل كذا».
\end{keyp}

\subsection{يعني إيه برنامج ويب؟ الصالة والمطبخ والمخزن}

أيّ موقع ويب تتعامل معه (فيسبوك، بنكك، \en{ScholarPath}) يتكوّن في الحقيقة من
\textbf{ثلاث قطع كبيرة} تتعاون. أسهل طريقة لفهمها هي تشبيه \textbf{المطعم}:

\begin{center}
\begin{tabular}{L{3.4cm} L{4.3cm} L{4.6cm}}
\toprule
\textbf{\color{spPrimary}في المطعم} & \textbf{\color{spPrimary}في النظام} & \textbf{\color{spPrimary}وظيفته} \\
\midrule
الصالة والجرسون & الواجهة \en{(Frontend)} & ما تراه وتلمسه في المتصفّح: الأزرار، الشاشات، النماذج. \\
المطبخ & الخلفية \en{(Backend)} & يستقبل «الطلبات»، ينفّذ قواعد العمل، ويجهّز النتيجة. \\
المخزن والثلّاجة & قاعدة البيانات \en{(Database)} & تخزّن كلّ شيء بشكل دائم: المستخدمون، المنح، الحجوزات. \\
\bottomrule
\end{tabular}
\end{center}

\begin{idea}[الفكرة: كيف تتعاون القطع الثلاث؟]
عندما تطلب أكلة، أنت لا تدخل المطبخ ولا المخزن بنفسك. تقول للجرسون (الواجهة) ما تريد،
الجرسون ينقل الطلب للمطبخ (الخلفية)، المطبخ يأخذ المكوّنات من المخزن (قاعدة البيانات)،
يطبخها، ثمّ يعيد الطبق للجرسون الذي يضعه أمامك. \textbf{الواجهة} تعيش داخل متصفّحك على
جهازك، بينما \textbf{الخلفية} و\textbf{قاعدة البيانات} تعيشان على حاسوب بعيد قويّ اسمه
\en{«server»} (الخادم).
\end{idea}

\begin{note}
لماذا نفصل هذا الفصل؟ لأنّ لكلّ قطعة أدوات مختلفة. واجهة \en{ScholarPath} مكتوبة بلغة
\en{React} وتعمل في المتصفّح، وخلفيّته مكتوبة بلغة \en{C\#} على منصّة \en{.NET}،
وقاعدة بياناته هي \en{SQL Server}. سنشرح كلّ اسم من هذه بعد قليل.
\end{note}

\subsection{يعني إيه \en{API}؟ ورحلة ضغطة زرّ واحدة}

طيب، بأيّ «لغة» يتكلّم الجرسون مع المطبخ؟ في عالم الويب، الواجهة والخلفية يتكلّمان
عبر ما يُسمّى \en{API}.

\begin{idea}[الفكرة: ما هو \en{API}؟]
\en{API} اختصار لـ \en{Application Programming Interface}، أي \textbf{«واجهة برمجية»}.
تخيّلها \textbf{قائمة الطعام} في المطعم: قائمة محدَّدة من الطلبات المسموح بها، ولكلّ طلب
شكل معروف. الواجهة لا تكلّم قاعدة البيانات مباشرة أبدًا؛ هي فقط تطلب من الخلفية عبر هذه
القائمة. مثلًا: «هات لي المنح صفحة رقم 1»، أو «احجز الجلسة الفلانية».
\end{idea}

وكلّ طلب-وردّ بينهما يمشي على بروتوكول اسمه \en{HTTP}. الطلب اسمه \en{request}،
والردّ اسمه \en{response}. لِنتتبَّع معًا \textbf{رحلة ضغطة زرّ واحدة}: الطالب فتح صفحة
المنح وضغط «بحث».

\begin{enumerate}
\item \textbf{المتصفّح (الواجهة)} يجهّز \en{request} ويرسله للخلفية: «\en{GET} المنح».
\item \textbf{الخلفية} تستقبل الطلب، وتتأكّد أنّ المستخدم مسجَّل الدخول (صلاحيّات).
\item الخلفية تسأل \textbf{قاعدة البيانات}: «هات المنح المنشورة، صفحة 1».
\item قاعدة البيانات ترجع الصفوف، والخلفية تحوّلها إلى \en{response} بصيغة \en{JSON}.
\item المتصفّح يستقبل الردّ ويرسم بطاقات المنح على الشاشة. كلّ ذلك في أقلّ من ثانية.
\end{enumerate}

\begin{note}[ما هو \en{JSON}؟]
\en{JSON} هي طريقة قياسية لكتابة البيانات كنصّ منظَّم بين أقواس، يفهمها الطرفان بسهولة.
هي «لغة التغليف» التي يسافر بها الردّ من الخلفية إلى الواجهة.
\end{note}

خلّينا نشوف الشكل التوضيحيّ للطلب والردّ. \textbf{المقتطف التالي توضيحيّ فقط} (مبسَّط
لنفهم الشكل، وليس منسوخًا من ملفّ):

\begin{english}
\begin{lstlisting}[language={}]
// Request — the browser asks the server for scholarships
GET /api/scholarships?page=1&fundedOnly=true
Accept-Language: ar

// Response — the server replies (200 OK) with JSON
{
  "items": [
    { "id": "a1b2-3c4d", "title": "Chevening", "fundingType": "FullyFunded" }
  ],
  "totalCount": 42
}
\end{lstlisting}
\end{english}

هنشوف إيه في المقتطف ده؟ السطر الأوّل \code{GET /api/scholarships} هو \emph{الطلب}: كلمة
\en{GET} تعني «هات لي» (قراءة بيانات)، والباقي هو العنوان مع فلاتر البحث. السطر
\code{Accept-Language: ar} يقول للخلفية «أنا عايز الردّ بالعربي». ثمّ يأتي \emph{الردّ}:
الرقم \code{200} معناه «تمام، نجح الطلب»، وبعده كتلة \en{JSON} فيها قائمة \code{items}
(المنح) وعددها الكلّي \code{totalCount}.

الآن نرى \textbf{الكود الحقيقيّ} من المشروع الذي يقوم بهذا الطلب. أوّلًا من جهة
\textbf{الواجهة}: هذا المقتطف من ملفّ الخدمة \code{client/src/services/api/scholarships.ts}
الذي يرسل الطلب فعلًا:

\begin{english}
\begin{lstlisting}[language=Java]
const { data } = await apiClient.get<PaginatedListWireDto<ScholarshipWireDto>>(
  "/api/scholarships",
  { params },
);
return {
  items: data.items.map(toListItem),
  total: data.totalCount,
  totalPages: data.totalPages,
};
\end{lstlisting}
\end{english}

هنشوف إيه هنا؟ الدالة \code{apiClient.get} هي «الجرسون» الذي يذهب للخلفية على العنوان
\code{"/api/scholarships"} ومعه فلاتر البحث \code{params}. كلمة \code{await} معناها «استنّى
الردّ يرجع». عندما يعود الردّ نأخذ منه \code{data.items} (قائمة المنح) ونجهّزها لعرضها.

وفي \textbf{الجهة المقابلة} على الخلفية، هذا هو الجزء الذي \emph{يستقبل} الطلب، من ملفّ
\code{ScholarshipsController.cs}:

\begin{english}
\begin{lstlisting}[language={[Sharp]C}]
[HttpGet]
public async Task<ActionResult<PaginatedList<ScholarshipDto>>> Get(
    [FromQuery] GetScholarshipsQuery query)
{
    // read the requested language from the Accept-Language header
    var headerLang = Request.Headers["Accept-Language"].ToString()
        .Split(',').FirstOrDefault() ?? "en";
    var lang = headerLang.StartsWith("ar",
        StringComparison.OrdinalIgnoreCase) ? "ar" : "en";

    var updatedQuery = query with { Language = lang };
    return await mediator.Send(updatedQuery);
}
\end{lstlisting}
\end{english}

هنشوف إيه هنا؟ السطر \code{[HttpGet]} يقول «هذه الدالة تردّ على طلبات \en{GET}». الدالة
اسمها \code{Get} وتستقبل شروط البحث في \code{query}. السطرين في الوسط يقرآن ترويسة
\code{Accept-Language} ليعرفا هل يردّ بالعربي أم الإنجليزي. أخيرًا \code{mediator.Send}
يمرّر الطلب للطبقة الداخلية التي تجلب البيانات وترجعها. لاحظ التطابق الجميل: الواجهة
نادت \code{"/api/scholarships"}، وهنا \code{[HttpGet]} على نفس العنوان يردّ عليها.

\begin{secnote}
لاحظ سطر \code{Request.Headers} وقبله التحقّق من تسجيل الدخول (الصفة \en{[Authorize]} على
هذا المتحكِّم). في \en{ScholarPath} \emph{لا يُقدَّم أيّ شيء دون التأكّد من هوية المستخدم
وصلاحيّته} — قاعدة أمان أساسية اسمها «البوّابة المغلقة» \en{(gated)}.
\end{secnote}

\subsection{يعني إيه «طبقة» \en{(layer)}؟ ولماذا نقسّم الكود؟}

لاحظت في المقتطف السابق أنّ المتحكِّم لم يجلب البيانات بنفسه، بل «مرّرها للطبقة
الداخلية». هذه فكرة محورية في \en{ScholarPath} اسمها \textbf{الطبقات} \en{(layers)}.

\begin{idea}[الفكرة: لماذا طبقات؟]
بدل أن نكتب كلّ شيء (شكل الشاشة + قواعد العمل + كلام قاعدة البيانات) في مكان واحد
متشابك، نقسّم الكود إلى \textbf{طبقات} لكلٍّ منها مسؤولية واحدة واضحة. تخيّل المطعم مرة
أخرى: الجرسون لا يطبخ، والطبّاخ لا يزرع الخضار. هذا الفصل يجعل الكود \textbf{أسهل في
الفهم، والاختبار، والتعديل} دون كسر باقي أجزاء النظام.
\end{idea}

خلفيّة \en{ScholarPath} مقسَّمة إلى \textbf{أربع طبقات}، والقلب فيها هو طبقة \code{Domain}.
هذا مقتطف حقيقيّ من قلب النظام، من ملفّ \code{ScholarPath.Domain/Entities/Scholarships.cs}،
يمثّل «المنحة» ككائن أعمال خالص:

\begin{english}
\begin{lstlisting}[language={[Sharp]C}]
public class Scholarship : AuditableEntity, ISoftDeletable
{
    public string TitleEn { get; set; } = default!;
    public string TitleAr { get; set; } = default!;   // bilingual title
    public ScholarshipStatus Status { get; set; } = ScholarshipStatus.Draft;
    public DateTimeOffset Deadline { get; set; }
    public FundingType FundingType { get; set; }
    public decimal? FundingAmountUsd { get; set; }
    // ... more properties omitted
}
\end{lstlisting}
\end{english}

هنشوف إيه هنا؟ هذا الكائن \code{Scholarship} يصف «ما هي المنحة»: لها عنوان بالعربي
\code{TitleAr} وبالإنجليزي \code{TitleEn}، وحالة \code{Status} (مسودّة، منشورة...)، وميعاد
نهائيّ \code{Deadline}، ونوع تمويل \code{FundingType}. المهمّ: هذا الكائن \textbf{لا يعرف
أيّ شيء} عن قاعدة البيانات ولا عن الشاشات — هو قواعد عمل خالصة. سنتعمّق في هذا التصميم
(اسمه \en{Clean Architecture}) في الفصل التالي؛ يكفيك الآن أن تعرف \emph{أنّ} الكود
مقسَّم لطبقات، وأنّ \code{Domain} هي القلب.

\subsection{يعني إيه سحابة \en{(cloud)} و\en{Azure}؟}

قلنا إنّ الخلفية وقاعدة البيانات تعيشان على «خادم بعيد». لكن مَن يملك هذا الخادم؟
هنا تأتي فكرة \textbf{السحابة} \en{(cloud)}.

\begin{idea}[الفكرة: ما هي السحابة؟]
«السحابة» ببساطة تعني \textbf{حواسيب شركة كبيرة نستأجرها عبر الإنترنت}، بدل أن نشتري
خوادم ونضعها في غرفة عندنا. ندفع مقابل ما نستخدمه فقط، وهم يتولّون الكهرباء والتبريد
والصيانة. \en{ScholarPath} يعيش على سحابة \textbf{مايكروسوفت} اسمها \en{Azure}.
\end{idea}

نستخدم من \en{Azure} عدّة خدمات، كلمة واحدة لكلٍّ منها: \en{App Service} (يشغّل
الخلفية)، \en{SQL Database} (قاعدة البيانات)، \en{Blob Storage} (تخزين الملفّات
المرفوعة)، \en{Azure OpenAI} (محرّك الذكاء الاصطناعيّ)، و\en{Key Vault} (خزنة كلمات
السرّ). لا داعي لحفظها الآن؛ ستقابلها تباعًا في فصول لاحقة.

\subsection{مكوّنات المنصّة: مَن يعمل ماذا}

جمعنا في الجدول التالي \textbf{أهمّ التقنيات} التي يقوم عليها \en{ScholarPath}، بسطر
مبسَّط لكلّ واحدة. ستراها مرارًا في الكتاب، فلا بأس أن تعود لهذا الجدول كلّما نسيت:

\begin{center}
\begin{tabular}{L{3cm} L{9.2cm}}
\toprule
\textbf{\color{spPrimary}التقنية} & \textbf{\color{spPrimary}بتعمل إيه} \\
\midrule
\en{.NET 10} & المنصّة (بلغة \en{C\#}) التي تُبنى عليها الخلفية كلّها. \\
\en{React 19} & المكتبة (بلغة \en{TypeScript}) التي تُبنى بها الواجهة في المتصفّح. \\
\en{SQL Server} & قاعدة البيانات التي تخزّن كلّ شيء في جداول. \\
\en{Azure OpenAI} & محرّك الذكاء الاصطناعيّ الذي يشغّل المساعد الذكيّ ويجاوب الأسئلة. \\
\en{Stripe} & بوّابة الدفع الآمنة التي تُحصِّل ثمن جلسات الاستشاريّ بالبطاقة. \\
\en{SignalR} & القناة الفورية للتحديثات اللحظية: الدردشة والإشعارات دون تحديث الصفحة. \\
\en{Hangfire} & مدبّر المهام الخلفية المجدولة (تذكيرات، تنظيف، معالجة مؤجَّلة). \\
\bottomrule
\end{tabular}
\end{center}

\subsection{بنية المشروع على القرص}

كفاية نظريّة، تعال نفتح المشروع فعلًا. الكود مقسَّم إلى جزأين كبيرين على القرص:
مجلَّد \code{server/} للخلفية، ومجلَّد \code{client/} للواجهة. هذه هي شجرة \textbf{الخلفية}
الحقيقية، وفيها الطبقات الأربع التي تكلّمنا عنها:

\begin{english}
\begin{lstlisting}[language={}]
server/src/
|-- ScholarPath.Domain/           # the core: entities + business rules
|   |-- Entities/  Enums/  Events/  Interfaces/  Exceptions/
|-- ScholarPath.Application/       # use cases: Commands / Queries
|   |-- Scholarships/  Applications/  ConsultantBookings/  AI/  ...
|-- ScholarPath.Infrastructure/    # database + external services
|   |-- Persistence/  Migrations/  Services/  Hubs/  Jobs/
`-- ScholarPath.API/               # entry point: Controllers + Middleware
    |-- Controllers/  Middleware/  Program.cs  appsettings.json
\end{lstlisting}
\end{english}

هنشوف إيه هنا؟ القراءة من الداخل للخارج: \code{Domain} هي القلب (كائنات مثل
\code{Scholarship} الذي رأيناه). \code{Application} فيها «حالات الاستخدام» مقسَّمة حسب
الموديول (مجلَّد \code{Scholarships}، مجلَّد \code{ConsultantBookings}...). \code{Infrastructure}
تتكلّم مع قاعدة البيانات والخدمات الخارجية. و\code{API} هي المدخل الذي تصل إليه طلبات
الويب — وفيها مجلَّد \code{Controllers} حيث يعيش المتحكِّم الذي رأيناه.

وهذه شجرة \textbf{الواجهة} (\code{client/src}) باختصار:

\begin{english}
\begin{lstlisting}[language={}]
client/src/
|-- pages/         # one folder per role: student/ consultant/ company/ admin/
|-- components/    # reusable UI pieces (buttons, cards, forms)
|-- services/      # api/ + signalR/  (these talk to the backend)
|-- hooks/  stores/  routes/  locales/  theme/  types/
`-- App.tsx  main.tsx
\end{lstlisting}
\end{english}

هنشوف إيه هنا؟ مجلَّد \code{pages} فيه مجلَّد لكلّ دور (\code{student}، \code{consultant}...).
مجلَّد \code{components} فيه القطع القابلة لإعادة الاستخدام. والأهمّ لموضوعنا: مجلَّد
\code{services} فيه \code{api} — ومنه جاء المقتطف \code{apiClient.get} الذي يكلّم الخلفية.
مجلَّد \code{locales} يحمل ترجمات العربية والإنجليزية.

\begin{qa}
\textbf{سؤال متوقَّع:} ليه الواجهة والخلفية في مجلَّدين منفصلين وبلغتين مختلفتين؟
\newline
\textbf{الإجابة:} لأنّهما يعملان في بيئتين مختلفتين — الواجهة داخل متصفّح المستخدم
والخلفية على الخادم — ولكلٍّ منهما أدواته المثلى. الفصل يسمح لفريق الواجهة وفريق
الخلفية بالعمل بالتوازي، ويسمح باستبدال أيّ جهة دون كسر الأخرى، طالما بقيت لغة الحوار
بينهما (\en{API}) ثابتة.
\end{qa}

\subsection{إزاي تقرأ الكتاب ده — خريطة الـ 3 ساعات}

الكتاب 10 فصول، ومصمَّم بحيث تفهم النظام كاملًا في نحو \textbf{3 ساعات}. لست مضطرًّا
لحفظ أيّ شيء؛ فقط اقرأ بالترتيب، وتوقّف عند الصناديق الملوّنة — فهي تلخّص أهمّ الأفكار.
إليك خريطة تقريبية للوقت:

\begin{center}
\begin{tabular}{L{4.5cm} L{5cm} C{2.4cm}}
\toprule
\textbf{\color{spPrimary}المرحلة} & \textbf{\color{spPrimary}تتعلّم فيها} & \textbf{\color{spPrimary}وقت تقريبيّ} \\
\midrule
هذا الفصل (0) & الصورة الكبيرة والمصطلحات & 20 دقيقة \\
البنية والطبقات & \en{Clean Architecture} والتدفّق & 40 دقيقة \\
الموديولات الأساسية & المنح، التقديم، الحجز، الدفع & 70 دقيقة \\
الميزات المتقدّمة & الذكاء الاصطناعيّ، الفوريّة، التحليلات & 40 دقيقة \\
مراجعة وربط & الأمان والصورة النهائية & 10 دقائق \\
\bottomrule
\end{tabular}
\end{center}

\begin{note}[كيف تستفيد من الصناديق]
الصندوق \textbf{الأزرق} يشرح فكرة، و\textbf{الأخضر} يبرّر قرار تصميم، و\textbf{البنفسجيّ}
نقطة أساسية أو سؤال مناقشة، و\textbf{البرتقاليّ} زاوية أمان. لو كنت مستعجلًا، اقرأ
العناوين والصناديق فقط — تكفي لأخذ فكرة سريعة.
\end{note}

\begin{summary}[الخلاصة في سطور]
\begin{itemize}
\item \en{ScholarPath} منصّة ويب ثنائية اللغة تحلّ تِيه الطالب في عالم المنح: اكتشاف،
متابعة تقديم، حجز جلسات استشاريّ بالفيديو (مدفوعة)، مجتمع، ومساعد ذكاء اصطناعيّ.
\item أربعة أدوار: الطالب، الاستشاريّ، الجهة المانحة، والأدمن — كلٌّ يرى ما يخصّه.
\item أيّ نظام ويب = \textbf{واجهة} في المتصفّح + \textbf{خلفية} على الخادم +
\textbf{قاعدة بيانات} للتخزين (الصالة + المطبخ + المخزن).
\item الواجهة والخلفية يتحاوران عبر \en{API} برسائل \en{HTTP}: \en{request} يخرج،
و\en{response} بصيغة \en{JSON} يعود — كلّ ذلك في أقلّ من ثانية عند ضغطة زرّ.
\item الكود مقسَّم إلى \textbf{طبقات}؛ قلب الخلفية هو طبقة \code{Domain} التي لا تعرف
شيئًا عن الشاشات أو قاعدة البيانات. تفاصيل ذلك في الفصل التالي.
\item يعيش النظام على سحابة \en{Azure}، ويقوم على \en{.NET 10} و\en{React 19}
و\en{SQL Server}، بمساعدة \en{Stripe} و\en{SignalR} و\en{Hangfire} و\en{Azure OpenAI}.
\item على القرص: \code{server/src} فيه الطبقات الأربع، و\code{client/src} فيه الواجهة.
\end{itemize}
\end{summary}

\clearpage
% ch01-clean-architecture — المعماريّة النظيفة (Clean Architecture)
\section{المعماريّة النظيفة (Clean Architecture)}

في الفصل السابق اتّفقنا على المصطلحات الأساسية: الخلفية \en{(backend)}، والواجهة
\en{(frontend)}، وقاعدة البيانات، والـ \en{API}، وفكرة «الطبقة» بشكل عام. في هذا
الفصل نفتح صندوق الخلفية ونشوف كيف اتبنى من الداخل. الخلفية في \en{ScholarPath}
مكتوبة بـ \en{.NET 10}، ومنظَّمة بأسلوب اسمه \en{Clean Architecture} — أو
«المعماريّة النظيفة». وده أهم قرار هندسيّ في المشروع كلّه، فخلّينا نفهمه واحدة واحدة.

الفكرة ببساطة: بدل ما نكتب الكود كلّه مكوَّمًا في مكان واحد، نقسّمه إلى \textbf{أربع
طبقات متداخلة} زيّ حلقات دائرة. كل طبقة لها مسؤولية واحدة واضحة، والأهمّ: قاعدة صارمة
واحدة تحكم الاعتماد بينهم.

\begin{idea}[الفكرة: الاعتماد يروح للداخل فقط]
تخيّل أربع دوائر متداخلة. في القلب توجد \textbf{قواعد العمل} (منطق المِنَح
والحجوزات والدفع)، وفي الخارج توجد \textbf{التفاصيل التقنية} (قاعدة البيانات،
\en{Azure}، \en{Stripe}). القاعدة الذهبية: \textbf{الاعتماد يتّجه للداخل دائمًا}.
الطبقة الداخلية لا تعرف شيئًا عن الطبقة الخارجية. القلب (منطق العمل) لا يعرف أنّ
قاعدة البيانات اسمها \en{SQL Server}، ولا أنّ مزوّد الذكاء الاصطناعي اسمه
\en{Azure OpenAI}. هو يعرف \emph{ماذا} يحتاج، لا \emph{كيف} يُنفَّذ.
\end{idea}

\subsection{الطبقات الأربع ومسؤولية كل واحدة}

النظام مبنيّ على أربعة مشاريع \en{(.csproj)} منفصلة، كل واحد منها طبقة. ترتيبها من
الداخل للخارج هكذا: \en{Domain} ثم \en{Application} ثم \en{Infrastructure} ثم
\en{API}. اتجاه الاعتماد يمشي في العكس — من الخارج للداخل:

\begin{center}\begin{LTR}
\en{API} $\longrightarrow$ \en{Infrastructure} $\longrightarrow$ \en{Application} $\longrightarrow$ \en{Domain}
\end{LTR}\end{center}

يعني \en{API} يعتمد على مَن تحته، و\en{Domain} في القلب لا يعتمد على أحد إطلاقًا.
ده مضمون فعليًّا في الكود عن طريق مراجع المشاريع \en{(project references)}: مشروع
\code{ScholarPath.Domain} لا يحمل أي مرجع لأي مشروع آخر، ومشروع
\code{ScholarPath.Application} يشير إلى \code{ScholarPath.Domain} فقط. المُصرِّف
\en{(compiler)} نفسه يرفض أي محاولة لكسر هذا الترتيب.

\begin{center}
\begin{tabular}{L{2.3cm} L{3.7cm} L{8.8cm}}
\toprule
\textbf{\color{spPrimary}الطبقة} & \textbf{\color{spPrimary}المشروع \en{(csproj)}} & \textbf{\color{spPrimary}المحتوى الأساسي} \\
\midrule
\en{Domain} & \code{ScholarPath.Domain} & الكيانات \en{(entities)} $\approx$55، والـ \en{enums}، وأحداث الدومين، والمنفذان \code{ICurrentUserService} و\code{IDateTimeService}. بلا أي اعتماد على \en{framework}. \\
\en{Application} & \code{ScholarPath.Application} & حالات الاستخدام كأوامر واستعلامات \en{CQRS} عبر \en{MediatR}، والـ \en{validators}، وقوالب \en{AutoMapper}، و$\approx$33 \en{port} (واجهة) يجب على العالم الخارجي تنفيذها. \\
\en{Infrastructure} & \code{ScholarPath.Infrastructure} & $\approx$46 \en{adapter} تنفّذ الـ \en{ports}: \en{EF Core}، \en{Stripe}، \en{Azure OpenAI}، \en{ACS}، \en{Blob}، \en{Key Vault}، \en{SignalR}، \en{Hangfire}. \\
\en{API} & \code{ScholarPath.API} & الـ \en{controllers} الرفيعة، ومستقبِل \en{Stripe webhook}، والهُويّة \en{(ASP.NET Identity + JWT)}. \\
\bottomrule
\end{tabular}
\end{center}

خلّينا ندخل على كل طبقة بالتفصيل، من القلب للخارج.

\subsection{طبقة \en{Domain} — القلب النقيّ}

هذه أعمق طبقة، وأهم قاعدة فيها: \textbf{صفر اعتماد على أي \en{framework}}. لا تعرف
\en{Entity Framework}، ولا \en{ASP.NET}، ولا \en{Azure}. هي مجرّد \en{C\#} خالص
يصف عالم المشكلة: كيان الطالب، كيان المِنحة، كيان الحجز، حالات الدفع كـ \en{enum}،
وأحداث الدومين \en{(domain events)} التي تُطلَق عند حدوث شيء مهمّ.

الاستثناء الوحيد أنّ الدومين يُعرِّف \textbf{منفذين} \en{(ports)} فقط، لأنّهما مفهومان
ينتميان لقلب العمل نفسه: مَن هو المستخدم الحاليّ؟ وكم الساعة الآن؟ خلّينا نشوفهما.

هنشوف إيه: الملف الحقيقيّ الذي يُعرِّف المنفذين داخل الدومين.

\begin{english}
\begin{lstlisting}[language={[Sharp]C}]
namespace ScholarPath.Domain.Interfaces;

public interface ICurrentUserService
{
    Guid? UserId { get; }
    string? Email { get; }
    string? ActiveRole { get; }
    IReadOnlyCollection<string> Roles { get; }
    bool IsAuthenticated { get; }
    string? IpAddress { get; }
    // ...
}

public interface IDateTimeService
{
    DateTimeOffset UtcNow { get; }
    DateOnly Today { get; }
}
\end{lstlisting}
\end{english}

الشرح سطرًا بسطر لأهم الأسطر:
\begin{itemize}
\item \code{ICurrentUserService}: منفذ يجيب على سؤال «مَن يُنفّذ هذا الطلب الآن؟» —
  مُعرِّفه \code{UserId}، بريده، ودوره النشِط \code{ActiveRole}. الدومين يحتاج هذه
  المعلومة لتطبيق قواعد الملكية (مثلًا: هل هذا الطالب يملك هذا الطلب؟)، لكنّه لا
  يعرف \emph{من أين} تأتي — من \en{JWT}؟ من \en{cookie}؟ لا يهمّه.
\item \code{IDateTimeService}: منفذ بسيط يعطي \code{UtcNow} و\code{Today}. لماذا لا
  نستدعي \code{DateTime.Now} مباشرة؟ لأنّ تجميد الوقت خلف منفذ يجعل الكود
  \textbf{قابلًا للاختبار}: في الاختبار نُثبّت «الآن» على لحظة معيّنة ونتأكّد أنّ
  منطق المواعيد النهائية يتصرّف صحّ.
\end{itemize}

\begin{why}
لماذا نُبقي الدومين نقيًّا بلا \en{framework}؟ لأنّ قواعد العمل هي أثمن ما نملك،
وأطول ما يعيش. الأُطُر التقنية \en{(frameworks)} تتغيّر: قد نستبدل \en{SQL Server}
بغيره، أو نُرقّي \en{.NET}. لو كانت قواعد العمل مشتبكة مع التفاصيل التقنية، فأي
تغيير تقنيّ يهدّد المنطق. عزلُها يعني أنّ منطق المِنَح يبقى ثابتًا مهما تبدّل ما
حوله — وده يجعله قابلًا للاختبار بلا قاعدة بيانات ولا شبكة.
\end{why}

\subsection{طبقة \en{Application} — حالات الاستخدام والمنافذ}

فوق الدومين تأتي طبقة \en{Application}. مسؤوليتها: تنسيق \textbf{حالات الاستخدام}
\en{(use cases)} — «الطالب يحجز جلسة»، «المستشار يقبل الحجز»، «أرشد لي منَحًا». كل
حالة استخدام مكتوبة على نمط \en{CQRS}.

\begin{keyp}[نقطة مهمة: \en{CQRS}]
\en{CQRS} اختصار لـ \en{Command Query Responsibility Segregation}. الفكرة: نفصل كل
عملية إلى نوعين. \textbf{أمر} \en{(Command)} يُغيّر الحالة (كتابة) مثل
\code{AcceptBookingCommand}. و\textbf{استعلام} \en{(Query)} يقرأ فقط بلا أي أثر
جانبيّ مثل \code{GetMyRecommendationsQuery}. كل أمر أو استعلام له مُعالِج
\en{(handler)} واحد مسؤول عنه، فيصبح كل مِلفّ صغيرًا ومملوكًا لمطوّر واحد.
\end{keyp}

هذه الطبقة تحتوي أيضًا على الـ \en{validators} (تتحقّق من صحّة المُدخَلات قبل أي
عمل) عبر \en{FluentValidation}، وقوالب \en{AutoMapper} (تحوّل الكيانات إلى
\en{DTOs} للعرض). لكن الجزء الأهمّ معماريًّا: هنا تُعرَّف الـ \textbf{$\approx$33
منفذًا} \en{(ports)}. المنفذ هو مجرّد \en{interface} — وعدٌ بخدمة يحتاجها المنطق دون
أن يعرف تنفيذها. مثال حيّ: منفذ الذكاء الاصطناعي.

هنشوف إيه: منفذ \code{IAiService} الحقيقيّ — عقد الذكاء الاصطناعي كما تراه طبقة العمل.

\begin{english}
\begin{lstlisting}[language={[Sharp]C}]
namespace ScholarPath.Application.Common.Interfaces;

public interface IAiService
{
    Task<AiRecommendationResult> GenerateRecommendationsAsync(
        Guid userId, int topN, CancellationToken ct);

    Task<AiEligibilityResult> CheckEligibilityAsync(
        Guid userId, Guid scholarshipId, CancellationToken ct);

    Task<AiChatResponse> AskAsync(
        Guid userId, string sessionId, string message,
        IReadOnlyList<AiChatHistoryTurn> history, CancellationToken ct);
}
\end{lstlisting}
\end{english}

الشرح: هذا الـ \en{interface} يصف \emph{ماذا} يقدّم الذكاء الاصطناعي في ثلاث
عمليات فقط: \code{GenerateRecommendationsAsync} للترشيحات،
و\code{CheckEligibilityAsync} لفحص الأهليّة، و\code{AskAsync} للمحادثة. لاحِظ أنّه
\textbf{لا يوجد ذِكر لـ \en{Azure} ولا \en{OpenAI} ولا أي مِفتاح}. طبقة العمل تعرف
أنّها تحتاج «خدمة ذكاء اصطناعي» بهذا الشكل، ولا تهتمّ من يقف خلفها. هذا هو معنى
«المنفذ»: عقدٌ نظيف يعيش في الداخل.

\subsection{طبقة \en{Infrastructure} — المحوّلات التي تُنفّذ الوعود}

كل منفذ في طبقة \en{Application} لا بدّ أن يجد مَن يُنفّذه في الواقع. هنا يأتي دور
طبقة \en{Infrastructure}: هي بيت الـ \textbf{$\approx$46 محوّلًا} \en{(adapters)}.
المحوّل هو الصف \en{(class)} الذي يقول «أنا أُنفّذ هذا المنفذ باستخدام تقنية
حقيقية». هنا يعيش \en{EF Core} (قاعدة البيانات)، و\en{Stripe} (الدفع)، و\en{Azure
OpenAI} (الذكاء الاصطناعي)، و\en{ACS} (الفيديو)، و\en{Blob} (تخزين الملفات)،
و\en{Key Vault} (المفاتيح السرّية)، و\en{SignalR} (الوقت الحقيقيّ)، و\en{Hangfire}
(المهام الخلفية).

هنشوف إيه: محوّل \code{AzureOpenAiService} الحقيقيّ الذي يُنفّذ منفذ \code{IAiService}
الذي رأيناه للتوّ.

\begin{english}
\begin{lstlisting}[language={[Sharp]C}]
public sealed class AzureOpenAiService(
    IHttpClientFactory httpFactory,
    IOptions<AiOptions> opts,
    IKnowledgeRetriever retriever,
    LocalAiService local,
    ILogger<AzureOpenAiService> logger) : IAiService
{
    // Recommendations & eligibility are deterministic -> delegate to the local engine
    public Task<AiRecommendationResult> GenerateRecommendationsAsync(
        Guid userId, int topN, CancellationToken ct)
        => local.GenerateRecommendationsAsync(userId, topN, ct);

    public async Task<AiChatResponse> AskAsync(...)
    {
        // If Azure endpoint/key are not set, answer with the local RAG router
        if (!IsAzureChatConfigured(ai))
            return await local.AskAsync(userId, sessionId, msg, history, ct);
        // ... retrieve (RAG) -> augment -> call Azure OpenAI -> on failure fall back to local
    }
}
\end{lstlisting}
\end{english}

الشرح لأهمّ الأسطر:
\begin{itemize}
\item السطر \code{... : IAiService}: إعلان صريح بأنّ هذا الصفّ \textbf{يُنفّذ}
  المنفذ. المُصرِّف يُلزمه بتقديم العمليات الثلاث كلّها.
\item السطر \code{=> local.GenerateRecommendationsAsync(...)}: الترشيحات ليست نداءً
  لنموذج لغويّ، بل حساب حتميّ. لذلك يُفوّض المحوّل السحابيّ هذه العملية لمحوّل محلّي
  \code{LocalAiService} — نفس المنفذ، تنفيذ مختلف حسب الحاجة.
\item كتلة الـ \code{if} والتعليقات: لو لم يُضبَط مِفتاح \en{Azure}، أو فشل النداء
  الشبكيّ، يتراجع المحوّل بأمان إلى المسار المحلّي بدل أن يُفشل الطلب بخطأ 500.
  المستخدم يحصل على إجابة دائمًا.
\end{itemize}

\begin{note}
لاحِظ اتجاه الاعتماد: محوّل \code{AzureOpenAiService} في الطبقة الخارجية يعرف
المنفذ \code{IAiService} الموجود في الداخل ويُنفّذه. العكس مستحيل: طبقة
\en{Application} لا تعرف اسم \code{AzureOpenAiService} أبدًا. هذا هو «قلب الاعتماد»
\en{(Dependency Inversion)} في الممارسة.
\end{note}

\subsection{طبقة \en{API} — الواجهة الرفيعة}

أعلى طبقة هي \en{API}. مهمّتها الوحيدة: استقبال طلب \en{HTTP}، ترجمته إلى أمر أو
استعلام، وتمريره إلى \en{MediatR}، ثم إعادة النتيجة. الـ \en{controllers} هنا
\textbf{رفيعة جدًّا} (سطر أو ثلاثة) — لا منطق عمل فيها إطلاقًا. تضمّ هذه الطبقة
أيضًا مستقبِل \en{Stripe webhook} (الذي يتحقّق من توقيع \en{Stripe} ويطابق
المدفوعات)، وإعداد الهُويّة \en{(ASP.NET Identity + JWT)}، وربط كل شيء معًا في
\code{Program.cs}.

\subsection{\en{Ports \& Adapters} — أهمّ قرار في النظام}

وصلنا لجوهر الفصل. الفكرة التي تربط كل ما سبق اسمها \en{Ports \& Adapters}
(المنافذ والمحوّلات): طبقة العمل تُعلن \textbf{منافذ} (واجهات نظيفة)، وطبقة
\en{Infrastructure} تقدّم \textbf{محوّلات} تُنفّذها. اختيار أيّ محوّل يُستخدَم يحدث
عند تشغيل التطبيق في مكان واحد اسمه \code{DependencyInjection.cs}، وبناءً على
\textbf{الإعداد فقط}، بلا تغيير أي سطر كود.

هنشوف إيه: مقتطف حقيقيّ من \code{DependencyInjection.cs} يختار المحوّل حسب الإعداد.

\begin{english}
\begin{lstlisting}[language={[Sharp]C}]
// Stripe: real client only with a genuine secret key; else the dev stub.
var stripeSecretKey = config.GetValue<string>($"{StripeOptions.SectionName}:SecretKey");
if (stripeSecretKey is not null && stripeSecretKey.StartsWith("sk_"))
    services.AddSingleton<IStripeService, StripeService>();
else
    services.AddSingleton<IStripeService, StubStripeService>();

// AI provider is config-driven (Ai:Provider): AzureOpenAi | OpenAi | Local
var aiProvider = config.GetValue<string>($"{AiOptions.SectionName}:Provider") ?? "Stub";
if (string.Equals(aiProvider, "AzureOpenAi", StringComparison.OrdinalIgnoreCase))
    services.AddScoped<IAiService, AzureOpenAiService>();
else
    services.AddScoped<IAiService>(sp => sp.GetRequiredService<LocalAiService>());

// Antivirus: real ClamAV only when FileScanning:Enabled = true.
if (config.GetValue<bool>($"{FileScanningOptions.SectionName}:Enabled"))
    services.AddSingleton<IFileScanService, ClamAvFileScanService>();
else
    services.AddSingleton<IFileScanService, NoOpFileScanService>();
\end{lstlisting}
\end{english}

الشرح: كل كتلة \code{if/else} تختار محوّلًا حقيقيًّا أو بديلًا للتطوير
\en{(stub)} بناءً على قيمة في الإعداد:
\begin{itemize}
\item مِفتاح \code{Stripe:SecretKey} يبدأ بـ \code{sk} حقيقيّ؟ استخدم
  \code{StripeService}. وإلّا فالبديل \code{StubStripeService} لبيئة التطوير.
\item مِفتاح \code{Ai:Provider} يساوي \code{AzureOpenAi}؟ استخدم محوّل \en{Azure}.
  وإلّا فالمحوّل المحلّي \code{LocalAiService}.
\item مِفتاح \code{FileScanning:Enabled} مفعَّل؟ استخدم فحص \en{ClamAV} الحقيقيّ.
  وإلّا فمحوّل لا يفعل شيئًا.
\end{itemize}

\begin{why}
لماذا هذا أهمّ قرار في النظام؟ لسببين عمليّين ضخمين. \textbf{أولًا: القابلية
للاختبار.} نستطيع في الاختبارات استبدال \en{Stripe} و\en{Azure} و\en{ACS} ببدائل
حتميّة، فنختبر منطق العمل كلّه بلا إنترنت ولا تكلفة ولا بطء. \textbf{ثانيًا:
القابلية للتبديل.} نبدّل مزوّدًا كاملًا بتغيير سطر في الإعداد فقط — \code{Ai:Provider}،
\code{Storage:Provider}، \code{Email:Provider}، \code{FileScanning:Enabled}،
\code{Hangfire:Enabled} — بلا لمس كود المنطق. في التطوير كل شيء بدائل محلّية مجّانية،
وفي الإنتاج تُبدَّل بمزوّدات حقيقية عبر الإعداد. نفس الكود، بيئتان.
\end{why}

\begin{secnote}
هذا الفصل بين المنفذ والمحوّل هو أيضًا حصنٌ أمنيّ. المفاتيح السرّية والاتّصالات
الخارجية كلّها محبوسة في طبقة \en{Infrastructure} وحدها. طبقة العمل — قلب النظام —
لا تلمس مِفتاحًا سرّيًّا ولا تعرف عنوان خادم. لو أراد مُدقّق أمنيّ معرفة «أين تُستخدَم
المفاتيح؟»، فالإجابة في مكان واحد محصور، لا منثورة في كل مِلفّ.
\end{secnote}

\begin{qa}
«لماذا \en{Clean Architecture} بدل معماريّة طبقيّة بسيطة؟» — لأنّها تفرض قلب
الاعتماد \en{(Dependency Inversion)}: الدومين لا يعرف شيئًا عن \en{EF Core} أو
\en{Azure}، فنستطيع تبديل قاعدة البيانات أو أي خدمة دون لمس منطق العمل. أثبتنا هذا
عمليًّا بحقن بدائل \en{(stubs)} لكل الخدمات الخارجية في الاختبارات، وبتبديل مزوّد
الذكاء الاصطناعي عبر مِفتاح \code{Ai:Provider} وحده. في المعماريّة الطبقيّة
التقليدية تعتمد الطبقة العليا على الدنيا مباشرةً، فيصبح المنطق مرهونًا بالتفاصيل
التقنية — وهذا بالضبط ما نتجنّبه.
\end{qa}

\begin{summary}[الخلاصة في سطور]
\begin{itemize}
\item \en{ScholarPath} خلفيّتها مبنيّة على \en{Clean Architecture}: أربع طبقات
  متداخلة، والاعتماد يتّجه \textbf{للداخل فقط}.
\item الطبقات هي: \en{Domain} (القلب النقيّ) ثم \en{Application} (حالات الاستخدام)
  ثم \en{Infrastructure} (التقنيات) ثم \en{API} (الواجهة الرفيعة). أربعة مشاريع
  \en{.csproj} منفصلة، والترتيب مضمون بمراجع المشاريع.
\item \en{Domain} يحوي $\approx$55 كيانًا والـ \en{enums} وأحداث الدومين ومنفذين
  فقط \en{(}\code{ICurrentUserService}\en{, }\code{IDateTimeService}\en{)}، بلا أي
  اعتماد على \en{framework}.
\item \en{Application} يكتب حالات الاستخدام كـ \en{CQRS}، ومعها الـ \en{validators}
  و\en{AutoMapper}، ويُعلن $\approx$33 منفذًا يصف \emph{ماذا} يحتاج النظام.
\item \en{Infrastructure} يقدّم $\approx$46 محوّلًا تُنفّذ المنافذ بتقنيات حقيقية
  \en{(EF Core, Stripe, Azure OpenAI, ACS, Blob, Key Vault, SignalR, Hangfire)}.
\item القرار الأهمّ هو \en{Ports \& Adapters}: يجعل النظام \textbf{قابلًا للاختبار}
  و\textbf{قابلًا لتبديل المزوّد بالإعداد وحده} — \code{Ai:Provider}،
  \code{Storage:Provider}، \code{Email:Provider}، \code{FileScanning:Enabled}،
  \code{Hangfire:Enabled} — دون تغيير أي كود.
\end{itemize}
\end{summary}

\clearpage
% ch02-cqrs-integrity
\section{نمط CQRS وخطّ المعالجة وسلامة البيانات}

في الفصل السابق عرفنا الطبقات الأربع، وكيف أنّ الطبقة الداخلية لا تعرف شيئًا عن
الخارجية، وكيف تتكلّم الطبقات مع العالم الخارجيّ عبر \en{ports \& adapters}. في هذا
الفصل ننزل درجة أعمق: كيف \emph{تسافر} أيّ عملية من لحظة وصولها للـ \en{API} حتى
تُحفَظ في قاعدة البيانات؟ ومَن الذي يحرس هذه الرحلة حتى لا تُكتَب بيانات ناقصة أو
مكرّرة أو متعارضة؟

الإجابة مبنيّة على فكرتين متكاملتين. الأولى تنظيميّة: نمط \en{CQRS} الذي يفصل عمليّات
\emph{الكتابة} عن عمليّات \emph{القراءة}، ويمرّر كلّ عملية عبر خطّ معالجة موحّد. والثانية
دفاعيّة: مجموعة ضوابط في طبقة البيانات تضمن أنّ الصفّ الذي وصل للقاعدة صحيح ومتّسق. لو
فهمت هذين المحورين تكون فهمت العمود الفقريّ للنظام كلّه.

% ====================================================================
\subsection{نمط CQRS: كل عملية إمّا أمر أو استعلام}

\en{CQRS} اختصار لـ \en{Command Query Responsibility Segregation}، أي «فصل مسؤوليّات
الأمر عن الاستعلام». الفكرة بسيطة جدًّا: أيّ شيء يفعله النظام يقع في واحدة من خانتين
فقط، لا ثالث لهما.

\begin{idea}[الفكرة: أمر أم استعلام؟]
\textbf{الأمر (\en{Command})} يُغيّر حالة النظام — يكتب، يعدّل، أو يحذف. أمثلة:
\code{AcceptBookingCommand} (قبول حجز)، \code{CreatePostCommand} (إنشاء منشور).
عادةً يُرجِع لا شيء أو نتيجة صغيرة.\\[4pt]
\textbf{الاستعلام (\en{Query})} يقرأ الحالة \emph{دون أيّ أثر جانبيّ} — لا يكتب حرفًا
واحدًا. أمثلة: \code{GetMyRecommendationsQuery}، \code{GetBookingRecordingsQuery}.
كل طرف من الطرفين له نوع طلب مستقلّ، ومعالِج (\en{handler}) مستقلّ، وقواعد تحقّق
مستقلّة.
\end{idea}

كل هذه الطلبات — أوامر واستعلامات — تُرسَل عبر مكتبة واحدة اسمها \en{MediatR}. أنت
تبني كائن الطلب وتناديه بسطر واحد \code{mediator.Send(request)}، وهي تتكفّل بالعثور على
المعالِج الصحيح وتشغيله. النتيجة المباشرة لهذا التصميم أنّ الـ \en{controller} يصبح
\emph{رفيعًا جدًّا}: لا منطق عمل بداخله، فقط يستقبل الطلب ويحوّله.

\textbf{هنشوف إيه:} الـ \en{endpoint} المسؤول عن قبول الحجز في \code{BookingsController}.
لاحظ أنّ كل عمله الحقيقيّ سطرٌ واحد.

\begin{english}
\begin{lstlisting}[language={[Sharp]C}]
[HttpPost("{id:guid}/accept")]
public async Task<IActionResult> AcceptBooking(
    Guid id,
    [FromBody] AcceptBookingCommand command,
    CancellationToken cancellationToken)
{
    if (id != command.BookingId)
        return BadRequest("Route booking id does not match body booking id.");

    await _sender.Send(command, cancellationToken);   // one line
    return NoContent();
}
\end{lstlisting}
\end{english}

شرح سريع سطرًا بسطر لأهمّ ما في المقتطف:
\begin{itemize}
\item السمة \code{[HttpPost(...)]} تربط المسار \en{/accept} بهذه الدالة.
\item الطلب القادم من العميل يُربَط مباشرةً بكائن \code{AcceptBookingCommand} — أي أنّ
  «شكل الطلب» و«الأمر» شيء واحد.
\item السطر \code{await \_sender.Send(command, ...)} هو كل المنطق: أرسِل الأمر، و\en{MediatR}
  ستجد المعالِج وتشغّله. الـ \en{controller} لا يعرف \emph{كيف} يُقبَل الحجز إطلاقًا.
\item يتحقّق فقط من اتّساق بسيط (تطابق المُعرّف في المسار والجسد) ثم يُعيد \en{204 NoContent}.
\end{itemize}

\begin{why}
لماذا نُبقي الـ \en{controller} رفيعًا؟ لأنّ كل منطق العمل يصير في صنف معالِج مستقلّ
قابل للاختبار وحده دون تشغيل خادم \en{HTTP}. كل ميزة جديدة تعني إضافة \emph{معالِج جديد}،
لا تعديل صنف خدمة ضخم قائم. وهذا يفتح الباب لخطّ معالجة موحّد يلفّ \emph{كل} معالِج
تلقائيًّا — وهو موضوعنا التالي.
\end{why}

% ====================================================================
\subsection{مثال حيّ: أمر قبول الحجز}

لنتتبّع أمرًا حقيقيًّا كاملًا. عندما يقبل المستشار حجزًا مدفوعًا، يجب أن تحدث عدّة أشياء
بترتيب دقيق: تحميل الحجز، سحب المبلغ المحجوز من بطاقة الطالب عبر \en{Stripe}، تجهيز غرفة
اجتماع مرئيّ عبر \en{ACS}، حفظ كل ذلك، ثم إعلان أنّ الحجز تأكّد. كل هذا يعيش في معالِج
واحد.

أوّلًا، الأمر نفسه — لاحظ أنّه \code{record} صغير جدًّا يحمل مُعرّف الحجز فقط، ومُزخرَف
بالسمة \code{[Auditable]}:

\textbf{هنشوف إيه:} تعريف \code{AcceptBookingCommand} والسمة التي تجعله قابلًا للتدقيق.

\begin{english}
\begin{lstlisting}[language={[Sharp]C}]
[Auditable(
    AuditAction.BookingAccepted,
    "ConsultantBooking",
    TargetIdProperty = nameof(BookingId),
    SummaryTemplate = "Booking accepted: {TargetId}")]
public sealed record AcceptBookingCommand(Guid BookingId) : IRequest;
\end{lstlisting}
\end{english}

الأمر مجرّد بيانات: مُعرّف حجز واحد. أمّا السمة \code{[Auditable]} فتقول للنظام «سجّل هذا
الأمر تلقائيًّا في سجلّ التدقيق باسم الفعل \en{BookingAccepted}» — وسنرى آليّتها بعد قليل.
الآن المعالِج (مختصَرًا مع الحفاظ على الخطوات الحقيقيّة):

\textbf{هنشوف إيه:} قلب \code{AcceptBookingCommandHandler}: تحميل، سحب الدفع، تجهيز الغرفة،
حفظ، ثم نشر الحدث.

\begin{english}
\begin{lstlisting}[language={[Sharp]C}]
public async Task Handle(AcceptBookingCommand request, CancellationToken ct)
{
    // 1. Load the aggregate (with its Payment)
    var booking = await _context.Bookings
        .Include(b => b.Payment)
        .FirstOrDefaultAsync(b => b.Id == request.BookingId, ct);
    // ... ownership + state guards omitted ...

    // 2. Stripe capture - take the money that was only *held* at booking time
    var capture = await _stripeService.CapturePaymentIntentAsync(
        booking.StripePaymentIntentId!, amountCents, idempotencyKey, ct);

    // 3. Provision the ACS video room for the confirmed session
    var room = await _meetingService.CreateRoomAsync(booking.Id, ct);
    booking.MeetingRoomId = room.RoomId;
    booking.Status = BookingStatus.Confirmed;

    // 4. Persist, THEN publish the domain event
    await _context.SaveChangesAsync(ct);
    await _publisher.Publish(
        new BookingConfirmedEvent(
            booking.Id, booking.StudentId, booking.ConsultantId), ct);
}
\end{lstlisting}
\end{english}

اقرأ الخطوات كقصّة قصيرة:
\begin{itemize}
\item \textbf{تحميل}: يجلب الحجز ومعه سجلّ الدفع \code{Payment} في استعلام واحد بفضل
  \code{Include}. (حُذفت هنا فحوص الملكيّة والحالة لِقِصَر المساحة، وهي موجودة في الكود
  الأصليّ.)
\item \textbf{سحب الدفع}: عند الحجز كان المبلغ \emph{محجوزًا} فقط على البطاقة لا مسحوبًا.
  الآن يُسحَب فعليًّا عبر \code{CapturePaymentIntentAsync} — سنشرح نموذج «احجز ثمّ اسحب»
  في فصل المدفوعات.
\item \textbf{تجهيز الغرفة}: يُنشئ غرفة اجتماع عبر خدمة \en{ACS} ويحفظ مُعرّفها على الحجز.
\item \textbf{حفظ ثم نشر}: يُستدعى \code{SaveChangesAsync} أوّلًا، وبعد نجاح الحفظ فقط
  يُنشَر \code{BookingConfirmedEvent}. هذا الترتيب مقصود تمامًا — نُوضّحه في قسم المعترِض.
\end{itemize}

\begin{keyp}[نقطة مهمة: المعالِج هو موطن منطق العمل]
كل القرارات — مَن يملك الحجز، أهو في الحالة الصحيحة، متى نسحب، متى ننشر الحدث — تعيش
هنا، في صنف واحد صغير قابل للاختبار. الـ \en{controller} لم يعرف شيئًا من ذلك. هذا هو
جوهر \en{CQRS} عمليًّا.
\end{keyp}

% ====================================================================
\subsection{خطّ المعالجة وترتيبه المقصود}

ذكرنا أنّ \en{MediatR} تلفّ كل معالِج. الأداة التي تفعل ذلك تُسمّى \en{pipeline behavior}:
حلقة تُغلِّف المعالِج فتنفّذ شيئًا \emph{قبله} و\emph{بعده}، تمامًا كطبقات البصلة. عندنا
أربع حلقات، وترتيب تسجيلها = ترتيب تنفيذها (الأبعد أوّلًا، الأقرب للمعالِج أخيرًا).

\textbf{هنشوف إيه:} تسجيل الحلقات الأربع في \code{DependencyInjection.cs} — الترتيب هنا هو
عين الترتيب الفعليّ عند التنفيذ.

\begin{english}
\begin{lstlisting}[language={[Sharp]C}]
services.AddMediatR(cfg =>
{
    cfg.RegisterServicesFromAssembly(assembly);
    cfg.AddOpenBehavior(typeof(LoggingBehavior<,>));      // 1 outermost
    cfg.AddOpenBehavior(typeof(ValidationBehavior<,>));   // 2
    cfg.AddOpenBehavior(typeof(PerformanceBehavior<,>));  // 3
    cfg.AddOpenBehavior(typeof(AuditBehavior<,>));        // 4 closest to handler
});
\end{lstlisting}
\end{english}

الجدول التالي يلخّص وظيفة كل حلقة وسبب موضعها:

\begin{center}
\begin{tabular}{L{3.2cm} L{4.4cm} L{5cm}}
\toprule
\textbf{\color{spPrimary}الحلقة} & \textbf{\color{spPrimary}ماذا تفعل} & \textbf{\color{spPrimary}لماذا هذا الموضع} \\
\midrule
\en{Logging} & تسجّل بداية ونهاية كل طلب مع مُعرّف المستخدم & الأبعد، لتلتقط \emph{كل} طلب حتى لو رُفِض لاحقًا \\
\en{Validation} & تشغّل كل مُتحقّقات \en{FluentValidation} وتقصر الطلب بـ 400 عند الفشل & مبكّرة عمدًا: لا نُضيّع وقتًا ولا نُدقّق طلبًا معطوبًا \\
\en{Performance} & تقيس زمن المعالِج وتُنبّه للبطيء & بعد التحقّق، لنقيس زمن العمل الحقيقيّ فقط \\
\en{Audit} & تكتب صفًّا في \code{AuditLog} للأوامر المُزخرَفة & الأقرب للمعالِج، لِتُسجّل فقط ما نُفِّذ فعلًا \\
\bottomrule
\end{tabular}
\end{center}

\begin{why}
الترتيب ليس اعتباطيًّا. \en{Logging} يجب أن يرى \emph{كل} طلب، لذا هو الأبعد.
\en{Validation} يجب أن يرفض المدخل السيّئ \emph{قبل} أن نصرف وقتًا في القياس أو نكتب صفوف
تدقيق لطلب لن يُنفَّذ أصلًا. \en{Audit} يجب أن يجلس أقرب ما يكون للمعالِج، فلا يُسجّل إلّا
فعلًا نُفِّذ حقًّا. اقلب الترتيب وستدقّق أفعالًا لم تحدث، أو تقيس زمن طلبات فاسدة.
\end{why}

ولنرَ حلقة التحقّق من الداخل — كيف «تقصر» الطلب:

\textbf{هنشوف إيه:} قلب \code{ValidationBehavior}: يجمع الأخطاء، وإن وُجدت يرمي استثناءً
قبل أن يصل الطلب للمعالِج.

\begin{english}
\begin{lstlisting}[language={[Sharp]C}]
var failures = validationResults
    .Where(r => !r.IsValid)
    .SelectMany(r => r.Errors)
    .ToList();

if (failures.Count > 0)
    throw new ValidationException(failures);   // short-circuit -> 400

return await next(cancellationToken);          // otherwise, continue inward
\end{lstlisting}
\end{english}

الفكرة أنّ \code{next(...)} هي «الحلقة التالية إلى الداخل» وصولًا للمعالِج. إن وُجد خطأ
واحد، نرمي \code{ValidationException} فورًا ولا نستدعي \code{next} أبدًا — فلا يعمل
المعالِج، ولا تُقاس مدّة، ولا يُكتَب تدقيق. هذا معنى «القصر» (\en{short-circuit}).

\begin{note}
المُتحقّقات ومُلفّات \en{AutoMapper} تُكتَشَف تلقائيًّا من التجميعة (\en{assembly}) عند
الإقلاع. أي أنّ إضافة ميزة جديدة بمُتحقّقها لا تحتاج أيّ ربط يدويّ في حاوية الاعتماديّات.
\end{note}

% ====================================================================
\subsection{التدقيق بالسمة Auditable}

الآن نفهم كيف يعمل التدقيق التلقائيّ. حلقة \en{Audit} — الأقرب للمعالِج — تنتظر انتهاء
المعالِج بنجاح، ثم تسأل: هل نوع هذا الطلب مُزخرَف بالسمة \code{[Auditable]}؟ إن كان
كذلك، تكتب صفًّا في جدول \code{AuditLog}. إن لم يكن، لا تفعل شيئًا.

\textbf{هنشوف إيه:} قلب \code{AuditBehavior}: يشغّل المعالِج، يقرأ السمة، وإن وُجدت يكتب
سجلّ التدقيق.

\begin{english}
\begin{lstlisting}[language={[Sharp]C}]
var response = await next(cancellationToken);   // run the handler first

var attr = (AuditableAttribute?)Attribute.GetCustomAttribute(
    typeof(TRequest), typeof(AuditableAttribute));
if (attr is null)
    return response;                            // no attribute -> no audit row

var targetId = ResolveTargetId(response, request, attr.TargetIdProperty);
var summary  = BuildSummary(attr.SummaryTemplate, targetId, response, request);

await audit.WriteAsync(attr.Action, attr.TargetType, targetId,
    beforeJson: null, afterJson: null, summary, cancellationToken);
return response;
\end{lstlisting}
\end{english}

نقاط جديرة بالانتباه:
\begin{itemize}
\item \code{await next(...)} أوّلًا — التدقيق يحدث \emph{بعد} نجاح المعالِج، فلا نُسجّل
  فعلًا فشل.
\item \code{GetCustomAttribute} يقرأ السمة عن نوع الطلب. لو غابت، يخرج فورًا — فالحلقة
  «لا تفعل شيئًا» لكل الأوامر غير المُعلَّمة.
\item \code{ResolveTargetId} و\code{BuildSummary} يستخرجان مُعرّف الهدف ونصًّا وصفيًّا من
  السمة (تذكّر \code{TargetIdProperty = nameof(BookingId)} في أمر قبول الحجز).
\item لو فشلت الكتابة في التدقيق، تُلتقَط ولا تُفشِل الأمر الأصليّ — فنجاح العمل لا يُلغى
  بسبب سجلّ.
\end{itemize}

\begin{keyp}[نقطة مهمة: التدقيق بالتزيين لا بالتكرار]
لتجعل أمرًا قابلًا للتدقيق تضع سطرًا واحدًا فوقه: \code{[Auditable(...)]}. لا كود مكرّر
داخل كل معالِج. عشرات الأوامر الحسّاسة في النظام (قبول حجز، موافقة على منحة، تغيير دور
مستخدم، ...) مُعلّمة هكذا، فتُسجَّل كلها بنفس الآليّة المركزيّة.
\end{keyp}

% ====================================================================
\subsection{سلامة البيانات: الحذف الناعم}

ننتقل الآن من «رحلة الطلب» إلى «حراسة البيانات». أوّل حارس هو \emph{الحذف الناعم}
(\en{soft delete}). الصفوف الحسّاسة أو ذات القيمة التاريخيّة لا تُحذَف حذفًا صلبًا أبدًا.
بدلًا من ذلك تُنفّذ الواجهة \code{ISoftDeletable} (تحمل \code{IsDeleted} و\code{DeletedAt}
و\code{DeletedByUserId})، و«الحذف» ما هو إلّا قلب \code{IsDeleted} إلى \en{true}.

لكن كيف نضمن أنّ الصفّ «المحذوف» يختفي من \emph{كل} استعلام دون أن يتذكّر كل مطوّر ذلك؟
عبر \emph{مُرشِّح استعلام عموميّ} في \en{EF Core} على كل كيان من هذا النوع:

\textbf{هنشوف إيه:} سطر المُرشِّح العموميّ — يُطبَّق على قرابة 15 كيانًا.

\begin{english}
\begin{lstlisting}[language={[Sharp]C}]
b.HasQueryFilter(a => !a.IsDeleted);   // ~15 entities carry this
\end{lstlisting}
\end{english}

\begin{idea}[الفكرة: الصفّ يختفي دون أن يُدمَّر]
بمجرّد وضع هذا المُرشِّح على الكيان، \emph{كل} استعلام \en{LINQ} يولّده \en{EF} لهذا
الكيان يُضاف إليه ضمنيًّا \code{WHERE IsDeleted = 0}. المطوّر لا يستطيع أن ينسى إضافته —
فالصفوف المحذوفة ببساطة لا تظهر في أيّ استعلام أو ربط أو عدّ، إلّا إن اختار الكود صراحةً
تجاوز المُرشِّح بـ \code{IgnoreQueryFilters()} (وهذا محصور في أدوات الإدارة و\en{GDPR}).
\end{idea}

\begin{why}
لماذا نتجنّب الحذف الصلب؟ لأنّه يكسر التاريخ المترابط (مراجعات، مدفوعات، آثار تدقيق)،
ويجعل سؤال «مَن حذف ماذا ومتى؟» مستحيلًا. الحذف الناعم يحفظ ترابط التاريخ، ويدعم مسار
تراجع/استرجاع، ومع سجلّ التدقيق يعطي تتبّعًا كاملًا. وهو أيضًا الخيار الآمن الافتراضيّ
لمنصّة ملزَمة بقواعد حفظ البيانات وحلّ النزاعات.
\end{why}

% ====================================================================
\subsection{معترِض SaveChangesAsync: مكان واحد، ضمانات كثيرة}

الحارس الثاني — والأقوى — هو تجاوُز الدالة \code{SaveChangesAsync} داخل \code{DbContext}.
هذه الدالة هي «الاختناق» الوحيد الذي يمرّ منه كل حفظ للبيانات، فجعلناها تفرض عدّة قواعد
عرضيّة (\en{cross-cutting}) دفعةً واحدة.

\textbf{هنشوف إيه:} قلب \code{SaveChangesAsync}: ختم التواريخ، نشر الأحداث بعد الحفظ،
وترجمة تعارض الفهرس الفريد.

\begin{english}
\begin{lstlisting}[language={[Sharp]C}]
foreach (var entry in ChangeTracker.Entries<AuditableEntity>())
{
    if (entry.State == EntityState.Added && entry.Entity.CreatedAt == default)
        entry.Entity.CreatedAt = auditNow;
    else if (entry.State == EntityState.Modified)
        entry.Entity.UpdatedAt = auditNow;
}

try
{
    var result = await base.SaveChangesAsync(cancellationToken);
    await DispatchDomainEventsAsync(cancellationToken);   // only after commit
    return result;
}
catch (DbUpdateException ex) when (ex.InnerException is SqlException sqlEx
       && (sqlEx.Number == 2601 || sqlEx.Number == 2627))
{
    throw new ConflictException("This action conflicts with an existing record.");
}
\end{lstlisting}
\end{english}

هذا المقتطف الصغير يفرض ثلاث ضمانات (والرابعة نشرحها بعده نصًّا):
\begin{itemize}
\item \textbf{ختم التواريخ}: قبل الحفظ يمرّ على الكيانات المتغيّرة، فيضع \code{CreatedAt}
  على الصفوف الجديدة (فقط إن لم تُضبَط، حتى تحتفظ بيانات البذر بتواريخها المقصودة)،
  و\code{UpdatedAt} على كل تعديل. بدون هذا تُحفَظ صفوف بتاريخ صفريّ يكسر العرض والترتيب.
\item \textbf{أحداث الدومين بعد الحفظ فقط}: يُستدعى \code{DispatchDomainEventsAsync}
  \emph{بعد} أن ينجح \code{base.SaveChangesAsync}. لماذا؟ حتى لا يُطلَق أثر جانبيّ
  (إشعار، بريد) لتغيير تراجَع ولم يُحفَظ أصلًا. (هذا يفسّر ترتيب «احفظ ثمّ انشر» في معالِج
  قبول الحجز.)
\item \textbf{ترجمة التعارض}: لو ضرب الحفظ فهرسًا فريدًا (رقما الخطأ \en{2601}/\en{2627}
  في \en{SQL Server})، نلتقطه ونعيد رميه كـ \code{ConflictException} ودّيّ (مثل «هذا
  الموعد حُجِز للتوّ») بدل خطأ 500 غامض.
\end{itemize}

\begin{idea}[الفكرة: تحويل كل FK نحو المستخدم إلى Restrict]
الضمانة الرابعة تعيش في \code{OnModelCreating} لا في الحفظ، لكنّها من نفس فلسفة «قاعدة
مركزيّة واحدة». يرفض \en{SQL Server} وجود أكثر من مسار \en{CASCADE} إلى الجدول نفسه —
وهو ما يحدث لأيّ كيان له مفتاحان أجنبيّان أو أكثر نحو \code{ApplicationUser}. فبدل إضافة
\code{Restrict} في كل إعداد على حدة، نمرّ على \emph{كل} مفتاح أجنبيّ يشير للمستخدم ونحوّله
إلى \code{DeleteBehavior.Restrict} دفعةً واحدة — فنتجنّب الخطأ ونحمي تاريخ المستخدم.
\end{idea}

\begin{secnote}
تجميع هذه القواعد في نقطة واحدة أمانٌ بحدّ ذاته: لا يمكن لمطوّر أن «ينسى» ختم تاريخ أو
يترك تعارضًا يتسرّب كخطأ 500 يكشف تفاصيل داخليّة. القاعدة تُفرَض على الجميع لأنّها تعيش في
الطريق الذي لا مفرّ منه.
\end{secnote}

% ====================================================================
\subsection{ضوابط أخرى: التزامن والتكرار والفهارس الفريدة}

يبقى ثلاثة حرّاس صغار لكن حاسمين:

\subsubsection{التزامن المتفائل عبر RowVersion}
الكيانات القابلة للتدقيق تحمل عمودًا \code{RowVersion} يتغيّر مع كل تعديل. لو حاول
مستخدمان تعديل الصفّ نفسه في آنٍ، ينجح الأوّل ويفشل الثاني لأنّ نسخته صارت قديمة — فنتجنّب
«التحديث الضائع» دون أقفال ثقيلة. يُعرَّف بسطر:
\code{b.Property(p => p.RowVersion).IsRowVersion();}

\subsubsection{مفاتيح عدم التكرار IdempotencyKey}
المال والإشعارات يجب أن يحدثا \emph{مرّةً واحدة بالضبط} حتى مع إعادة المحاولة. لذلك يحمل
\code{Payment} و\code{Notification} حقل \code{IdempotencyKey} فريدًا. لو أُعيد إرسال نفس
الطلب، يصطدم بالمفتاح الموجود فلا يتكرّر خصمٌ ولا إشعار. (تذكّر \code{idempotencyKey} في
مقتطف سحب الدفع سابقًا — نفس الفكرة تُرسَل إلى \en{Stripe} كذلك.)

\subsubsection{الفهارس الفريدة المُصفّاة}
هذه أذكاها. أحيانًا نريد قاعدة مثل «تطبيق \emph{فعّال} واحد لكل طالب لكل منحة» — لكن نسمح
بعدّة تطبيقات مسحوبة أو مرفوضة. الفهرس الفريد العاديّ لا يكفي، فنستخدم \emph{فهرسًا فريدًا
مُصفّى}: يُطبّق شرط التفرّد على صفوف بعينها فقط.

\textbf{هنشوف إيه:} الفهرس الذي يفرض تطبيقًا فعّالًا واحدًا لكل طالب/منحة.

\begin{english}
\begin{lstlisting}[language={[Sharp]C}]
b.HasIndex(a => new { a.StudentId, a.ScholarshipId })
    .IsUnique()
    .HasFilter(
        "[ScholarshipId] IS NOT NULL " +
        "AND [Status] <> 'Withdrawn' " +
        "AND [Status] <> 'Rejected' " +
        "AND [Status] <> 'Accepted'")
    .HasDatabaseName("UX_Applications_Student_Scholarship_Active");
\end{lstlisting}
\end{english}

الشرط في \code{HasFilter} يستثني الحالات المنتهية (مسحوب/مرفوض/مقبول) والتطبيقات الخارجيّة
البحتة (بلا رابط منحة). فيبقى التفرّد مفروضًا على التطبيقات \emph{النشطة} فقط. النتيجة:
قاعدة العمل صارت محفورةً في القاعدة نفسها، لا مجرّد فحص في الكود يمكن تجاوزه.

\begin{qa}
\textbf{سؤال متوقّع:} «لو الكود فيه فحص قبل الإدراج، فلماذا نحتاج الفهرس الفريد أيضًا؟»
لأنّ الفحص في الكود عرضة لسباق التزامن (\en{race condition}): طلبان متزامنان قد يمرّان
كلاهما من الفحص ثم يُدرِجان. الفهرس الفريد في القاعدة هو خطّ الدفاع الأخير الذي لا
يُخترَق: أحد الطلبين يفشل حتمًا، ونحن نترجم فشله إلى \code{ConflictException} ودّيّ كما
رأينا في المعترِض. الكود يعطي رسالة لطيفة، والقاعدة تعطي الضمان.
\end{qa}

% ====================================================================
\begin{summary}[الخلاصة في سطور]
\begin{itemize}
\item \textbf{\en{CQRS}}: كل عملية إمّا \en{Command} (كتابة) أو \en{Query} (قراءة)،
  تُرسَل عبر \en{MediatR}، والـ \en{controller} رفيع بسطر واحد \code{mediator.Send}.
\item \textbf{المعالِج} هو موطن منطق العمل: رأينا \code{AcceptBookingCommand} يحمّل ثم
  يسحب الدفع ثم يجهّز غرفة \en{ACS} ثم يحفظ ثم ينشر حدث الدومين.
\item \textbf{خطّ المعالجة} أربع حلقات بترتيب مقصود: \en{Logging} ثم \en{Validation} ثم
  \en{Performance} ثم \en{Audit} (الأقرب للمعالِج) — كلٌّ في موضعه لسبب.
\item \textbf{التدقيق} تلقائيّ بسطر \code{[Auditable]} فوق الأمر؛ حلقة \en{Audit} تكتب صفّ
  \code{AuditLog} بعد نجاح المعالِج فقط.
\item \textbf{الحذف الناعم} عبر \code{ISoftDeletable} و\code{HasQueryFilter} على قرابة 15
  كيانًا يجعل الصفوف المحذوفة تختفي من كل استعلام دون تدميرها.
\item \textbf{معترِض \code{SaveChangesAsync}} يختم التواريخ، وينشر الأحداث بعد الحفظ فقط،
  ويترجم تعارض الفهرس الفريد إلى \code{ConflictException}، ويحوّل مفاتيح المستخدم إلى
  \code{Restrict}.
\item \textbf{ضوابط أخرى}: \code{RowVersion} للتزامن، \code{IdempotencyKey} لعدم التكرار،
  وفهارس فريدة مُصفّاة تحفر قواعد العمل داخل القاعدة نفسها.
\end{itemize}
\end{summary}

\clearpage
% ch03-auth
\section{التوثيق والصلاحيّات}

في الفصلين السابقين عرفنا الطبقات و\en{CQRS} والمنافذ (\en{ports}). في هذا الفصل
ننتقل إلى أول ما يقابله أي مستخدم عند فتح \en{ScholarPath}: كيف يُثبت أنه \emph{هو}،
وكيف يقرّر النظام بعد ذلك ما \emph{يُسمح} له بعمله. الموضوعان متلازمان لكنهما مختلفان
تمامًا، ومن المهم التفريق بينهما من البداية.

\begin{idea}[الفكرة: توثيق مقابل تفويض]
\textbf{التوثيق (\en{Authentication})} يجيب على سؤال: «مَن أنت؟» — أي التأكّد من
هويّة المستخدم عبر كلمة مرور أو حساب \en{Google} أو رمز.
أمّا \textbf{التفويض (\en{Authorization})} فيجيب على سؤال آخر: «طيّب، وأنت مين، تقدر
تعمل إيه؟» — أي هل يحقّ لهذا المستخدم رؤية هذا المورد أو تعديله؟
النظام ينفّذ الأول عند الدخول (ويصدر رمزين)، ثم ينفّذ الثاني في \emph{كل} طلب لاحق.
\end{idea}

كل ما سنشرحه هنا مأخوذ من الكود الحيّ داخل \code{server/src}، تحديدًا من طبقة
\en{Infrastructure} (الخدمات) وطبقة \en{Application} (المُعالِجات / \en{handlers}).

% ====================================================================
\subsection{الهويّة عبر ASP.NET Core Identity}

بدلًا من كتابة نظام مستخدمين من الصفر (وهو أمر خطير أمنيًّا)، يبني \en{ScholarPath}
على \en{ASP.NET Core Identity}: مكتبة مايكروسوفت الناضجة لإدارة الحسابات وكلمات المرور
والأدوار. نقطة الوصل هي الصنف \code{ApplicationUser} الذي \emph{يرث} من
\code{IdentityUser<Guid>}، فيأخذ كل ما توفّره \en{Identity} (كلمة المرور المُجزّأة،
عدّاد المحاولات الفاشلة، تأكيد البريد) ويضيف فوقه حقول العمل الخاصّة بنا.

\begin{note}
اخترنا \code{Guid} كنوع للمُعرِّف بدل الرقم المتسلسل: المُعرِّفات العشوائية لا تُخمَّن
ولا تكشف عدد المستخدمين، وتناسب التوزيع دون تصادم.
\end{note}

هنشوف تعريف \code{ApplicationUser} (مختصرًا) من ملف \code{Identity.cs}:

\begin{english}
\begin{lstlisting}[language={[Sharp]C}]
public class ApplicationUser : IdentityUser<Guid>, ISoftDeletable
{
    public string FirstName { get; set; } = default!;
    public string LastName  { get; set; } = default!;
    public AccountStatus AccountStatus { get; set; } = AccountStatus.Unassigned;
    public bool IsOnboardingComplete { get; set; }
    public DateTimeOffset? LastLoginAt { get; set; }

    // Dual-role support: the role active for the current session
    public string? ActiveRole { get; set; }

    public ICollection<RefreshToken> RefreshTokens { get; } = [];
    public ICollection<LoginAttempt> LoginAttempts { get; } = [];

    public string FullName => $"{FirstName} {LastName}".Trim();
}
\end{lstlisting}
\end{english}

نقرأ الأسطر المهمّة واحدة واحدة:
\begin{itemize}
\item \code{: IdentityUser<Guid>} — الوراثة التي تجلب البنية التحتية للهوية جاهزة،
  فلا نُخزّن كلمة المرور بأنفسنا أبدًا.
\item \code{AccountStatus} و\code{IsOnboardingComplete} — حقول عمل تُقرّر مسار
  المستخدم بعد الدخول (مثلًا: هل أكمل خطوات التسجيل؟).
\item \code{ActiveRole} — لأن المستخدم قد يملك أكثر من دور، نحفظ الدور الفعّال للجلسة.
\item \code{RefreshTokens} — مجموعة رموز التحديث المرتبطة بالمستخدم؛ سنعود إليها بعد قليل.
\end{itemize}

\subsubsection{تجزئة كلمة المرور}
النظام لا يخزّن كلمة المرور نصًّا صريحًا، بل \emph{تجزئتها} فقط عبر المنفذ
\code{IPasswordHasher}، وتنفيذه \code{IdentityPasswordHasher} يفوّض العمل إلى
\code{PasswordHasher} من \en{Identity} (خوارزمية \en{PBKDF2} مع مِلح
عشوائي لكل تجزئة). عند الدخول نستدعي \code{Verify} الذي يعيد النتيجة
دون أن يفكّ التجزئة أبدًا (التجزئة عملية باتجاه واحد).

\begin{secnote}
لماذا مِلح (\en{salt}) لكل مستخدم؟ لأنه يجعل تجزئتين لكلمتَي مرور متطابقتين
\emph{مختلفتين} في القاعدة، فيُبطل هجمات جداول \en{rainbow} الجاهزة، ويمنع
معرفة أن مستخدِمَين يستعملان الكلمة نفسها. ولأن التجزئة خلف المنفذ
\code{IPasswordHasher}، يمكن استبدال الخوارزمية مستقبلًا دون لمس بقية الكود.
\end{secnote}

\subsubsection{القفل بعد المحاولات الفاشلة}
لصدّ هجمات التخمين، يعدّ \code{LoginCommandHandler} المحاولات الفاشلة
\emph{المتتالية} داخل نافذة زمنية؛ فإذا بلغت الحدّ يُقفَل الحساب مؤقّتًا. هنشوف
قلب هذا المنطق:

\begin{english}
\begin{lstlisting}[language={[Sharp]C}]
private const int MaxFailedAttempts = 5;
private static readonly TimeSpan LockoutDuration = TimeSpan.FromMinutes(30);
// ...
var priorFailures = await db.LoginAttempts
    .CountAsync(a => a.Email == normalizedEmail && !a.Succeeded
                     && a.OccurredAt >= windowStart, ct);

if (priorFailures + 1 >= MaxFailedAttempts)
{
    user.LockoutEnd = now + LockoutDuration;   // lock for 30 minutes
    user.AccessFailedCount = priorFailures + 1;
}
\end{lstlisting}
\end{english}

الشرح:
\begin{itemize}
\item \code{MaxFailedAttempts = 5} — بعد خمس محاولات فاشلة متتالية يُقفَل الحساب.
\item \code{CountAsync(...)} — يحسب الفشل \emph{غير الناجح} منذ بداية النافذة فقط،
  فأي دخول ناجح يصفّر العدّاد عمليًّا.
\item \code{user.LockoutEnd = now + LockoutDuration} — يضبط لحظة انتهاء القفل (30 دقيقة)؛
  وأي محاولة قبلها تُرفَض مباشرة برسالة «الحساب مقفول مؤقّتًا».
\end{itemize}

كل محاولة — ناجحة أو فاشلة — تُسجَّل كصفٍّ في \code{LoginAttempts} مع الوقت وعنوان
\en{IP}، فيصبح لدينا سِجلّ يمكن تدقيقه لاحقًا. أمّا \textbf{تأكيد البريد} فتديره
\en{Identity} عبر خاصية \code{EmailConfirmed}؛ فحساب لم يؤكّد بريده يظل محدود
الصلاحيات حتى يضغط رابط التفعيل.

% ====================================================================
\subsection{الرمزان: access JWT و refresh token}

بعد نجاح التوثيق يُصدر \code{TokenService} \emph{زوجًا} من الرموز، لكل منهما دور مختلف:

\begin{center}
\begin{tabular}{L{3.2cm} L{11.4cm}}
\toprule
\textbf{\color{spPrimary}الرمز} & \textbf{\color{spPrimary}دوره وعمره} \\
\midrule
\en{Access JWT} & رمز حامل (\en{bearer}) قصير العمر جدًّا — \textbf{15 دقيقة} —
يُرفَق مع كل طلب في ترويسة \en{Authorization}. يحمل هويّة المستخدم وأدواره موقّعةً. \\
\en{Refresh token} & رمز عشوائي أطول عمرًا (7 أيام، أو 30 مع «تذكُّرني») وظيفته
\emph{الوحيدة} استبدال رمز الوصول المنتهي دون إعادة إدخال كلمة المرور. \\
\bottomrule
\end{tabular}
\end{center}

\begin{why}
لماذا رمز وصول قصير جدًّا؟ لأن الـ\en{JWT} — بمجرد إصداره — لا يمكن سحبه فعليًّا حتى
ينتهي (فحصه يتم دون الرجوع للقاعدة). فلو سُرِّب، نريد نافذة خطر ضيّقة: 15 دقيقة فقط.
والثمن (تجديد متكرّر) رخيص لأن التجديد يتم بصمت عبر رمز التحديث دون إزعاج المستخدم.
\end{why}

\subsubsection{كيف يُولَّد رمز التحديث ويُخزَّن؟}
النقطة الحسّاسة: رمز التحديث \emph{لا} يُخزَّن كما هو في القاعدة، بل تُخزَّن
\emph{تجزئته} فقط في الحقل \code{RefreshToken.TokenHash}. هنشوف التوليد والتجزئة من
\code{TokenService}:

\begin{english}
\begin{lstlisting}[language={[Sharp]C}]
private static string GenerateSecureRefreshToken()
{
    var bytes = RandomNumberGenerator.GetBytes(64);   // 64 crypto-random bytes
    return Convert.ToBase64String(bytes);
}

private static string HashRefreshToken(string token)
{
    var bytes = SHA256.HashData(Encoding.UTF8.GetBytes(token));
    return Convert.ToHexString(bytes);                // store the hash, not the token
}
\end{lstlisting}
\end{english}

الشرح:
\begin{itemize}
\item \code{RandomNumberGenerator.GetBytes(64)} — مولّد عشوائي \emph{تشفيري} (لا
  \code{Random} العادي)، فالرمز غير قابل للتخمين.
\item الرمز الخام يُرسَل للعميل فقط. أمّا ما يُحفَظ في القاعدة فهو ناتج
  \code{HashRefreshToken} — تجزئة \en{SHA-256} بصيغة \en{hex}.
\end{itemize}

\begin{keyp}[نقطة مهمة: تسريب القاعدة لا يعطي رمزًا صالحًا]
لو تسرّبت قاعدة البيانات كاملة، لن يجد المهاجم إلا \emph{تجزئات} لا يمكن عكسها إلى
الرمز الأصلي. وبما أن النظام يقارن بتجزئة ما يُقدَّم إليه، فالتجزئة المسروقة عديمة
الفائدة. نفس المبدأ نطبّقه على رمز إعادة تعيين كلمة المرور (سيأتي بعد قليل).
\end{keyp}

\subsubsection{التدوير عند كل استخدام}
رمز التحديث \textbf{يُدوَّر}: كل مرة يُستخدَم فيها، يُبطَل القديم فورًا ويُصدَر بديل
جديد. وإذا عاد رمزٌ مُبطَل للظهور مرّة أخرى، فهذه إشارة سرقة، فنُلغي \emph{كل} جلسات
المستخدم احتياطًا. هنشوف \code{RotateRefreshTokenAsync} (مختصرًا):

\begin{english}
\begin{lstlisting}[language={[Sharp]C}]
// Reuse detection: a revoked token presented again means it was already
// rotated -- a replay attack or a stale client. Revoke the whole family
// so a stolen token can never be exchanged.
if (stored.IsRevoked)
{
    await RevokeAllForUserAsync(stored.UserId, "Refresh token reuse detected", ct);
    return null;
}

// Otherwise: mark the old one revoked and issue a fresh replacement pair.
stored.IsRevoked = true;
stored.RevokedAt = clock.UtcNow;
stored.RevokedReason = "Rotated";
var pair = IssueTokens(stored.User, roles, stored.User.ActiveRole, rememberMe);
stored.ReplacedByTokenId = newToken.Id;   // chain old -> new for traceability
\end{lstlisting}
\end{english}

الشرح:
\begin{itemize}
\item \code{if (stored.IsRevoked)} — لو الرمز المُقدَّم مُبطَل سلفًا فهذا «إعادة
  استخدام»؛ نستدعي \code{RevokeAllForUserAsync} فنقطع كل جلسات المستخدم.
\item \code{stored.IsRevoked = true} — نُبطِل القديم أوّلًا، ثم نُصدر زوجًا جديدًا،
  فلا يبقى رمزان صالحان في وقت واحد.
\item \code{ReplacedByTokenId} — نربط القديم بالجديد، فيصير لدينا سلسلة يمكن تتبّعها.
\end{itemize}

\subsubsection{نفس الحماية لرمز إعادة تعيين كلمة المرور}
حقل \code{PasswordResetToken.TokenHash} يتبع نفس الفلسفة: يُولَّد رمز خام، يُرسَل
بالبريد، ويُخزَّن \emph{مُجزَّأً} فقط، صالحًا لساعة واحدة ولمرّة واحدة:

\begin{english}
\begin{lstlisting}[language={[Sharp]C}]
var rawToken = Convert.ToHexString(RandomNumberGenerator.GetBytes(32));
db.PasswordResetTokens.Add(new PasswordResetToken
{
    UserId    = user.Id,
    TokenHash = Hash(rawToken),        // only the hash is persisted
    ExpiresAt = now + TokenLifetime,   // 1 hour, single use (UsedAt)
});
// the raw token travels only inside the emailed link, never to the DB
\end{lstlisting}
\end{english}

لاحظ أيضًا أن معالج «نسيت كلمة المرور» \emph{لا} يكشف إن كان البريد مسجّلًا أم لا —
يعيد نجاحًا صامتًا دائمًا — حتى لا يُستغَل لاكتشاف الحسابات الموجودة.

% ====================================================================
\subsection{خوارزميّة التوقيع: RS256 مقابل HS256}

الـ\en{JWT} يُوقَّع كي لا يستطيع أحد تزييف محتواه. أمامنا خياران شائعان، واختار
\en{ScholarPath} الأأمن منهما:

\begin{center}
\begin{tabular}{L{3.4cm} L{5.5cm} L{5.5cm}}
\toprule
\textbf{\color{spPrimary}الوجه} & \textbf{\color{spPrimary}\en{HS256} (متماثل)} & \textbf{\color{spPrimary}\en{RS256} (غير متماثل)} \\
\midrule
المفتاح & سرّ واحد مشترك يُوقّع ويتحقّق & زوج: مفتاح خاصّ يوقّع، وعامّ يتحقّق \\
مَن يحتاج السرّ؟ & كل طرف يتحقّق يحتاج السرّ نفسه & المتحقّق يحتاج المفتاح العامّ فقط \\
خطر التسريب & مَن يعرف السرّ يستطيع \emph{تزوير} الرموز & المفتاح العامّ يتحقّق ولا يوقّع \\
التدوير & صعب (السرّ منتشر في كل مكان) & سهل عبر منفذ مفاتيح مستقلّ \\
\bottomrule
\end{tabular}
\end{center}

في \code{TokenService} نوقّع بالمفتاح الخاصّ الآتي من منفذ المفاتيح
\code{IJwtKeyProvider}:

\begin{english}
\begin{lstlisting}[language={[Sharp]C}]
// RS256: sign with the RSA private key supplied by the key provider
// (Azure Key Vault in production, a local key in development).
var signingKey = keyProvider.GetSigningKey();
var creds = new SigningCredentials(signingKey, SecurityAlgorithms.RsaSha256);

var token = new JwtSecurityToken(
    issuer:   _opts.Issuer,
    audience: _opts.Audience,
    claims:   claims,
    notBefore: clock.UtcNow.UtcDateTime,
    expires:   expires.UtcDateTime,
    signingCredentials: creds);

return new JwtSecurityTokenHandler().WriteToken(token);
\end{lstlisting}
\end{english}

الشرح:
\begin{itemize}
\item \code{keyProvider.GetSigningKey()} — يجلب المفتاح \emph{الخاصّ}
  (\en{RSA private key})؛ في الإنتاج من \en{Azure Key Vault}، وفي التطوير مفتاح محلّي.
\item \code{SecurityAlgorithms.RsaSha256} — هذا هو \en{RS256} بالضبط.
\item \code{claims} — تحمل هويّة المستخدم وأدواره؛ لأنها موقّعة، أي تلاعب بها يُبطِل
  التوقيع فيُرفَض الرمز.
\end{itemize}

\begin{why}
لماذا \en{RS256}؟ لأن الخدمات التي \emph{تتحقّق} من الرمز تحتاج المفتاح
\emph{العامّ} فقط، والمفتاح العامّ لا يستطيع التوقيع؛ فحتى لو تسرّب لا يستطيع أحد
تزوير رموز. أمّا \en{HS256} فسرّه الواحد يوقّع \emph{ويتحقّق}، فمن يقرأه يقدر يزوّر.
وبفصل المفتاح خلف المنفذ \code{IJwtKeyProvider} صار \textbf{تدوير المفاتيح} عملية
تشغيلية بلا إعادة نشر للكود.
\end{why}

% ====================================================================
\subsection{الدخول الاجتماعي (SSO) والتحقّق بخطوتين (2FA)}

إلى جانب كلمة المرور، يدعم النظام الدخول عبر \en{Google} و\en{Microsoft} بمعيار
\en{OAuth 2.0} (تدفّق كود التفويض)، ويديره \code{SsoService}. الفكرة أن المستخدم
يوثّق نفسه عند \en{Google}، فيعود متصفّحه إلينا حاملًا \emph{كودًا} مؤقّتًا، ونحن
نبادل هذا الكود بمعلومات المستخدم في مكالمتين خلفيتين.

هنشوف قلب المبادلة من \code{SsoService.ExchangeCodeForTokenAsync} (مختصرًا):

\begin{english}
\begin{lstlisting}[language={[Sharp]C}]
var form = new Dictionary<string, string>
{
    ["code"]          = code,
    ["client_id"]     = provider.ClientId!,
    ["client_secret"] = provider.ClientSecret!,
    ["redirect_uri"]  = redirectUri,
    ["grant_type"]    = "authorization_code",
};

var response = await client.PostAsync(tokenEndpoint, content, ct);
// on success -> read the access_token, then GET the userinfo endpoint
// (email, name, picture) to build our SsoUserInfo
\end{lstlisting}
\end{english}

الشرح:
\begin{itemize}
\item \code{["code"]} — الكود المؤقّت العائد من المزوّد؛ يُبادَل مرّة واحدة فقط.
\item \en{client secret} — سرّ تطبيقنا لدى المزوّد، يأتي من الإعدادات لا من
  الكود، فلا يظهر أبدًا للمتصفّح.
\item بعد الحصول على \en{access token} نطلب نقطة \en{userinfo} لجلب البريد والاسم
  والصورة، ونبني منها \code{SsoUserInfo} لإنشاء الحساب أو ربطه.
\end{itemize}

\begin{secnote}
عنوان العودة (\en{redirect URI}) لا يُعكَس كما جاء من العميل؛ يفضّل النظام عنوانًا
ثابتًا من الإعدادات، أو يطابقه بقائمة سماح صريحة. هذا يسدّ ثغرة إعادة توجيه الكود إلى
موقع مهاجم — نقطة ضعف كلاسيكية في \en{OAuth}.
\end{secnote}

\begin{note}
للحسابات التي تريد حماية أقوى، يدعم النظام \textbf{التحقّق بخطوتين (\en{2FA})} عبر
\en{TOTP} — رمز لمرّة واحدة يتغيّر كل 30 ثانية من تطبيق مصادقة (مثل \en{Google
Authenticator}) — مبنيًّا على تدفّق \en{TwoFactor} في \en{ASP.NET Identity}. حتى لو
سُرِّبت كلمة المرور، يبقى الرمز الدوري حاجزًا ثانيًا.
\end{note}

% ====================================================================
\subsection{التفويض: الأدوار وفحوص الملكيّة}

بعد أن عرفنا \emph{مَن} المستخدم، يبدأ التفويض. يعمل على مستويين متكاملين:

\textbf{أولًا: الأدوار (\en{roles}).} لكل مستخدم دور أو أكثر، ويُحدَّد المسموح لكل دور:

\begin{center}
\begin{tabular}{L{3.4cm} L{11.3cm}}
\toprule
\textbf{\color{spPrimary}الدور} & \textbf{\color{spPrimary}صلاحيّته العامّة} \\
\midrule
\en{SuperAdmin} & أعلى صلاحية؛ إدارة المسؤولين والإعدادات الحسّاسة للمنصّة. \\
\en{Admin} & إدارة المحتوى والمستخدمين والمراجعات ومهامّ التشغيل اليومية. \\
\en{Student} & البحث عن المنح، التقديم، حجز الاستشارات، المشاركة في المجتمع. \\
\en{Consultant} & استقبال الحجوزات، عقد الجلسات، إدارة ملفّه المهني. \\
\en{Company} & نشر المنح ومراجعة الطلبات (مزوّد المنح). \\
\bottomrule
\end{tabular}
\end{center}

الأدوار تحمى نقاط النهاية على مستوى واسع (مثلًا: لوحة الإدارة لِـ\en{Admin} فقط).
لكن الدور وحده لا يكفي: كوْن المستخدم \en{Student} لا يعني أن له حقّ تعديل طلب
\emph{طالب آخر}.

\textbf{ثانيًا: فحوص الملكيّة (\en{ownership checks}).} داخل كل \en{handler} يتحقّق
النظام أن المورد المطلوب يخصّ صاحب الطلب فعلًا — بمقارنة مُعرِّف المستخدم الحالي
(من المنفذ \code{ICurrentUserService}) بمالك المورد — وإلّا يرفض العملية.

\begin{keyp}[نقطة مهمة: الدور طبقة أولى، والملكيّة طبقة ثانية]
الدور يجيب «هل يحقّ لهذا \emph{النوع} من المستخدمين استعمال هذه الوظيفة؟»، والملكيّة
تجيب «وهل يملك \emph{هذا} المستخدم \emph{هذا} المورد بالذات؟». الطبقتان معًا تمنعان
ثغرة الوصول المباشر لمورد غير مملوك (\en{IDOR}) — من أخطر أخطاء تطبيقات الويب.
\end{keyp}

% ====================================================================
\subsection{مخطّط تدفّق تسجيل الدخول}

نجمع كل ما سبق في تدفّق واحد: من وصول الطلب حتى إصدار الرمزين. اقرأ الجدول من أعلى
إلى أسفل:

\begin{center}
\begin{tabular}{L{1.1cm} L{13.4cm}}
\toprule
\textbf{\color{spPrimary}\#} & \textbf{\color{spPrimary}ما يحدث} \\
\midrule
1 & يصل طلب \en{POST} إلى نقطة الدخول ومعه البريد وكلمة المرور. \\
2 & يبحث \code{LoginCommandHandler} عن المستخدم بالبريد المُطبَّع؛ إن لم يوجد يُسجَّل فشل ويُرفَض برسالة عامّة. \\
3 & إن كان الحساب مقفولًا (\code{LockoutEnd} في المستقبل) يُرفَض فورًا. \\
4 & تُتحقّق كلمة المرور عبر \code{IPasswordHasher.Verify}؛ عند الفشل يُزاد العدّاد وقد يُقفَل الحساب (5 محاولات $\rightarrow$ 30 دقيقة). \\
5 & إن كان الحساب موقوفًا أو مُعطَّلًا يُرفَض الدخول رغم صحّة كلمة المرور. \\
6 & عند النجاح: تُصفَّر المحاولات، ويُحدَّث \code{LastLoginAt}، وتُسجَّل محاولة ناجحة. \\
7 & يجلب النظام أدوار المستخدم، ثم يستدعي \code{TokenService.IssueTokens}. \\
8 & يُنشأ \en{access JWT} موقّع بـ\en{RS256} (15 دقيقة) $+$ \en{refresh token} عشوائي. \\
9 & تُخزَّن \emph{تجزئة} رمز التحديث في \code{RefreshToken.TokenHash}، ويُعاد الزوج للعميل. \\
\bottomrule
\end{tabular}
\end{center}

\begin{qa}
\textbf{سؤال متوقّع:} «لو انتهى رمز الوصول والمستخدم لسّه شغّال، هيحصل إيه؟»\\
\textbf{الإجابة:} العميل يرسل رمز التحديث إلى نقطة التجديد، فيدوّره
\code{RotateRefreshTokenAsync}: يُبطِل القديم ويصدر زوجًا جديدًا، فتستمرّ الجلسة بلا
إدخال كلمة مرور. ولو ظهر رمز مُبطَل ثانيةً، تُلغى كل الجلسات كإجراء ضدّ السرقة.
\end{qa}

% ====================================================================
\begin{summary}[الخلاصة في سطور]
\begin{itemize}
\item \textbf{التوثيق} يجيب «مَن أنت؟»، و\textbf{التفويض} يجيب «تقدر تعمل إيه؟» —
  الأول عند الدخول، والثاني في كل طلب.
\item \code{ApplicationUser : IdentityUser<Guid>} يبني على \en{ASP.NET Core Identity}:
  كلمة مرور مُجزّأة بـ\en{PBKDF2}$+$مِلح، قفل بعد 5 محاولات، وتأكيد بريد.
\item النظام يُصدر \textbf{رمزين}: \en{access JWT} قصير (15 دقيقة) و\en{refresh token}
  يُدوَّر عند كل استخدام، ويُخزَّن \emph{مُجزَّأً} في \code{RefreshToken.TokenHash}.
\item نفس التجزئة تحمي \code{PasswordResetToken}: تسريب القاعدة لا يعطي رمزًا صالحًا.
\item التوقيع بـ\en{RS256} (غير متماثل) أأمن من \en{HS256}: المتحقّق يحتاج المفتاح
  العامّ فقط، والتدوير سهل عبر المنفذ \code{IJwtKeyProvider}.
\item \code{SsoService} يبادل أكواد \en{OAuth} مع \en{Google} و\en{Microsoft}،
  ويُتاح \en{2FA} عبر \en{TOTP}.
\item التفويض طبقتان: \textbf{الأدوار} (\en{SuperAdmin/Admin/Student/Consultant/Company})
  ثم \textbf{فحوص الملكيّة} داخل الـ\en{handlers} لسدّ ثغرات \en{IDOR}.
\end{itemize}
\end{summary}

\clearpage
% ch04-security
\section{الأمان بعُمق (Defense in Depth)}

في الفصل السابق شرحنا كيف يعرف النظام \emph{مَن أنت} (التوثيق والرموز). في هذا
الفصل نجيب عن سؤال مختلف: ماذا يحدث \emph{بعد} أن نعرفك؟ كيف نحمي بياناتك وبيانات
غيرك حتى لو نجح مهاجم في اختراق طبقة واحدة؟ الجواب هو فلسفة اسمها \en{Defense in
Depth}، أو «الأمان بعُمق». وسنسير فيها \emph{واحدة واحدة}.

\begin{idea}[الفكرة: طبقات مستقلّة لا حائط واحد]
تخيّل قلعة قديمة: خندق، ثم سور، ثم بوابة، ثم حرّاس بالداخل. لو قفز أحدهم فوق
الخندق، ما زال السور يوقفه. الأمان بعُمق نفس الفكرة في البرمجيات: بدلًا من الاعتماد
على ضابط أمان \emph{واحد}، نضع عدّة ضوابط \textbf{مستقلّة}، كلٌّ منها يعمل وحده. فشل
طبقة لا يعني سقوط النظام، لأن الطبقة التالية ما زالت قائمة. هذا يحوّل «ثغرة واحدة =
كارثة» إلى «ثغرة واحدة = إزعاج محدود».
\end{idea}

\subsection{خريطة الطبقات كاملةً}

قبل الغوص في التفاصيل، هذه صورة عامّة لكل طبقات الحماية في \en{ScholarPath}. كل صفّ
في الجدول ضابط مستقلّ؛ وباقي الفصل يشرح الصفوف الحسّاسة منها بالتفصيل (التشفير،
الرفع، الفحص، الخصوصيّة، والتدقيق).

\begin{center}
\begin{tabular}{L{3.4cm} L{9.4cm}}
\toprule
\textbf{\color{spPrimary}الطبقة} & \textbf{\color{spPrimary}الضابط} \\
\midrule
النقل \en{(Transport)} & \en{HTTPS} في كل مكان؛ رمز \en{JWT bearer} على الـ \en{API}؛ محاور \en{SignalR} موثَّقة \\
الهويّة \en{(Identity)} & \en{ASP.NET Core Identity} — كلمات سرّ مُجزّأة \en{(hashed)}، قفل عند التخمين، تأكيد بريد \\
الرموز \en{(Tokens)} & رمز وصول \en{JWT} قصير العمر + رموز تحديث \en{(refresh)} دوّارة تُخزَّن \emph{مُجزّأة} \\
التخويل \en{(Authorization)} & أدوار (\en{Student}/\en{ScholarshipProvider}/\en{Consultant}/\en{Admin}) + ملكيّة لكلّ مورد \\
الأسرار والمفاتيح & مزوّدات \en{Azure Key Vault} (مفتاح تشفير الحقول، مفتاح توقيع \en{JWT}) \\
البيانات المخزَّنة & تشفير حقول \en{AES-256-GCM} على الأعمدة الحسّاسة \\
الرفع \en{(Uploads)} & فحص-قبل-التخزين بـ \en{ClamAV}، بسياسة \en{fail-closed} \\
الخصوصيّة/\en{PII} & تنقية النصّ قبل الذكاء الاصطناعي؛ مراجع فضفاضة؛ تصدير/حذف \en{GDPR} \\
التكامل \en{(Integrity)} & مفاتيح \en{idempotency}، فهارس فريدة مُرشَّحة، \code{RowVersion} \\
قابليّة التدقيق & سجلّ \code{AuditLog} إلحاقيّ فقط (قبل/بعد بصيغة \en{JSON}) \\
الحذف الناعم & مُرشِّحات استعلام عامّة تُخفي الصفوف المحذوفة من كلّ استعلام \\
\bottomrule
\end{tabular}
\end{center}

\begin{why}
لاحظ أنّ الطبقات \emph{لا تتكلّم مع بعضها}. التشفير لا يعرف شيئًا عن فحص الفيروسات،
والتدقيق لا يعرف شيئًا عن الأدوار. هذا الاستقلال مقصود: لو ارتبطت الطبقات ببعضها،
لأصبح خطأ في واحدة قادرًا على تعطيل الباقي. الاستقلال هو ما يمنح «العُمق» معناه.
\end{why}

\subsection{التشفير عند التخزين: AES-256-GCM}

بعض الأعمدة في قاعدة البيانات حسّاسة جدًّا — مثل نبذة المستخدم \code{Biography} أو
الملاحظات الشخصيّة على الطلب \code{PersonalNotes}. لو تسرّبت نسخة احتياطيّة من قاعدة
البيانات، لا نريد أن يقرأ أحدٌ هذه الحقول. الحلّ: نشفّرها \emph{قبل} أن تُكتَب، على
مستوى التطبيق نفسه.

\begin{idea}[الفكرة: تشفير مُصادَق (AEAD) بمغلّف ذاتيّ الوصف]
نستخدم خوارزمية \en{AES-256-GCM}. كلمة \en{GCM} مهمّة: إنها خوارزمية
\textbf{مُصادَقة} \en{(AEAD)} تعطيك \emph{السرّيّة والسلامة معًا} في خطوة واحدة.
المفتاح 32 بايت (256 بت)، ولكلّ عمليّة تشفير \en{nonce} عشوائيّ طازج طوله 12 بايت،
ووسم مُصادَقة \en{(tag)} طوله 16 بايت. القيمة المخزَّنة «مغلّف ذاتيّ الوصف»:
البادئة الحرفيّة \code{enc:v1:} متبوعةً بـ \en{Base64} للسلسلة
(\en{nonce} ثم \en{tag} ثم \en{ciphertext}). لأنّ الـ \en{nonce} عشوائيّ، فإنّ تشفير
النصّ نفسه مرّتين يعطي ناتجًا مختلفًا في كلّ مرّة.
\end{idea}

المنفذ \code{IFieldEncryptionService} يُنفّذه الصنف \code{AesGcmFieldEncryptionService}.
هنشوف إزاي بيبني المغلّف عند التشفير:

\begin{english}
\begin{lstlisting}[language={[Sharp]C}]
public const string EnvelopePrefix = "enc:v1:";
private const int NonceSizeBytes = 12;  // AES-GCM standard nonce
private const int TagSizeBytes   = 16;  // authentication tag

public string? Encrypt(string? plaintext)
{
    if (plaintext is null) return null;                 // null -> null
    var plaintextBytes = Encoding.UTF8.GetBytes(plaintext);

    var nonce = new byte[NonceSizeBytes];
    RandomNumberGenerator.Fill(nonce);                  // fresh random nonce

    var tag = new byte[TagSizeBytes];
    var ciphertext = new byte[plaintextBytes.Length];
    using (var aes = new AesGcm(_key, TagSizeBytes))
        aes.Encrypt(nonce, plaintextBytes, ciphertext, tag);

    // Envelope payload: nonce (12) + tag (16) + ciphertext
    var payload = new byte[NonceSizeBytes + TagSizeBytes + ciphertext.Length];
    // ... copy nonce, then tag, then ciphertext into payload ...
    return EnvelopePrefix + Convert.ToBase64String(payload);
}
\end{lstlisting}
\end{english}

نشرح أهمّ الأسطر:
\begin{itemize}
\item \code{if (plaintext is null) return null} — القيمة الفارغة تمرّ كما هي؛ عمود
  فارغ يبقى فارغًا ولا نضيّع تشفيرًا عليه.
\item \code{RandomNumberGenerator.Fill(nonce)} — هنا يُولَد الـ \en{nonce} العشوائيّ
  الجديد. هذا السطر تحديدًا هو سرّ أنّ نفس النصّ لا يعطي أبدًا نفس التشفير.
\item \code{aes.Encrypt(nonce, plaintextBytes, ciphertext, tag)} — خطوة واحدة تُنتج
  النصّ المشفَّر \emph{و} وسم المُصادَقة معًا.
\item السطر الأخير يجمّع (\en{nonce} + \en{tag} + \en{ciphertext}) ويضع أمامها
  البادئة \code{enc:v1:} — فتصير القيمة تصف نفسها بنفسها.
\end{itemize}

وعلى الجانب الآخر، عند فكّ التشفير، يُتحقَّق من الوسم \en{(tag)} أوّلًا: أيّ عبثٍ
ببِتّ واحد، أو اقتطاع للقيمة، أو مفتاح خاطئ — يجعل الدالّة \code{Decrypt}
\textbf{ترمي استثناءً} بدلًا من أن تُرجع قمامة. هذا هو معنى «مُصادَق».

\begin{keyp}[نقطة مهمة: versioned + طرح آمن]
المقطع \code{v1} في البادئة ليس زينة — إنه \textbf{رقم إصدار}. لو دوّرنا المفتاح
مستقبلًا، نُدخل \code{enc:v2:} للقيم الجديدة بينما نظلّ قادرين على فكّ القيم القديمة
\code{v1}. وأهمّ من ذلك: أيّ قيمة \emph{لا} تبدأ بالبادئة (نصّ قديم كُتب قبل تفعيل
الميزة) تُرجَع كما هي دون تعديل، ثم يُعاد تشفيرها تلقائيًّا عند أوّل كتابة لها. هذا
«الطرح الآمن» يعني تفعيل التشفير على نظام حيّ \emph{بلا توقّف ولا هجرة بيانات}.
\end{keyp}

والأجمل أنّ بقيّة الكود لا تشعر بشيء من هذا. نربط التشفير بالعمود عبر \en{value
converter} في \en{EF Core}، فيصير النصّ صريحًا في الذاكرة ومشفَّرًا في القاعدة:

\begin{english}
\begin{lstlisting}[language={[Sharp]C}]
public sealed class EncryptedStringConverter : ValueConverter<string?, string?>
{
    public EncryptedStringConverter(IFieldEncryptionService encryption)
        : base(
            plaintext => encryption.Encrypt(plaintext),  // to the database
            stored    => encryption.Decrypt(stored))     // from the database
    { }
}
\end{lstlisting}
\end{english}

بمجرّد إسناد هذا المحوّل لعمود، تظلّ معالِجات الأوامر والاستعلامات تقرأ وتكتب نصًّا
عاديًّا، ولا تعلم أبدًا أنّ العمود مشفَّر عند التخزين. طبقة واحدة، ومسؤوليّة واحدة.

\begin{secnote}
لماذا \en{GCM} وليس تشفيرًا عاديًّا؟ التشفير العاديّ يحمي السرّيّة لكنه \emph{لا}
يكشف العبث. \en{GCM} يعطيك الاثنين. وميزته على تشفير القاعدة الشفّاف \en{(TDE)} أنّ
المفتاح لا يعيش أبدًا في الكود ولا في ملفّات الإعداد — بل يأتي من \en{Key Vault} —
وحتى نسخة احتياطيّة مسروقة أو مدير قاعدة بيانات فضوليّ لن يرى إلّا نصًّا مشفَّرًا
لهذه الأعمدة، وهو ما لا يضمنه تشفير القاعدة وحده بمجرّد أن يُرجِع استعلامٌ صفًّا.
\end{secnote}

\subsubsection{إدارة المفاتيح}

المفتاح يأتي من المنفذ \code{IFieldEncryptionKeyProvider}، وله مُنفِّذان: واحد يقرأ
من \en{Key Vault} في الإنتاج، وآخر محلّيّ للتطوير. وعند الإقلاع نتحقّق أنّ طوله
\textbf{32 بايت بالضبط}، وإلّا يرفض النظام أن يبدأ أصلًا. وهناك منفذ مستقلّ آخر هو
\code{IJwtKeyProvider} لمفتاح توقيع \en{JWT}. وضع كلّ مفتاح خلف منفذه الخاصّ هو ما
يجعل \textbf{تدوير المفاتيح} عمليّة تشغيليّة بسيطة، لا إعادة نشرٍ للكود.

\begin{qa}
«ماذا لو تسرّب المفتاح نفسه؟» — عندها يفشل هذا الضابط، لكن تذكّر فلسفة العُمق:
التشفير طبقة واحدة فقط. ما زالت هناك طبقات الهويّة والتخويل والتدقيق. ثم إنّ المفتاح
لا يعيش في الكود ولا في \en{git}؛ يعيش في \en{Key Vault} بصلاحيّات محدودة، ويمكن
تدويره فورًا إلى \code{enc:v2:} دون أن نفقد القيم القديمة.
\end{qa}

\subsection{رفع الملفات: مسار مُعتِم ومزوّدان}

المستخدمون يرفعون مرفقات: مستندات الطلبات، ملفّات الترقية، مرفقات المنتدى،
وتسجيلات الجلسات. المنفذ \code{IBlobStorageService} يُنفّذه \code{FileStorageService}،
وله مزوّدان يُختاران من الإعداد \code{Storage:Provider}: \en{Local} (قرص، للتطوير)
أو \en{AzureBlob} (تخزين \en{Azure}).

هنشوف إزاي بيتخزّن الملفّ وبيرجّع مساره:

\begin{english}
\begin{lstlisting}[language={[Sharp]C}]
public async Task<string> UploadAsync(
    Stream content, string fileName, string contentType,
    string container, CancellationToken ct)
{
    // Unique key keeps two same-named uploads from colliding.
    var key = $"{Guid.NewGuid():N}/{fileName}";

    if (UseAzure)
    {
        var blob = ContainerClient(container).GetBlobClient(key);
        await blob.UploadAsync(content, /* ...ContentType... */, ct);
        return $"azure:{container}/{key}";     // opaque path
    }
    // ... Local provider writes under LocalRoot()/container/key ...
    return $"local:{container}/{key}";
}
\end{lstlisting}
\end{english}

نشرح أهمّ الأسطر:
\begin{itemize}
\item \code{var key = \$"\{Guid.NewGuid():N\}/\{fileName\}"} — يُسبَق اسم الملفّ
  بـ \en{GUID} عشوائيّ، فملفّان بنفس الاسم لن يتصادما أبدًا، ولا يستطيع مستخدم أن
  يخمّن مسار ملفّ آخر.
\item القيمة المُرجَعة \code{azure:container/key} أو \code{local:container/key} هي
  \textbf{مسار مُعتِم} \en{(opaque)} يحمل اسم المزوّد بداخله — فيعمل التنزيل والحذف
  للمزوّدَين معًا دون لبس.
\end{itemize}

وحين نُنشئ حاوية \en{Azure}، نُنشئها بـ \code{PublicAccessType.None} — أي \emph{غير
عامّة أبدًا}. لا يوجد رابط عامّ للملفّ؛ كلّ تنزيل يمرّ عبر الـ \en{API} الموثَّق الذي
يتحقّق من صلاحيّة الطالب أوّلًا.

\begin{note}
المسار المُعتِم فكرة صغيرة بأثر كبير: الكود الأعلى يتعامل مع «سلسلة» واحدة يخزّنها في
عمود، ولا يهتمّ إن كان الملفّ على قرصٍ محلّيّ أو في \en{Azure}. تبديل المزوّد يصبح
تغيير إعدادٍ واحد، لا تغيير كود.
\end{note}

\subsection{الفحص قبل التخزين: ClamAV بسياسة fail-closed}

الملفّ المرفوع قد يكون خبيثًا. لذلك نفحصه \emph{قبل} أن نخزّنه، \textbf{وهو ما زال
تيّارًا في الذاكرة}. المنفذ \code{IFileScanService} يُنفّذه \code{ClamAvFileScanService}،
الذي يرسل البايتات إلى خادم \en{ClamAV} عبر بروتوكول \en{INSTREAM}. الأحكام ثلاثة:
\textbf{Clean} (نظيف، نكمل)، و\textbf{Infected} (مُصاب، نرفض ونسجّل التوقيع)،
و\textbf{ScanUnavailable} (تعذّر الفحص، نرفض).

هنشوف قلب الدالّة: كيف يترجم ناتج الفاحص إلى حُكم، وكيف يتصرّف عند الفشل:

\begin{english}
\begin{lstlisting}[language={[Sharp]C}]
switch (result.Result)
{
    case ClamScanResults.Clean:
        return new FileScanResult(FileScanVerdict.Clean, null);

    case ClamScanResults.VirusDetected:
        var signature = result.InfectedFiles?.FirstOrDefault()?.VirusName
                        ?? "unknown";
        return new FileScanResult(FileScanVerdict.Infected, signature);

    default: // Error / Unknown -> treat as unavailable
        return new FileScanResult(FileScanVerdict.ScanUnavailable, ...);
}
// ... and if the daemon is unreachable:
catch (Exception ex)   // socket timeout, protocol error, no daemon
{
    // Fail closed: an unscanned file is never stored.
    return new FileScanResult(FileScanVerdict.ScanUnavailable, "...");
}
\end{lstlisting}
\end{english}

نشرح المنطق:
\begin{itemize}
\item \code{Clean} فقط هو الذي يمرّ ويُخزَّن؛ التيّار يُعاد لبدايته لتُكتَب نفس
  البايتات المفحوصة.
\item \code{VirusDetected} يُرجِع الحُكم \code{Infected} مع توقيع الفيروس ويُسجَّل
  تحذيرًا؛ الملفّ يُرفَض.
\item السطر \code{catch (Exception ex)} هو جوهرة الفصل: أيّ تعذّرٍ في الوصول للخادم
  أو مهلة أو خطأ بروتوكول يُترجَم إلى \code{ScanUnavailable}، ومسار الرفع يرفض هذا
  الحُكم.
\end{itemize}

\begin{secnote}
هذه هي سياسة \en{fail-closed} (الفشل إلى الإغلاق). القاعدة: عند الشكّ، \emph{ارفض}.
لو تعطّل الفاحص، لا نقول «لا بأس، خزّن الملفّ»؛ بل نرفضه. لأنّ البديل — تخزين ملفّ
\emph{غير مفحوص} — يفتح نافذة زمنيّة قد يصبح فيها ملفّ خبيث في متناول اليد. الفحص
أوّلًا على التيّار، والرفض عند العجز، يغلقان هذه النافذة تمامًا. (والميزة مُتحكَّم بها
عبر \code{FileScanning:Enabled}، وفي التطوير يُستخدَم \code{NoOpFileScanService}.)
\end{secnote}

\subsection{بيانات الأشخاص (PII)}

بيانات الأشخاص \en{(Personally Identifiable Information)} — الإيميل، رقم الكارت، رقم
التليفون — نتعامل معها بحذر خاصّ، عبر عدّة طبقات في مسار الذكاء الاصطناعي وفي نموذج
البيانات وفي دورة الحياة.

\begin{idea}[الفكرة: تنقية قبل التخزين وقبل المزوّد]
حين يكتب المستخدم رسالة للمساعد الذكيّ، تمرّ الرسالة أوّلًا على دالّة
\code{RedactPii()} \textbf{قبل} أن نخزّنها \emph{وقبل} أن نرسلها لأيّ مزوّد خارجيّ.
فالنصّ الخام لا يغادر العمليّة أبدًا؛ ما نخزّنه (\code{AiInteraction.PromptText})
وما نرسله إلى \en{OpenAI} هو النسخة \emph{المُنقّاة}.
\end{idea}

هنشوف الدالّة الحقيقيّة وترتيب أنماطها:

\begin{english}
\begin{lstlisting}[language={[Sharp]C}]
internal static string RedactPii(string input)
{
    if (string.IsNullOrWhiteSpace(input)) return string.Empty;
    // Order matters: email, then card, then the most permissive phone last.
    var s = EmailRegex().Replace(input, "[redacted-email]");
    s     = CreditCardRegex().Replace(s, "[redacted-card]");
    s     = PhoneRegex().Replace(s, "[redacted-phone]");
    return s;
}

[GeneratedRegex(@"[\w\.-]+@[\w\.-]+\.\w{2,}")]      // email
private static partial Regex EmailRegex();
[GeneratedRegex(@"\b(?:\d[ -]*?){13,19}\b")]        // 13-19 digit card
private static partial Regex CreditCardRegex();
[GeneratedRegex(@"(?:\+?\d[\d\s\-\.]{8,}\d)")]      // 10+ digit phone
private static partial Regex PhoneRegex();
\end{lstlisting}
\end{english}

لماذا هذا الترتيب بالذات؟
\begin{itemize}
\item الإيميلات أوّلًا لأنها تتمركز حول علامة \code{@}، فيسهل تمييزها.
\item ثم الكروت: تسلسل من 13 إلى 19 رقمًا.
\item ثم التليفونات \textbf{أخيرًا}، لأنّ نمطها هو \emph{الأوسع} (أيّ 10 أرقام أو
  أكثر). لو جاء أوّلًا، لابتلع أرقام الكروت قبل أن يميّزها نمط الكارت. الترتيب هنا
  ليس تفصيلًا تجميليًّا؛ إنه صحّة.
\end{itemize}

وهناك طبقات مكمّلة حول الـ \en{PII}:
\begin{itemize}
\item \textbf{تخصيص بلا بيانات تواصُل} — الكتلة التعريفيّة التي نحقنها في المحادثة
  لتخصيص الإجابات تحمل حقائق أكاديميّة فقط (الجنسيّة، المستوى، المجال، المعدّل، الدول
  المفضّلة)، ولا تحمل إيميلًا ولا تليفونًا أبدًا.
\item \textbf{ذاكرة \en{AI} لكلّ مستخدم} — عند إعادة تشغيل تاريخ المحادثة نُرشّحه
  للمستخدم صاحب الطلب فقط، فحتى لو تسرّب \code{SessionId} لا يمكن أن يتسرّب نصّ
  محادثة مستخدمٍ آخر.
\item \textbf{عيّنات تدقيق التنقية} — الكيان \code{AiRedactionAuditSample} يخزّن
  عيّنات شهريّة من الرسائل المُنقّاة (حتى 50 عيّنة من الشهر المنصرم) لمراجعة بشريّة —
  حتى نُثبت أنّ المُنقّي يعمل فعلًا.
\item \textbf{مراجع فضفاضة \en{(loose FK)}} — الجداول عالية الحجم/التدقيق/التحليلات
  (المدفوعات، الإشعارات، الأصوات، الدردشة، سجلّ التدقيق، تفاعلات \en{AI}) تحمل
  \en{Guid} المستخدم \emph{بلا قيد مفتاح أجنبيّ}، فيمكن حذف مستخدم أو إخفاء هويّته
  دون أن ينكسر التاريخ تتابعيًّا.
\item \textbf{\en{GDPR}} — الكيان \code{UserDataRequest} يقود التصدير والحذف عبر
  وظيفتَي \code{DataExportJob} و\code{DataDeleteJob}؛ والنصّ الحرّ الحسّاس مشفَّر عند
  التخزين كما رأينا.
\end{itemize}

\begin{secnote}
لاحظ الفلسفة المتكرّرة: نقلّل ما نحتفظ به من الأساس. التنقية تحذف الـ \en{PII} قبل
أن يُكتَب، والمراجع الفضفاضة تجعل حذف المستخدم ممكنًا دون تدمير السجلّات المرتبطة.
أقلّ بيانات حسّاسة مخزَّنة = سطح هجوم أصغر.
\end{secnote}

\subsection{سجلّ التدقيق (AuditLog): إلحاقيّ فقط}

آخر طبقة نتناولها هي \textbf{قابليّة التدقيق}. كلّ فعل تجاريّ حسّاس يُسجَّل في جدول
\code{AuditLog} بصورة \en{append-only} (إلحاقيّ فقط) — نضيف صفوفًا ولا نعدّل ولا نحذف.
الخدمة \code{AuditService} تكتب الصفّ، وتحفظ الحالة \emph{قبل} و\emph{بعد} بصيغة
\en{JSON}:

\begin{english}
\begin{lstlisting}[language={[Sharp]C}]
var entry = new AuditLog
{
    ActorUserId  = currentUser.UserId, // null for system actions
    Action       = action,
    TargetType   = targetType,         // polymorphic: any entity kind
    TargetId     = targetId,
    BeforeJson   = beforeJson,          // state before the change
    AfterJson    = afterJson,           // state after the change
    IpAddress    = Trim(currentUser.IpAddress, 64),
    CorrelationId= Trim(currentUser.CorrelationId, 128),
    OccurredAt   = clock.UtcNow,
};
db.AuditLogs.Add(entry);
\end{lstlisting}
\end{english}

نشرح الحقول المهمّة:
\begin{itemize}
\item \code{ActorUserId} — مَن فعلها؛ ويكون فارغًا \en{(null)} إذا كان الفاعل هو
  النظام نفسه (وظيفة خلفيّة مثلًا).
\item \code{TargetType} و\code{TargetId} — مرجع \en{polymorphic}: يستطيع الصفّ أن
  يشير إلى \emph{أيّ} نوع كيان، فسجلٌّ واحد يخدم كلّ الكيانات.
\item \code{BeforeJson} و\code{AfterJson} — لقطتا الحالة قبل التغيير وبعده، فنعرف
  بالضبط \emph{ماذا} تغيّر، لا مجرّد أنّ شيئًا تغيّر.
\end{itemize}

\begin{keyp}[نقطة مهمة: لماذا «إلحاقيّ فقط»؟]
سجلّ التدقيق لا قيمة له إن أمكن تزويره. جعله \en{append-only} يعني أنّ التاريخ لا
يُعاد كتابته: مَن فعل ماذا، ومتى، وما القيمة قبل وبعد — كلّها محفوظة إلى الأبد. وحين
يجتمع هذا مع الحذف الناعم (الصفوف لا تُدمَّر بل تُخفى)، تحصل على أثرٍ كامل قابل
للتتبّع لأيّ نزاع أو تحقيق لاحق.
\end{keyp}

\begin{summary}[الخلاصة في سطور]
\begin{itemize}
\item \textbf{الأمان بعُمق} = طبقات ضوابط \emph{مستقلّة}؛ فشل طبقة لا يُسقط النظام،
  لأنّ الطبقة التالية ما زالت قائمة.
\item \textbf{التشفير عند التخزين} بـ \en{AES-256-GCM}: مفتاح 32 بايت، \en{nonce}
  عشوائيّ 12 بايت، وسم 16 بايت، ومغلّف ذاتيّ الوصف \code{enc:v1:}؛ مُصادَق فأيّ عبثٍ
  يُفشِل فكّ التشفير، ومُصدَّر \en{(versioned)} فالطرح آمن بلا توقّف.
\item \textbf{رفع الملفّات} عبر منفذٍ بمزوّدَين ومسارٍ مُعتِم بادئته \en{GUID}،
  والحاويات \code{PublicAccessType.None} — لا رابط عامّ أبدًا.
\item \textbf{الفحص قبل التخزين} بـ \en{ClamAV} على التيّار، بسياسة \en{fail-closed}:
  عند الشكّ ارفض؛ لا نُخزّن ملفًّا غير مفحوص.
\item \textbf{الـ \en{PII}} تُنقَّى قبل التخزين وقبل المزوّد (بترتيب: إيميل ثم كارت
  ثم تليفون)، مع تخصيصٍ بلا بيانات تواصُل، ومراجع فضفاضة، وتصدير/حذف \en{GDPR}.
\item \textbf{سجلّ التدقيق} \code{AuditLog} إلحاقيّ فقط، يحفظ الحالة قبل/بعد
  بـ \en{JSON} — أثرٌ لا يُزوَّر لكلّ فعلٍ حسّاس.
\end{itemize}
\end{summary}

\clearpage
% ch05-ai-concepts
\section{الذكاء الاصطناعيّ (1): المفاهيم الأساسية}

في هذا الفصل نبني الأساس النظريّ الذي يقوم عليه كلّ ما يخصّ الذكاء الاصطناعيّ في
\en{ScholarPath}. لن نكتب كودًا كثيرًا هنا؛ هدفنا أن نفهم \emph{الأفكار} أولًا:
كيف يحوّل الحاسوب النصّ إلى أرقام؟ وكيف يقيس أنّ جملتين متقاربتان في المعنى؟ وما
\en{RAG} ولماذا نحتاجه أصلًا؟ ثمّ نُنهي الفصل بالصورة الكاملة: عندنا ثلاث ميزات
ذكاء اصطناعيّ، لكنّ \textbf{واحدة فقط} منها تستعمل نموذجًا لُغويًّا حقيقيًّا. نمشي
واحدة واحدة، وكلّ فكرة نربطها بسطر كود حقيقيّ من المشروع.

% ─────────────────────────────────────────────────────────────────────────────
\subsection{الـ Embeddings — تحويل النصّ إلى أرقام}

\begin{idea}[الفكرة: \en{Embedding}]
الـ \en{Embedding} يأخذ أيّ نصّ — كلمة أو جملة أو فقرة — ويحوّله إلى مصفوفة أرقام
عشريّة ثابتة الطول اسمها \en{vector} (مُتّجه). الخاصّيّة الجوهريّة: \textbf{النصوص
المتقاربة في المعنى تُنتج مُتّجهات متقاربة عدديًّا}، بغضّ النظر عن الصياغة الحرفيّة،
بل وحتى عبر اللغات. سؤال بالإنجليزيّة ونفس السؤال بالعربيّة يُنتجان مُتّجهين يشيران
إلى نفس الاتّجاه تقريبًا داخل فضاء ذي 1536 بُعدًا.
\end{idea}

تخيّل خريطة ضخمة، لكن بدل بُعدين (طول وعرض) فيها 1536 بُعدًا. كلّ نصّ يصير نقطةً على
هذه الخريطة. المنح الدراسيّة تتجمّع في ناحية، وأسئلة الأمور الشائعة في ناحية، وملفّات
المستشارين في ناحية ثالثة. لمّا يسأل الطالب سؤالًا، نحوّل سؤاله إلى نقطة أيضًا، ثمّ
نبحث عن أقرب النقاط إليه — وتلك النقاط هي المستندات الأوثق صِلةً بسؤاله.

في \en{ScholarPath} نموذج الـ \en{Embeddings} هو \en{Azure OpenAI}
\code{text-embedding-3-small}، وهو يُنتج مُتّجهات بطول \textbf{1536} رقمًا. لكن هنا
تفصيلة مهمّة: لو خزّنّا المُتّجهات في القاعدة دون أن نُسجّل \emph{أيّ نموذج} ولّدها،
فسنقع في كارثة صامتة يوم نغيّر النموذج — لأنّ مُتّجهات نموذجٍ لا تُقارَن أبدًا
بمُتّجهات نموذجٍ آخر. الحلّ بسيط: نضع بادئة \code{sp:} أمام اسم النموذج، فتصير هويّة
المُتّجه واضحة ومحجوزة.

\textbf{هنشوف إزّاي الكود بيبني اسم النموذج بالبادئة \code{sp:}} — المقتطف من الملفّ
\code{AzureOpenAiEmbeddingService.cs}:

\begin{english}
\begin{lstlisting}[language={[Sharp]C}]
// Provider-neutral model tag. The prefix namespaces the vector space so
// documents embedded with a different model are never matched here.
public string ModelName =>
    IsAzureEmbeddingConfigured()
        ? $"sp:{opts.Value.AzureOpenAi.EmbeddingDeploymentName}"
        // -> "sp:text-embedding-3-small"
        : local.ModelName;   // fallback when Azure isn't configured

public int Dimensions => opts.Value.AzureOpenAi.EmbeddingDimensions; // 1536
\end{lstlisting}
\end{english}

نقرأ الكود سطرًا سطرًا:
\begin{itemize}
  \item \code{ModelName} خاصّيّة محسوبة: لو \en{Azure} مُهيّأ، تُرجِع الاسم مسبوقًا
    بـ \code{sp:} (فيصير \code{sp:text-embedding-3-small}).
  \item لو \en{Azure} غير مُهيّأ، ترجع لاسم النموذج المحلّيّ عبر \code{local.ModelName}
    — هذا مبدأ الاختيار بالإعداد الذي عرفناه سابقًا.
  \item \code{Dimensions} تُرجِع 1536، وهو طول المُتّجه الذي نلتزم به في القاعدة.
\end{itemize}

\begin{note}
البادئة \code{sp:} ليست جزءًا من اسم النموذج عند \en{Azure}؛ هي مجرّد «مساحة أسماء»
(\en{namespace}) داخل قاعدتنا. فائدتها الوحيدة: تمييز أيّ نموذجٍ ولّد كلّ مُتّجه،
حتى لا نخلط مُتّجهات نموذجين مختلفين في بحثٍ واحد.
\end{note}

% ─────────────────────────────────────────────────────────────────────────────
\subsection{الـ Cosine Similarity — قياس التقارب المعنويّ}

عرفنا أنّ النصوص المتقاربة تُنتج مُتّجهات متقاربة الاتّجاه. لكن كيف نقيس هذا التقارب
بالضبط؟ الجواب: نقيس \emph{الزاوية} بين المُتّجهين، لا المسافة بينهما ولا طولهما.
هذا هو \en{Cosine Similarity} (تشابه جيب التمام). لو كان لدينا مُتّجهان
$\mathbf{a}$ و $\mathbf{b}$، فالمعادلة:

\[
\cos(\theta)
= \frac{\mathbf{a}\cdot\mathbf{b}}{\|\mathbf{a}\|\;\|\mathbf{b}\|}
= \frac{\displaystyle\sum_{i=1}^{n} a_i\, b_i}
       {\displaystyle\sqrt{\sum_{i=1}^{n} a_i^2}\;\sqrt{\sum_{i=1}^{n} b_i^2}}
\quad\in[0,\,1]
\]

\begin{idea}[الفكرة: الزاوية الصغيرة = تشابه عالٍ]
البَسط $\mathbf{a}\cdot\mathbf{b}$ هو الضرب النقطيّ، والمقام يقسم على طولَي
المُتّجهين حتى يُلغي أثر الطول تمامًا ويُبقي الاتّجاه وحده. النتيجة رقم بين $0$
و $1$: كلّما \textbf{صغُرت الزاوية} بين المُتّجهين اقتربت النتيجة من $1$ (تشابه
عالٍ)، وكلّما كبُرت الزاوية اقتربت من $0$ (لا علاقة). ببساطة: مُتّجهان يشيران لنفس
الاتّجاه $\Rightarrow$ نصّان يتحدّثان عن نفس الموضوع.
\end{idea}

لكن ليس كلّ تقاربٍ يُقبَل. نحتاج \emph{عتبة} (\en{threshold}) نقول عندها: «هذا
المستند قريب بما يكفي لِنُرفقه للنموذج». هذه العتبة مضبوطة في الإعداد
\code{RagMinScore} وقيمتها \textbf{0.6}. الجدول التالي يوضّح كيف نُفسّر الدرجات:

\begin{center}
\begin{tabular}{L{2.8cm} L{4.2cm} L{6.6cm}}
\toprule
\textbf{\color{spPrimary}الدرجة} & \textbf{\color{spPrimary}التفسير} & \textbf{\color{spPrimary}السلوك في \en{ScholarPath}}\\
\midrule
$\geq 0.85$ & معنى شبه مطابق & يُسترجَع ويُستشهَد به في الإجابة\\
$0.70$–$0.84$ & متشابه جدًّا & يُسترجَع ويُستشهَد به\\
$0.60$–$0.69$ & مرتبط بوضوح & يُسترجَع (عند حدّ \code{RagMinScore})\\
$< 0.60$ & ارتباط ضعيف & يُستبعَد قبل إرسال البرومبت للنموذج\\
\bottomrule
\end{tabular}
\end{center}

\begin{keyp}[نقطة مهمة: العتبة تحمي الجودة]
لو أرسلنا للنموذج مستندات ضعيفة الصِّلة، سيبني إجابته على معلومات لا تخصّ السؤال —
وهذا مصدرٌ رئيسيّ للهلوسة. العتبة \code{RagMinScore = 0.6} هي صمّام الأمان: أيّ مستند
تحت 0.6 يُرمى قبل أن يرى النموذجُ سؤالَ المستخدم أصلًا.
\end{keyp}

% ─────────────────────────────────────────────────────────────────────────────
\subsection{ما هو الـ RAG؟}

\begin{idea}[الفكرة: \en{Retrieval-Augmented Generation}]
النموذج اللُّغويّ (\en{LLM}) يعرف الكثير، لكن معرفته مُجمّدة عند لحظة تدريبه، ولا
يعرف شيئًا عن بيانات \en{ScholarPath} الخاصّة (منحنا، مستشارونا، مقالاتنا). لو سألناه
مباشرةً «ما المنح المتاحة لمصر؟» فسيُخمّن من ذاكرته — وقد يخترع منحًا لا وجود لها.
الـ \en{RAG} يحلّ هذه \emph{فجوة المعرفة}: بدل التخمين، نبحث أوّلًا في قاعدة معرفتنا
الحيّة عن المستندات الأوثق صِلةً بالسؤال، ثمّ نُرفقها للنموذج \textbf{كسياق} مع
السؤال، فتُبنى الإجابة على بياناتٍ حقيقيّة ومحدَّثة.
\end{idea}

الاسم نفسه يشرح المكوّنات الثلاثة: \en{Retrieval} (استرجاع المستندات)، ثمّ
\en{Augmented} (تعزيز البرومبت بها)، ثمّ \en{Generation} (توليد الإجابة). التدفّق في
\en{ScholarPath} \textbf{خطّيّ ومرّة واحدة} — لا حلقات ولا تكرار:

\begin{center}
\begin{tabular}{L{1.4cm} L{12.4cm}}
\toprule
\textbf{\color{spPrimary}الخطوة} & \textbf{\color{spPrimary}ما يحدث}\\
\midrule
1 & يكتب المستخدم سؤالًا نصّيًّا.\\
2 & نحوّل السؤال إلى مُتّجه (\code{EmbedAsync} $\Rightarrow$ \code{float[1536]}).\\
3 & نحسب \en{cosine} بين مُتّجه السؤال وكلّ مستند، ونأخذ الأعلى تشابهًا.\\
4 & نُبقي المستندات فوق العتبة (0.6) فقط، ونصوغها في كتلة سياق.\\
5 & نُرسل: تعليمات النظام + السياق + السؤال إلى النموذج \en{GPT-4o-mini}.\\
6 & يُرجِع النموذج إجابةً مسنودةً، مع قائمة المصادر التي استُشهد بها.\\
\bottomrule
\end{tabular}
\end{center}

\begin{note}
لاحظ الفرق الجوهريّ: النموذج لا «يتذكّر» المنح من تدريبه، بل \emph{يقرؤها} من السياق
الذي حضّرناه له لحظتها. لو أضاف الأدمِن منحةً جديدة اليوم وأعاد بناء قاعدة المعرفة،
صار النموذج قادرًا على ذكرها فورًا — دون إعادة تدريبٍ ولا تعديل كود.
\end{note}

% ─────────────────────────────────────────────────────────────────────────────
\subsection{الـ Standard RAG مقابل الـ Agentic RAG}

التدفّق الذي شرحناه يُسمّى \en{Standard RAG}: خطّ أنابيب مستقيم، تمريرة واحدة، بلا
خطوة تقييم. هذا بسيط وسريع، لكن له ثلاثة أنماط فشل معروفة:

\begin{enumerate}
  \item \textbf{الأسئلة الغامضة}: «إزّاي أقدّم؟» سؤال ناقص السياق — هل يقصد التسجيل في
    المنصّة أم التقديم على منحة بعينها؟ الـ \en{Standard RAG} يسترجع دون أن يستوضح.
  \item \textbf{الأدلّة المبعثرة}: سؤال يحتاج تجميعًا من عدّة عناقيد مستندات (مثل «قارِن
    \en{Erasmus} بـ \en{Chevening} في التمويل») يحصل على إجابةٍ جزئيّة، لأنّ خطّ
    الأنابيب يستنفد تمريرة استرجاعٍ واحدة فقط.
  \item \textbf{الثقة الزائفة}: مستندات سجّلت تشابهًا عاليًا لكنّها لا تحوي الإجابة فعلًا،
    فيملأ النموذج الفراغ بتفاصيل مُقنعة لكنّها مُختلَقة.
\end{enumerate}

الـ \en{Agentic RAG} يعالج هذا باستبدال الخطّ المستقيم بـ \textbf{حلقة تحكّم}
(\en{control loop}) يقودها وكيلٌ (\en{agent}) ذاتيّ القرار: يستوضح لو كان السؤال
غامضًا، ويختار مصدر البحث المناسب، ويُقيّم هل النتائج كافية، ولو لم تكن كافية يُعيد
صياغة السؤال ويبحث مجدّدًا — قبل أن يُولّد الإجابة النهائيّة. أي أنّه يكرّر ويصحّح
نفسه، لكن بثمن. الجدول التالي يوازن بينهما:

\begin{center}
\begin{tabularx}{\linewidth}{L{3.6cm} Y Y}
\toprule
\textbf{\color{spPrimary}المعيار} & \textbf{\color{spPrimary}\en{Standard RAG}} & \textbf{\color{spPrimary}\en{Agentic RAG}}\\
\midrule
زمن الاستجابة & 1–3 ثوانٍ & 5–15+ ثانية (نداءات نموذج متعدّدة)\\
التكلفة لكلّ استعلام & $\approx \$0.0002$ & $\approx \$0.001$–$\$0.002$ ($3$–$10$ أضعاف)\\
سؤال بسيط أحاديّ المصدر & ممتاز & بلا تحسّن يُذكَر\\
سؤال معقّد متعدّد المصادر & متوسّط & أفضل بوضوح\\
تصحيح الأخطاء & حتميّ وقابل للاختبار & غير حتميّ، يصعُب تكراره\\
حجم المتون الأمثل & حتى $\sim$10 آلاف مستند & 10 آلاف+ ومجالات متعدّدة\\
\bottomrule
\end{tabularx}
\end{center}

\begin{why}
اخترنا \en{Standard RAG} للشات بوت. قاعدة معرفتنا فيها \textbf{692} مستندًا فقط،
والأسئلة موضوعيّة ومحدّدة (منح، مستشارون، نصائح تقديم). تمريرة استرجاعٍ واحدة تجد
المستند الصحيح باستمرار في \textbf{أقلّ من 10 ملّي ثانية}. إضافة حلقة \en{agentic}
كانت ستضرب الزمن في $5$ والتكلفة في $10$ دون أيّ مكسبٍ ملموس في الدقّة على مجموعة
أسئلتنا. عتبة الانتقال إلى \en{Agentic RAG} تقارب 10 آلاف مستندٍ بمجالاتٍ متغايرة —
ونحن بعيدون عنها.
\end{why}

\begin{qa}
\textbf{س: أليس \en{Agentic RAG} «أذكى» دائمًا؟ لماذا لا نستعمله لِنبدو متطوّرين؟}\\
لأنّ الهندسة الجيّدة تختار الأبسط الذي يفي بالغرض. «الأذكى» هنا يعني أبطأ بخمسة
أضعاف، وأغلى بعشرة، وأصعب في الاختبار وتتبّع الأخطاء — مقابل صفر تحسّنٍ قابلٍ للقياس
على 692 مستندًا. القرار الهندسيّ السليم أن نضيف التعقيد يوم تفرضه البيانات، لا قبله.
\end{qa}

% ─────────────────────────────────────────────────────────────────────────────
\subsection{الصورة الكاملة — ثلاث ميزات، مسارٌ واحد فقط يستعمل GPT}

\begin{keyp}[نقطة مهمة: تمييزٌ جوهريّ]
في \en{ScholarPath} \textbf{ثلاث} ميزات ذكاء اصطناعيّ، لكن \textbf{واحدة فقط}
تستعمل \en{RAG} + \en{GPT} حقيقيًّا (وهي الشات بوت). الميزتان الأخريان — التوصيات
والأهليّة — حسابيّتان بالكامل: خوارزميّة تسجيل نقاط حتميّة تعمل محلّيًّا داخل الخادم،
بلا أيّ نداء شبكة ولا نموذجٍ لُغويّ.
\end{keyp}

الجدول التالي يلخّص التوجيه (\en{routing}):

\begin{center}
\begin{tabular}{L{3.4cm} C{3.4cm} C{2.6cm} L{4.0cm}}
\toprule
\textbf{\color{spPrimary}الميزة} & \textbf{\color{spPrimary}تستعمل \en{GPT}؟} & \textbf{\color{spPrimary}تستعمل \en{RAG}؟} & \textbf{\color{spPrimary}الأسلوب}\\
\midrule
الشات بوت & نعم — \en{GPT-4o-mini} & نعم (بحث \en{cosine}) & \en{Standard RAG}\\
التوصيات & لا & لا & تسجيل نقاط محلّيّ\\
الأهليّة & لا & لا & مطابقة ملفٍّ محلّيّة\\
\bottomrule
\end{tabular}
\end{center}

كيف يظهر هذا في الكود؟ الخدمة \code{AzureOpenAiService} هي الواجهة الرسميّة حين يكون
\code{Ai:Provider = AzureOpenAi}، لكنّها لا تُنفّذ التوصيات والأهليّة بنفسها — بل
\textbf{تفوّضها} فورًا إلى \code{LocalAiService}. فقط \code{AskAsync} (الشات) هي التي
تشغّل مسار \en{RAG} الحقيقيّ.

\textbf{هنشوف التفويض في الكود} — المقتطف من \code{AzureOpenAiService.cs}:

\begin{english}
\begin{lstlisting}[language={[Sharp]C}]
// Recommendations & eligibility are structured profile-matching tasks:
// deterministic scoring beats LLM guesswork, so delegate to the local service.
public Task<AiRecommendationResult> GenerateRecommendationsAsync(
    Guid userId, int topN, CancellationToken ct)
    => local.GenerateRecommendationsAsync(userId, topN, ct);

public Task<AiEligibilityResult> CheckEligibilityAsync(
    Guid userId, Guid scholarshipId, CancellationToken ct)
    => local.CheckEligibilityAsync(userId, scholarshipId, ct);

// Only AskAsync runs the real RAG pipeline (retrieve -> augment -> generate).
public async Task<AiChatResponse> AskAsync(...) { /* ... */ }
\end{lstlisting}
\end{english}

نقرأ المقتطف:
\begin{itemize}
  \item الدالّتان \code{GenerateRecommendationsAsync} و \code{CheckEligibilityAsync}
    ليستا إلا سطر تمرير: \code{=> local.\dots} — أي أنّهما تُعيدان النداء حرفيًّا إلى
    \code{LocalAiService} بلا أيّ نداء \en{Azure}.
  \item الدالّة \code{AskAsync} وحدها هي التي تحمل جسدًا حقيقيًّا (\code{async \dots})
    فيه خطوات \en{RAG}. هذا هو الفرق العمليّ بين «ميزة تستعمل \en{GPT}» وأخرى لا تستعمله.
\end{itemize}

والآن ننظر داخل \code{AskAsync} نفسها لِنرى خطوتَي «الاسترجاع» و«الترشيح» — وهنا
يظهر الإعدادان اللذان تحدّثنا عنهما، \code{RagTopK} و \code{RagMinScore}:

\begin{english}
\begin{lstlisting}[language={[Sharp]C}]
// Retrieve (the "R" in RAG): the top-K most similar knowledge docs.
docs = await retriever.RetrieveAsync(msg, ai.RagTopK, ct);   // RagTopK = 5

// Keep only docs above the similarity threshold before building the prompt.
var relevant = docs.Where(d => d.Score >= ai.RagMinScore).ToList(); // 0.6
var context  = RagSupport.BuildContextBlock(relevant, arabic);
\end{lstlisting}
\end{english}

نقرأ:
\begin{itemize}
  \item \code{RetrieveAsync(msg, ai.RagTopK, ct)} تُرجِع أعلى \code{RagTopK} مستندًا
    تشابهًا مع السؤال — والقيمة \textbf{5}. أي نأخذ أفضل خمسة مرشّحين لا أكثر.
  \item \code{Where(d => d.Score >= ai.RagMinScore)} تُرشّح المرشّحين بحدّ \textbf{0.6}؛
    فلو لم يبلغ أيٌّ منها العتبة، تخرج كتلة سياقٍ فارغة، ويُجيب النموذج بحذرٍ بدل أن
    يخترع.
  \item \code{BuildContextBlock} تصوغ الناجين في نصٍّ مُرقّم يُلصَق داخل البرومبت.
\end{itemize}

\begin{secnote}
تفويض التوصيات والأهليّة محلّيًّا ليس اختصارًا فقط، بل قرار خصوصيّة وأمان: بيانات
الطالب الحسّاسة (مستواه الأكاديميّ، بلده، تخصّصه) تُطابَق داخل خادمنا ولا تغادره إلى
أيّ خدمة خارجيّة. لا يخرج إلى \en{Azure} إلا نصّ سؤال الشات وسياقٌ عامّ من قاعدة
المعرفة — لا ملفّات المستخدمين.
\end{secnote}

\begin{summary}[الخلاصة في سطور]
\begin{itemize}
  \item الـ \en{Embedding} يحوّل النصّ إلى مُتّجه ثابت الطول (1536 رقمًا عبر
    \code{text-embedding-3-small})؛ النصوص المتقاربة معنى تُنتج مُتّجهات متقاربة
    اتّجاهًا، حتى عبر اللغات.
  \item بادئة \code{sp:} تحجز «هويّة النموذج» في القاعدة حتى لا نخلط مُتّجهات نموذجين.
  \item الـ \en{Cosine Similarity} يقيس الزاوية بين مُتّجهين (رقم بين 0 و1)؛ الزاوية
    الصغيرة تعني تشابهًا عاليًا، والعتبة \code{RagMinScore = 0.6} تُبعد الضعيف قبل
    النموذج.
  \item الـ \en{RAG} يسدّ فجوة معرفة النموذج: نسترجع سياقًا حيًّا ونُرفقه بالسؤال بدل
    التخمين من الذاكرة — تدفّق خطّيّ واحد.
  \item اخترنا \en{Standard RAG} لا \en{Agentic}: 692 مستندًا واسترجاع أقلّ من 10
    ملّي ثانية لا يبرّران مضاعفة الزمن $5\times$ والتكلفة $10\times$.
  \item ثلاث ميزات \en{AI}، لكن الشات بوت وحده يستعمل \en{RAG}+\en{GPT}؛ التوصيات
    والأهليّة حسابيّتان محلّيًّا عبر \code{LocalAiService}.
\end{itemize}
\end{summary}

\clearpage
% ch06-ai-build
\section{الذكاء الاصطناعيّ (2): البناء والتشغيل}

في الفصل السابق فهمنا \emph{المفاهيم}: كيف يحوّل \en{Embedding} أي نصّ إلى متّجه
أرقام، وكيف يقيس \en{Cosine} قربَ متّجهين، وما معنى \en{RAG}، وكيف يوجّه النظام
الأسئلةَ بين ميزاته الثلاث (الشات، التوصيات، الأهليّة). في هذا الفصل ننزل إلى
\emph{المطبخ}: كيف بُنِي هذا كلّه فعليًّا في الكود، صفًّا صفًّا. هنمشي واحدة واحدة —
من بناء قاعدة المعرفة، إلى تخزين المتّجهات، إلى خطّ الشات الكامل، ثمّ التوصيات
والأهليّة الحسابيّتين، وأخيرًا شبكة الأمان والتكلفة.

\begin{idea}[الفكرة: ثلاث ميزات، محرّكان]
عندنا ثلاث ميزات ذكاء اصطناعيّ، لكن محرّكين فقط. \textbf{الشات} وحده يستخدم
\en{RAG} حقيقيًّا مع \en{GPT-4o-mini}. أمّا \textbf{التوصيات} و\textbf{الأهليّة}
فتعملان بخوارزميّة حسابيّة حتميّة (\en{deterministic}) بلا أيّ نداء شبكة ولا
\en{GPT} إطلاقًا. تذكّر هذا التقسيم؛ سيفسّر لك كثيرًا ممّا يأتي.
\end{idea}

% ═══════════════════════════════════════════════════════════════════════════
\subsection{قاعدة المعرفة: من خمسة مصادر إلى صفوف قابلة للبحث}

قبل أن يجيب الشات عن أيّ سؤال، لازم يكون عنده «كتاب» يقرأ منه. هذا الكتاب اسمه
\emph{قاعدة المعرفة} (\en{Knowledge Base})، ويبنيه صفٌّ واحد اسمه
\code{KnowledgeBaseIndexer}. مهمّته: يجمع البيانات من خمسة مصادر مختلفة، ويحوّل كلّ
عنصر إلى صفٍّ من نوع \code{KnowledgeDocument} في قاعدة \en{SQL Server}.

\begin{center}
\begin{tabular}{L{4.2cm} L{8.4cm}}
\toprule
\textbf{\color{spPrimary}المصدر} & \textbf{\color{spPrimary}ماذا يُفهرَس ومن أين} \\
\midrule
المنح المفتوحة & كلّ منحة حالتها \code{Open} — استعلام حيّ من القاعدة \\
الأسئلة الشائعة (\en{FAQ}) & مجموعة بيانات ثابتة من \code{KnowledgeBodies.cs} \\
الاستشاريّون الموثّقون & كلّ استشاريّ مُوثَّق وغير موقوف — استعلام حيّ \\
الموارد المنشورة & كلّ مورد حالته \code{Published} وغير محذوف \\
أفضل منشورات المجتمع & المواضيع التي صافي تقييمها 3 فأكثر (حتّى 200 موضوع) \\
\bottomrule
\end{tabular}
\end{center}

\begin{note}[ملاحظة: منشورات المجتمع مُطفأة افتراضيًّا]
المصدر الخامس (منشورات المجتمع) \textbf{لا} يدخل القاعدة إلّا لو فُعِّل المفتاح
\code{Ai:IncludeCommunityPostsInKb} صراحةً. السبب: هذه محتوى مستخدمين غير مُدقَّق،
ولا نريد للمساعد أن يُقدّم كلامَ أشخاص عشوائيّين وكأنّه معلومة رسميّة موثّقة، قبل أن
توجد آليّة مراجعة. لهذا عدد المستندات الفعليّ في الإنتاج يدور حول 692 مستندًا.
\end{note}

\subsubsection{الخوارزميّة: تحديث ذكيّ لا إعادة بناء كاملة}

الدالة \code{RebuildAsync} لا تمسح كلّ شيء وتبدأ من الصفر. بدلًا من ذلك تمشي في
أربع خطوات ذكيّة:

\begin{enumerate}
  \item \textbf{الإدراج والتحديث (\en{upsert})}: لكلّ مستند مطلوب، تحسب بصمة
    \en{SHA-256} على محتواه (\en{ContentEn + ContentAr}). لو المستند جديد تُدرِجه؛
    لو موجود تُحدّثه.
  \item \textbf{إزالة الأيتام (\en{orphans})}: أيّ مستند في القاعدة اختفى مصدره
    الأصليّ (مثلًا منحة أُغلقت) يُحذَف.
  \item \textbf{إعادة تضمين انتقائيّة}: المستند يُعاد تضمينه \emph{فقط} لو تغيّرت
    بصمته، أو كان تضمينه فارغًا، أو أُنشئ بنموذج مختلف. على قاعدة «دافئة» هذا صفرُ عمل.
  \item \textbf{التضمين على دُفعات}: المستندات المُعلَّقة تُرسَل إلى \en{Azure} في
    دُفعات من 16 مستندًا في المرّة.
\end{enumerate}

\begin{keyp}[نقطة مهمة: بصمة المحتوى هي مفتاح الرخص]
البصمة \en{SHA-256} هي حجر الزاوية. لو غيّرت وصف منحة واحدة، بصمتها فقط هي التي
تتغيّر، فيُعاد تضمين مستند واحد بتكلفة تكاد لا تُذكَر. أمّا الـ 691 مستندًا الباقون
فبصماتهم كما هي، فيُتخطَّون تمامًا. لهذا تكون عمليّة إعادة البناء الليليّة شبه مجّانيّة.
\end{keyp}

الآن نشوف الكود الحقيقيّ. أوّل مقتطف: خطوة الـ \en{upsert} مع فحص البصمة. لاحظ كيف
أنّ تغيّر البصمة \emph{يُفرِّغ} التضمين القديم حتّى يُعاد حسابه لاحقًا:

\begin{english}
\begin{lstlisting}[language={[Sharp]C}]
var hash = Sha256(d.EmbedText);              // SHA-256(ContentEn + ContentAr)
if (byKey.TryGetValue((d.SourceType, d.SourceKey), out var doc))
{
    doc.ContentEn = d.ContentEn;
    doc.ContentAr = d.ContentAr;
    // ...
    if (doc.ContentHash != hash)             // content actually changed?
    {
        doc.ContentHash = hash;
        doc.Embedding = [];                  // drop the stale vector
        doc.EmbeddingDimensions = 0;
        doc.EmbeddingModel = null;           // -> flagged for re-embedding
    }
}
\end{lstlisting}
\end{english}

شرح سطرًا بسطر: السطر الأوّل يحسب البصمة الجديدة. لو المستند موجود مسبقًا نُحدّث
حقوله. الشرط \code{doc.ContentHash != hash} يعني «هل تغيّر المحتوى حقًّا؟» — لو نعم،
نُفرِّغ حقل \code{Embedding} ونُصفّر النموذج، فيصير المستند «مُعلَّقًا» بانتظار إعادة
التضمين. لو المحتوى لم يتغيّر، لا نلمس التضمين، فلا تكلفة.

المقتطف الثاني: حلقة التضمين على دُفعات من 16. هنا يحدث الاتّصال الوحيد المدفوع مع
\en{Azure}، ثمّ تُحزَم النتيجة عبر \code{VectorMath.Pack}:

\begin{english}
\begin{lstlisting}[language={[Sharp]C}]
private const int EmbedBatchSize = 16;
// ...
for (var i = 0; i < pending.Count; i += EmbedBatchSize)
{
    var batch = pending.Skip(i).Take(EmbedBatchSize).ToList();
    var texts = batch.Select(EmbedText).ToList();

    var vectors = await embeddings.EmbedBatchAsync(texts, ct);   // Azure REST

    for (var j = 0; j < batch.Count && j < vectors.Count; j++)
    {
        batch[j].Embedding = VectorMath.Pack(vectors[j]);  // float[1536] -> byte[6144]
        batch[j].EmbeddingModel = model;
    }
}
\end{lstlisting}
\end{english}

الحلقة تتقدّم 16 عنصرًا في كلّ دورة (\code{i += EmbedBatchSize}). نأخذ دفعةً،
نُجهّز نصوصها، ونرسلها مرّةً واحدة إلى \code{EmbedBatchAsync}. الردّ متّجهات
\en{float}، ونحزم كلّ متّجه في عمود \en{VARBINARY} ونسجّل اسم النموذج الذي أنتجه.
إرسال 16 معًا أوفر بكثير من 16 نداءً منفصلًا.

\begin{secnote}[زاوية الأمان: بوّابة تمنع الصرف المزدوج]
عمليّة إعادة البناء تقرأ كلّ المستندات وتُضمِّن المُعلَّق منها مقابل مزوّد
\emph{مدفوع}. لو تشغّلت عمليّتان معًا (مثلًا تهيئة بدء التشغيل + طلب أدمن)، لضُمِّنت
المستنداتُ مرّتين — أي صرفٌ مضاعف. لهذا يحرس الكودُ العمليّةَ بقُفل ثابت على مستوى
العمليّة (\code{SemaphoreSlim RebuildGate}) يُسلسِل الطلبات فلا تتداخل أبدًا.
\end{secnote}

% ═══════════════════════════════════════════════════════════════════════════
\subsection{تخزين المتّجهات في \en{SQL Server} لا في قاعدة متّجهات}

قرارٌ معماريّ مهمّ: لم نستخدم قاعدة متّجهات متخصّصة (مثل \en{Pinecone} أو
\en{Qdrant})، بل خزّنّا المتّجهات داخل \en{SQL Server} نفسها. لكن \en{SQL} لا يعرف
نوع «متّجه»؛ يعرف \en{VARBINARY} فقط. فكيف نحوّل مصفوفة \en{float} إلى بايتات
والعكس؟ هذه مهمّة الصفّ \code{VectorMath}.

\begin{idea}[الفكرة: حسبة البايتات]
كلّ متّجه فيه 1536 رقمًا من نوع \en{float}. كلّ \en{float} يحتلّ 4 بايتات. إذن كلّ
مستند يحتاج $1536 \times 4 = 6144$ بايتًا فقط في عمود \en{VARBINARY}. رقمٌ صغير جدًّا،
ولهذا لا نحتاج بنيةً خاصّة لتخزينه.
\end{idea}

المقتطف التالي يُظهر التحويل في الاتّجاهين. لاحظ \code{Buffer.BlockCopy} — نسخٌ خامّ
للذاكرة بلا حلقات، سريعٌ للغاية:

\begin{english}
\begin{lstlisting}[language={[Sharp]C}]
public static byte[] Pack(float[] vector)
{
    if (vector.Length == 0) return [];
    var bytes = new byte[vector.Length * sizeof(float)];   // 1536 * 4 = 6144
    Buffer.BlockCopy(vector, 0, bytes, 0, bytes.Length);
    return bytes;
}

public static float[] Unpack(byte[] bytes)
{
    if (bytes.Length == 0) return [];
    var vector = new float[bytes.Length / sizeof(float)];
    Buffer.BlockCopy(bytes, 0, vector, 0, vector.Length * sizeof(float));
    return vector;
}
\end{lstlisting}
\end{english}

الدالة \code{Pack} تحجز مصفوفة بايتات بطول $\text{count} \times 4$، ثمّ تنسخ
الذاكرة نسخًا مباشرًا. \code{Unpack} تفعل العكس تمامًا: تقسم عدد البايتات على 4 لتعرف
عدد الأرقام، ثمّ تنسخها إلى مصفوفة \en{float} جاهزة لحساب \en{Cosine}.

\begin{why}
عند 692 مستندًا، مسحُ القاعدة كلّها في الذاكرة وحسابُ \en{Cosine} لكلّ مستند يكتمل في
أقلّ من 10 مِلّي ثانية. أيّ قاعدة متّجهات خارجيّة كانت ستضيف خدمةً جديدة تُنشَر
وتُدار، وقفزةَ شبكة إضافيّة، وتكلفةً تشغيليّة — كلّ ذلك لِيصير أبطأ لا أسرع عند هذا
الحجم. عتبةُ الانتقال إلى قاعدة متّجهات حقيقيّة تبدأ عند عشرات الآلاف من المستندات.
\end{why}

% ═══════════════════════════════════════════════════════════════════════════
\subsection{خطّ الشات الكامل: تسع خطوات}

الآن الجزء الأهمّ: ماذا يحدث بالضبط عندما يكتب الطالب سؤالًا في الشات؟ الرحلة تمرّ
بتسع خطوات مرتّبة. الجدول التالي يلخّصها، ثمّ نشرح المحطّات الحسّاسة:

\begin{center}
\begin{tabular}{C{0.7cm} L{4.4cm} L{7.2cm}}
\toprule
\textbf{\color{spPrimary}\#} & \textbf{\color{spPrimary}الخطوة} & \textbf{\color{spPrimary}ماذا يحدث} \\
\midrule
1 & استقبال الرسالة & الطالب يرسل سؤاله النصّيّ \\
2 & كشف اللغة & \code{RagSupport.IsArabic()} يقرّر: عربيّ أم إنجليزيّ \\
3 & فحص الإعداد & هل \en{Azure} مُهيّأ؟ لو لا $\rightarrow$ سقوط محلّيّ فورًا \\
4 & تضمين السؤال & \code{EmbedAsync()} يحوّل السؤال إلى \en{float[1536]} \\
5 & الاسترجاع & حساب \en{Cosine} على 692 مستندًا، وأخذ أعلى 5 \\
6 & التصفية & إبقاء المستندات التي درجتها 0.6 فأكثر فقط \\
7 & بناء السياق & \code{BuildContextBlock()} ينسّق المستندات المختارة \\
8 & نداء \en{GPT} & إرسال: التعليمات + التاريخ + السياق + السؤال \\
9 & بناء الاستجابة & نصّ الجواب + قائمة المصادر + التكلفة المقدّرة \\
\bottomrule
\end{tabular}
\end{center}

الخطوات 4 و5 و6 هي قلب \en{RAG}. المقتطف التالي — من صفّ الاسترجاع
\code{KnowledgeRetriever} — يُظهر كيف نحمّل المستندات المُضمَّنة \emph{بالنموذج
النشِط فقط}، ثمّ نرتّبها بالدرجة ونأخذ أعلى \en{topK}:

\begin{english}
\begin{lstlisting}[language={[Sharp]C}]
var docs = await db.KnowledgeDocuments.AsNoTracking()
    .Where(d => d.EmbeddingModel == model && d.EmbeddingDimensions > 0)
    .Select(d => new { d.Id, d.TitleEn, d.ContentEn, d.Embedding /* ... */ })
    .ToListAsync(ct);

var queryVector = await embeddings.EmbedAsync(query, ct);

return [.. docs
    .Select(d => new {
        Doc = d,
        Score = VectorMath.CosineSimilarity(queryVector, VectorMath.Unpack(d.Embedding)),
    })
    .OrderByDescending(x => x.Score)
    .Take(topK)];                            // top-5 by default
\end{lstlisting}
\end{english}

الشرط \code{EmbeddingModel == model} مهمّ: متّجهاتُ نموذجٍ مختلف تعيش في «فضاء» آخر
ولا يصحّ مقارنتها، فنتجاهلها. نحسب \en{Cosine} بين متّجه السؤال وكلّ مستند بعد فكّ
حزمته بـ \code{Unpack}، نرتّب تنازليًّا، ونأخذ الأعلى. القيمة الافتراضيّة لـ
\code{Ai:RagTopK} هي 5، ولـ \code{Ai:RagMinScore} هي 0.6.

\subsubsection{تصميم تعليمات النظام (\en{System Prompt})}

قبل نداء \en{GPT}، نُرسِل «تعليمات نظام» تُشكّل شخصيّة المساعد (بنسختين: عربيّة
وإنجليزيّة). أهمّها أن يُطلَب منه:

\begin{itemize}
  \item أن يذكر المنحَ والاستشاريّين والموارد \textbf{بالاسم} حين تَرِد في كتلة السياق.
  \item أن يساعد في صياغة خطاب الغرض وخطابات التوصية والسيرة الذاتيّة عند الطلب.
  \item ألّا \textbf{يخترع} أبدًا اسمَ منحةٍ أو موعدًا أو مبلغًا غير موجود في السياق.
  \item ألّا \textbf{يرفض} طلبًا معقولًا من المستخدم.
\end{itemize}

\subsubsection{ذاكرة الجلسة (\en{Session Memory})}

كيف «يتذكّر» الشات ما قيل قبل قليل بلا مخزنٍ على الخادم؟ ببساطة: قبل كلّ نداء نُضيف
أدوار المحادثة السابقة (آخر ~20 دورًا) بين تعليمات النظام والسؤال الحاليّ:

\begin{english}
\begin{lstlisting}[language={[Sharp]C}]
var messages = new List<object>
{
    new { role = "system", content = system },
};
foreach (var turn in history)                     // prior turns this session
    messages.Add(new { role = turn.Role, content = turn.Content });
messages.Add(new { role = "user", content = user });   // current question
\end{lstlisting}
\end{english}

القائمة تبدأ بتعليمات النظام، ثمّ نُضيف كلّ دورٍ سابق بدوره (\code{user} أو
\code{assistant})، ثمّ نُلحِق السؤالَ الحاليّ. النموذج يقرأ السياق كاملًا فيبدو
«متذكّرًا». نداء \en{GPT-4o-mini} يمرّ بـ \code{temperature=0.3} — قيمة منخفضة تجعل
الإجابات ثابتةً ومنضبطة لا خياليّة.

% ═══════════════════════════════════════════════════════════════════════════
\subsection{التوصيات: درجاتٌ حسابيّة بلا \en{GPT}}

هنا ننتقل إلى المحرّك الثاني. خلافًا لما قد يظنّه كثيرون، توليدُ التوصيات
\textbf{ليس} فيه أيّ ذكاء اصطناعيّ توليديّ. إنّه صفٌّ حسابيّ اسمه
\code{LocalAiService} يُعطي كلّ منحةٍ مفتوحة \emph{درجةً} من 100 بناءً على مطابقتها
لملفّ الطالب.

\begin{secnote}[زاوية الأمان: لا شبكة، لا \en{GPT}، لا مفاجآت]
لا يوجد نداء \en{GPT}، ولا \en{RAG}، ولا أيّ طلب \en{HTTP} في خطّ التوصيات إطلاقًا.
كلّه يعمل داخل الخادم بشكل حتميّ. النتيجة نفسها للمدخلات نفسها دائمًا — قابلة للاختبار
ولا تكلّف شيئًا ولا تتعطّل بسبب مزوّد خارجيّ.
\end{secnote}

معادلة الدرجة تجمع أربعة مكوّنات بأوزان مختلفة:

\[
\text{Score} \;=\; \underbrace{S_{\text{level}}}_{\le 30}
\;+\; \underbrace{S_{\text{field}}}_{\le 40}
\;+\; \underbrace{S_{\text{country}}}_{\le 25}
\;+\; \underbrace{B_{\text{funding}}}_{\le 5}
\;\in\; [0,\,100]
\]

مجال الدراسة هو الوزن الأكبر (حتّى 40 نقطة) لأنّ توافق التخصّص أهمّ من الجغرافيا.
المقتطف الحقيقيّ من دالة التوليد:

\begin{english}
\begin{lstlisting}[language={[Sharp]C}]
var score = 0;
if (userLevel.HasValue && userLevel.Value == c.TargetLevel) score += 30;
score += OverlapScore(preferredFields, fosFields, maxPoints: 40);
// countries compared on a canonical ISO key ("Egypt"/"masr" -> "EG")
score += OverlapScore(NormalizeCountries(preferredCountries),
                      NormalizeCountries(countries), maxPoints: 25);
if (c.FundingAmountUsd >= 10_000) score += 5;      // funding bonus
scored.Add((c.Id, Math.Min(score, 100), c.TitleEn, c.TitleAr));
\end{lstlisting}
\end{english}

المستوى الأكاديميّ: تطابقٌ تامّ يعطي 30. المجال والدولة: عبر دالّة
\code{OverlapScore} التي تحسب نسبة التقاطع مضروبةً في السقف. مكافأةُ 5 نقاط لو كان
التمويل 10000 دولار فأكثر. أخيرًا \code{Math.Min(score, 100)} تضمن ألّا تتجاوز
الدرجةُ 100. ثمّ تُترجَم الدرجة إلى تسمية مفهومة للطالب:

\begin{center}
\begin{tabular}{C{2.6cm} L{3.0cm} L{6.6cm}}
\toprule
\textbf{\color{spPrimary}الدرجة} & \textbf{\color{spPrimary}التسمية} & \textbf{\color{spPrimary}ما يظهر للطالب} \\
\midrule
80--100 & توافق قويّ  & قريبة من اهتماماتك ومستواك الأكاديميّ \\
50--79  & توافق جزئيّ & تتقاطع مع بعض تفضيلاتك \\
20--49  & توافق محدود & تلمس مجالاتك، تستحقّ نظرةً سريعة \\
0--19   & توافق ضعيف  & لا تتماشى حاليًّا مع ملفّك \\
\bottomrule
\end{tabular}
\end{center}

\subsubsection{التخزين المؤقّت وإعادة الترطيب (\en{Re-hydration})}

حسابُ الدرجات يجري مرّةً وتُخزَّن النتيجة \emph{مؤقّتًا} كـ \en{JSON} داخل الحقل
\code{ResponseText} من كيان \code{AiInteraction}. لكن هنا حيلةٌ ذكيّة: الـ \en{JSON}
المخزَّن يحفظ فقط الحقولَ التي أنتجها الذكاء (المعرّف، الدرجة، الشرح). أمّا
\emph{الموعد النهائيّ} و\emph{مبلغ التمويل} فيُجلَبان \textbf{حيًّا} من القاعدة عند
كلّ قراءة.

\begin{keyp}[نقطة مهمة: لماذا نُعيد الترطيب؟]
لو خزّنّا الموعد النهائيّ داخل الكاش، ثمّ غيّره الأدمن، لظهر للطالب موعدٌ قديمٌ خاطئ.
بإعادة الترطيب من القاعدة حيًّا، الدرجة المحفوظة تبقى (فهي لا تتغيّر كثيرًا)، بينما
البياناتُ الحسّاسة للوقت تكون دائمًا حديثة. الكاش صالح 24 ساعة؛ بعدها تُولَّد توصياتٌ
جديدة.
\end{keyp}

% ═══════════════════════════════════════════════════════════════════════════
\subsection{الأهليّة: مطابقة الملفّ الشخصيّ}

المحرّك الحسابيّ نفسه يخدم الميزة الثالثة: فحص الأهليّة. هنا نقارن ملفّ الطالب
بمتطلّبات منحةٍ واحدة عبر ثلاثة معايير:

\begin{center}
\begin{tabular}{L{3.2cm} L{4.4cm} L{4.6cm}}
\toprule
\textbf{\color{spPrimary}المعيار} & \textbf{\color{spPrimary}مصدر الطالب} & \textbf{\color{spPrimary}منطق المطابقة} \\
\midrule
المستوى الأكاديميّ & \code{AcademicLevel} & تطابق تامّ \\
الدولة & الدول المفضّلة + الجنسيّة & عبر \code{CountryNormalizer} \\
مجال الدراسة & \code{FieldOfStudy} & احتواء غير حسّاس للحالة \\
\bottomrule
\end{tabular}
\end{center}

كلّ معيار يأخذ إحدى القيم: \code{yes} أو \code{partial} أو \code{no} أو
\code{unknown} (حين يكون حقلُ الطالب فارغًا). ثمّ تُجمَع هذه القيم في حكمٍ واحد عبر
دالّة \code{DeriveVerdict}:

\begin{english}
\begin{lstlisting}[language={[Sharp]C}]
internal static EligibilityVerdict DeriveVerdict(
    IReadOnlyList<AiEligibilityCriterion> criteria)
{
    if (criteria.Count == 0) return EligibilityVerdict.NotEligible;

    if (criteria.Any(c => c.Match == "no"))            // any hard fail
        return EligibilityVerdict.NotEligible;

    if (criteria.Any(c => c.Match == "partial" || c.Match == "unknown"))
        return EligibilityVerdict.PartiallyEligible;   // any uncertainty

    return EligibilityVerdict.Eligible;                // all clear
}
\end{lstlisting}
\end{english}

المنطق بسيطٌ ومتحفّظ: أيُّ معيار قيمته \code{no} كافٍ وحده ليكون الحكم
\code{NotEligible}. لو لا يوجد رفضٌ صريح لكن يوجد \code{partial} أو \code{unknown}،
فالحكم \code{PartiallyEligible}. ولا نقول \code{Eligible} إلّا لو كانت كلّ المعايير
\code{yes}.

\begin{note}[ملاحظة: تطبيع أسماء الدول ليس بديهيًّا]
منحةٌ تكتب دولتها «\en{France}» يجب أن تطابق طالبًا كتب «\en{FR}» أو «\en{france}»
أو «فرنسا». الدالّة \code{CountryNormalizer.Matches} تُرجِع كلّ الصيغ إلى مفتاحٍ
معياريّ واحد قبل المقارنة. بدون هذا، لقرأ طالبٌ مصريّ يفضّل ألمانيا كلمة «\en{no}»
على منحةٍ كتبت الدولة «\en{DE}».
\end{note}

% ═══════════════════════════════════════════════════════════════════════════
\subsection{شبكة الأمان: السقوط الآمن (\en{Fallback}) ثلاثيّ المستوى}

ماذا لو تعطّل \en{Azure} تمامًا؟ لا يجوز أن يرى المستخدمُ خطأً. لهذا كلّ نداء شات
مُحاطٌ بسلسلة سقوطٍ من ثلاثة مستويات، أيُّ فشلٍ فيها يوجّه الطلبَ إلى
\code{LocalAiService.AskAsync}:

\begin{center}
\begin{tabular}{C{0.7cm} L{5.4cm} L{6.2cm}}
\toprule
\textbf{\color{spPrimary}\#} & \textbf{\color{spPrimary}نقطة الفحص} & \textbf{\color{spPrimary}عند الفشل} \\
\midrule
1 & هل النقطة والمفتاح مُهيّآن؟ & سقوطٌ محلّيّ فوريّ \\
2 & هل نجح نداء الاسترجاع؟ & سقوطٌ محلّيّ (خطأ \en{HTTP}) \\
3 & هل نجح نداء \en{GPT}؟ & سقوطٌ محلّيّ (خطأ أو مهلة) \\
\bottomrule
\end{tabular}
\end{center}

والجميل أنّ المستوى المحلّيّ ليس ردًّا أخرقَ ثابتًا. \code{AskAsync} المحلّيّة تُجرِّب
\emph{مسارين}: أوّلًا نفس استرجاع \en{RAG} (بنفس \code{KnowledgeRetriever}) لكن بجواب
\emph{استخلاصيّ} — تعرض أقربَ مستند مفهرَس بدل توليد نصّ. ولو لم تجد شيئًا فوق العتبة
0.6، تسقط إلى موجّهِ كلماتٍ مفتاحيّة (\en{keyword router}) يفهم نيّة السؤال (موعد،
أهليّة، توصية...) ويردّ ردًّا مفيدًا. هكذا يبقى \en{RAG} قابلًا للعرض حتّى بلا أيّ
مفتاح \en{API}.

% ═══════════════════════════════════════════════════════════════════════════
\subsection{التكلفة والضبط الدقيق وكيان \code{AiInteraction}}

\subsubsection{نموذج التكلفة وسقف الإنفاق}

نموذج \en{GPT-4o-mini} رخيصٌ جدًّا. التكلفة تُحسَب من عدد الرموز (\en{tokens}) داخلًا
وخارجًا:

\[
\text{Cost} \;=\; N_{\text{prompt}} \times \frac{0.15}{10^6}
\;+\; N_{\text{completion}} \times \frac{0.60}{10^6}
\]

(أي 0.15 دولار لكلّ مليون رمز دخْل، و0.60 دولار لكلّ مليون رمز خرْج). الجدول يوضّح
حجم الأرقام العمليّ:

\begin{center}
\begin{tabular}{L{5.0cm} C{2.2cm} C{2.2cm} C{2.6cm}}
\toprule
\textbf{\color{spPrimary}السيناريو} & \textbf{الدخل} & \textbf{الخرج} & \textbf{التكلفة} \\
\midrule
دورة شات بسيطة & 300 & 150 & \$0.000135 \\
سياق غنيّ (5 مستندات) & 600 & 300 & \$0.000270 \\
1000 دورة يوميًّا & --- & --- & \$0.27 تقريبًا \\
سقف \code{DailyCostLimitUsd} & \multicolumn{3}{L{7.0cm}}{\$5.00 --- حدّ يوميّ صارم في \code{AiOptions}} \\
\bottomrule
\end{tabular}
\end{center}

\begin{why}
حتّى لو خرجت الأمور عن السيطرة، المفتاح \code{Ai:DailyCostLimitUsd} يضع سقفًا صلبًا
عند خمسة دولارات في اليوم. عند بلوغه يتوقّف النظام عن نداء \en{Azure} ويسقط تلقائيًّا
إلى المحرّك المحلّيّ المجّانيّ. هكذا لا يمكن لخطأٍ أو إساءة استخدام أن يُنتِج فاتورةً
مفاجئة.
\end{why}

\subsubsection{مسار الضبط الدقيق (\en{Fine-Tuning}) داخل لوحة الأدمن}

جُهِّزت مجموعةُ بياناتٍ للضبط الدقيق (نحو 332 مثالًا بصيغة محادثات \en{OpenAI}).
والمسار كلّه يُدار من لوحة الأدمن في أربع خطوات: \emph{تصدير} الـ \en{JSONL}، ثمّ
\emph{رفعه}، ثمّ \emph{تدريب} المهمّة على \en{Azure}، ثمّ \emph{تفعيل} النموذج
الناتج. وأجمل ما فيه أنّ تفعيله لا يحتاج تغييرَ كود؛ يكفي أن يختار الكودُ اسمَ
النشرِ بأولويّةٍ ثلاثيّة:

\begin{english}
\begin{lstlisting}[language={[Sharp]C}]
// Priority: DB platform setting -> appsettings override -> base model
var deployment = !string.IsNullOrWhiteSpace(dbDeployment)
    ? dbDeployment                       // e.g. "scholarpath-ft-001"
    : !string.IsNullOrWhiteSpace(az.FineTunedDeploymentName)
        ? az.FineTunedDeploymentName
        : az.DeploymentName;             // "gpt-4o-mini"  (base)
\end{lstlisting}
\end{english}

أوّلًا يُقرأ اسمُ النشر من إعدادات المنصّة في القاعدة؛ فإن كان فارغًا يُجرَّب إعداد
الملفّ؛ وإلّا يُستخدَم النموذجُ الأساس \en{gpt-4o-mini}. الأدمن يستطيع «الرجوع للنموذج
الأساس» بضغطة واحدة بلا إعادة تشغيل.

\subsubsection{كيان \code{AiInteraction}: سجلٌّ وكاشٌ في جدولٍ واحد}

كلّ عمليّات الذكاء الاصطناعيّ (شات، توصية، أهليّة) تُسجَّل في جدولٍ واحد هو
\code{AiInteraction}. إنّه في آنٍ معًا \textbf{سجلّ تدقيق} و\textbf{طبقة تخزين مؤقّت}:

\begin{english}
\begin{lstlisting}[language={[Sharp]C}]
public class AiInteraction : AuditableEntity
{
    public Guid UserId { get; set; }
    public AiFeature Feature { get; set; }   // Recommendation | Chat | Eligibility
    public AiProvider Provider { get; set; } // AzureOpenAi | Local
    public string? SessionId { get; set; }   // groups chat turns into a session

    public string PromptText { get; set; }   // what was asked
    public string ResponseText { get; set; } // answer text OR cached JSON
    public string? MetadataJson { get; set; }// eligibility breakdown, etc.

    public int PromptTokens { get; set; }
    public int CompletionTokens { get; set; }
    public decimal CostUsd { get; set; }     // billing
}
\end{lstlisting}
\end{english}

المفتاح لفهم هذا الكيان أنّ معنى الحقل \code{ResponseText} \emph{يتغيّر} حسب
\code{Feature}: في الشات هو نصّ الجواب؛ في التوصيات هو الـ \en{JSON} المخزَّن مؤقّتًا؛
في الأهليّة هو الحكم، وتفاصيله في \code{MetadataJson}. جدولٌ واحد، ثلاثة استعمالات —
بسيطٌ وأنيق.

% ═══════════════════════════════════════════════════════════════════════════
\begin{summary}[الخلاصة في سطور]
\begin{itemize}
  \item \textbf{قاعدة المعرفة} تُبنى من خمسة مصادر (منح، أسئلة شائعة، استشاريّون،
    موارد، منشورات مجتمع) إلى صفوف \code{KnowledgeDocument}؛ بصمة \en{SHA-256} تجعل
    إعادة البناء انتقائيّة ورخيصة، والتضمين يجري في دُفعاتٍ من 16.
  \item \textbf{المتّجهات} تُخزَّن في \en{SQL Server} كـ \en{VARBINARY} بحجم
    $1536 \times 4 = 6144$ بايتًا لكلّ مستند؛ كافٍ تمامًا عند 692 مستندًا، والمسح
    يكتمل في أقلّ من 10 مِلّي ثانية بلا قاعدة متّجهات.
  \item \textbf{خطّ الشات} تسع خطوات: كشف لغة، فحص إعداد، تضمين، استرجاع \en{Cosine}،
    تصفية عند 0.6، بناء سياق، نداء \en{GPT-4o-mini} بـ \code{temperature=0.3}، ثمّ
    بناء الجواب والمصادر؛ مع ذاكرة جلسةٍ من آخر ~20 دورًا.
  \item \textbf{التوصيات والأهليّة} حسابيّتان بلا \en{GPT}: درجةٌ من
    (مستوى 30 + مجال 40 + دولة 25 + مكافأة 5)، وحكمُ أهليّةٍ من ثلاثة معايير مع تطبيع
    الدول؛ التوصيات تُخزَّن مؤقّتًا وتُعاد ترطيبُ موعدها وتمويلها حيًّا من القاعدة.
  \item \textbf{شبكة الأمان} ثلاثيّة المستوى تسقط دائمًا إلى \code{LocalAiService}؛
    وسقفُ \code{DailyCostLimitUsd} عند خمسة دولارات يمنع أيّ فاتورةٍ مفاجئة؛ وكلّ
    عمليّة تُسجَّل في كيان \code{AiInteraction} الذي يجمع بين السجلّ والكاش.
\end{itemize}
\end{summary}

\clearpage
% ch07-meetings-realtime
\section{الجلسات المرئيّة والزمن الحقيقيّ}

في الفصول السابقة رأينا كيف يُقسَّم النظام إلى طبقات، وكيف يمرّ كل طلب عبر
\en{CQRS}، وكيف تُنشِر العملياتُ أحداثًا دوميّة، وكيف نعزل العالَم الخارجيّ خلف
منافذ ومحوّلات (\en{ports/adapters}). في هذا الفصل نطبّق كل ذلك على ميزتين
«حيّتين» بطبيعتهما: \textbf{مكالمة الفيديو} بين الطالب والمستشار، و\textbf{الدفع
اللحظيّ} للإشعارات والمحادثة عبر \en{SignalR}. الميزتان مختلفتان في التقنية لكنهما
تتشاركان الفكرة نفسها: النظام لا «يستضيف» الوسائط بنفسه، بل ينسّق بين المستخدم
وخدمة سحابية متخصّصة، ويحتفظ هو فقط بالحقائق المهمّة (مَن حضر؟ متى؟ أين التسجيل؟).

\begin{idea}[الفكرة: النظام مُنسِّق لا خادم وسائط]
مكالمة الفيديو تجري بين متصفّحَي المشاركَين وخدمة \en{Azure Communication Services}
(نسمّيها اختصارًا \en{ACS}). خادمُنا \en{.NET} لا يمرّر بايتات الصوت أو الصورة أبدًا؛
كل ما يفعله هو: إنشاء «غرفة»، إصدار رمز دخول قصير العمر لكل طرف، وطلب بدء التسجيل.
هذا يجعل الحِمل على خادمنا خفيفًا جدًّا، ويترك أمور الترميز والشبكة لِمن يُتقِنها.
\end{idea}

\subsection{الغرفة التي بلا خادم: منفذ \en{IMeetingService}}

كل تعامل مع \en{ACS} يمرّ عبر منفذ واحد اسمه \code{IMeetingService}. في الإنتاج
يُختار المحوّل \code{AzureCommunicationMeetingService} تلقائيًّا حين يكون مفتاح
الإعداد \code{Acs:ConnectionString} مضبوطًا؛ وفي التطوير — حيث لا نملك اشتراك
\en{ACS} — يُحقَن بديل بسيط اسمه \code{StubMeetingService} يُرجِع قيمًا وهميّة حتى
تعمل الشاشات بلا اتصال حقيقيّ. هذا هو نفس نمط المنافذ الذي رأيناه من قبل: تبديل
المزوّد قرارُ \emph{إعداد} لا قرار \emph{كود}.

النقطة المفاجئة أوّل مرّة: \textbf{«الغرفة» ليست كائنًا على الخادم}. الغرفة مجرّد
مُعرِّف مكالمة جماعية في \en{ACS}، أي \en{group-call id} وهو ببساطة \en{GUID}. حين
يقبل المستشار الحجز، نولّد هذا المُعرِّف ونحفظه على صفّ الحجز، ولا شيء غير ذلك:

\begin{english}
\begin{lstlisting}[language={[Sharp]C}]
// A booking's "room" is just an ACS group-call id (a GUID).
public Task<MeetingRoom> CreateRoomAsync(Guid bookingId, CancellationToken ct)
    => Task.FromResult(new MeetingRoom(Guid.NewGuid().ToString()));
\end{lstlisting}
\end{english}

هنشوف إيه في السطر ده: الدالة \code{CreateRoomAsync} لا تتّصل بأي خدمة إطلاقًا — إنّها
فقط تُنشئ \en{GUID} جديدًا (\code{Guid.NewGuid}) وتغلّفه في كائن \code{MeetingRoom}
وتُعيده فورًا. لا يوجد نداء شبكة، لا كائن غرفة يُحجَز، لا موارد تُستهلَك مسبقًا.
الغرفة «توجد» بمجرّد أن يعرف الطرفان مُعرِّفها.

\begin{why}
لماذا غرفة بلا خادم؟ لأنّ إنشاء كائنات غرف مسبقًا لكل حجز يعني موارد محجوزة قد لا
تُستخدَم أبدًا (المستشار قد يرفض، أو ينسى الطالب). المُعرِّف الكسول (\en{lazy}) يعني
أنّنا لا ندفع شيئًا حتى لحظة الانضمام الفعليّ. وإن كان المُعرِّف مفقودًا لسبب ما
(سجلّ قديم)، تُنشئ الواجهةُ الخلفية غرفةً جديدة عند الانضمام — بلا تدخّل يدويّ.
\end{why}

\subsection{الانضمام: رمز وصول لكل مشارك ونافذة \en{$\pm$15} دقيقة}

معرفة مُعرِّف الغرفة وحدها لا تكفي للدخول. عند ضغط «انضمّ»، يتحقّق النظام أوّلًا من
\textbf{نافذة الانضمام}: يُسمَح بالدخول فقط ضمن \en{$\pm$15} دقيقة حول موعد الجلسة
المجدوَل (القيمتان \code{JoinOpensBefore} و\code{JoinClosesAfter} كلاهما 15 دقيقة).
قبل ذلك تظهر رسالة «لم تُفتَح بعد»، وبعده «أُغلِقت الغرفة». ثمّ — والأهمّ — يُصدِر
النظام \textbf{رمز وصول مؤقّتًا لكل مستخدم على حدة} من هويّة \en{ACS}:

\begin{english}
\begin{lstlisting}[language={[Sharp]C}]
public async Task<MeetingAccessToken> IssueAccessTokenAsync(
    string roomId, Guid userId, CancellationToken ct)
{
    var result = await _identityClient
        .CreateUserAndTokenAsync(new[] { CommunicationTokenScope.VoIP }, ct)
        .ConfigureAwait(false);

    return new MeetingAccessToken(
        RoomId: roomId,
        Token: result.Value.AccessToken.Token,
        AcsUserId: result.Value.User.Id,
        ExpiresAt: result.Value.AccessToken.ExpiresOn);
}
\end{lstlisting}
\end{english}

هنشوف إيه هنا سطرًا-بسطر:
\begin{itemize}
\item \code{CreateUserAndTokenAsync} تُنشئ هويّة \en{ACS} مؤقّتة \emph{وتُصدِر لها رمزًا}
  في نداء واحد. النطاق المطلوب هو \code{CommunicationTokenScope.VoIP} فقط — أي «صوت
  وفيديو»، لا صلاحيّة أوسع.
\item الرمز الناتج (\code{AccessToken.Token}) هو ما تأخذه الواجهةُ لتدخل به المكالمة،
  ومعه \code{AcsUserId} الذي يميّز المشارك داخل \en{ACS}.
\item \code{ExpiresOn} تحمل لحظة انتهاء الصلاحيّة: الرمز قصير العمر (مدّته نحو ساعتين)،
  وهي أكثر من كافية لجلسة أقصاها 60 دقيقة.
\end{itemize}

\begin{secnote}
\textbf{رمز الفيديو = بيانات اعتماد.} أي شخص يحمل الرمز يستطيع الدخول للمكالمة، ولهذا
نُصدِر رمزًا \emph{لكل مشارك} مقصورًا على \en{VoIP} وينتهي سريعًا. النتيجة: لو تسرّب
رمزٌ، فهو عديم الفائدة بعد انتهاء الجلسة، ولا يمكن «إعادة تشغيله» في مكالمة أخرى.
لاحظ أيضًا أنّنا لا نسلّم المُعرِّف والرمز إلّا للطرفين المصرَّح لهما عبر نقطة
الانضمام المحميّة — لا عبر أي رابط عامّ.
\end{secnote}

\begin{qa}
\textbf{سؤال متوقّع:} «لماذا لا تُصدِرون رمزًا واحدًا للغرفة وتشاركونه؟» لأنّ الرمز
المشترَك لا يمكن إبطاله لطرفٍ دون الآخر، ولا يميّز مَن دخل. الرمز لكل مستخدم يربط
الدخول بهويّة، ويجعل تتبّع الحضور دقيقًا، ويقلّص أثر أي تسريب إلى مستخدم واحد ولمدّة
قصيرة فقط.
\end{qa}

\subsection{التسجيل والحضور و«الغياب»}

بمجرّد أن تتّصل المكالمة فعليًّا، تطلب الواجهةُ بدءَ التسجيل. التسجيل في \en{ACS}
يتمّ عبر خدمة \en{Call Automation}، ونادينا لها بسيط:

\begin{english}
\begin{lstlisting}[language={[Sharp]C}]
public async Task<string> StartRecordingAsync(string serverCallId, CancellationToken ct)
{
    var result = await _callAutomationClient.GetCallRecording()
        .StartAsync(new StartRecordingOptions(new ServerCallLocator(serverCallId)), ct)
        .ConfigureAwait(false);
    return result.Value.RecordingId;   // handle to the running recording
}
\end{lstlisting}
\end{english}

هنشوف إيه: نمرّر \code{serverCallId} — وهو مُعرِّف المكالمة الجارية كما يراه الخادم —
إلى \code{StartAsync}، فيبدأ \en{ACS} بتسجيل الجلسة إلى تخزين \en{Blob}، ويُعيد لنا
\code{RecordingId} وهو «مقبض» للتسجيل الجاري نحتفظ به. بعد انتهاء الجلسة، نحفظ صفّ
\code{SessionRecording} يحمل مسار التخزين وحجم الملفّ، ونستطيع لاحقًا تنزيل التسجيل
بثًّا عبر \code{DownloadRecordingAsync}، فيظهر للطالب زرّ تنزيل \en{MP4}.

أمّا \textbf{الحضور}، فهو حقائق نحرص على تسجيلها على صفّ الحجز نفسه:
\begin{itemize}
\item \code{StudentJoinedAt} و\code{ConsultantJoinedAt}: لحظة أوّل دخول لكل طرف
  (متطلَّب \en{FR-217}).
\item \code{RecordingStartedAt}: لحظة بدء التسجيل (\en{PB-006}).
\item \code{NoShowMarkedAt}: تُضبَط حين تكتشف مهمّةٌ خلفيّة أنّ طرفًا لم يحضر.
\end{itemize}

\begin{keyp}[نقطة مهمة: مَن يكتشف الغياب؟]
مهمّة خلفيّة دوريّة اسمها \code{MeetingNoShowSweepJob} تمرّ على الحجوزات المنتهية،
فتقارن أوقات الدخول: إن لم يحضر طرفٌ ضِمن النافذة، تُعلِّم الحجز بـ\en{no-show}
(تضبط \code{NoShowMarkedAt}) وتُطلِق منطق الاسترداد المناسب. أي أنّ «الغياب» ليس
حُكمًا يدويًّا بل نتيجة قياس آليّ على أوقات دخول مسجَّلة بدقّة.
\end{keyp}

\subsection{التدفّق الكامل في 12 خطوة}

الجدول التالي يجمع كل ما سبق في رحلة واحدة، من لحظة قبول الحجز حتى تنزيل ملفّ
\en{MP4}. المراحل ملوّنة ذهنيًّا: \textbf{القبول} (خطوات 1–3)، \textbf{التحقّق
والرمز} (4–7)، \textbf{الواجهة و\en{SDK}} (8–10)، ثمّ \textbf{التسجيل والحفظ}
(11–12).

\begin{center}
\begin{tabular}{C{0.9cm} L{13.0cm}}
\toprule
\textbf{\color{spPrimary}\#} & \textbf{\color{spPrimary}الخطوة} \\
\midrule
1 & المستشار يقبل الحجز، فيعمل الأمر \code{AcceptBookingCommand}. \\
2 & المعالِج ينادي \code{IMeetingService.CreateRoomAsync} فيتولّد \code{roomId} من نوع \en{GUID}. \\
3 & يُحفَظ \code{MeetingRoomId} في \en{SQL}، ويُنشَر \code{BookingConfirmedEvent} فيصل إشعار للطالب. \\
4 & ضمن نافذة \en{$\pm$15} دقيقة من الموعد، يضغط الطالب أو المستشار «انضمّ». \\
5 & الأمر \code{RecordMeetingJoinCommand} يتحقّق من النافذة، ثمّ ينادي \code{IssueAccessTokenAsync}. \\
6 & \en{ACS} عبر \code{CreateUserAndTokenAsync} بنطاق \en{VoIP} يُصدِر رمزًا قصير العمر (نحو ساعتين). \\
7 & الواجهةُ تستقبل \code{roomId} و\code{accessToken} و\code{acsUserId}. \\
8 & \en{ACS Calling SDK}: إنشاء \code{CallClient} ثمّ \code{createCallAgent} ثمّ \code{agent.join} بـ\code{groupId}. \\
9 & طلب أذونات الأجهزة، تعداد الكاميرات، إنشاء \code{LocalVideoStream} وعرضه. \\
10 & عند \code{Connected}، تنادي الواجهةُ نقطةَ بدء التسجيل مرسِلةً \code{serverCallId}. \\
11 & \code{StartRecordingAsync} يبدأ التسجيل إلى \en{Blob} ويُعيد \code{RecordingId}. \\
12 & بعد الجلسة: يُحفَظ صفّ \code{SessionRecording}، ويستطيع الطالب تنزيل \en{MP4}. \\
\bottomrule
\end{tabular}
\end{center}

\begin{note}
لاحظ أنّ التسجيل يبدأ من \emph{الواجهة} عند لحظة الاتصال، لا من الخادم مسبقًا. هذا
مقصود: لا معنى لتسجيل غرفة فارغة، والبدء عند \code{Connected} يضمن أنّ هناك مكالمة
حيّة فعلًا. والنداء مُصمَّم ليكون آمنًا للتكرار (\en{idempotent})، فلو أطلقه الطرفان
معًا لا يبدأ تسجيلان.
\end{note}

\subsection{الواجهة: تكامل \en{ACS Calling SDK} في \en{Meeting.tsx}}

الآن نرى الجانب الآخر من نقطة الانضمام: ملفّ \code{Meeting.tsx} في الواجهة. بعد أن
يوافق المستخدم على إشعار التسجيل (بوّابة الموافقة \en{CONSENT-01})، يبدأ الكود
الحقيقيّ التالي (مختصَر):

\begin{english}
\begin{lstlisting}[language=Java]
const join = await meetingsApi.join(bookingId!);
// When ACS isn't configured the server issues a stub token.
if (join.provider === "Stub") {
  throw new Error(t("bookings:meeting.errorNotConfigured"));
}

const callClient = new CallClient();
const credential = new AzureCommunicationTokenCredential(join.accessToken);
const agent = await callClient.createCallAgent(credential, { displayName });
// ... ask device permissions, pick a camera, build localStream ...

const call = agent.join(
  { groupId: join.roomId },
  { videoOptions: { localVideoStreams: localStream ? [localStream] : [] },
    audioOptions: { muted: false } },
);

// Start the recording once the call is actually connected (idempotent).
call.on("stateChanged", () => {
  if (call.state === "Connected") void ensureRecording();
});
\end{lstlisting}
\end{english}

هنشوف إيه في هذا المقتطف:
\begin{itemize}
\item \code{meetingsApi.join} يطلب من الخادم بيانات الدخول؛ لو كان المزوّد
  \code{Stub} (بيئة تطوير) نرمي رسالة صريحة «الفيديو غير مُهيَّأ» بدل خطأ غامض من
  الـ\en{SDK}.
\item ننشئ \code{CallClient} ونغلّف الرمز في \code{AzureCommunicationTokenCredential}،
  ثمّ نصنع «وكيلًا» (\code{createCallAgent}) هو مَن يجري المكالمات باسمنا.
\item الانضمام هو \code{agent.join} مع \code{groupId: join.roomId} — هنا يظهر بوضوح
  أنّ الغرفة مجرّد مُعرِّف مجموعة، ونمرّر خيارات الفيديو والصوت.
\item نستمع لتغيّر الحالة، وعند \code{Connected} فقط نستدعي \code{ensureRecording}
  التي تجلب \code{serverCallId} وتطلب بدء التسجيل من الخادم.
\end{itemize}

\begin{secnote}
لاحظ بوّابة \en{CONSENT-01}: لا ينضمّ المستخدم — ومن ثمّ لا يبدأ أي تسجيل — قبل أن
يوافق صراحةً على أنّ الجلسة مسجَّلة. المؤثّر (\en{effect}) الذي يشغّل الانضمام لا
يعمل إلّا بعد قلب علامة الموافقة إلى «نعم»، فلا شيء يدخل الغرفة قبل الموافقة. هذه
حماية للخصوصيّة مبنيّة في التصميم لا مضافة لاحقًا.
\end{secnote}

\subsection{الزمن الحقيقيّ: \en{SignalR} وثلاثة \en{hubs}}

الجزء الثاني من الفصل يخصّ الدفع اللحظيّ للأحداث النصّية: رسالة محادثة جديدة، إشعار
حجز، تفاعل مجتمعيّ. هنا نستخدم \en{SignalR} عبر ثلاثة \en{hubs}، كلّها ترث قاعدةً
واحدة مُوثَّقة (\code{AuthenticatedHub} المعلَّمة بـ\code{[Authorize]})، فلا يتّصل
بها إلّا مستخدم يحمل \en{JWT} صالحًا. الجدول يلخّص مَن يفعل ماذا:

\begin{center}
\begin{tabular}{L{3.0cm} L{6.6cm} L{4.0cm}}
\toprule
\textbf{\color{spPrimary}الـ\en{hub}} & \textbf{\color{spPrimary}الأحداث المدفوعة} & \textbf{\color{spPrimary}من يستخدمه} \\
\midrule
\code{ChatHub} & رسالة جديدة، مؤشّر الكتابة، الحضور (\en{online/offline}) & محادثات المجتمع الخاصّة \\
\code{NotificationHub} & تأكيد حجز، دفعة معالَجة، تذكير موعد & كل الأدوار \\
\code{CommunityHub} & ردود وتفاعلات جديدة، انضمام للأقسام والمواضيع & موجز المجتمع \\
\bottomrule
\end{tabular}
\end{center}

كيف يصل الحدث إلى مستخدم بعينه؟ الفكرة أنّ كل اتّصال جديد يُضاف إلى «مجموعة» خاصّة
بصاحبه. في \code{NotificationHub} مثلًا، عند الاتّصال نضيف الاتّصال إلى مجموعة اسمها
\code{user:<userId>}:

\begin{english}
\begin{lstlisting}[language={[Sharp]C}]
public override async Task OnConnectedAsync()
{
    // Await the group-add: a notification fired right after connect must not
    // race past the join and miss the recipient.
    if (Context.UserIdentifier is not null)
    {
        await Groups.AddToGroupAsync(
            Context.ConnectionId, $"user:{Context.UserIdentifier}");
    }
    await base.OnConnectedAsync();
}
\end{lstlisting}
\end{english}

هنشوف إيه: \code{Context.UserIdentifier} هو مُعرِّف المستخدم المستخرَج من الـ\en{JWT}.
نضيف اتّصاله الحاليّ إلى مجموعة \code{user:<id>}، \emph{وننتظر} إتمام الإضافة
(\code{await}) لا نطلقها ونمضي؛ لأنّ إشعارًا يُدفَع بعد الاتّصال مباشرةً قد «يسبق»
انضمام المجموعة فيضيع على صاحبه ثمّ يظهر لاحقًا كأنّه مكرَّر.

بعد ذلك، الدفعُ نفسه سطرٌ واحد من أي خدمة تحمل \code{IHubContext}. هذا مثال حقيقيّ
من موزّع الإشعارات، يدفع للمجموعة الخاصّة بالمستخدم المستهدَف:

\begin{english}
\begin{lstlisting}[language={[Sharp]C}]
await hub.Clients.Group($"user:{row.RecipientUserId}")
    .SendAsync("notification", new
    {
        id = row.Id,
        type = row.Type.ToString(),
        row.TitleEn,
        row.TitleAr,
        row.DeepLink,
        createdAt = DateTimeOffset.UtcNow,
    }, ct);
\end{lstlisting}
\end{english}

هنشوف إيه: نأخذ \code{hub} (من نوع \code{IHubContext<NotificationHub>})، نستهدف
مجموعة المستخدم، ونرسل حدثًا باسم \code{notification} يحمل حمولة \en{JSON} بسيطة.
على الواجهة، خطّاف \en{React} يستمع لهذا الاسم فيحدّث الحالة فورًا — فيظهر الإشعار
دون أن يُحدِّث المستخدم الصفحة. لاحظ أنّنا نرسل العنوان بالعربية والإنجليزية معًا
(\code{TitleAr} و\code{TitleEn})، فتختار الواجهةُ لغة العرض.

\begin{why}
لماذا مجموعة \code{user:<id>} بدل استهداف اتّصال بعينه؟ لأنّ للمستخدم الواحد قد تكون
عدّة تبويبات أو أجهزة متّصلة في آنٍ. الدفع للمجموعة يصل إلى \emph{كل} جلساته الحيّة
مرّة واحدة، وهو المكافئ العمليّ لِنمط \code{Clients.User(id).SendAsync}. وحين نحتاج
دفعًا لطرفَي محادثة أو لمشتركي موضوع، نستخدم مجموعات أخرى مثل
\code{conversation:<id>} و\code{forum-post:<id>}.
\end{why}

\begin{secnote}
الاشتراك في المجموعات ليس مفتوحًا: قبل أن ينضمّ عميلٌ لمجموعة محادثة أو موضوع، يتحقّق
الـ\en{hub} أنّ صاحب الـ\en{JWT} طرفٌ فعليّ فيها (بوّابة \en{SEC-02})، وإلّا يرمي
\en{HubException}. بغير هذا التحقّق كان أي مستخدم مصادَق يستطيع الاشتراك في
\code{conversation:<id>} والتنصّت على رسائل خاصّة — لذا الاشتراك نفسه محميٌّ لا الدفع
فقط.
\end{secnote}

\begin{summary}[الخلاصة في سطور]
\begin{itemize}
\item مكالمة الفيديو تمرّ عبر منفذ \code{IMeetingService} (محوّل
  \code{AzureCommunicationMeetingService} في الإنتاج، \code{StubMeetingService} في
  التطوير)، ويُختار بالإعداد \code{Acs:ConnectionString}.
\item «الغرفة» = \en{ACS group-call id} أي \en{GUID}، بلا كائن خادم؛ تُنشأ كسولًا
  عند القبول.
\item عند الانضمام: نافذة \en{$\pm$15} دقيقة، ورمز وصول \en{VoIP} قصير العمر (نحو ساعتين)
  \emph{لكل مشارك} عبر \code{CreateUserAndTokenAsync} — لأنّ رمز الفيديو بيانات
  اعتماد.
\item التسجيل عبر \en{Call Automation} (\code{StartRecordingAsync} \en{$\rightarrow$}
  \code{RecordingId}) يُحفَظ كـ\code{SessionRecording} + \en{Blob}؛ والحضور
  (\code{StudentJoinedAt}/\code{ConsultantJoinedAt}) يقوده \code{MeetingNoShowSweepJob}
  لاكتشاف \en{no-show}.
\item الواجهة \code{Meeting.tsx} تنضمّ بـ\code{groupId}، تعرض الكاميرا، وتبدأ التسجيل
  عند \code{Connected} خلف بوّابة موافقة \en{CONSENT-01}.
\item الزمن الحقيقيّ عبر ثلاثة \en{hubs} (\code{ChatHub}/\code{NotificationHub}/%
  \code{CommunityHub}) مصادَقة بـ\en{JWT}؛ الدفع سطرٌ واحد عبر \code{IHubContext}
  إلى مجموعة \code{user:<id>}، والاشتراك في المجموعات محميّ (\en{SEC-02}).
\end{itemize}
\end{summary}

\clearpage
% ch07b-notifications-chat
\section{الإشعارات والشات: إزاي بيشتغلوا فعلاً}

في الفصل السابق فكّكنا \en{SignalR}: ثلاثة \en{hubs} ترث قاعدةً واحدة مصادَقة
بـ\en{JWT}، ومجموعات مثل \code{user:<id>} و\code{conversation:<id>}، والدفع اللحظيّ
سطرٌ واحد عبر \code{IHubContext}. الآن نبني \emph{فوق} تلك الأساسات ميزتين
يستعملهما كلّ مستخدم كل يوم: \textbf{جرس الإشعارات} الذي يومض دون تحديث الصفحة،
و\textbf{الشات الخاصّ 1:1} بين طرفين. لن نُعيد شرح \en{SignalR} نفسه؛ سنتتبّع
بدلًا من ذلك رحلة \emph{الحدث} — من اللحظة التي يقع فيها شيء في النظام حتى تظهر
النتيجة على شاشة الطرف الآخر.

\begin{idea}[الفكرة: ميزتان حيّتان فوق طبقةٍ واحدة]
الجرس والشات يبدوان مختلفَين، لكنّهما يتشاركان الهيكل نفسه: \textbf{يحفظ النظام
الحقيقة أوّلًا في قاعدة البيانات}، ثمّ \textbf{يدفعها حيّةً} عبر \en{SignalR} لِمن
هو متّصل الآن. الحفظ يضمن ألّا يضيع شيء (المستخدم غير المتّصل يجده حين يعود)،
والدفع يضمن الفوريّة لِمن هو حاضر. كل ما يلي مجرّد تفصيل لهاتين الخطوتين.
\end{idea}

\subsection{نقطة الدخول: \code{DispatchAsync} خطوة بخطوة}

كل إشعار في النظام — تأكيد حجز، دفعة نجحت، رسالة جديدة — يمرّ عبر باب واحد اسمه
\code{INotificationDispatcher.DispatchAsync}. المُرسِل لا يعرف شيئًا عن اللغات أو
القنوات أو ساعات الهدوء؛ يقول فقط «أبلِغ هذا المستخدم بهذا النوع»، والموزّع يتكفّل
بالباقي بترتيب ثابت. هذا هو قلب الدالة (مقتصّ):

\begin{english}
\begin{lstlisting}[language={[Sharp]C}]
public async Task DispatchAsync(
    Guid recipientUserId, NotificationType type, NotificationParams parameters,
    string? deepLink, string? idempotencyKey, CancellationToken ct)
{
    // 1) Idempotency — Stripe/event replays must not double-insert.
    if (!string.IsNullOrEmpty(idempotencyKey)
        && await db.Notifications.AnyAsync(n => n.IdempotencyKey == idempotencyKey, ct))
        return;

    var content  = catalog.Render(type, parameters);                  // 2) EN + AR
    var silenced = await IsSilencedAsync(recipientUserId, type, ct);   // 3) DND?
    var channels = await ResolveChannelsAsync(recipientUserId, type, silenced, ct);
    // 4) build a Notification row per channel, SaveChanges, then DeliverAsync each.
}
\end{lstlisting}
\end{english}

هنشوف إيه سطرًا-بسطر:
\begin{itemize}
\item الخطوة \textbf{الأولى} حارس ضدّ التكرار: إن وُجد \code{idempotencyKey} وسبق أن
  خُزِّن إشعارٌ بنفس المفتاح، نخرج فورًا (\code{return}). هذا يمنع الإشعار المكرَّر
  حين تعيد \en{Stripe} إرسال الحدث نفسه مرّتين.
\item الخطوة \textbf{الثانية} \code{catalog.Render}: تحوّل النوعَ والمعاملات إلى نصّ
  فعليّ \textbf{بالعربية والإنجليزية معًا} (عنوان + نصّ لكل لغة). المُرسِل لا يكتب
  نصوصًا، بل يمرّر نوعًا وقيمًا فقط.
\item الخطوة \textbf{الثالثة} \code{IsSilencedAsync}: هل هذا المستخدم في وضع «عدم
  إزعاج» الآن؟ نتيجة \code{silenced} تسري على كل ما بعدها.
\item الخطوة \textbf{الرابعة} \code{ResolveChannelsAsync}: تقرّر القنوات (جرس داخل
  التطبيق؟ بريد؟)، ثمّ يُبنى صفّ \code{Notification} لكل قناة، يُحفَظ، ثمّ يُسلَّم.
\end{itemize}

\subsubsection{منع التكرار: الفحص المسبق وسباق الفهرس الفريد}

الحارس المسبق (\code{AnyAsync}) يكفي في الحالة العاديّة، لكنّه لا يكفي وحده لو وصل
حدثان بنفس المفتاح في اللحظة ذاتها: قد يمرّ كلاهما من الفحص قبل أن يحفظ أيٌّ منهما.
لذلك توجد شبكة أمان ثانية على مستوى قاعدة البيانات: فهرس فريد على
\code{(IdempotencyKey, Channel)}. من يصل ثانيًا يصطدم بالفهرس، فنترجم الاصطدام إلى
«لا عمل» (\en{no-op}) بدل خطأ:

\begin{english}
\begin{lstlisting}[language={[Sharp]C}]
try { await db.SaveChangesAsync(ct); }
catch (ConflictException) when (!string.IsNullOrEmpty(idempotencyKey))
{
    // A concurrent dispatch with the same key won the race; the unique
    // (IdempotencyKey, Channel) index rejected our rows — detach and no-op.
    db.Notifications.RemoveRange(rows);
    return;
}
\end{lstlisting}
\end{english}

هنشوف إيه: لو نجح حفظ الطرف الآخر أوّلًا، يرفض الفهرسُ صفوفَنا فتُرمى
\code{ConflictException}. نلتقطها \textbf{بشرط} وجود مفتاح تكرار (\code{when})، ثمّ
نفصل صفوفنا المرفوضة (\code{RemoveRange}) ونخرج بهدوء. النتيجة النهائيّة واحدة سواء
وصل الحدث مرّة أو ألف مرّة: إشعارٌ واحد فقط.

\begin{why}
لماذا حارسان لا حارس واحد؟ الفحص المسبق يوفّر عملًا في الحالة الشائعة (لا نحاول
الحفظ أصلًا)، لكنّه ليس ذرّيًّا. الفهرس الفريد ذرّيّ ومضمون من قاعدة البيانات لكنّه
يكلّف استثناءً. الجمع بينهما يعطي الأداء الجيّد \textbf{و} الصحّة تحت التزامن معًا —
وهو نمط الحماية بطبقتين نفسه الذي رأيناه في المدفوعات.
\end{why}

\subsection{«عدم الإزعاج»: مكتوم، ساعات هدوء، وسياسة \en{fail-open}}

الدالة \code{IsSilencedAsync} تجيب عن سؤال واحد: هل نكتم هذا المستخدم الآن؟ الجواب
«نعم» في حالتين: إمّا أنّه فعّل الكتم الكامل (\code{NotificationsMuted})، أو أنّ
اللحظة الحاليّة تقع داخل \textbf{ساعات الهدوء} (\code{QuietHours}) التي حدّدها.
وفيها ثلاث دقائق مهمّة:

\begin{itemize}
\item \textbf{بمنطقة المستخدم الزمنيّة}: نافذة الهدوء تُحسَب بتوقيت المستخدم لا
  بتوقيت الخادم، فـ«العاشرة مساءً» تعني مساءه هو.
\item \textbf{قد تلفّ حول منتصف الليل}: نافذة مثل \en{22:00 $\rightarrow$ 08:00}
  تبدأ في يوم وتنتهي في اليوم التالي.
\item \textbf{اختبار المستخدم يمرّ دائمًا}: النوع \code{SystemTest} يُعيد «غير مكتوم»
  فورًا — لأنّه الوسيلة التي يتحقّق بها المستخدم من إعداداته نفسها.
\end{itemize}

منطق «اللفّ حول منتصف الليل» هو النقطة الذكيّة، وهذا هو الكود:

\begin{english}
\begin{lstlisting}[language={[Sharp]C}]
private static bool InQuietWindow(TimeOnly start, TimeOnly end, string? tzId)
{
    var localNow = TimeOnly.FromDateTime(
        TimeZoneInfo.ConvertTime(DateTimeOffset.UtcNow, ResolveTimeZone(tzId)).DateTime);
    // The window may wrap past midnight (e.g. 22:00 -> 08:00).
    return start <= end
        ? localNow >= start && localNow < end    // same day
        : localNow >= start || localNow < end;   // 22:00 -> 08:00
}
\end{lstlisting}
\end{english}

هنشوف إيه:
\begin{itemize}
\item نحوّل «الآن» \en{UTC} إلى الوقت المحلّي للمستخدم عبر \code{ConvertTime} ثمّ
  نأخذ الجزء الزمنيّ فقط (\code{localNow}).
\item إن كانت \code{start <= end} فالنافذة داخل اليوم نفسه، فنطلب أن يقع «الآن»
  بينهما (\code{and}).
\item وإلّا فالنافذة ملتفّة حول منتصف الليل، فنطلب أن يقع «الآن» \textbf{بعد} البداية
  \textbf{أو قبل} النهاية (\code{or}) — فتشمل ليلة كاملة تعبر منتصف الليل.
\end{itemize}

\begin{secnote}
لاحظ فلسفة \code{IsSilencedAsync} كلّها: \textbf{تفشل مفتوحةً} (\en{fail-open}). لو
كان ملفّ المستخدم مفقودًا، أو مُعرِّف المنطقة الزمنيّة تالفًا، أو أيّ استثناء آخر —
تُرجِع «غير مكتوم» فيصل الإشعار عاديًّا. القاعدة: خطأٌ في منطق الكتم يجب ألّا
\textbf{يبتلع} إشعارًا مهمًّا. الأسوأ المقبول هو إزعاجٌ عابر، لا معلومة ضائعة.
\end{secnote}

\subsection{القنوات والدفع الحيّ للجرس}

بعد معرفة حالة الكتم، تقرّر \code{ResolveChannelsAsync} أين يذهب الإشعار. القاعدة
بسيطة: \code{InApp} (الجرس) \textbf{دائمًا}، و\code{Email} حسب تفضيلات المستخدم لكل
\emph{نوع} إشعار على حدة. أمّا تحت «عدم الإزعاج» فالاثنان يُكتَمان — \textbf{لكن}
صفّ \code{InApp} يظلّ يُخزَّن في قاعدة البيانات. الجدول يلخّص السلوك:

\begin{center}
\begin{tabular}{L{4.4cm} C{2.4cm} C{2.9cm} C{2.0cm}}
\toprule
\textbf{\color{spPrimary}الحالة} & \textbf{\color{spPrimary}صفّ \code{InApp}} & \textbf{\color{spPrimary}دفعة حيّة (\en{toast})} & \textbf{\color{spPrimary}\code{Email}} \\
\midrule
عاديّ، وتفضيل البريد مفعّل   & يُخزَّن & يُدفَع & يُرسَل \\
عاديّ، وتفضيل البريد معطَّل  & يُخزَّن & يُدفَع & لا \\
تحت «عدم الإزعاج»           & يُخزَّن & لا     & لا \\
\bottomrule
\end{tabular}
\end{center}

المفتاح في السطر الأخير: تحت الكتم تُرجِع \code{ResolveChannelsAsync} قناة
\code{InApp} وحدها (تُسقِط \code{Email})، وفي داخل التسليم تُكتَم أيضًا الدفعةُ الحيّة
لتلك القناة. النتيجة: المستخدم \emph{يجد} الإشعار في قائمته لاحقًا، لكنّه لا يُزعَج
به لحظتَها. وهذا هو سطر الدفع الحيّ داخل \code{DeliverAsync} لقناة \code{InApp}:

\begin{english}
\begin{lstlisting}[language={[Sharp]C}]
// The in-app row is always stored; the live toast is suppressed under DND.
if (!silenced)
    await hub.Clients.Group($"user:{row.RecipientUserId}")
        .SendAsync("notification", new {
            id = row.Id, type = row.Type.ToString(),
            row.TitleEn, row.TitleAr, row.BodyEn, row.BodyAr,
            row.DeepLink, createdAt = DateTimeOffset.UtcNow }, ct);
\end{lstlisting}
\end{english}

هنشوف إيه: الشرط \code{if (!silenced)} هو ما يكتم الدفعة الحيّة فقط (الصفّ خُزِّن قبل
هذا السطر أصلًا). نستهدف مجموعة \code{user:<id>} — التي انضمّ إليها المستخدم عند
الاتّصال — ونرسل حدثًا اسمه \code{notification} يحمل العنوان والنصّ باللغتين مع
رابط عميق (\code{DeepLink}). الواجهة تستمع لهذا الاسم فيومض الجرس فورًا.

\begin{qa}
\textbf{سؤال متوقّع:} «تحت عدم الإزعاج، لماذا نخزّن صفّ الجرس أصلًا بدل تجاهله؟» لأنّ
الكتم يعني «لا تُزعِجني الآن»، لا «احذف الخبر». المستخدم قد يفتح قائمة الإشعارات بعد
ساعات الهدوء ويريد أن يجد كل ما فاته. الكتم يقمع \emph{المقاطعة} لا \emph{المعلومة}.
\end{qa}

\begin{keyp}[نقطة مهمة: لماذا \code{await} في \code{OnConnectedAsync}؟]
رأينا في الفصل السابق أنّ \code{NotificationHub} يضيف الاتّصال إلى مجموعة
\code{user:<id>} عند الاتّصال، \textbf{وينتظر} الإضافة (\code{await}) لا يطلقها
ويمضي. السبب حرِج هنا تحديدًا: لو دُفِع إشعارٌ في اللحظة التالية للاتّصال مباشرةً،
ومجموعةُ المستخدم لم تكتمل بعد، لَسَبَق الإشعارُ الانضمامَ وضاع على صاحبه — ثمّ ظهر
لاحقًا كأنّه «مكرَّر» بعد إعادة الإرسال. الانتظار يقفل هذه الفجوة الزمنيّة.
\end{keyp}

\subsection{الشات 1:1: الحضور والبوّابة قبل أي رسالة}

ننتقل للميزة الثانية. \code{ChatHub} يرث \code{AuthenticatedHub} المعلَّم
بـ\code{[Authorize]}، فلا يتّصل به إلّا حاملُ \en{JWT} صالح. فوق ذلك يتتبّع
\textbf{الحضور} (\en{online/offline}) بعدّاد مرجعيّ: أوّل اتّصال للمستخدم يعلنه
«متّصلًا» للآخرين، وآخر اتّصال يسقط (وصل العدّاد صفرًا) يعلنه «غير متّصل». والعميلُ
حديثُ الاتّصال يستدعي \code{GetOnlineUsers} ليهيّئ حالته دفعةً واحدة بدل انتظار
الأحداث تتقاطر.

لكن قبل أن ينضمّ أي عميل إلى مجموعة محادثة أو يبثّ مؤشّر كتابة، تمرّ كل عمليّة عبر
بوّابة واحدة، \code{EnsureParticipantAsync}:

\begin{english}
\begin{lstlisting}[language={[Sharp]C}]
private async Task EnsureParticipantAsync(string conversationId)
{
    if (!Guid.TryParse(Context.UserIdentifier, out var userId))
        throw new HubException("Unauthorized.");
    if (!Guid.TryParse(conversationId, out var convId))
        throw new HubException("Invalid conversation id.");

    var isParticipant = await db.Conversations.AsNoTracking()
        .AnyAsync(c => c.Id == convId
            && (c.ParticipantOneId == userId || c.ParticipantTwoId == userId));
    if (!isParticipant)
        throw new HubException("You are not a participant in this conversation.");
}
\end{lstlisting}
\end{english}

هنشوف إيه: نأخذ هويّة المتّصل من \code{Context.UserIdentifier} (من الـ\en{JWT} لا من
طلب العميل)، ثمّ نسأل قاعدة البيانات: هل هذا المستخدم أحد طرفَي المحادثة
(\code{ParticipantOneId} أو \code{ParticipantTwoId})؟ إن لم يكن، نرمي
\code{HubException} فورًا فلا انضمام ولا كتابة. الدوالّ \code{JoinConversation}
و\code{TypingStart} و\code{TypingStop} كلّها تستدعي هذه البوّابة أوّلًا.

\begin{secnote}
هذه بوّابة \en{SEC-02 IDOR}. بدونها كان \textbf{أي} مستخدم مصادَق يستطيع الانضمام
إلى \code{conversation:<id>} لأي محادثة — بمجرّد تخمين المُعرِّف — والتنصّت على
رسائل خاصّة تُدفَع لتلك المجموعة. الدرس: المصادقة (\en{authentication}) تُثبِت
\emph{مَن أنت}، لكنّ التفويض (\en{authorization}) — أنت طرفٌ في \emph{هذه}
المحادثة تحديدًا — يجب فحصه على الخادم قبل كل اشتراك، لا الاكتفاء بأنّك مسجَّل
دخولك.
\end{secnote}

مؤشّرات الكتابة تتبع النمط نفسه: \code{TypingStart}/\code{TypingStop} تبثّان — بعد
البوّابة — إلى \code{OthersInGroup} أي «كل مَن في المجموعة عداي». فيرى الطرف الآخر
«يكتب الآن…» دون أن يرى المرسِلُ إشعارَ نفسه.

\subsection{من الرسالة إلى الطرف الآخر: \code{SendMessageCommandHandler}}

إرسال رسالة يمرّ بستّ خطوات مرتّبة داخل معالِج واحد. المقتطف التالي يجمع أهمّها
(مقتصّ):

\begin{english}
\begin{lstlisting}[language={[Sharp]C}]
// 1) canonical order (smaller GUID first) -> one conversation per pair
var (participantOne, participantTwo) = senderId.CompareTo(request.RecipientId) <= 0
    ? (senderId, request.RecipientId)
    : (request.RecipientId, senderId);
// ... find-or-create the conversation with that ordered pair ...

// 2) block check (either direction)
if (isBlocked) throw new ConflictException("Cannot send message. User is blocked.");

// 3) SEC-07 defense-in-depth: strip HTML before it hits any sink
var sanitizedBody = new Ganss.Xss.HtmlSanitizer().Sanitize(request.Body);

// 4) persist the message + touch the conversation
db.Messages.Add(new ChatMessage { ConversationId = conversation.Id,
    SenderId = senderId, Body = sanitizedBody, SentAt = DateTimeOffset.UtcNow });
await db.SaveChangesAsync(ct);

// 5) live push to whoever has the conversation open
await chatNotifier.NotifyNewMessageAsync(
    conversation.Id, message.Id, senderId, sanitizedBody, message.SentAt, ct);
// 6) bell + email for the absent party (elsewhere in the app)
await publisher.Publish(new ChatMessageReceivedEvent(/* ... */), ct);
\end{lstlisting}
\end{english}

هنشوف إيه خطوةً-خطوة:
\begin{itemize}
\item \textbf{(1) الترتيب \en{canonical}}: نرتّب الطرفين بحيث يكون الأصغر
  (\code{CompareTo}) أوّلًا قبل إنشاء المحادثة. هذا يجعل الفهرس الفريد على الزوج
  يضمن \textbf{محادثة واحدة لكل ثنائيّ}، ويمنع سباق رسالتين متزامنتين من الطرفين
  يُنشئان صفّين بأعمدة مقلوبة.
\item \textbf{(2) فحص الحظر}: إن كان أيٌّ من الطرفين حاظرًا للآخر (في أيّ اتّجاه)
  نرمي \code{ConflictException} فلا رسالة تُرسَل.
\item \textbf{(3) تنقية \en{HTML}} عبر \code{HtmlSanitizer}: نجرّد أي وسوم أو سكربتات
  من النصّ \emph{قبل} تخزينه.
\item \textbf{(4) الحفظ}: نضيف \code{ChatMessage} ونحدّث المحادثة (آخر رسالة
  ووقتها)، ثمّ \code{SaveChanges}.
\item \textbf{(5) الدفع الحيّ}: \code{NotifyNewMessageAsync} يبثّ لمجموعة المحادثة.
\item \textbf{(6) الحدث}: ننشر \code{ChatMessageReceivedEvent} ليتكفّل بجرس الطرف
  الغائب وبريده.
\end{itemize}

\begin{secnote}
لماذا ننقّي \en{HTML} والواجهةُ \en{React} تعرض النصّ آمنًا أصلًا؟ لأنّ نفس النصّ
المخزَّن يذهب إلى \textbf{ثلاثة مصارف} (\en{sinks}): عقدة نصّ في \en{React}
(آمنة)، ودفعة \en{SignalR} حيّة، ومعاينة داخل بريد \en{HTML} — والأخيران ليسا
\en{React}. التنقية \textbf{عند التخزين} تجعل كل مصرف آمنًا مرّة واحدة بدل الاعتماد
على أن يتذكّر كلُّ مستهلِك أن يهرب النصّ. هذا معنى \en{SEC-07}: دفاعٌ في العمق.
\end{secnote}

الدفع الحيّ نفسه سطرٌ واحد في \code{ChatRealtimeNotifier}، يبثّ لمجموعة المحادثة:

\begin{english}
\begin{lstlisting}[language={[Sharp]C}]
hub.Clients.Group($"conversation:{conversationId}")
   .SendAsync("MessageReceived",
       new { conversationId, messageId, senderId, body, sentAt }, ct);
\end{lstlisting}
\end{english}

هنشوف إيه: نستهدف مجموعة \code{conversation:<id>} — التي لا ينضمّ إليها إلّا من
اجتاز بوّابة \code{EnsureParticipantAsync} — ونرسل حدثًا اسمه \code{MessageReceived}.
كلّ عميل يشاهد هذه المحادثة الآن يستقبله فتظهر الرسالة فورًا في الخيط.

\begin{keyp}[نقطة مهمة: المسار المزدوج]
لماذا خطوتان (5) و(6) معًا؟ لأنّهما تصلان إلى شخصين مختلفين. الدفع الحيّ
(\code{NotifyNewMessageAsync}) يصل \textbf{فقط} لِمن يفتح المحادثة على شاشته الآن —
هو مشترِكٌ في مجموعة \code{conversation:<id>}. أمّا الطرف الموجود في مكان آخر من
التطبيق (صفحة المنح مثلًا) فليس في تلك المجموعة، فيتكفّل به الحدثُ
\code{ChatMessageReceivedEvent}: يُطلِق إشعار جرسٍ عبر المسار الذي شرحناه في أوّل
الفصل. \textbf{الخلاصة}: المفتوحة أمامك تصلك حيّةً، والمغلقة يذكّرك بها الجرس — فلا
رسالة تفوت.
\end{keyp}

\begin{note}
لاحظ التناظر بين الميزتين: كلتاهما \textbf{تحفظ ثمّ تدفع}. الإشعار يحفظ صفّ
\code{Notification} ثمّ يدفع لمجموعة \code{user:<id>}؛ والرسالة تحفظ \code{ChatMessage}
ثمّ تدفع لمجموعة \code{conversation:<id>} \emph{و}تنشر حدثًا للجرس. الدوام أوّلًا،
الفوريّة ثانيًا — هو النمط الذي يجعل الميزتين موثوقتين لا مجرّد «حيّتين».
\end{note}

\begin{summary}[الخلاصة في سطور]
\begin{itemize}
\item كل إشعار يمرّ عبر \code{DispatchAsync} بترتيب ثابت: منع تكرار
  \en{$\rightarrow$} تصيير النصّ باللغتين (\code{catalog.Render}) \en{$\rightarrow$}
  فحص الكتم \en{$\rightarrow$} تحديد القنوات \en{$\rightarrow$} حفظ \en{$\rightarrow$}
  تسليم.
\item منع التكرار بطبقتين: فحص \code{IdempotencyKey} المسبق + فهرس فريد على
  \code{(IdempotencyKey, Channel)} يترجم السباق إلى \en{no-op}.
\item «عدم الإزعاج» = مكتوم أو داخل ساعات الهدوء (بمنطقة المستخدم، وقد تلفّ حول منتصف
  الليل)؛ السياسة \en{fail-open}، و\code{SystemTest} يمرّ دائمًا. تحت الكتم يُخزَّن
  صفّ \code{InApp} لكن لا \en{toast} ولا \code{Email}.
\item الشات محميّ ببوّابة \code{EnsureParticipantAsync} (\en{SEC-02 IDOR}) قبل أي
  انضمام أو كتابة — المتّصل من الـ\en{JWT} يجب أن يكون طرفًا في المحادثة.
\item \code{SendMessageCommandHandler}: ترتيب \en{canonical} (محادثة واحدة لكل زوج)
  \en{$\rightarrow$} فحص حظر \en{$\rightarrow$} تنقية \en{HTML} (\en{SEC-07}، لثلاثة
  مصارف) \en{$\rightarrow$} حفظ \en{$\rightarrow$} دفع حيّ \en{$\rightarrow$} حدثٌ
  للغائب.
\item المسار المزدوج: الدفع الحيّ لِمن يفتح المحادثة، وحدث الجرس لِمن هو في مكانٍ
  آخر — فلا رسالة تُفوَّت.
\end{itemize}
\end{summary}

\clearpage
% ch08-payments-jobs — الحجز والدفع (Money Flow) والمهامّ الخلفيّة
\section{الحجز والدفع (Money Flow) والمهامّ الخلفيّة}

في هذا الفصل نمشي مع «الفلوس» من أوّل ما الطالب يحجز جلسة مع مستشار، لحدّ ما
تصل صافي أرباحها لجيب المستشار. الفكرة الأساسية بسيطة جدًّا في جملة واحدة:
\textbf{فلوس الطالب تتحجز ولا تتاخد}، ولا تُخصَم فعليًّا إلا لمّا المستشار يقبل
الجلسة. ده نموذج شبيه بالـ \en{escrow} (الوسيط الأمين) مبنيّ على \en{Stripe}.
وفي آخر الفصل هنشوف «المهامّ الخلفيّة» (\en{background jobs}) اللي بتشتغل لوحدها
في الخلفية: تذكيرات، صرف أرباح، تنظيف، والتزامات \en{GDPR}.

كل الكلام اللي جاي مقتبس من الكود الحقيقي للنظام في \code{server/src/...}؛
هنعرض مقتطفات قصيرة حقيقية ونشرحها سطرًا بسطر.

\subsection{لماذا نموذج authorize-then-capture؟}

تخيّل إنّ الطالب دفع كامل المبلغ لحظة الحجز، وبعدين المستشار رفض أو مردّش. دلوقتي
إحنا ماسكين فلوس الطالب ولازم نعمل عملية استرداد كاملة. الأفضل من كده إنّنا
\textbf{نحجز} المبلغ على البطاقة (نأخذ «تفويضًا» فقط) من غير ما نخصمه، ونستنّى.

\begin{idea}[الفكرة: التفويض ثم القبض (authorize then capture)]
عملية الدفع على البطاقة بتتقسّم لخطوتين منفصلتين:
\begin{itemize}
\item \textbf{التفويض} (\en{authorize}): البنك بيحجز المبلغ على بطاقة الطالب
  ويضمن إنّه موجود، لكن الفلوس \emph{لسه ما اتحوّلتش} لحسابنا.
\item \textbf{القبض} (\en{capture}): دي الخطوة اللي بتحوّل الفلوس فعليًّا. إحنا
  بنأجّلها لحدّ ما المستشار يقبل الجلسة.
\end{itemize}
ما بين الخطوتين، فلوس الطالب «محجوزة» في مكان آمن: لا هو خسرها، ولا إحنا أخدناها.
\end{idea}

المنفذ (\en{port}) اللي بيمثّل كل تعاملاتنا مع \en{Stripe} اسمه \code{IStripeService}،
وله محوّلان (\en{adapters}): \code{StripeService} (الحقيقي، بيتكلّم مع \en{Stripe}
فعلًا) و\code{StubStripeService} (بديل وهميّ للتطوير والاختبار، بيردّ نتائج جاهزة
من غير أي اتّصال بالإنترنت). زيّ ما اتكلّمنا في مبدأ الـ \en{ports} قبل كده، ده
بيخلّي منطق الحجز نفسه ما يعرفش حاجة عن \en{Stripe}، ونقدر نبدّل المزوّد أو نختبر
من غير ما نلمس قواعد العمل.

\begin{why}
ليه التفويض-ثم-القبض ومش الخصم المباشر؟ لأنّه بيحمي الطرفين. الطالب مش هيتخصم منه
قرش إلا لو الخدمة اتأكّدت فعلًا (المستشار قبل)، والمنصّة مش محتاجة تعمل عمليات
استرداد كتير ومعقّدة لو الحجز اتلغى قبل ما يبدأ. ده كمان بيقلّل نزاعات البطاقات
(\en{chargebacks}) لأنّ المبلغ ما بيتقبضش غير مقابل خدمة حقيقية.
\end{why}

\subsection{دورة الحياة خطوةً بخطوة}

خلّينا نتفرّج على رحلة الدفعة الواحدة من أوّلها لآخرها. الجدول ده بيلخّص المحطّات
الخمس، وبعده هنفصّل كل واحدة بالكود:

\begin{center}
\begin{tabular}{L{3.6cm} L{4.2cm} L{5.2cm}}
\toprule
\textbf{\color{spPrimary}المحطّة} & \textbf{\color{spPrimary}الفعل التقني} & \textbf{\color{spPrimary}مصير الفلوس} \\
\midrule
حجز + تفويض & \en{PaymentIntent} بـ \en{manual capture} & محجوزة على البطاقة (لا مخصومة) \\
المستشار يقبل & \code{CapturePaymentIntentAsync} & تُقبَض فعليًّا وتدخل النظام \\
رفض / انتهاء المهلة & \code{CancelPaymentIntentAsync} & يُفَكّ الحجز، الطالب ما اتخصمش \\
\en{no-show} / إلغاء & \code{RefundPaymentAsync} + \code{RefundCalculatorService} & استرداد كامل أو جزئيّ حسب السياسة \\
صرف للمستشار & \code{StripePayoutJob} + \en{Stripe Connect} & صافي الأرباح يوصل للمستشار \\
\bottomrule
\end{tabular}
\end{center}

\subsubsection{المحطّة الأولى: الحجز والتفويض}

هنشوف إزاي بننشئ الـ \en{PaymentIntent} بنمط \en{manual capture} لحظة الحجز.
ده مقتطف من إنشاء نيّة الدفع في \code{StripeService}:

\begin{english}
\begin{lstlisting}[language={[Sharp]C}]
var createOptions = new PaymentIntentCreateOptions
{
    Amount = amountCents,           // integer cents — no floating point
    Currency = currency,            // "usd"
    CaptureMethod = captureMethod,  // "manual" -> authorize, don't charge
    Metadata = metadata?.ToDictionary(k => k.Key, v => v.Value),
    PaymentMethodTypes = ["card"],  // card only
};

var requestOptions = new RequestOptions
{
    ApiKey = _opts.SecretKey,
    IdempotencyKey = idempotencyKey, // retry-safe: never double-charge
};

var intent = await service.CreateAsync(createOptions, requestOptions, ct);
\end{lstlisting}
\end{english}

نفصّص أهمّ الأسطر:
\begin{itemize}
\item \code{Amount = amountCents}: المبلغ دايمًا بالسنتات كعدد صحيح
  (\en{integer cents}). يعني 10.50 دولار بتتخزّن 1050. ليه؟ عشان نتجنّب أخطاء
  الكسور العشرية في أرقام الفاصلة العائمة (\en{floating point}). قرش واحد
  ضايع في حسابات الفلوس مشكلة كبيرة.
\item \code{CaptureMethod = "manual"}: ده قلب الموضوع كلّه. «\en{manual}» معناها
  «فوّض دلوقتي، واقبض لمّا أنا أقولك». من غير السطر ده، \en{Stripe} كان هيخصم
  فورًا.
\item \code{PaymentMethodTypes} بقيمة \code{card}: نقصر الدفع على الكروت بس. من
  غير كده \en{Stripe} بيورّي كل المحافظ المفعّلة في لوحة التحكّم وده كان بيلخبط
  المستخدمين.
\item \code{IdempotencyKey}: مفتاح «عدم التكرار». لو الطلب اتبعت مرّتين (بسبب
  ضغطة زرّ مكرّرة أو إعادة محاولة الشبكة)، \en{Stripe} بيتعرّف على المفتاح ده
  ويرجّع نفس النتيجة \emph{من غير ما يعمل حجزًا تاني}. الطالب ما بيتخصمش مرّتين.
\end{itemize}

الطريف إنّ صفّ \code{Payment} في قاعدة بياناتنا بيبدأ بحالة \code{Pending} (مش
\code{Held})، لأنّ نيّة الدفع اليدوية «مش محجوزة فعلًا» إلا بعد ما الطالب يفوّض
البطاقة في المتصفّح. ساعتها \en{Stripe} بيبعت لنا \en{webhook} يقول «تمّ التفويض»،
فنقلب الحالة لـ \code{Held}. (هنرجع للـ \en{webhooks} بعد شوية.)

\subsubsection{المحطّة الثانية: القبول يعني القبض}

لمّا المستشار يقبل الحجز، بنستدعي \code{CapturePaymentIntentAsync} عشان نقبض
المبلغ المحجوز. لو النيّة مش في حالة «جاهزة للقبض» (يعني الطالب ما فوّضش أصلًا)،
بنترجم الخطأ لرسالة واضحة بدل خطأ 500 غامض:

\begin{english}
\begin{lstlisting}[language={[Sharp]C}]
catch (StripeException ex)
{
    // A capture failure is almost always a payment-state problem — the
    // student never authorized the card, so the intent is not in
    // requires_capture. Surface it as a clean 422, not an opaque 500.
    throw new BookingDomainException(
        "The session fee could not be captured. The student may not have "
        + "completed the card authorization for this booking yet — ask them "
        + "to finish payment, then try accepting again.",
        ex);
}
\end{lstlisting}
\end{english}

هنا الفكرة: لو القبض فشل، غالبًا السبب إنّ الطالب ما كمّلش التفويض. فبدل ما
نرمي خطأ سيرفر مبهم، بنرمي استثناء عمل (\en{domain exception}) بيتحوّل لردّ 422
برسالة مفهومة للمستشار: «كمّل الطالب الدفع الأوّل، وبعدين اقبل تاني». تجربة أفضل
بكتير من شاشة خطأ حمرا.

\subsubsection{المحطّة الثالثة والرابعة: الإفراج والاسترداد}

لو المستشار رفض الحجز أو انتهت مهلة ردّه، بنفكّ الحجز عبر
\code{CancelPaymentIntentAsync}، والطالب ما بيتخصمش أصلًا. أما لو حصل إلغاء بعد
القبول، أو حصل \en{no-show} (حد ما حضرش)، فالموضوع محتاج حساب: كام نرجّع للطالب؟
ده شغل \code{RefundCalculatorService}. خلّينا نشوف منطق الطالب اللي بيلغي:

\begin{english}
\begin{lstlisting}[language={[Sharp]C}]
var moreThan24HoursBefore = scheduledStartAt.ToUniversalTime() > nowUtc.AddHours(24);

if (moreThan24HoursBefore)
{
    return new RefundCalculationResult
    {
        RefundPercentage = 100,
        RefundAmountCents = amountCents,
        CancellationReason = CancellationReason.StudentCancelledMoreThan24HoursBefore
    };
}

return new RefundCalculationResult
{
    RefundPercentage = 50,
    RefundAmountCents = amountCents / 2,
    CancellationReason = CancellationReason.StudentCancelledLessThan24HoursBefore
};
\end{lstlisting}
\end{english}

القاعدة زيّ ما بتشوف مباشرة وعادلة:
\begin{itemize}
\item لو الطالب لغى قبل الموعد بأكتر من 24 ساعة: استرداد كامل 100\%.
\item لو لغى في آخر 24 ساعة: استرداد 50\% بس (لأنّ المستشار حجز وقته).
\item لو \emph{المستشار} هو اللي لغى: الطالب بياخد 100\% دايمًا (ولو كان في آخر
  24 ساعة، بيتّطبق كمان خصم على تقييم المستشار من قِبَل المُعالِج).
\end{itemize}

\begin{keyp}[نقطة مهمة: أساس الاسترداد هو المبلغ المخزَّن]
لاحظ إنّ الحساب بيستخدم \code{amountCents} المخزَّن على صفّ الـ \code{Payment}،
مش سعر الحجز الحاليّ. ليه؟ لأنّ السعر ممكن يتغيّر بعد كده (خصم، عملة، إعادة
تسعير)، فاللقطة (\en{snapshot}) المخزَّنة وقت الدفع هي المرجع الموثوق. ده مبدأ
مهمّ: في المحاسبة، بنعتمد على القيمة وقت المعاملة مش القيمة الحالية.
\end{keyp}

بعد ما الحاسبة تقرّر المبلغ، بيتنفّذ الاسترداد فعليًّا عبر \code{RefundPaymentAsync}
على \en{Stripe}. ولو الاسترداد بسبب \en{no-show}، فالمهمّة الخلفيّة
\code{MeetingNoShowSweepJob} هي اللي بترصد الموقف وتشغّل المسار ده تلقائيًّا.

\subsection{تقسيم الأرباح}

كل صفّ \code{Payment} بيخزّن ثلاثة أرقام رئيسية بالسنتات:

\begin{center}
\begin{tabular}{L{4.4cm} L{8.4cm}}
\toprule
\textbf{\color{spPrimary}الحقل} & \textbf{\color{spPrimary}معناه} \\
\midrule
\code{AmountCents} & إجمالي المبلغ اللي دفعه الطالب (الجُملة). \\
\code{ProfitShareAmountCents} & حصّة المنصّة (الرسوم + نسبة الربح). \\
\code{PayeeAmountCents} & صافي نصيب المستشار بعد خصم حصّة المنصّة. \\
\bottomrule
\end{tabular}
\end{center}

النِّسَب نفسها مش مكتوبة في الكود؛ هي بتيجي من قاعدة ماليّة قابلة للتهيئة اسمها
\code{FinancialConfigRule} (ومعاها \code{ProfitShareConfig}). الحساب نفسه دالة
رياضية نقيّة (\en{pure function}) بتشتغل بالسنتات الصحيحة عشان المجموع يرجع
مظبوط بالظبط:

\begin{english}
\begin{lstlisting}[language={[Sharp]C}]
var feeCents = feeKind switch
{
    FeeKind.Percentage => (long)Math.Round(
        grossAmountCents * (feePercentage ?? 0m), MidpointRounding.AwayFromZero),
    FeeKind.FixedAmount => feeAmountCents ?? 0L,
    _ => 0L,
};

var profitShareCents = (long)Math.Round(
    grossAmountCents * profitSharePercentage, MidpointRounding.AwayFromZero);

var platformTotalCents = feeCents + profitShareCents;
// Payee takes the remainder so the split always re-sums to gross exactly.
var payeeNetCents = grossAmountCents - platformTotalCents;
\end{lstlisting}
\end{english}

السطر الأخير هو الحيلة الذكية: نصيب المستشار = الإجمالي ناقص حصّة المنصّة.
بكده مهما حصل تقريب في الكسور، \textbf{مجموع الحصّتين بيساوي الإجمالي بالظبط}
ولا يضيع قرش. الرسوم ممكن تكون نسبة مئوية أو مبلغًا ثابتًا (حسب \code{FeeKind})،
لكن نسبة الربح دايمًا نسبة مئوية.

\begin{note}
النِّسَب بتتحدّد عند القبول، مش عند الحجز. عند الحجز بيتخزّن
\code{ProfitShareAmountCents} بصفر و\code{PayeeAmountCents} بكامل المبلغ (كأنّ
كل المبلغ للمستشار مؤقّتًا)، وبس لمّا المستشار يقبل والقبض ينجح، بيتعاد حساب
الحصّتين من القاعدة الماليّة السارية ساعتها. كده لو الأدمن غيّر النِّسَب، الحجوزات
القديمة المدفوعة ما تتأثّرش.
\end{note}

\subsubsection{قاعدة فعّالة واحدة لكل نوع}

طب لو في أكتر من قاعدة ماليّة لنفس نوع الدفع، نستخدم أنهي واحدة؟ الحلّ أنيق:
\textbf{فهرس فريد مُصفّى} (\en{filtered unique index}) في قاعدة البيانات بيضمن إنّه
لا يمكن وجود أكتر من قاعدة \emph{فعّالة} واحدة لكل نوع:

\begin{english}
\begin{lstlisting}[language={[Sharp]C}]
// FR-170: at most one Active rule per payment type.
b.HasIndex(r => r.PaymentType)
    .IsUnique()
    .HasFilter("[Status] = 'Active'")
    .HasDatabaseName("UX_FinancialConfigRule_ActivePerType");
\end{lstlisting}
\end{english}

الشرط \code{HasFilter} بقيمة \code{[Status] = 'Active'} معناه إنّ قيد التفرّد
بيتطبّق بس على الصفوف الفعّالة. يعني ممكن يكون عندك عشر قواعد قديمة بحالة غير
فعّالة، لكن \emph{واحدة بس} فعّالة لكل نوع. لو حد حاول يفعّل تانية، قاعدة البيانات
نفسها بترفض. القاعدة دي محفورة في الداتابيز مش في كود التطبيق، فمستحيل تتنسى.

\subsection{الصرف للاستشاري عبر Stripe Connect}

لحدّ دلوقتي، الفلوس المقبوضة موجودة في حساب المنصّة على \en{Stripe}. إزاي توصل
للمستشار؟ عن طريق \en{Stripe Connect}: كل مستشار بيعمل حساب \en{Express} خاص بيه
(عبر \code{CreateConnectAccountAsync} + رابط تسجيل)، وبعدين مهمّة خلفيّة اسمها
\code{StripePayoutJob} بتشتغل كل ليلة وتجمّع كل الدفعات المقبوضة اللي لسه ما
اتصرفتش لكل مستشار، وتبعتله دفعة واحدة (\en{payout}) بمجموع \code{PayeeAmountCents}
بتاعته.

\begin{idea}[الفكرة: المطالبة المُسبقة قبل استدعاء Stripe]
مهمّة الصرف بتعمل حاجة ذكية: قبل ما تنادي \en{Stripe} خالص، بتسجّل صفّ الـ
\code{Payout} في قاعدة البيانات وتربط الدفعات بيه (\code{PayoutId}). ليه الأول؟
عشان لو المهمّة اتشغّلت مرّتين بالغلط، النسخة التانية تلاقي الدفعات «متطالَب بيها»
فما تصرفهاش تاني. ولو استدعاء \en{Stripe} فشل بعد كده، بيتشال الربط
(بترجع \code{PayoutId} فاضية) عشان الليلة الجاية تحاول من جديد. وكمان الصرف نفسه
بيتبعت بمفتاح عدم تكرار مشتقّ من معرّف الـ \code{Payout}، فمفيش صرف مزدوج.
\end{idea}

\subsection{مطابقة الـ webhook والضمانات الثلاث}

\en{Stripe} بيكلّمنا رجوعًا عن طريق \en{webhooks}: رسائل \en{HTTP} بيبعتها لخادمنا
لمّا يحصل حدث (اكتمل تفويض، حصل استرداد، اكتمل صرف...). المشكلة: أي حد يقدر يبعت
طلب \en{HTTP} مزيّف يدّعي إنّه من \en{Stripe}. عشان كده لازم \textbf{نتحقّق من
التوقيع} قبل ما نصدّق أي رسالة. ده بيحصل في \code{ParseWebhook}:

\begin{english}
\begin{lstlisting}[language={[Sharp]C}]
var stripeEvent = EventUtility.ConstructEvent(
    payload,
    signatureHeader,
    webhookSecret,
    throwOnApiVersionMismatch: false);
// throws StripeException if the signature is invalid -> we reject the request

switch (stripeEvent.Data?.Object)
{
    case PaymentIntent pi: paymentIntentId = pi.Id; amountCents = pi.Amount; break;
    case Charge ch:        chargeId = ch.Id; paymentIntentId = ch.PaymentIntentId; break;
    case Payout po:        payoutId = po.Id; break;
    // Account / Dispute cases ...
}
\end{lstlisting}
\end{english}

الدالة \code{EventUtility.ConstructEvent} بتاخد الحمولة الخام والتوقيع والسرّ
المشترك (\code{webhookSecret})، وتتأكّد إنّ التوقيع مطابق. لو مش مطابق، بترمي
استثناء ونرفض الطلب فورًا. بعد التحقّق، بنستخرج الكائن المناسب (نيّة دفع، فاتورة،
صرف...) عشان المُعالِج يقدر يطابقه مع صفّ الـ \code{Payment} الصحيح.

وعشان \en{Stripe} بيعيد إرسال نفس الحدث لحدّ 3 أيّام لو ما ردّيناش بنجاح، لازم
\textbf{نتجنّب معالجته مرّتين}. ده بيحصل بإزالة التكرار عبر
\code{StripeWebhookEvent.StripeEventId}:

\begin{english}
\begin{lstlisting}[language={[Sharp]C}]
// Idempotency: Stripe retries the same event for up to 3 days.
if (await db.StripeWebhookEvents.AnyAsync(e => e.StripeEventId == request.EventId, ct))
{
    logger.LogInformation("Webhook {EventId} already processed, skipping.", request.EventId);
    return true;
}
\end{lstlisting}
\end{english}

قبل ما نعمل أي حاجة، بنسأل: «الحدث ده اتعالج قبل كده؟». لو أيوة، بنتجاهله
ونرجّع نجاح. حقل \code{StripeEventId} عليه فهرس فريد في الداتابيز، فحتى لو
اتبعت نفس الحدث مرّتين في نفس اللحظة، واحد بس بينجح والتاني بيتعامل معاه كنجاح
مكرّر.

\begin{secnote}[زاوية الأمان: ثلاث ضمانات تحمي مسار الفلوس]
\begin{itemize}
\item \textbf{مفاتيح عدم التكرار} (\en{idempotency keys}): ما فيش خصم مزدوج ولا
  صرف مزدوج مهما اتكرّرت المحاولات.
\item \textbf{القبض اليدويّ} (\en{manual capture}): الفلوس ما بتتاخدش إلا لمّا
  الخدمة تتأكّد فعلًا (المستشار قبل).
\item \textbf{\en{webhooks} موقَّعة + إزالة تكرار الأحداث}: ما فيش حدث دفع مزيّف
  أو معاد إرساله يقدر يخدعنا.
\end{itemize}
وفوق ده كلّه: \textbf{بيانات الكروت لا تلمس خوادمنا أبدًا}. المتصفّح بيكلّم
\en{Stripe} مباشرة عن طريق الـ \en{client secret}، وإحنا بس بنشوف معرّف نيّة
الدفع (\en{intent id}). فحتى لو تسرّبت قاعدة بياناتنا، ما فيش أرقام كروت تتسرّب.
\end{secnote}

\begin{qa}
\textbf{سؤال متوقّع: «فين الفلوس ما بين الحجز والجلسة؟»} \\
الإجابة: مُفوَّضة ومحجوزة على بطاقة الطالب عبر نيّة دفع بنمط \en{manual capture}.
بنقبضها بس لمّا المستشار يقبل؛ ولو رفض أو انتهت المهلة بنفكّ الحجز؛ ولو حصل
\en{no-show} أو إلغاء مؤهَّل، \code{RefundCalculatorService} بيقرّر الاسترداد.
وصافي الأرباح بيوصل للمستشار عبر صرف \en{Stripe Connect}.
\end{qa}

\subsection{المهامّ الخلفيّة (Hangfire)}

كتير من الشغل في النظام مش بيحصل لحظة طلب المستخدم، لكن بيحصل «لوحده» في الخلفية
على فترات منتظمة. إحنا بنستخدم \en{Hangfire} لجدولة المهامّ دي. كل المهامّ
متحكَّم فيها بمفتاح إعداد واحد: \code{Hangfire:Enabled}؛ لو اتقفل، ما بيتشغّلش أي
جدول (مفيد في بيئات التطوير والاختبار).

الجدول ده بيلخّص المهامّ المُجدوَلة الفعليّة في النظام (13 مهمّة):

\begin{center}
\begin{tabular}{L{4.9cm} L{2.5cm} L{5.1cm}}
\toprule
\textbf{\color{spPrimary}المهمّة} & \textbf{\color{spPrimary}الجدولة} & \textbf{\color{spPrimary}الغرض} \\
\midrule
\code{StripePayoutJob} & يوميًّا 02:00 & صرف صافي الأرباح للمستشارين دفعة واحدة \\
\code{MeetingNoShowSweepJob} & كل 15 دقيقة & رصد غياب أحد الطرفين وفتح مسار الاسترداد \\
\code{SessionExpiryJob} & كل 15 دقيقة & إنهاء الحجوزات اللي فاتت مهلة ردّ المستشار \\
\code{CompletionJob} & كل 15 دقيقة & إغلاق الجلسات المكتملة تلقائيًّا \\
\code{BookingReminderJob} & كل 15 دقيقة & تذكير الطالب والمستشار قبل الجلسة \\
\code{DeadlineReminderJob} & يوميًّا 09:00 & تنبيه الطلاب بقرب انتهاء مواعيد المنح \\
\code{NotificationDispatcherJob} & يوميًّا 10:00 & تذكير بالمسودّات غير المُرسَلة \\
\code{ScholarshipAutoCloseJob} & كل ساعة & إغلاق المنح اللي فات موعدها تلقائيًّا \\
\code{ReviewTimeoutRefundJob} & يوميًّا 05:00 & استرداد رسوم مراجعة تأخّرت عن مهلتها \\
\code{DataExportJob} & كل ساعة & تجهيز تصدير بيانات المستخدم (\en{GDPR}) \\
\code{DataDeleteJob} & يوميًّا 03:00 & حذف/تعمية بيانات المستخدم (\en{GDPR}) \\
\code{RedactionAuditSamplingJob} & شهريًّا & عيّنات تنقية للتحقّق من حجب الـ \en{PII} \\
\code{IntegrityCheckJob} & يوميًّا 04:00 & فحوص سلامة للصفوف اليتيمة والعالقة \\
\bottomrule
\end{tabular}
\end{center}

نلاحظ إنّ المهامّ دي مش بس تذكيرات؛ فيها منها اللي بيمسّ الفلوس مباشرة
(\code{StripePayoutJob}، واسترداد المراجعة)، ومنها اللي بيخدم الالتزام القانونيّ
(\en{GDPR})، ومنها اللي بيراقب صحّة النظام نفسه (\code{IntegrityCheckJob} بيعدّ
الدفعات اليتيمة والـ \en{webhooks} العالقة). كلها مبنيّة عشان تكون
\textbf{آمنة للتكرار}: لو المهمّة اتشغّلت مرّتين، مفيش ضرر — نفس مبدأ عدم التكرار
اللي شفناه في مسار الدفع.

\begin{why}
ليه نفصل الشغل ده في الخلفية بدل ما نعمله لحظة الطلب؟ عشان تجربة المستخدم تفضل
سريعة (ما نخلّيهوش يستنّى صرف فلوس أو تصدير بيانات)، وعشان الشغل الدوريّ (زيّ
إغلاق المنح المنتهية) ما ليهوش «طلب مستخدم» أصلًا يشغّله — لازم حد يفتكره ويشغّله،
والـ \en{Hangfire} هو ده «الحد».
\end{why}

\begin{summary}[الخلاصة في سطور]
\begin{itemize}
\item النظام بيستخدم \en{authorize-then-capture} (شبه \en{escrow}) على \en{Stripe}:
  فلوس الطالب \textbf{محجوزة لا مأخوذة} لحدّ ما المستشار يقبل.
\item المنفذ \code{IStripeService} له محوّلان: \code{StripeService} الحقيقي و
  \code{StubStripeService} للتطوير — نبدّل من غير ما نلمس منطق الحجز.
\item دورة الحياة: حجز + تفويض (\en{manual capture}، سنتات صحيحة، مفتاح عدم تكرار)
  ثمّ قبول = قبض ثمّ رفض/انتهاء = إفراج ثمّ \en{no-show}/إلغاء = استرداد بحساب
  \code{RefundCalculatorService} ثمّ صرف عبر \code{StripePayoutJob}.
\item تقسيم الأرباح بالسنتات: \code{AmountCents} تساوي \code{ProfitShareAmountCents}
  (المنصّة) + \code{PayeeAmountCents} (المستشار)، من قاعدة ماليّة قابلة للتهيئة،
  بفهرس فريد مُصفّى يضمن قاعدة فعّالة واحدة لكل نوع.
\item الـ \en{webhook} بيتحقّق من توقيع \en{Stripe} ويزيل التكرار عبر
  \code{StripeEventId}. ثلاث ضمانات: عدم التكرار + القبض اليدويّ + \en{webhooks}
  موقَّعة. وبيانات الكروت \textbf{ما بتلمسش خوادمنا}.
\item 13 مهمّة \en{Hangfire} خلفيّة (تذكيرات، صرف، \en{no-show}، إغلاق منح،
  \en{GDPR}، تنقية، فحوص سلامة) متحكَّم فيها بـ \code{Hangfire:Enabled}.
\end{itemize}
\end{summary}

\clearpage
% ch09-analytics-summary
\section{التحليلات والخلاصة وأسئلة المناقشة}

وصلنا للفصل الأخير. في الفصول اللي فاتت مشينا خطوة بخطوة: بنينا الطبقات
(\en{Clean Architecture})، وفهمنا إزاي كل طلب بيتحوّل لأمر أو استعلام
(\en{CQRS})، وأمّنّا الدخول، وركّبنا مزايا الـ\en{AI}، وشبكنا الفيديو والدفع
والمهام الخلفية. في الفصل ده هنعمل حاجتين: الأولى نشرح \textbf{خطّ التحليلات}
(إزاي بيانات النظام الحيّة بتتحوّل لتقارير ولوحات تحكّم)، والثانية \textbf{نلمّ
الكتاب كلّه} في مكان واحد، وبعدها نجهّز \textbf{أسئلة لجنة المناقشة} بإجاباتها
عشان تبقى جاهز في يوم الدفاع.

هدفنا هنا مش نضيف تعقيد جديد، لكن نربط الصورة. فخُد نفَس، وتعالَ نمشي واحدة واحدة.

% ═══════════════════════════════════════════════════════════════════════════
\subsection{خطّ التحليلات: من قاعدة البيانات الحيّة إلى لوحات \en{Power BI}}
% ═══════════════════════════════════════════════════════════════════════════

النظام بيولّد بيانات كتير كل يوم: طلبات منح، حجوزات استشاريين، مدفوعات، تقييمات.
قاعدة البيانات التشغيلية (\en{SQL Server}) شغلها الأساسي إنها تخدم المستخدمين بسرعة،
مش إنها تعمل تقارير تحليلية ثقيلة. عشان كده بنفصل عالَم \textbf{التشغيل} عن عالَم
\textbf{التحليل}، وده اللي بنسمّيه \textbf{خطّ التحليلات} (\en{Analytics Data
Pipeline})، وهو مسؤولية الوحدات \en{PB-015} لحد \en{PB-018}.

\begin{idea}[الفكرة: يعني إيه \en{Data Warehouse}؟]
مخزن البيانات (\en{data warehouse}) هو قاعدة بيانات \textbf{منفصلة} مبنيّة مخصوص
للتحليل والتقارير، مش للمعاملات اللحظية. بننقل نسخة من بيانات النظام إليه، ونعيد
ترتيبها بشكل يخلّي الأسئلة التحليلية (مثال: «كام طلب منحة اتقدّم الشهر ده لكل
دولة؟») تتجاوب بسرعة، من غير ما نضغط على قاعدة البيانات الحيّة اللي بيشتغل عليها
الطلاب في نفس اللحظة.
\end{idea}

\subsubsection{المراحل بالترتيب}

الخطّ عندنا بيمشي في خمس محطات أساسية، من الشمال لليمين:

\begin{idea}[الفكرة: يعني إيه \en{CDC}؟]
\en{Change Data Capture} (اختصارًا \en{CDC}) معناه «التقاط التغييرات فقط». بدل ما
ننسخ الجدول كامل كل ليلة (تبذير)، الـ\en{CDC} بيراقب قاعدة البيانات ويلتقط
\textbf{الصفوف اللي اتغيّرت بس} — أُضيفت أو اتعدّلت أو اتمسحت — وينقلها. ده بيخلّي
النقل خفيف وسريع.
\end{idea}

\begin{idea}[الفكرة: يعني إيه \en{Star Schema}؟]
المخطّط النجمي (\en{star schema}) هو طريقة لترتيب جداول التحليل. في النص جدول
واحد كبير اسمه \textbf{جدول الحقائق} (\en{fact table}) بيحتوي الأحداث القابلة
للقياس (مثال: كل عملية دفع بمبلغها وتاريخها)، وحواليه \textbf{جداول الأبعاد}
(\en{dimension tables}) اللي بتوصف السياق (الطالب، الدولة، نوع المنحة، الزمن).
الشكل بيطلع زي النجمة: حقيقة في النص وأبعاد على الأطراف. الترتيب ده بيخلّي التقارير
سهلة وسريعة جدًّا.
\end{idea}

نبدأ من \en{SQL Server} (المصدر التشغيلي). الـ\en{Azure Data Factory} بيسحب
التغييرات بتقنية \en{CDC} وينزّلها في طبقة اسمها \textbf{\en{Bronze}} — دي نسخة
خام طبق الأصل من البيانات كما هي. بعدها أداة \en{dbt} بتنضّف وتحوّل: من
\en{Bronze} لطبقة \textbf{\en{Silver}} (بيانات متنظّفة وموحّدة)، ومنها لطبقة
\textbf{\en{Gold}} (الـ\en{star schema} النهائي الجاهز للتحليل). بعد كده
\en{Synapse Serverless} بيعرض الـ\en{Gold} على هيئة \en{views} يقدر أي حد
يستعلم عنها بلغة \en{SQL} من غير سيرفر دائم شغّال. وأخيرًا \en{Power BI} بيرسم
اللوحات التفاعلية اللي المشرف بيبصّ عليها.

\begin{center}
\begin{tabular}{L{3.1cm} L{4.0cm} L{5.2cm}}
\toprule
\textbf{\color{spPrimary}المرحلة} & \textbf{\color{spPrimary}الأداة} & \textbf{\color{spPrimary}بتعمل إيه} \\
\midrule
المصدر التشغيلي & \en{SQL Server 2022} & قاعدة البيانات الحيّة اللي فيها كل المعاملات \\
الالتقاط والنقل & \en{Azure Data Factory} (\en{CDC}) & تلتقط الصفوف المتغيّرة فقط وتنقلها \\
طبقة \en{Bronze} & \en{Data Lake} (خام) & نسخة طبق الأصل من البيانات كما هي \\
التحويل & \en{dbt} (\en{Silver} ثم \en{Gold}) & تنضّف وتوحّد ثم تبني الـ\en{star schema} \\
طبقة الاستعلام & \en{Synapse Serverless} & تعرض الـ\en{Gold} كـ\en{views} بلا خادم دائم \\
العرض & \en{Power BI} & لوحات تحكّم تفاعلية للمشرفين \\
الزمن الحقيقي & \en{Event Hubs} ثم \en{Stream Analytics} & يبثّ الأحداث لحظيًّا ويكشف الشذوذ \\
التنبيه & \en{PagerDuty} & ينبّه الفريق فورًا عند اكتشاف شذوذ \\
\bottomrule
\end{tabular}
\end{center}

\subsubsection{الفرع اللحظي: كشف الشذوذ في الوقت الحقيقي}

الخطّ اللي فوق ده \textbf{دفعيّ} (\en{batch}): بيشتغل مرّة كل ليلة. بس فيه أحداث
لازم ننتبه لها \textbf{في اللحظة}، مش بكرة الصبح — زي طفرة مدفوعات مفاجئة ممكن
تكون احتيال. عشان كده فيه \textbf{فرع لحظي} موازي: النظام بيبثّ ثلاثة أحداث
عالية القيمة إلى \en{Azure Event Hubs} أول ما تحصل، و\en{Azure Stream Analytics}
بيقرأها ويكشف الشذوذ، ولو لقى حاجة غريبة بيبعت تنبيه لـ\en{PagerDuty} وكمان
بيغذّي بلاطة حيّة في \en{Power BI}.

\begin{keyp}[نقطة مهمة: التحليلات لا تعطّل التشغيل أبدًا]
مبدأ التصميم هنا إن خطّ التحليلات \textbf{انتهازيّ} (\en{opportunistic}): لو
وقع، النظام الأساسي بيفضل شغّال عادي. الطالب اللي بيدفع أو بيقدّم طلب مش هيتأثّر
لو الـ\en{Event Hubs} كان واقع. الحدث الضايع بيتعوّض بالحدث اللي بعده، مش بإيقاف
المستخدم.
\end{keyp}

تعالَ نشوف ده في الكود الحقيقي. الأحداث الثلاثة اللي بنبثّها هي:
\en{ApplicationSubmitted} (تقديم طلب)، و\en{PaymentCaptured} (تحصيل دفعة)،
و\en{BookingCompleted} (اكتمال جلسة). كل واحد ليه مُعالِج
(\en{MediatR notification handler}) رفيع بيبني ظرفًا صغيرًا ويسلّمه للناشر.

\textbf{هنشوف إيه:} مُعالِج حدث «تحصيل دفعة» وهو بيوجّه الحدث لخطّ التحليلات
(من ملف \code{DomainEventPublishingHandlers.cs}).

\begin{english}
\begin{lstlisting}[language={[Sharp]C}]
public sealed class PaymentCapturedStreamHandler(
    IEventPublisher publisher,
    ILogger<PaymentCapturedStreamHandler> logger)
    : INotificationHandler<PaymentCapturedEvent>
{
    public async Task Handle(PaymentCapturedEvent n, CancellationToken ct)
    {
        // Forward the domain event to the streaming pipeline
        await publisher.PublishAsync(
            "PaymentCaptured",
            new
            {
                paymentId   = n.PaymentId,
                type        = n.Type.ToString(),
                amountCents = n.AmountCents,
                payerUserId = n.PayerUserId,
                payeeUserId = n.PayeeUserId,
            },
            ct).ConfigureAwait(false);
    }
}
\end{lstlisting}
\end{english}

نلاحظ حاجات بسيطة: الصنف بيطبّق \code{INotificationHandler<PaymentCapturedEvent>}،
يعني \en{MediatR} بينده عليه أوتوماتيك أول ما الحدث يتنشر بعد نجاح المعاملة. جوّه
\code{Handle} مفيش أي منطق أعمال — بس نداء واحد لـ\code{publisher.PublishAsync}
بنبعت له اسم الحدث \code{"PaymentCaptured"} وكائن صغير فيه الأرقام المهمّة
(المبلغ بالسنت، ونوع الدفعة، والدافع، والمستفيد). المُعالِج مقصود إنه رفيع كده،
وكل التعامل مع الأخطاء متروك للناشر.

\textbf{هنشوف إيه:} الناشر الحقيقي (\en{Azure Event Hubs}) وإزاي بيبلع الأخطاء
عشان ميعطّلش المعاملة (مقتطف مختصر من \code{EventHubPublisher.cs}).

\begin{english}
\begin{lstlisting}[language={[Sharp]C}]
public async Task PublishAsync(string eventType, object payload, CancellationToken ct = default)
{
    try
    {
        var envelope = new
        {
            eventType,
            occurredAt = DateTimeOffset.UtcNow,
            source = "ScholarPath.API",
            data = payload,
        };
        var json = JsonSerializer.Serialize(envelope);
        var data = new EventData(Encoding.UTF8.GetBytes(json));

        await using var producer = new EventHubProducerClient(
            _opts.ConnectionString, _opts.HubName);
        using var batch = await producer.CreateBatchAsync(ct);
        batch.TryAdd(data);
        await producer.SendAsync(batch, ct);        // send to Azure Event Hubs
    }
    catch (Exception ex)
    {
        // A hub outage must NOT roll back the business transaction
        logger.LogError(ex, "EventHub: failed to publish {EventType}; continuing.", eventType);
    }
}
\end{lstlisting}
\end{english}

الفكرة المحوريّة هنا سطرين: أول حاجة بنغلّف الحمولة في \textbf{ظرف} (\en{envelope})
فيه نوع الحدث ووقت حدوثه ومصدره، ونحوّله \en{JSON} ونبعته لِـ\en{Event Hubs}.
تاني حاجة — وهي الأهم — الـ\code{catch} بيمسك أي استثناء \textbf{ويسجّله بس من
غير ما يعيد رميه}. ده اللي بيحقّق مبدأ «التحليلات ما تعطّلش التشغيل»: لو الـ\en{hub}
واقع، الدالة بتكمل بهدوء والمعاملة الأصلية بتنجح عادي.

\textbf{هنشوف إيه:} إزاي النظام بيختار بين الناشر الحقيقي والبديل الآمن حسب
الإعداد (مقتطف من \code{DependencyInjection.cs}).

\begin{english}
\begin{lstlisting}[language={[Sharp]C}]
// Real publisher only when a connection string is configured; otherwise a stub
var ehConnString = config.GetValue<string>($"{EventHubOptions.SectionName}:ConnectionString");
if (!string.IsNullOrWhiteSpace(ehConnString))
    services.AddSingleton<IEventPublisher, EventHubPublisher>();
else
    services.AddSingleton<IEventPublisher, StubEventPublisher>();
\end{lstlisting}
\end{english}

ده نمط بنكرّره في كل النظام: لو الإعداد موجود (سلسلة اتصال حقيقية) نركّب الخدمة
الحقيقية \code{EventHubPublisher}؛ ولو مش موجود (زي بيئة التطوير على جهاز مطوّر)
نركّب بديلًا آمنًا \code{StubEventPublisher} بيسجّل الأحداث ويكمّل من غير ما يرمي
خطأ. النتيجة: المطوّر يقدر يشغّل النظام كامل من غير أي حساب \en{Azure}.

\begin{secnote}
كل الأحداث المبثوثة \textbf{ما فيهاش} بيانات حسّاسة زايدة: بنبعت مُعرِّفات
(\en{IDs}) وأرقام مجمّعة زي المبلغ والنوع، مش أسماء ولا إيميلات. كده لو اتّسرب
الـ\en{stream} ما يكشفش هويّات المستخدمين — وده يتماشى مع نهج تقليل البيانات
(\en{data minimisation}) اللي شفناه في فصل الأمان.
\end{secnote}

% ═══════════════════════════════════════════════════════════════════════════
\subsection{ملخّص الكتاب كلّه: ماذا تعلّمنا فصلًا فصلًا}
% ═══════════════════════════════════════════════════════════════════════════

خلّينا نلمّ الرحلة. كل فقرة تحت بتلخّص فصلًا كامل في سطرين:

\textbf{الفصل التمهيدي —} عرّفنا \en{ScholarPath}: منصّة منح دراسية ثنائية اللغة
(عربي وإنجليزي بدعم \en{RTL} كامل)، بتخدم الطلاب والاستشاريين والجهات المانحة
والمشرفين، ومقسّمة إلى \textbf{ثماني عشرة وحدة} (\en{PB-001..PB-018}) بتغطّي دورة
الحياة كاملة من إنشاء الحساب لحد تخزين البيانات للتحليل.

\textbf{فصل \en{Clean Architecture} —} اتعلّمنا إن الكود مقسّم لطبقات متداخلة،
والاعتماديّات بتشاور \textbf{للداخل بس}: طبقة \en{Domain} (قلب النظام) ما تعرفش
حاجة عن قاعدة البيانات ولا الإطار ولا الواجهة. ده بيخلّي منطق الأعمال قابل
للاختبار والنقل، وبيسمح باستبدال أي خدمة خارجية ببديل من غير ما نلمس القلب.

\textbf{فصل \en{CQRS} والسلامة —} كل عمليّة بقت إمّا \textbf{أمرًا} (\en{Command})
بيغيّر الحالة، أو \textbf{استعلامًا} (\en{Query}) بيقرأ بس. \en{MediatR} بيوصّل
كل طلب لِمُعالِجه، والمتحكّمات (\en{Controllers}) بقت رفيعة جدًّا. وسلوكيّات
الأنبوب (\en{pipeline behaviours}) زي التحقّق والتدقيق بتلفّ كل مُعالِج تلقائيًّا.

\textbf{فصل المصادقة —} الدخول بيعتمد \en{ASP.NET Core Identity} مع رموز
\en{JWT} موقّعة بخوارزمية \en{RS256} غير المتماثلة، ومدّة الرمز خمس عشرة دقيقة،
مع رمز تجديد بيتدوّر كل استعمال. وفيه دخول موحّد (\en{Google} و\en{Microsoft})
ومصادقة ثنائية (\en{TOTP}).

\textbf{فصل الأمان —} اتفرّجنا على التدقيق الآليّ: أي أمر متعلّم بـ\code{[Auditable]}
بيتسجّل تلقائيًّا في \code{AuditLog}، وعلى تنقية البيانات الحسّاسة
(\en{PII redaction})، وتحديد المعدّل (\en{rate-limiting})، والتوقيع على أحداث
\en{Stripe}.

\textbf{فصلا الـ\en{AI} —} اتعلّمنا مفاهيم \en{RAG} (استرجاع ثم توليد)
والـ\en{embeddings} والبحث بجيب التمام (\en{cosine similarity})، وبعدها بنينا
المزايا فعليًّا: شات بوت يستند لقاعدة معرفة حيّة، وتوصيات بتحسب درجة حتميّة
(\en{deterministic scoring})، وفحص أهليّة، وكل ده بسقوط آمن ثلاثي المستويات.

\textbf{فصل الاجتماعات والزمن الحقيقي —} شفنا مكالمة الفيديو الكاملة داخل المتصفّح
عبر \en{Azure Communication Services}، من إنشاء الغرفة عند تأكيد الحجز، لإصدار رمز
قصير العمر، للانضمام، للتسجيل على \en{Blob Storage}. وشفنا الـ\en{SignalR} بيدفع
التحديثات لحظيًّا (رسائل، إشعارات، تفاعلات).

\textbf{فصل المدفوعات والمهام —} اتعلّمنا نيّات الدفع (\en{PaymentIntents}) بالتحصيل
اليدوي، وحصّة المنصّة، والمهام الخلفية عبر \en{Hangfire}: تذكير المواعيد، وتسوية
المدفوعات، وكشف الغياب، ومزامنة التحليلات الليلية.

\textbf{هذا الفصل —} لمّينا خطّ التحليلات، وبنلخّص الكتاب، وبنجهّز أسئلة المناقشة.

\begin{summary}[ما الذي يميّز \en{ScholarPath}؟]
لو حد سألك «إيه اللي يخلّي المشروع ده مميّز؟» — دي التلات نقط اللي تقولها بثقة:

\medskip
\textbf{١) تطوير مدفوع بالمواصفات (\en{spec-driven development}).} كل طلب دمج
(\en{PR}) بيشير لمتطلّب وظيفي (\en{FR-xxx}) أو قصّة مستخدم (\en{US-xxx}). كل كيان
له هجرة قاعدة بيانات (\en{migration})، وكل أمر مؤثّر متعلّم بـ\code{[Auditable]}.
المواصفات مش ورق على جنب — هي مصدر الحقيقة اللي الكود بيتبعه.

\medskip
\textbf{٢) ذكاء اصطناعيّ إنتاجيّ بتكلفة حديّة تقترب من الصفر ومع سقوط آمن.}
شات بوت \en{RAG} + توصيات حتميّة + فحص أهليّة، وكلّه وراه \textbf{سقوط ثلاثي
المستويات} بيخلّي النظام يتدهور بلُطف مهما حصل عطل في \en{Azure} — المستخدم دايمًا
بياخد ردّ. وسقف تكلفة يوميّ بيمنع أي فاتورة تطير.

\medskip
\textbf{٣) بنية فيديو متكاملة من الطرف للطرف.} إنشاء غرفة \en{ACS}، وإصدار رمز
قصير العمر، و\en{SDK} في الواجهة، وتسجيل من جهة الخادم على \en{Blob Storage}،
وتحميل \en{MP4} — كلّه مربوط في اثنتي عشرة خطوة مع تتبّع الحضور وكشف الغياب.
\end{summary}

% ═══════════════════════════════════════════════════════════════════════════
\subsection{أسئلة المناقشة المتوقّعة وإجاباتها}
% ═══════════════════════════════════════════════════════════════════════════

الجزء ده هو اللي هيفيدك يوم الدفاع. جمّعنا أكتر الأسئلة المتوقّعة من لجنة المناقشة،
وكل سؤال في صندوقه بإجابة جاهزة ومختصرة. اقراها بصوت عالٍ مرّة أو اتنين قبل الدفاع.

\begin{qa}[ليه \en{Clean Architecture} بدل معمارية طبقية بسيطة؟]
لأنها بتفرض \textbf{عكس الاعتماديّة} (\en{dependency inversion}): طبقة الـ\en{Domain}
ما تعرفش حاجة عن \en{EF Core} ولا \en{Azure}. يعني نقدر نستبدل قاعدة البيانات أو
أي خدمة من غير ما نلمس منطق الأعمال. وأثبتنا ده عمليًّا لمّا حقنّا بدائل
(\en{Stub}) لكل الخدمات الخارجية (\en{ACS}، \en{Stripe}، \en{Azure}) في الاختبارات
وشغّلنا النظام كامل من غيرها.
\end{qa}

\begin{qa}[ليه \en{CQRS} في مشروع تخرّج؟ مش ده أكتر من اللازم؟]
المشروع فيه مئات المتطلّبات الوظيفية موزّعة على خمسة مطوّرين. الـ\en{CQRS} فرض
\textbf{حدود ملكيّة واضحة}: كل مُعالِج ملف \en{C\#} واحد يملكه شخص واحد. من غيره
كان صنف خدمة مشترك هيسبّب تعارضات دمج (\en{merge conflicts}) في كل سبرنت. وكمان
\en{MediatR} إدّانا التدقيق (\en{audit}) كسلوك عرضيّ \textbf{ببلاش}.
\end{qa}

\begin{qa}[ليه ما استخدمتوش قاعدة بيانات متجهات (\en{vector database})؟]
عند حجم يقارب \textbf{692 مستندًا}، المسح الكامل في الذاكرة بجيب التمام
(\en{cosine scan}) بيخلص في أقلّ من عشرة أجزاء من الألف من الثانية. قاعدة المتجهات
هتضيف قفزة شبكة وتكلفة تشغيل من غير أي فايدة في الدقّة عند الحجم ده. ومسار الترقية
موجود: نستبدل \code{KnowledgeRetriever} بتطبيق يعتمد \en{Qdrant} أو \en{pgvector}
لو كبرنا.
\end{qa}

\begin{qa}[هل الـ\en{AI} بيهلوس (\en{hallucinate})؟]
موجّه النظام (\en{system prompt}) \textbf{بيمنع الاختلاق}، والإجابة لازم تستشهد
بكتلة السياق المرقّمة. سياق الـ\en{RAG} مستند لقاعدة البيانات الحيّة؛ ولو مفيش
مستند سجّل درجة فوق العتبة (\en{0.6})، البديل المحلّي بيدّي إجابة عامّة بدل ما
يخمّن. يعني النموذج بيتكلّم من مصدر، مش من خياله.
\end{qa}

\begin{qa}[ليه \en{Azure ACS} للفيديو ومش \en{Jitsi}؟]
\en{Jitsi} بيحتاج خادمًا ذاتيّ الاستضافة بعبء تشغيليّ، ومفيهوش \en{API} تسجيل
مدمج. \en{ACS} بيوفّر خطّ تسجيل مُدار ومتوافق مع خصوصيّة البيانات
(\en{PII-compliant})، وبيتكامل بشكل طبيعيّ مع باقي منظومة \en{Azure} عندنا.
\end{qa}

\begin{qa}[إزاي بتشتغل فوترة الاستشاريين؟]
عبر \en{Stripe PaymentIntents} بتحصيل يدويّ (\en{manual capture}). كارت الطالب
\textbf{يُصرَّح عليه} وقت الحجز، والتحصيل الفعليّ ما بيحصلش إلا لمّا الاستشاري يقبل
الجلسة. وحصّة المنصّة بتتحسب على مستوى المُعالِج عبر \en{Stripe Transfer}.
\end{qa}

\begin{qa}[لو \en{Azure OpenAI} سقط، يحصل إيه؟]
سقوط آمن على \textbf{ثلاثة مستويات}: (١) الإعداد ناقص فبنروح للمحلّي؛ (٢) خطأ
\en{HTTP} في نداء التضمين (\en{embedding}) فبنروح للمحلّي؛ (٣) نداء التوليد فشل
فبنروح للمحلّي. الخدمة المحلّية \code{LocalAiService.AskAsync} بتشغّل نفس استرجاع
الـ\en{RAG} بدون التوليد، وبعدها بترجع لموجّه بالكلمات المفتاحيّة. النتيجة:
\textbf{المستخدم دايمًا بياخد ردّ}.
\end{qa}

\begin{note}
لاحظ إن كل الإجابات دي بتربط بين \textbf{قرار تصميميّ} و\textbf{سبب عمليّ ملموس}.
اللجنة مش بتدوّر على «اخترنا الأحدث»، بتدوّر على «اخترنا الأنسب وليه». خلّي كل
إجابة تبدأ بالسبب.
\end{note}

% ═══════════════════════════════════════════════════════════════════════════
\subsection{الخاتمة}
% ═══════════════════════════════════════════════════════════════════════════

كده نكون قفلنا الدائرة. النظام بيبدأ من طالب بيعمل حساب، بيمرّ باكتشاف المنح
والتقديم والحجز والدفع والاستشارة بالفيديو، وبينتهي ببياناته وهي بتتحوّل لمعرفة في
لوحات التحليل. كل حلقة في السلسلة دي مبنيّة على نفس المبادئ: طبقات نظيفة، أوامر
واستعلامات واضحة، تدقيق آليّ، وسقوط آمن عند أي عطل. ده مش بس نظام شغّال، ده نظام
\textbf{مفهوم} — تقدر تشرح أي جزء منه، وده أهمّ حاجة في يوم الدفاع.

\begin{summary}[الخلاصة في سطور]
\begin{itemize}
  \item خطّ التحليلات بيفصل عالَم \textbf{التشغيل} عن عالَم \textbf{التحليل}:
    \en{SQL Server} ثم \en{ADF} (\en{CDC}, طبقة \en{Bronze}) ثم \en{dbt}
    (\en{Silver} ثم \en{Gold} \en{star schema}) ثم \en{Synapse} ثم \en{Power BI}.
  \item فيه \textbf{فرع لحظيّ} موازٍ: أحداث تُبثّ إلى \en{Event Hubs}، و\en{Stream
    Analytics} يكشف الشذوذ ويُنبّه \en{PagerDuty}.
  \item مبدأ حاكم: التحليلات \textbf{ما تعطّلش التشغيل} — الناشر بيبلع الأخطاء
    ولا يعيد رميها، والبديل الآمن (\en{Stub}) بيشتغل بدون \en{Azure}.
  \item الأحداث المبثوثة بتحمل \textbf{مُعرِّفات وأرقام مجمّعة فقط}، لا بيانات
    شخصيّة حسّاسة.
  \item ما يميّز \en{ScholarPath}: تطوير مدفوع بالمواصفات، \en{AI} إنتاجيّ بتكلفة
    حديّة تقترب من الصفر مع سقوط آمن، وبنية فيديو متكاملة من الطرف للطرف.
  \item في الدفاع: ابدأ كل إجابة بالسبب العمليّ، مش بالتقنية.
\end{itemize}
\end{summary}


% ── الجزء الثاني: تعمّق تقنيّ ───────────────────────────────────────
\clearpage
\begin{mdframed}[style=spbase,linewidth=0pt,roundcorner=4pt,
   backgroundcolor=spAccent,innertopmargin=14pt,innerbottommargin=14pt]
{\color{white}\Huge\bfseries الجزء الثاني}\\[4pt]
{\color{white}\LARGE تعمّق تقنيّ ورحلات كاملة}\\[6pt]
{\color{white}\normalsize للي عايز يغوص أعمق: إزاي الطلب بيمشي من الواجهة للقاعدة،
مستويات التحقّق والوسطاء، \en{WebSocket}\slash\en{SignalR} والمهامّ الخلفيّة،
المجتمع والحظر، الأحداث الدوميّة والتدقيق، الهُوية والتوكنات و\en{SSO}،
الاختبارات، وأخيرًا رحلة طالب كاملة عبر النظام كلّه.}
\end{mdframed}
\clearpage
% ch20-request-lifecycle
\section{رحلة الطلب الكاملة: من زرّ في الواجهة إلى صفّ في القاعدة}

في الجزء الأوّل عرفنا الطبقات الأربع، ونمط \en{CQRS} مع \en{MediatR}،
وفكرة \en{ports \& adapters}، وسلاسل المعالجة \en{(pipeline behaviors)}.
لكنّنا رأينا كلّ قطعة على حِدَة. في هذا الفصل نُركّب القطع معًا ونتتبّع
\emph{طلبًا واحدًا حقيقيًّا} من أوّله لآخره: مستخدم يفتح صفحة الإشعارات، تُرسِل
الواجهة نداءً للـ \en{API}، يعبُر النداء عدّة محطّات، تُترجَم قراءة البيانات إلى
جملة \en{SQL} مُبارَمة، ثمّ يعود الردّ ليُحدِّث الشاشة. هدفنا أن تعرف — بدقّة —
مَن يفعل ماذا في كلّ محطّة، ولماذا وُضِع بهذا الترتيب.

سنمشي «واحدة واحدة»: نبدأ من المتصفّح، ندخل الخادم، نمرّ بأنبوب الـ \en{middleware}،
ندخل \en{MediatR} وسلاسله، نصل للمعالِج ثمّ \en{EF Core}، ونرى كيف يتحوّل الكيان إلى
\en{DTO} فـ \en{JSON}. وفي الطريق نُجيب على سؤال يشغل كثيرين: لماذا يوجد أحيانًا
\en{Service} في طبقة الـ \en{Infrastructure} بدل وضع كلّ شيء داخل المعالِج؟ ونختم
بتمهيد للفرق بين مسار \en{HTTP} العاديّ ومسار الزمن الحقيقيّ \en{(SignalR)}.

% ====================================================================
\subsection{نظرة طائر: خريطة الرحلة كاملةً}

قبل التفاصيل، لِنَرَ الصورة الكبيرة. أيّ طلب قراءة يمرّ بسلسلة محطّات ثابتة، ثمّ
يرجع الردّ بنفس المسار عكسيًّا. الرسم التالي يلخّص الرحلة: الأسهم المتّصلة هي رحلة
\emph{الذهاب} (الطلب)، والسهم المتقطّع هو رحلة \emph{العودة} (الردّ).

\begin{center}\begin{LTR}
\begin{tikzpicture}[node distance=0.55cm and 0.7cm]
  \node[actor,text width=1.7cm] (br) {Browser\\(React)};
  \node[sbox,right=of br,text width=1.9cm] (mw) {Middleware\\pipeline};
  \node[pbox,right=of mw,text width=1.9cm] (ctl) {Controller};
  \node[gbox,right=of ctl,text width=1.7cm] (med) {MediatR};

  \node[tbox,below=1.15cm of med,text width=1.9cm] (beh) {Behaviors};
  \node[pbox,left=of beh,text width=1.7cm] (hnd) {Handler};
  \node[sbox,left=of hnd,text width=1.7cm] (ef) {EF Core};
  \node[dbbox,left=of ef,text width=1.7cm] (sql) {SQL\\Server};

  \draw[arr] (br) -- (mw);
  \draw[arr] (mw) -- (ctl);
  \draw[arr] (ctl) -- (med);
  \draw[arr] (med) -- (beh);
  \draw[arr] (beh) -- (hnd);
  \draw[arr] (hnd) -- (ef);
  \draw[arr] (ef) -- (sql);

  \draw[arrd] (sql.south) to[out=-90,in=-90]
    node[lbl,below] {response: entity $\rightarrow$ DTO $\rightarrow$ JSON} (br.south);
\end{tikzpicture}
\end{LTR}\end{center}

كلّ صندوق محطّة سنشرحها بالترتيب. الفكرة المحوريّة أنّ هذا المسار \emph{واحد ثابت}
لكلّ الطلبات تقريبًا؛ ما يتغيّر هو المعالِج في النهاية فقط.

\begin{spotlight}[ركّز هنا لإبهار اللجنة: مسار موحّد لكلّ طلب]
كلّ طلب في \en{ScholarPath} — قبول حجز، جلب إشعارات، نشر منشور — يمشي في
\emph{نفس} الأنبوب: \en{middleware} ثمّ \en{controller} رفيع ثمّ \en{MediatR}
ثمّ سلاسل المعالجة ثمّ معالِج واحد. إضافة ميزة جديدة تعني إضافة
\en{command/query} ومعالِجًا لها فقط — \textbf{بلا لمس البنية ولا الأنبوب}.
هذا الانتظام هو ما يجعل النظام قابلًا للنموّ دون فوضى، وهو أوّل ما تُبرِزه للجنة:
«عندنا شكل واحد لكلّ عملية، ونضيف الميزات كوحدات صغيرة مستقلّة».
\end{spotlight}

% ====================================================================
\subsection{المحطّة الأولى: الواجهة \en{(React)} تبني الطلب}

الواجهة تطبيق صفحة-واحدة \en{(SPA — Single Page Application)}. معناها أنّ المتصفّح
يُحمّل تطبيق \en{React} مرّة واحدة، وبعدها لا يُعيد تحميل الصفحة كاملةً؛ بل عند الحاجة
لبيانات جديدة \emph{يُنادي} الخادم عبر الشبكة ويستقبل \en{JSON}، ثمّ يُحدِّث جزءًا من
الشاشة. إذًا الواجهة والخادم برنامجان منفصلان يتكلّمان بلغة \en{HTTP}.

مَن يُنظّم هذه النداءات في الواجهة؟ مكتبة \en{TanStack Query}. هي طبقة فوق نداءات
الشبكة تتكفّل بثلاثة أشياء: تُطلِق النداء، وتُخزّن نتيجته في ذاكرة مؤقّتة
\en{(cache)}، وتعرف متى تُعيد الجلب. الصفحة تقول «أريد الإشعارات» فتُعطيها المكتبة
البيانات من الـ \en{cache} فورًا إن وُجدت، أو تُطلق نداءً جديدًا.

النداء الفعليّ يمرّ عبر عميل \en{axios} مشترك اسمه \code{apiClient}. أهمّ ما فيه شيءٌ
اسمه \en{interceptor}: قطعة تلتفّ حول كلّ نداء صادر وتتدخّل قبله. عندنا
\en{interceptor} على الطلب يُلصِق رأس \en{Authorization} تلقائيًّا بكلّ نداء، فلا
تحتاج كلّ صفحة أن تتذكّر ذلك.

هنشوف إيه: الـ \en{request interceptor} الحقيقيّ في \code{client.ts} — يقرأ الرمز من
مخزن الحالة ويُلصِقه كرأس \en{Bearer}.

\begin{english}
\begin{lstlisting}[language=Java]
apiClient.interceptors.request.use((config) => {
  const tokens = useAuthStore.getState().tokens;
  if (tokens?.accessToken) {
    config.headers["Authorization"] = `Bearer ${tokens.accessToken}`;
  }
  // send the active UI language so the server localises the response
  config.headers["Accept-Language"] = i18n.language || "ar";
  return config;
});
\end{lstlisting}
\end{english}

شرح الأسطر المهمّة:
\begin{itemize}
\item \code{interceptors.request.use(...)}: سجِّل دالة تُشغَّل قبل \emph{كلّ} نداء صادر.
\item نقرأ \code{tokens.accessToken} من مخزن الحالة \en{(auth store)} — رمز الدخول الحاليّ.
\item نُلصِقه في رأس \code{Authorization} بصيغة \code{Bearer <token>}؛ هذا ما يُثبِت
  للخادم هُويّة المستخدم.
\item نُرسِل أيضًا \code{Accept-Language} لتأتي الردود بلغة الواجهة (عربي أو إنجليزي).
\end{itemize}

\begin{keyp}[نقطة مهمة: الرمز يُلصَق مرّة واحدة في مكان واحد]
لولا الـ \en{interceptor} لاضطرّت كلّ صفحة أن تكتب رأس \en{Authorization} يدويًّا،
وهو تكرار خطير: نسيان واحد يعني نداءً بلا هُويّة. المركزة في مكان واحد تجعل الأمن
افتراضيًّا لا اختياريًّا.
\end{keyp}

% ────────────────────────────────────────────────────────────────
\subsubsection{ماذا لو انتهت صلاحيّة الرمز؟ الردّ \en{401} والتجديد}

رمز الدخول قصير العمر. حين ينتهي، يردّ الخادم بحالة \code{401 Unauthorized}. هنا يعمل
\en{interceptor} ثانٍ — على \emph{الردّ} هذه المرّة: إن رأى \code{401} يُحاول تجديد
الرمز عبر \code{/api/auth/refresh}، ثمّ \emph{يُعيد} النداء الأصليّ كأنّ شيئًا لم يكن.
المستخدم لا يشعر بشيء.

هنشوف إيه: قلب الـ \en{response interceptor} — التقاط \code{401}، تجديد الرمز مرّة
واحدة، ثمّ إعادة تشغيل الطلب الأصليّ.

\begin{english}
\begin{lstlisting}[language=Java]
apiClient.interceptors.response.use(
  (response) => response,
  async (error) => {
    const status = error.response?.status;
    const original = error.config;
    if (status === 401 && original && !original._retry) {
      original._retry = true;                       // guard: try refresh once
      const pending = refreshing ?? (refreshing = refreshTokens());
      const newToken = await pending;               // share ONE in-flight refresh
      if (newToken) {
        original.headers["Authorization"] = `Bearer ${newToken}`;
        return apiClient(original);                 // replay the original request
      }
      useAuthStore.getState().clear();              // refresh failed -> log out
    }
    throw new ApiError(/* ... */);
  },
);
\end{lstlisting}
\end{english}

الفكرة سطرًا بسطر:
\begin{itemize}
\item الدالة الأولى تمرّر أيّ ردّ ناجح كما هو؛ الثانية تُعالج الأخطاء.
\item \code{status === 401} مع علامة \code{\_retry}: نتأكّد أن نُجرِّب التجديد
  \emph{مرّة واحدة} فقط لا حلقة لا نهائيّة.
\item \code{refreshing ?? (...)}: لو وصلت عدّة نداءات \code{401} معًا، نتشارك عمليّة
  تجديد \emph{واحدة} بدل إطلاق عدّة عمليّات تتصادم.
\item عند نجاح التجديد نُحدِّث الرأس ونُعيد \code{apiClient(original)} — أي نُعيد
  إرسال النداء الأصليّ بالرمز الجديد.
\item عند الفشل \code{clear()} — تسجيل خروج نظيف.
\end{itemize}

\begin{note}
هذا التجديد الصامت هو سبب أنّ المستخدم يبقى داخل النظام ساعاتٍ دون أن يُطلَب منه
تسجيل الدخول ثانيةً، رغم أنّ رمز الدخول نفسه ينتهي كلّ بضع دقائق. تفاصيل الهُويّة
والرموز شرحناها في فصل التوثيق؛ هنا يهمّنا فقط كيف تتصرّف الواجهة تلقائيًّا.
\end{note}

% ====================================================================
\subsection{المحطّة الثانية: الدخول للخادم والتوجيه \en{(Routing)}}

وصل النداء الآن إلى الخادم. أوّل قرار: \emph{أيّ دالة} في أيّ \en{controller} تُعالج
هذا المسار؟ هذا هو \en{Routing} — مطابقة عنوان \en{URL} والفعل \en{(GET/POST/...)}
مع الدالة الصحيحة. في \en{ASP.NET Core} تتمّ المطابقة عبر سمات \en{(attributes)}
مكتوبة فوق الصنف والدوال.

هنشوف إيه: \code{NotificationController} — لاحظ السمات التي تُعرّف المسار، وأنّ الدالة
سطرٌ واحد يستدعي \code{mediator.Send}.

\begin{english}
\begin{lstlisting}[language={[Sharp]C}]
[ApiController]
[Route("api/notifications")]
[Authorize]
public sealed class NotificationController(IMediator mediator) : ControllerBase
{
    [HttpGet]
    public async Task<ActionResult<NotificationsPageDto>> GetMine(
        [FromQuery] int page = 1,
        [FromQuery] int pageSize = 20,
        CancellationToken ct = default)
        => Ok(await mediator.Send(new GetNotificationsQuery(page, pageSize), ct));
}
\end{lstlisting}
\end{english}

كيف يقرأ الإطار هذه السمات؟
\begin{itemize}
\item \code{[ApiController]}: يقول «هذا \en{controller} لـ \en{API}»، ويُفعّل سلوكيّات
  مريحة (ربط تلقائيّ للبيانات، ردود أخطاء موحّدة).
\item \code{[Route("api/notifications")]}: البادئة المشتركة لكلّ مسارات هذا الصنف.
\item \code{[HttpGet]} فوق \code{GetMine}: مع البادئة تعني «الفعل \en{GET} على
  \code{/api/notifications} يُنفِّذ هذه الدالة».
\item \code{[Authorize]}: لا يدخل أحد دون رمز صالح — وهنا يظهر أثر رأس \en{Bearer}
  الذي ألصقته الواجهة.
\item \code{[FromQuery]}: خُذ \code{page} و\code{pageSize} من نصّ الاستعلام في العنوان.
\end{itemize}

المهمّ أنّ الدالة \emph{رفيعة}: لا منطق عمل بداخلها. تبني كائن \code{GetNotificationsQuery}
وتُسلِّمه لـ \en{MediatR} عبر \code{mediator.Send}، وتلفّ النتيجة في \code{Ok(...)}
(أي حالة \code{200}). كلّ «كيف نجلب الإشعارات» يعيش في مكان آخر.

\begin{why}
لماذا \en{controller} رفيع لهذه الدرجة؟ لأنّ فصل «استقبال الطلب» عن «تنفيذ العمل»
يجعل المنطق قابلًا للاختبار وحده دون تشغيل خادم \en{HTTP}، ويسمح بأنبوب موحّد يلفّ
كلّ معالِج تلقائيًّا. الـ \en{controller} مجرّد بوّابة ترجمة: \en{HTTP} داخلًا، رسالة
\en{MediatR} خارجًا.
\end{why}

% ====================================================================
\subsection{المحطّة الثالثة: أنبوب الـ \en{Middleware}}

قبل أن تصل الدالة أصلًا، يمرّ الطلب عبر سلسلة من الطبقات اسمها \en{middleware}.
الـ \en{middleware} قطعة تلتفّ حول كلّ طلب: تتدخّل قبل أن يصل للـ \en{controller}،
وقد تتدخّل بعد أن يعود الردّ. تخيّلها طبقات بصلة: الطلب يدخل من الخارج للداخل حتى
المعالِج، ثمّ يخرج الردّ من الداخل للخارج عبر نفس الطبقات عكسيًّا.

\begin{center}\begin{LTR}
\begin{tikzpicture}[node distance=0.24cm]
  \node[wbox] (m1) {UseSerilogRequestLogging  —  time \& log every request};
  \node[wbox,below=of m1] (m2) {SecurityHeadersMiddleware  —  add safe HTTP headers};
  \node[wbox,below=of m2] (m3) {ExceptionHandlerMiddleware  —  errors into ProblemDetails};
  \node[wbox,below=of m3] (m4) {MaintenanceMiddleware  —  short-circuit in maintenance mode};
  \node[wbox,below=of m4] (m5) {UseAuthentication / UseAuthorization  —  who? allowed?};
  \node[pbox,below=0.34cm of m5,text width=6.6cm] (core)
    {Controller  $\rightarrow$  MediatR  $\rightarrow$  Handler};

  \draw[arr] (m1) -- (m2);
  \draw[arr] (m2) -- (m3);
  \draw[arr] (m3) -- (m4);
  \draw[arr] (m4) -- (m5);
  \draw[arr] (m5) -- (core);
  \draw[arrd] (core.east) to[out=0,in=0]
    node[lbl,right] {response travels back up} (m1.east);
\end{tikzpicture}
\end{LTR}\end{center}

الترتيب ليس عشوائيًّا. مثلًا \code{ExceptionHandlerMiddleware} موضوع \emph{مبكّرًا} حتى
يلتقط أيّ خطأ يحدث في أيّ طبقة تحته ويحوّله لردّ موحّد نظيف
\en{(ProblemDetails)}. والصيانة \code{MaintenanceMiddleware} تُقصِّر المسار إن كان
النظام في وضع صيانة فلا يُتعِب القاعدة أصلًا. والتوثيق \code{UseAuthentication} يأتي
قبل التخويل \code{UseAuthorization} لأنّه لا معنى للسؤال «هل يُسمح له؟» قبل معرفة
«مَن هو؟». سنُفرِد الفصل التالي لتفصيل هذا الأنبوب والتحقّق؛ يكفينا الآن أنّه الطبقة
التي تلفّ كلّ شيء.

\begin{idea}[الفكرة: الترتيب هو السلوك]
في أنبوب الـ \en{middleware} \emph{ترتيب} التسجيل هو نفسه ترتيب التنفيذ. تحريك سطر
واحد لأعلى أو لأسفل يُغيّر السلوك: لو وضعت معالج الأخطاء \emph{بعد} التوثيق لَما
التقط أخطاء التوثيق. القاعدة: ضَع ما يجب أن يلفّ الجميع في الخارج (مبكّرًا).
\end{idea}

% ====================================================================
\subsection{المحطّة الرابعة: \en{MediatR} وسلاسل المعالجة \en{(Behaviors)}}

سلّم الـ \en{controller} الرسالة لـ \en{MediatR}. مهمّتها إيجاد المعالِج المناسب
\en{(handler)} وتشغيله. لكن قبل أن تصل للمعالِج، تمرّ الرسالة بسلسلة أخرى داخليّة اسمها
\en{pipeline behaviors}. وهذه بالضبط فكرة \en{interceptor} لكن على مستوى الرسائل:
كلّ \en{behavior} يلتفّ حول المعالِج، يتدخّل قبله وبعده.

\begin{keyp}[نقطة مهمة: الـ \en{behavior} هو \en{interceptor} للرسائل]
عندنا نوعان من الالتفاف في النظام: \en{middleware} يلفّ طلب \en{HTTP}،
و\en{behavior} يلفّ رسالة \en{MediatR}. نفس النمط، مستوًى مختلف. الـ \en{behaviors}
عندنا تفعل أربعة أشياء عرضيّة: \textbf{التسجيل} (كم استغرق الطلب)،
\textbf{التحقّق} \en{(FluentValidation)}، \textbf{الأداء} (إنذار لو تجاوز حدًّا)،
و\textbf{التدقيق} \en{(audit)} للأوامر المُزخرَفة بـ \code{[Auditable]}.
\end{keyp}

الجميل أنّ المعالِج \emph{لا يعرف} بوجود هذه السلاسل إطلاقًا. هو يكتب منطق العمل فقط،
والتحقّق والتسجيل والتدقيق تُضاف حوله تلقائيًّا. لو أردت التحقّق من صحّة
\code{GetNotificationsQuery} تكتب صنف \en{validator} مستقلًّا، وسلسلة التحقّق تلتقطه
دون تعديل المعالِج. هذا الفصل بين «منطق العمل» و«الاهتمامات العرضيّة» هو جوهر قوّة
النمط.

% ====================================================================
\subsection{المحطّة الخامسة: المعالِج ثمّ \en{EF Core}}

وصلنا للقلب. المعالِج \en{(handler)} هو المكان الوحيد الذي يعرف \emph{كيف} تُنفَّذ هذه
العمليّة تحديدًا. لكنّه لا يكتب \en{SQL} بيده؛ يكتب استعلامًا بلغة \en{C\#} باستخدام
\en{LINQ}، و\en{EF Core} يترجمه إلى \en{SQL}. لفهم ذلك نحتاج مفهومين.

\begin{idea}[الفكرة: \en{IQueryable} والتنفيذ المؤجَّل]
\code{IQueryable} ليس بيانات؛ هو \emph{وصف} لاستعلام لم يُنفَّذ بعد. حين تكتب
\code{.Where(...)} ثمّ \code{.OrderBy(...)} أنت تبني \emph{شجرة تعبير} تصف ما تريد،
لا أكثر. لا يذهب شيء للقاعدة حتى تطلب النتيجة فعليًّا. هذا اسمه \en{deferred
execution} — التنفيذ المؤجَّل. اللحظة الحاسمة هي نداء مثل \code{ToListAsync} أو
\code{FirstOrDefaultAsync} أو \code{CountAsync}: عندها \emph{فقط} تُترجَم الشجرة
لـ \en{SQL} وتُرسَل وتُنفَّذ.
\end{idea}

هنشوف إيه: معالِج \code{GetNotificationsQueryHandler} الحقيقيّ. راقب أين يُبنى الاستعلام
وأين يُنفَّذ فعلًا.

\begin{english}
\begin{lstlisting}[language={[Sharp]C}]
public async Task<NotificationsPageDto> Handle(
    GetNotificationsQuery request, CancellationToken ct)
{
    var userId = currentUser.UserId
        ?? throw new ForbiddenAccessException("Not authenticated.");

    // IQueryable: nothing has run yet, just a description of the query
    var baseQuery = db.Notifications.Where(
        n => n.RecipientUserId == userId && n.Channel == NotificationChannel.InApp);

    var total = await baseQuery.CountAsync(ct);       // runs: SELECT COUNT(*)

    var items = await baseQuery
        .OrderByDescending(n => n.CreatedAt)
        .Skip((request.Page - 1) * pageSize)
        .Take(pageSize)
        .Select(n => new NotificationDto(n.Id, n.TitleAr, n.IsRead, n.CreatedAt))
        .ToListAsync(ct);                             // <-- SQL runs HERE

    return new NotificationsPageDto(items, request.Page, pageSize, total, unread);
}
\end{lstlisting}
\end{english}

النقاط المفتاحيّة:
\begin{itemize}
\item \code{db.Notifications.Where(...)} يُنتِج \code{baseQuery} من نوع
  \code{IQueryable} — \emph{لم يُنفَّذ بعد}، ولا استعلام ذهب للقاعدة.
\item \code{.Where(n => n.RecipientUserId == userId ...)}: مرشّح — نُريد إشعارات هذا
  المستخدم فقط، وقناة العرض داخل التطبيق فقط.
\item \code{.OrderByDescending(...).Skip(...).Take(...)}: ترتيب وتقسيم صفحات — كلّها
  ما زالت \emph{وصفًا} يُضاف للشجرة.
\item \code{.Select(n => new NotificationDto(...))}: إسقاط \en{(projection)} —
  نطلب من القاعدة الأعمدة التي نحتاجها فقط لا الصفّ كاملًا.
\item \code{await ...ToListAsync(ct)}: \emph{هنا} تُترجَم الشجرة كلّها لجملة \en{SELECT}
  واحدة، وتُرسَل، وتعود النتائج. سطر \code{CountAsync} قبله نفّذ \code{SELECT COUNT(*)}
  منفصلًا.
\end{itemize}

كيف تبدو الترجمة؟ \en{EF Core} يمتلك «مزوِّدًا» \en{(provider)} لـ \en{SQL Server}
يمشي في شجرة التعبير ويبني الجملة المقابلة. والأهمّ: أيّ قيمة متغيّرة (مثل
\code{userId}) لا تُلصَق في نصّ الجملة، بل تُمرَّر كـ \emph{بارامتر} منفصل.

\begin{center}\begin{LTR}
\begin{tikzpicture}[node distance=0.7cm]
  \node[srcbox,text width=5.3cm,align=left,minimum height=1.6cm] (cs)
    {\textbf{C\# / LINQ}\\[2pt]
     db.Notifications\\
     \ \ .Where(n => n.RecipientUserId == userId)\\
     \ \ .OrderByDescending(n => n.CreatedAt)};
  \node[dbbox,right=1.4cm of cs,text width=5.3cm,align=left,minimum height=1.6cm] (sql)
    {\textbf{SQL (parameterized)}\\[2pt]
     SELECT Id, TitleAr, IsRead, CreatedAt\\
     FROM Notifications\\
     WHERE RecipientUserId = @p0\\
     ORDER BY CreatedAt DESC};
  \draw[arr] (cs) -- node[lbl,above,align=center]
    {EF Core provider\\translates} (sql);
  \node[dbox,below=0.5cm of sql,text width=5.0cm]
    {@p0 = userId  (a bound parameter, never string-concatenated)};
\end{tikzpicture}
\end{LTR}\end{center}

القيمة \code{userId} تُرسَل في خانة البارامتر \code{@p0} (اسمها الفعليّ عند
\en{EF Core} أطول قليلًا، لكن الفكرة واحدة). النصّ لا يحوي القيمة أبدًا. لهذا أثرٌ أمنيّ
كبير جدًّا.

\begin{spotlight}[ركّز هنا لإبهار اللجنة: الترجمة الآمنة \en{LINQ} إلى \en{SQL}]
لأنّ \en{EF Core} يُمرِّر كلّ القيم كبارامترات منفصلة، فإنّ \textbf{حقن \en{SQL}
\en{(SQL Injection)} يصبح مستحيلًا افتراضيًّا} — دون أن يكتب المطوّر سطرًا واحدًا
للحماية. لو أدخل مستخدمٌ خبيثٌ نصًّا مثل \code{'; DROP TABLE ...} كقيمة، فإنّه يُعامَل
\emph{بيانةً} في خانة البارامتر لا \emph{أمرًا}، فلا يُنفَّذ. هذه نقطة قويّة أمام
اللجنة: «أماننا ضدّ الحقن ليس حِيلة أضفناها، بل خاصيّة بنيويّة نحصل عليها مجّانًا لأنّنا
لا نبني \en{SQL} نصيًّا أبدًا».
\end{spotlight}

\begin{secnote}
القاعدة الذهبيّة: لا تبنِ استعلامًا بدمج نصوص \en{(string concatenation)}. طالما بقيت
داخل \en{LINQ}، فالبَرمَتة تلقائيّة. الخطر الوحيد يظهر لو لجأ أحدٌ لـ \en{SQL} خام
مكتوب يدويًّا بدمج مدخلات المستخدم — وهذا ممنوع في مشروعنا.
\end{secnote}

% ====================================================================
\subsection{المحطّة السادسة: الرجوع — كيان إلى \en{DTO} إلى \en{JSON}}

عادت النتائج من القاعدة على شكل كائنات في الذاكرة. لكنّنا لا نُرسِل الكيان
\en{(entity)} الداخليّ للواجهة مباشرةً؛ نُرسِل \en{DTO} — كائن نقل بيانات يحمل فقط ما
تحتاجه الشاشة. في مثالنا رأينا الإسقاط \code{Select(n => new NotificationDto(...))}
يبني الـ \en{DTO} مباشرةً داخل الاستعلام. في حالات أخرى نستخدم \en{AutoMapper} ليُحوِّل
الكيان إلى \en{DTO} بقاعدة تحويل مُعرَّفة مسبقًا، فلا نكتب النسخ حقلًا حقلًا.

بعدها يتولّى الإطار تسلسل الـ \en{DTO} إلى نصّ \en{JSON} ويضعه في جسم الردّ مع حالة
\code{200}. يعود هذا الردّ صعودًا عبر أنبوب الـ \en{middleware} (مسجِّلًا زمنه في
الطريق) حتى يخرج للشبكة.

\begin{why}
لماذا \en{DTO} وليس الكيان مباشرةً؟ لثلاثة أسباب: لا نُسرِّب حقولًا داخليّة أو حسّاسة
للواجهة؛ ولا نربط شكل الـ \en{API} بشكل جداول القاعدة (نستطيع تغيير أحدهما دون كسر
الآخر)؛ ونُرسِل بيانات أقلّ على الشبكة لأنّ الـ \en{DTO} يحمل ما يلزم فقط.
\end{why}

في الواجهة تستقبل \en{TanStack Query} الـ \en{JSON}، تحوّله لكائنات، وتُخزّنه في
الـ \en{cache} تحت مفتاح الاستعلام. أيّ مكوّن يطلب نفس البيانات يحصل عليها فورًا من
الذاكرة. وعند حدوث تغيير (مثل تعليم إشعار مقروءًا) نُبطِل \en{(invalidate)} المفتاح،
فتُعيد المكتبة الجلب تلقائيًّا وتُحدِّث الشاشة. بهذا تُغلَق الدائرة: من زرٍّ في الواجهة،
إلى صفٍّ في القاعدة، ورجوعًا إلى بكسل على الشاشة.

% ====================================================================
\subsection{سؤالك الصريح: لماذا يوجد \en{Service} في الـ \en{Infrastructure} ولا نضع كلّ شيء في المعالِج؟}

سؤال ممتاز ويستحقّ إجابةً صريحة. رأينا أنّ المعالِج يقرأ من القاعدة عبر \en{EF Core}
مباشرةً، فلماذا لا نفعل \emph{كلّ} شيء داخله؟ الجواب: قراءة القاعدة عمليّة نظيفة
ومملوكة لنا، أمّا بعض الأعمال فتحتاج التعامل مع \emph{عالم خارجيّ} لا نتحكّم فيه —
وهذه نعزلها في \en{Service} داخل الـ \en{Infrastructure} (وهو تنفيذ لـ \en{port}
مُعرَّف في طبقة \en{Application}). أربعة أسباب:

\begin{center}
\begin{tabular}{L{4.2cm} L{8.6cm}}
\toprule
\textbf{\color{spPrimary}السبب} & \textbf{\color{spPrimary}الشرح} \\
\midrule
اتّصال \en{I/O} خارجيّ &
  الدفع عبر \en{Stripe}، إرسال البريد، رفع ملفّ إلى \en{Blob}، غرفة \en{ACS} — كلّها
  نداءات شبكة لطرف ثالث، بمنطق إعادة محاولة ومفاتيح سرّيّة. لا يليق حشوها في معالِج. \\
إعادة الاستخدام &
  \code{StripeService} أو \code{IEmailService} يُستدعى من معالِجات كثيرة. نكتب المنطق
  مرّة ونحقنه حيث نشاء بدل تكراره. \\
تبديل المزوّد بالإعداد &
  المعالِج يعتمد على \en{port} (واجهة) لا على تنفيذ بعينه. نستبدل التنفيذ من الإعداد:
  \code{StubMeetingService} محليًّا و\code{AzureCommunicationMeetingService} إنتاجيًّا،
  دون لمس المعالِج. \\
قابليّة الاختبار &
  في الاختبار نحقن \en{stub} مزيّفًا للـ \en{port} فنختبر منطق المعالِج دون لمس
  \en{Stripe} أو الشبكة إطلاقًا — سريع وحتميّ. \\
\bottomrule
\end{tabular}
\end{center}

القاعدة العمليّة: \textbf{القراءة/الكتابة البسيطة للقاعدة تبقى في المعالِج عبر
\en{EF Core}؛ وأيّ تعامل مع نظام خارجيّ يُعزَل خلف \en{port} ويُنفَّذ كـ \en{Service}
في الـ \en{Infrastructure}}. هكذا يبقى المعالِج مركّزًا على «قصّة العمل»، وتبقى
تفاصيل الشبكة والمزوّدين قابلةً للتبديل والاختبار.

\begin{qa}[سؤال متوقّع: أليس هذا طبقات زائدة بلا داعٍ؟]
لا، لأنّ العزل يشتري لنا شيئين ملموسين: نُشغّل النظام محليًّا بلا حساب \en{Stripe}
أو \en{Azure} (عبر تنفيذات \en{stub})، ونكتب اختبارات لا تلمس الشبكة فتكون سريعة
وموثوقة. الكلفة سطر واجهة واحد؛ العائد استقلالٌ عن مزوّدين نتحكّم فيهم بالإعداد لا
بالكود.
\end{qa}

% ====================================================================
\subsection{تمهيد: مسار \en{HTTP} مقابل مسار الزمن الحقيقيّ \en{(SignalR)}}

كلّ ما سبق كان مسار \en{HTTP}: العميل \emph{يسأل} والخادم \emph{يجيب}، ثمّ يُغلَق
الاتّصال. هذا مثاليّ لطلبٍ صريح مثل «أعطني الإشعارات». لكن ماذا لو أردنا العكس: أن
\emph{يدفع} الخادم خبرًا للعميل دون أن يسأل — رسالة دردشة وصلت الآن، أو إشعار جديد
ظهر؟ في \en{HTTP} وحده لا يستطيع الخادم بدء الكلام.

هنا يأتي \en{SignalR}: قناة \emph{مفتوحة ثنائيّة الاتّجاه} (عبر \en{WebSockets}) تبقى
حيّةً بين العميل والخادم، فيستطيع أيّ طرف الإرسال في أيّ لحظة. الفرق الجوهريّ:

\begin{center}
\begin{tabular}{L{4.0cm} L{4.3cm} L{4.3cm}}
\toprule
 & \textbf{\color{spPrimary}\en{HTTP} (طلب/ردّ)} & \textbf{\color{spPrimary}\en{SignalR} (زمن حقيقيّ)} \\
\midrule
مَن يبدأ الكلام؟ & العميل دائمًا & أيّ طرف \\
عمر الاتّصال & لحظيّ (يُغلَق بعد الردّ) & مفتوح ودائم \\
الاتّجاه & باتّجاه واحد لكلّ نداء & ثنائيّ \\
مناسب لِـ & جلب/حفظ بيانات & إشعارات ودردشة مباشرة \\
\bottomrule
\end{tabular}
\end{center}

في \code{Program.cs} نربط الاثنين جنبًا إلى جنب: \code{app.MapControllers()} لمسارات
\en{HTTP}، و\code{app.MapHub(...)} لقنوات \en{SignalR} مثل \code{/hubs/notifications}
و\code{/hubs/chat}. نفس التطبيق يخدم المسارين. تفاصيل هذا المسار الحيّ — كيف يُبنى
الاتّصال، وكيف يُبثّ الإشعار للمستخدم المعنيّ فقط — موضوع فصلٍ قادمٍ مستقلّ؛ هنا فقط
غرسنا البذرة: طلبٌ يسأل، وقناةٌ تدفع.

% ====================================================================
\begin{summary}[الخلاصة في سطور]
\begin{itemize}
\item الواجهة \en{(React + TanStack Query)} تبني النداء، و\en{interceptor} على العميل
  يُلصِق \en{Authorization: Bearer} تلقائيًّا، ويجدّد الرمز عند \code{401} ويُعيد الطلب.
\item التوجيه \en{(routing)} عبر \code{[ApiController]} و\code{[HttpGet]} يوصّل المسار
  للدالة الصحيحة؛ والـ \en{controller} رفيع: سطرٌ واحد \code{mediator.Send}.
\item الـ \en{middleware} سلسلة تلفّ كلّ طلب (تسجيل، أمان، أخطاء، صيانة، توثيق)؛
  ترتيبها هو سلوكها.
\item الـ \en{MediatR pipeline behaviors} هي \en{interceptor} للرسائل: تحقّق وتسجيل
  وأداء وتدقيق حول المعالِج دون أن يعلم.
\item المعالِج يكتب \en{LINQ} على \code{IQueryable} (تنفيذ مؤجَّل)؛ لا شيء يُنفَّذ حتى
  \code{ToListAsync}؛ و\en{EF Core} يترجم لـ \en{SQL} مُبارَم يمنع حقن \en{SQL}.
\item الردّ: كيان إلى \en{DTO} إلى \en{JSON}، والواجهة تُحدِّث \en{cache}‎ها؛ فتُغلَق
  الدائرة من الزرّ إلى الصفّ ورجوعًا.
\item ما يتعامل مع عالم خارجيّ (دفع/بريد/تخزين) يُعزَل في \en{Service} خلف \en{port}
  للإعادة والتبديل والاختبار.
\item \en{HTTP} للطلب/الردّ، و\en{SignalR} للدفع الحيّ ثنائيّ الاتّجاه — وهو فصل قادم.
\end{itemize}
\end{summary}

\clearpage
% ch21 — مستويات التحقّق (Validation) والوسطاء (Middleware / Filters / Behaviors)
\section{مستويات التحقّق والوسطاء}

في الفصل السابق مشينا مع الطلب من ساعة ما يدخل الخادم لحدّ ما يرجع الردّ، وشُفنا الـ\en{pipeline behaviors} اللي بتلفّ كل حالة استخدام. في الفصل ده هنركّز على سؤالين عمليّين جدًّا: أين نتحقّق من المدخلات؟ وكيف نعترض الطلب في المكان الصحّ؟

هنفهم إيه بالظبط؟ أولًا: ليه بنكرّر التحقّق أربع مرّات (عند العميل، ثم \en{FluentValidation}، ثم قواعد العمل داخل الـ\en{handler}، ثم قاعدة البيانات) وكل طبقة بتحرس إيه. ثانيًا: الفرق بين ثلاث آليّات اعتراض مختلفة \en{middleware} و\en{filters} و\en{behaviors}، وكل واحدة بتشتغل على أيّ مستوى. ثالثًا: الوسطاء الحقيقية عندنا بمقتطفات كود فعليّة. رابعًا: إزّاي مجرّد سمة (\en{attribute}) فوق كلاس بتغيّر سلوك الطلب. خلّينا نبدأ واحدة واحدة.

\begin{idea}[الفكرة: دفاع في العمق]
التحقّق مش خطوة واحدة، ده أربع حلقات متتالية. كل حلقة أرخص وأسرع من اللي بعدها لكن أقلّ ثقة، والحلقة الأخيرة (قاعدة البيانات) أبطأ لكن لا تُخترَق حتى تحت التزامن. لو طبقة فشلت، اللي بعدها تمسك الكرة.
\end{idea}

\subsection{أربع طبقات تحقّق — وليه بنكرّرها}

تخيّل التحقّق كأنّه قِمع (\en{funnel}): المدخلات تدخل من فوق، وكل طبقة بتصفّي شيئًا معيّنًا، وما ينزل للقاع هو بس الطلب النظيف اللي يستحقّ إنه يُحفَظ. الرسمة دي بتوضّح الأربع طبقات ونقاط الرفض عند كل واحدة.

\begin{center}\begin{LTR}
\begin{tikzpicture}[node distance=0.5cm]
  \node[wbox,text width=8cm,fill=spLightBlue] (l1)
    {1 — Client (React form): fast UX, not trusted};
  \node[wbox,text width=8cm,fill=spLightTeal,below=of l1] (l2)
    {2 — FluentValidation: shape / limits};
  \node[wbox,text width=8cm,fill=spLightPurple,below=of l2] (l3)
    {3 — Handler: business rules, ownership};
  \node[wbox,text width=8cm,fill=spLightGreen,below=of l3] (l4)
    {4 — Database: unique filtered index, FK, RowVersion};
  \node[dbbox,below=of l4] (ok) {Persisted};
  \draw[arr] (l1) -- (l2);
  \draw[arr] (l2) -- (l3);
  \draw[arr] (l3) -- (l4);
  \draw[arr] (l4) -- (ok);
  \node[fbox,right=1.1cm of l2] (r2) {400 rejected};
  \node[fbox,right=1.1cm of l3] (r3) {403 / 409 / 422};
  \node[fbox,right=1.1cm of l4] (r4) {2601 -> 409};
  \draw[arrd] (l2) -- (r2);
  \draw[arrd] (l3) -- (r3);
  \draw[arrd] (l4) -- (r4);
\end{tikzpicture}
\end{LTR}\end{center}

\subsubsection{الطبقة الأولى — العميل (نموذج \en{React})}

أوّل تحقّق بيحصل في المتصفّح: النموذج بيمنعك تبعت رسالة فاضية، أو بيقولك «الحقل ده مطلوب» قبل ما تضغط إرسال. ده ممتاز لتجربة المستخدم لأنه فوريّ وبدون رحلة شبكة. لكنه غير موثوق أمنيًّا إطلاقًا: الكود ده بيشتغل على جهاز المستخدم، وأيّ حدّ يقدر يفتح أدوات المطوّر أو يبعت طلبًا مباشرًا بـ\en{curl} يتخطّى النموذج بالكامل. القاعدة الذهبية: تحقّق العميل للراحة، مش للأمان.

\begin{why}
لو اعتمدنا على تحقّق العميل وحده، أيّ مهاجم يبعت \en{JSON} يدويّ ويحقن قيمًا مستحيلة. عشان كده كل تحقّق مهمّ لازم يتكرّر على الخادم — العميل مجرّد طبقة تجميل سريعة.
\end{why}

\subsubsection{الطبقة الثانية — \en{FluentValidation} عبر \en{ValidationBehavior}}

أوّل خطّ دفاع حقيقيّ على الخادم. قبل ما أيّ \en{handler} يشتغل، بيمرّ الطلب على \en{ValidationBehavior} اللي بيجمع كل الـ\en{validators} المسجّلة للأمر ده وينفّذها. لو فيه فشل واحد، بيرمي \en{ValidationException} فورًا، والطلب يتوقّف قبل أيّ شغل أو أيّ لمسة لقاعدة البيانات. ده المقتطف الحقيقيّ من \code{ValidationBehavior}:

\begin{english}
\begin{lstlisting}[language={[Sharp]C}]
if (!_validators.Any())
    return await next(cancellationToken).ConfigureAwait(false);

var context = new ValidationContext<TRequest>(request);
var validationResults = await Task.WhenAll(
    _validators.Select(v => v.ValidateAsync(context, cancellationToken)))
    .ConfigureAwait(false);

var failures = validationResults
    .Where(r => !r.IsValid)
    .SelectMany(r => r.Errors)
    .ToList();

if (failures.Count > 0)
    throw new AppValidationException(failures);

return await next(cancellationToken).ConfigureAwait(false);
\end{lstlisting}
\end{english}

نقرأ الأسطر المهمّة: لو مفيش \en{validators} للأمر، بنعدّي مباشرة على \code{next} بدون أيّ تكلفة. لو فيه، بننفّذهم كلهم بالتوازي بـ\code{Task.WhenAll}. بعدين بنجمّع كل الأخطاء من كل النتائج في قائمة واحدة بـ\code{SelectMany}. لو القائمة فيها ولو خطأ واحد، بنرمي \code{AppValidationException} — ولاحظ إن الطلب ما وصلش أبدًا للـ\en{handler}. ده معناه إننا بنرفض بسرعة قبل أيّ عمل، وده رخيص جدًّا.

طيب الـ\en{validator} نفسه شكله إيه؟ ده مثال حقيقيّ لأمر إرسال رسالة في الشات:

\begin{english}
\begin{lstlisting}[language={[Sharp]C}]
public sealed class SendMessageCommandValidator
    : AbstractValidator<SendMessageCommand>
{
    public SendMessageCommandValidator()
    {
        RuleFor(v => v.RecipientId).NotEmpty();
        RuleFor(v => v.Body).NotEmpty().MaximumLength(2000);
    }
}
\end{lstlisting}
\end{english}

القواعد هنا «شكليّة» بحتة: المستلِم مش فاضي، ونصّ الرسالة مش فاضي وطوله ما يتعدّاش 2000 حرف. لاحظ إن الـ\en{validator} ما بيعرفش حاجة عن قواعد العمل: هل المستخدم محظور؟ هل المحادثة موجودة؟ دي أسئلة مش من شغله — دي شغل الطبقة اللي بعده. الـ\en{validator} بيحرس الشكل والحدود فقط.

\begin{keyp}[نقطة مهمة: حدّ الطبقة الثانية]
\en{FluentValidation} بيجاوب على سؤال «هل المدخلات لها شكل صحيح؟» — أنواع، مش فاضي، حدود دنيا وعليا، صيغة \en{email}. مش بيجاوب على «هل مسموح للمستخدم ده يعمل ده دلوقتي؟» — ده قرار عمل، محلّه الطبقة الثالثة.
\end{keyp}

\subsubsection{الطبقة الثالثة — قواعد العمل داخل الـ\en{handler}}

هنا بنطبّق ثوابت الدومين اللي محتاجة تقرأ من قاعدة البيانات أو تفهم السياق. أمثلة حقيقية: منع التقديم على منحة بعد انتهاء موعدها، فحص إن المستخدم فعلًا يملك المورد اللي بيعدّله (الملكيّة)، ومنع إرسال رسالة لمستخدم حظرك. في نفس أمر إرسال الرسالة، الـ\en{handler} بيعمل الفحص ده:

\begin{english}
\begin{lstlisting}[language={[Sharp]C}]
var isBlocked = await db.UserBlocks.AnyAsync(b =>
    (b.BlockerId == recipientId && b.BlockedUserId == currentUser.UserId) ||
    (b.BlockerId == currentUser.UserId && b.BlockedUserId == recipientId), ct);

if (isBlocked)
    throw new ConflictException("Cannot send message. User is blocked.");
\end{lstlisting}
\end{english}

الفحص ده مستحيل نعمله في \en{FluentValidation}، لأنه محتاج يستعلم قاعدة البيانات ويعرف مين المستخدم الحاليّ. فحوصات الملكيّة والحالة والمواعيد كلها من نوعية «قواعد عمل»، فمحلّها الطبيعيّ داخل الـ\en{handler} حيث السياق الكامل متاح. لو القاعدة اتكسرت، بنرمي استثناءً دلاليًّا (\code{ConflictException} أو \code{ForbiddenAccessException}) والوسيط بيترجمه لكود \en{HTTP} مناسب — هنشوف ده بعد شويّة.

\subsubsection{الطبقة الرابعة — قيود قاعدة البيانات (الحصن الأخير)}

الطبقات الثلاث السابقة كلها في الكود، والكود ممكن يتسابق. تخيّل مستخدمين، أو ضغطتين متتاليتين سريعتين، بيقدّموا على نفس المنحة في نفس اللحظة. الطبقة الثالثة بتقرأ «هل فيه تقديم سابق؟» وتلاقي لأ في الحالتين، فتسمح للاتنين — لأنهم قروا قبل ما أيّ واحد يكتب. ده اسمه \en{race condition}، وميقدرش يمسكه أيّ كود تطبيقيّ لوحده.

هنا بييجي دور قاعدة البيانات: الفهرس الفريد المصفّى (\en{unique filtered index}) بيرفض ذرّيًّا (\en{atomically}) الصفّ التاني على مستوى المحرّك، مهما كان التزامن. ده التعريف الحقيقيّ لتقديمات المنح عندنا:

\begin{english}
\begin{lstlisting}[language={[Sharp]C}]
b.HasIndex(a => new { a.StudentId, a.ScholarshipId })
    .IsUnique()
    .HasFilter(
        "[ScholarshipId] IS NOT NULL " +
        "AND [Status] <> 'Withdrawn' " +
        "AND [Status] <> 'Rejected' " +
        "AND [Status] <> 'Accepted'")
    .HasDatabaseName("UX_Applications_Student_Scholarship_Active");
\end{lstlisting}
\end{english}

نقرأه: الفهرس فريد على الزوج (الطالب، المنحة)، يعني ممنوع صفّين نشطين لنفس الزوج. كلمة «مصفّى» هي شرط الـ\code{HasFilter}: الفهرس بيطبّق قيد التفرّد فقط على الصفوف اللي حالتها نشطة — مش منسحبة ولا مرفوضة ولا مقبولة. ليه؟ عشان لو الطالب انسحب من منحة، نسمحله يقدّم عليها تاني؛ الفهرس بيتجاهل الصفوف القديمة غير النشطة. وشرط \code{[ScholarshipId] IS NOT NULL} بيستثني تتبّع المنح الخارجية (اللي مالهاش رابط منحة داخليّة) من قاعدة التفرّد. لمّا التزامن يحاول يدخل صفّ ثانٍ نشط، \en{SQL Server} بيرمي خطأ 2601، والطبقة بتترجمه لـ\en{HTTP} 409.

\begin{secnote}
لاحظ إن الـ\en{predicate} بتاع الفهرس المصفّى في \en{SQL Server} ما بيقبلش \code{NOT IN} ولا \code{OR} في بعض المواضع؛ عشان كده اتكتب كـ\code{<>} متسلسلة بـ\code{AND}. نفس المعنى، صيغة مقبولة. التفصيلة دي مهمّة: القيد لازم يكون صحيحًا على مستوى المحرّك، مش بس نيّة في الكود.
\end{secnote}

\begin{spotlight}[ركّز هنا لإبهار اللجنة: كل طبقة تحرس شيئًا مختلفًا]
لو سألتك اللجنة «ليه بتكرّروا التحقّق أربع مرّات؟» الإجابة القويّة: مش تكرار، ده دفاع في العمق. العميل يحرس تجربة المستخدم، \en{FluentValidation} يحرس شكل المدخلات، الـ\en{handler} يحرس قواعد العمل، وقاعدة البيانات تحرس التكامل تحت التزامن. الفهرس الفريد هو آخر حصن لا يُخترَق حتى لو وصل طلبان في نفس الميلّي-ثانية — وده اللي أيّ كود تطبيقيّ لوحده ميقدرش يضمنه.
\end{spotlight}

\subsection{ثلاث آليّات اعتراض — أين يعمل كلٌّ}

عندنا ثلاث أماكن نقدر نعترض فيها الطلب، وكل واحدة على مستوى مختلف. أسهل طريقة لفهمها إنك تتخيّلهم طبقات بصل (\en{onion}): الأبعد يلفّ الأقرب، والقلب هو الـ\en{handler} اللي فيه المنطق الفعليّ.

\begin{center}\begin{LTR}
\begin{tikzpicture}
  \draw[draw=spSec!65,fill=spLightOrange,rounded corners=6pt] (0,0) rectangle (11,7.2);
  \node[color=spSec!80,font=\small\bfseries] at (5.5,6.8) {Middleware — HTTP level};
  \draw[draw=spHow!65,fill=spLightBlue,rounded corners=6pt] (1,0.8) rectangle (10,6.1);
  \node[color=spHow!80,font=\small\bfseries] at (5.5,5.7) {Filter / Attribute — MVC action};
  \draw[draw=spQA!65,fill=spLightPurple,rounded corners=6pt] (2,1.7) rectangle (9,5.0);
  \node[color=spQA!80,font=\small\bfseries] at (5.5,4.6) {Pipeline Behavior — use case};
  \draw[draw=spAccent!70,fill=spLightTeal,rounded corners=6pt] (3.2,2.6) rectangle (7.8,3.9);
  \node[color=spAccent!80,font=\small\bfseries,align=center] at (5.5,3.25) {Handler\\core logic};
\end{tikzpicture}
\end{LTR}\end{center}

\en{Middleware} بيشتغل على مستوى \en{HTTP} الخام: بيلفّ كل طلب داخل ولا خارج، حتى الطلبات اللي مالهاش أكشن أصلًا (زيّ \code{/health}). \en{Filters} و\en{attributes} بيشتغلوا على مستوى الـ\en{MVC}: بعد ما \en{ASP.NET} يحدّد الأكشن، وبيقدروا يقروا سماته زيّ \code{[Authorize]}. أخيرًا \en{MediatR pipeline behaviors} بتشتغل على مستوى حالة الاستخدام: جوّه الـ\en{handler} تقريبًا، وبتعرف نوع الأمر بالظبط. الجدول ده بيلخّص متى تختار أيًّا:

\begin{center}
\begin{tabular}{L{3.4cm} L{5.4cm} L{5.2cm}}
\toprule
الآليّة & المستوى & تختارها لِـ \\
\midrule
\en{Middleware} & \en{HTTP} — يلفّ كل الطلبات & رؤوس أمان، صيانة، ترجمة استثناءات، \en{CORS} \\
\en{Filter / Attribute} & الأكشن (\en{MVC}) & مصادقة وتفويض، تقييد الأدوار، علامات مثل \code{[Auditable]} \\
\en{Pipeline Behavior} & حالة الاستخدام (\en{MediatR}) & تحقّق، تدقيق، معاملات، سجلّ لكل أمر \\
\bottomrule
\end{tabular}
\end{center}

\begin{note}
القاعدة العمليّة: لو الاهتمام يخصّ كل طلب \en{HTTP} بلا استثناء (زيّ رؤوس الأمان)، حطّه في \en{middleware}. لو يخصّ الأوامر والاستعلامات فقط (زيّ التحقّق والتدقيق)، حطّه في \en{behavior}. لو يخصّ سلوك الأكشن نفسه (زيّ التفويض)، استخدم \en{attribute/filter}.
\end{note}

\subsection{الوسطاء الكاستم عندنا}

عندنا كذا \en{middleware} مكتوبة بإيدينا. خلّينا ناخد أهمّ ثلاثة.

\subsubsection{\en{ExceptionHandlerMiddleware} — ترجمة الاستثناءات}

ده أهمّهم. مكانه بيلفّ باقي الـ\en{pipeline} كله في \code{try/catch} واحد. الفكرة: بدل ما أيّ استثناء يطلع للمستخدم كـ500 مكشوف (مع أثر مكدّس محرج وربّما تسريب معلومات)، بنترجم كل نوع استثناء لكود \en{HTTP} نظيف وجسم \en{ProblemDetails} موحّد. مقتطف من قلب الوسيط:

\begin{english}
\begin{lstlisting}[language={[Sharp]C}]
catch (ValidationException vex)
{
    await WriteProblemAsync(context, HttpStatusCode.UnprocessableEntity,
        "Validation failed", detail: vex.Message, errors: vex.Errors);
}
catch (NotFoundException nex)
{
    await WriteProblemAsync(context, HttpStatusCode.NotFound,
        "Not found", nex.Message);
}
// ... ConflictException -> 409, ForbiddenAccessException -> 403 ...
catch (Exception ex)
{
    logger.LogError(ex, "Unhandled exception at {Path}", context.Request.Path);
    var detail = env.IsDevelopment() ? ex.ToString()
        : "An error occurred. Please try again.";
    await WriteProblemAsync(context, HttpStatusCode.InternalServerError,
        "Internal server error", detail);
}
\end{lstlisting}
\end{english}

كل \code{catch} بيربط نوع استثناء دلاليّ بكود \en{HTTP} صحيح: فشل التحقّق يبقى 422، «غير موجود» يبقى 404، التعارض يبقى 409، الممنوع يبقى 403. لاحظ الـ\code{catch} الأخير: أيّ استثناء غير متوقّع بيتسجّل في اللوج (عشان نشوفه)، لكن التفاصيل الخام ما بتظهرش للمستخدم في الإنتاج — بيشوف رسالة عامّة نظيفة. دالة \code{WriteProblemAsync} بتكتب \en{JSON} بنوع \code{application/problem+json} فيه \code{title} و\code{status} و\code{detail} و\code{traceId} — والأخير بيربط شكوى المستخدم بسطر اللوج المحدّد.

الرسمة دي بتوضّح الفكرة: كل الأنواع بتتجمّع في مكان واحد وتتفرّع لأكواد \en{HTTP} نظيفة (\en{fan-in} ثم \en{fan-out}):

\begin{center}\begin{LTR}
\begin{tikzpicture}[node distance=0.45cm and 2.3cm]
  \node[gbox,text width=3.2cm] (h) {ExceptionHandler};
  \node[srcbox,above left=0.3cm and 2.3cm of h] (e1) {ValidationException};
  \node[srcbox,left=2.3cm of h] (e2) {NotFoundException};
  \node[srcbox,below left=0.3cm and 2.3cm of h] (e3) {ForbiddenAccess};
  \node[srcbox,below left=1.5cm and 2.3cm of h] (e4) {Exception (unknown)};
  \node[fbox,above right=0.3cm and 2.3cm of h] (c1) {422};
  \node[fbox,right=2.3cm of h] (c2) {404};
  \node[fbox,below right=0.3cm and 2.3cm of h] (c3) {403};
  \node[fbox,below right=1.5cm and 2.3cm of h] (c4) {500 + log};
  \draw[arr] (e1) -- (h);
  \draw[arr] (e2) -- (h);
  \draw[arr] (e3) -- (h);
  \draw[arr] (e4) -- (h);
  \draw[arr] (h) -- (c1);
  \draw[arr] (h) -- (c2);
  \draw[arr] (h) -- (c3);
  \draw[arr] (h) -- (c4);
\end{tikzpicture}
\end{LTR}\end{center}

\begin{spotlight}[ركّز هنا لإبهار اللجنة: الترجمة مرّة واحدة في مكان واحد]
النقطة اللي بتفرّق مشروع محترف: مفيش \code{try/catch} متناثر في كل \en{controller}. كل الأخطاء بتُترجَم في نقطة واحدة (\en{ExceptionHandler})، فالردود متّسقة، وما فيش 500 مكشوف بيسرّب تفاصيل داخلية، وكل خطأ 500 حقيقيّ بيتسجّل مع \code{traceId} يخلّي الدعم يلاقيه في ثانية. مركزيّة معالجة الأخطاء = اتّساق + أمان + قابليّة تشخيص.
\end{spotlight}

\subsubsection{\en{MaintenanceMiddleware} — وضع الصيانة}

لمّا نحتاج نوقف المنصّة مؤقّتًا (ترحيل قاعدة بيانات مثلًا)، بنقلب مفتاح \code{maintenance.enabled} في الإعدادات، والوسيط ده بيردّ 503 على كل طلب — إلا مسارات مستثناة بعناية عشان الأدمن نفسه ما يتقفلش برّه:

\begin{english}
\begin{lstlisting}[language={[Sharp]C}]
var path = context.Request.Path.Value ?? string.Empty;
if (BypassedPrefixes.Any(p =>
        path.StartsWith(p, StringComparison.OrdinalIgnoreCase)))
{
    await next(context).ConfigureAwait(false);
    return;
}
// ... read maintenance.enabled from cache (30s TTL) ...
if (isInMaintenance)
{
    context.Response.StatusCode = StatusCodes.Status503ServiceUnavailable;
    await context.Response.WriteAsJsonAsync(new { maintenance = true });
    return;
}
\end{lstlisting}
\end{english}

المسارات المستثناة بتشمل \code{/health} (فحص الخادم)، و\code{/api/auth} (عشان الأدمن يقدر يسجّل دخول ويطفّي الصيانة)، و\code{/api/admin/settings} (نفس المفتاح لازم يفضل قابل للوصول). ولاحظ إن قراءة الإعداد بتتخزّن مؤقّتًا 30 ثانية عشان ما نضربش قاعدة البيانات في كل طلب. وفي حال فشل القراءة، الوسيط «يفشل مفتوحًا» — ما بيحجبش الطلبات الشرعيّة بسبب عطل مؤقّت في القراءة.

\subsubsection{\en{SecurityHeadersMiddleware} — رؤوس الأمان}

وسيط بسيط لكن مهمّ: بيضيف رؤوس أمان على كل ردّ. بيستخدم \code{OnStarting} عشان الرؤوس تتكتب بالظبط قبل ما جسم الردّ يبدأ يتبعت:

\begin{english}
\begin{lstlisting}[language={[Sharp]C}]
context.Response.OnStarting(() =>
{
    var headers = context.Response.Headers;
    headers["X-Content-Type-Options"] = "nosniff";
    headers["X-Frame-Options"] = "DENY";
    headers["Referrer-Policy"] = "strict-origin-when-cross-origin";
    // ... Permissions-Policy, Cross-Origin-Opener-Policy ...
    if (context.Request.IsHttps)
        headers["Strict-Transport-Security"] =
            "max-age=63072000; includeSubDomains; preload";
    return Task.CompletedTask;
});
\end{lstlisting}
\end{english}

كل رأس بيسدّ ثغرة معروفة: \code{X-Content-Type-Options: nosniff} بيمنع المتصفّح من تخمين نوع المحتوى، \code{X-Frame-Options: DENY} بيمنع تضمين الموقع في \en{iframe} (حماية من \en{clickjacking})، و\code{Strict-Transport-Security} بيجبر المتصفّح يستخدم \en{HTTPS} لمدّة سنتين. النقطة إن ده اهتمام يخصّ كل ردّ بلا استثناء — فمحلّه الطبيعيّ \en{middleware}، مش \en{filter} ولا \en{behavior}.

\subsection{السمات كـ\en{filters} — يعني إيه \en{attribute} بيغيّر السلوك}

آخر قطعة في اللغز: إزّاي مجرّد سطر فوق كلاس بيغيّر مصير الطلب؟ السمة (\en{attribute}) هي وسم بيتحطّ على كود، وحاجة تانية بتقرأه وتتصرّف بناءً عليه.

أوضح مثال \code{[Authorize]}. لمّا تحطّ \code{[Authorize(Roles = "Admin")]} فوق \en{controller}، إنت مش بتكتب منطق فحص؛ إنت بتحطّ علامة. الـ\en{authorization filter} في \en{ASP.NET} بيقرا العلامة دي قبل ما الأكشن يشتغل، ولو المستخدم مش أدمن بيردّ 403 من غير ما الأكشن يتنفّذ أصلًا. السمة أعلنت النيّة؛ الـ\en{filter} نفّذها.

عندنا نفس النمط في التدقيق. الأمر بيتوسم بـ\code{[Auditable]}، والـ\en{AuditBehavior} (اللي شفناه في الفصل السابق) بيقرا العلامة دي ويسجّل الفعل تلقائيًّا:

\begin{english}
\begin{lstlisting}[language={[Sharp]C}]
[Auditable(AuditAction.Create, "ChatMessage",
    TargetIdProperty = nameof(RecipientId),
    SummaryTemplate = "Sent message to {RecipientId}")]
public sealed record SendMessageCommand(
    Guid RecipientId,
    string Body) : IRequest<Guid>;
\end{lstlisting}
\end{english}

الأمر ده ما فيهوش أيّ كود تدقيق. مجرّد سمة \code{[Auditable]} بتقول: «لمّا ينجح ده، سجّل فعل إنشاء على كيان \en{ChatMessage}، والهدف هو \code{RecipientId}». الـ\en{AuditBehavior} بيقرا السمة بالـ\en{reflection} ويكتب سجلّ التدقيق. الجمال هنا الفصل: الأمر بيعلن «أنا قابل للتدقيق»، والآليّة بتتكفّل بالتنفيذ في مكان واحد — نفس فلسفة \code{[Authorize]} بالظبط.

\begin{qa}[سؤال متوقّع: ليه سمة بدل استدعاء دالة؟]
لأن السمة تصريحيّة (\en{declarative}): بتقول «إيه» مش «إزّاي». ده بيخلّي الأمر نظيف ومركّز على منطقه، وبيضمن إن كل الأوامر القابلة للتدقيق تتعامل بنفس الطريقة تلقائيًّا — مستحيل حد ينسى يستدعي دالة التدقيق، لأنها مش استدعاء أصلًا، دي علامة بتتقرا مركزيًّا.
\end{qa}

\begin{summary}[الخلاصة في سطور]
\begin{itemize}
  \item التحقّق أربع طبقات، كل واحدة تحرس شيئًا مختلفًا: العميل (تجربة سريعة، غير موثوقة)، \en{FluentValidation} (شكل وحدود المدخلات، رفض 400)، الـ\en{handler} (قواعد العمل والملكيّة والمواعيد)، وقاعدة البيانات (تكامل تحت التزامن).
  \item الفهرس الفريد المصفّى هو الحصن الأخير: يرفض ذرّيًّا الصفّ المكرّر مهما كان التزامن — وده اللي الكود التطبيقيّ لوحده ميقدرش يضمنه.
  \item ثلاث آليّات اعتراض على مستويات مختلفة: \en{middleware} (كل طلب \en{HTTP})، \en{filter/attribute} (الأكشن)، \en{behavior} (حالة الاستخدام). اختر بحسب نطاق الاهتمام.
  \item وسطاؤنا الحقيقية: \en{ExceptionHandler} يترجم كل استثناء لكود \en{HTTP} نظيف في مكان واحد (لا 500 مكشوف)، \en{Maintenance} يردّ 503 مع مسارات مستثناة، \en{SecurityHeaders} يضيف رؤوس أمان لكل ردّ.
  \item السمة تصريحيّة: \code{[Authorize]} و\code{[Auditable]} علامات تقرأها آليّة مركزيّة وتتصرّف — تفصل «الإعلان عن النيّة» عن «تنفيذها».
\end{itemize}
\end{summary}

\clearpage
% الفصل 22 — الزمن الحقيقيّ بعمق: WebSocket و SignalR وفكرة الـ Hubs
\section{الزمن الحقيقيّ بعمق: \en{WebSocket} و\en{SignalR} وفكرة الـ \en{Hubs}}

في الجزء الأول رأينا الإشعارات والشات وهما يشتغلان: رسالة تظهر عند الطرف
الآخر لحظيًّا، وجرس الإشعارات يرنّ بلا تحديث للصفحة. هذا الفصل يفتح الغطاء ويجيب
عن السؤال الأصعب: كيف يصل التحديث من الخادم إلى المتصفّح في نفس اللحظة،
بلا أن يسأل المتصفّح أصلًا؟ هنمشي «واحدة واحدة»: نبدأ من مشكلة الـ \en{HTTP}
العاديّ، ثم نفهم الـ \en{WebSocket} كحلّ، ثم نرى \en{SignalR} كطبقة تجريد فوقه،
وأخيرًا نغوص في الكود الحقيقيّ لثلاثة \en{hubs} في \en{ScholarPath}: كيف يتوثّق
الاتصال، وكيف يُوجَّه الدفع بالمجموعات، وأين يجلس «المشتركون» على الطرفين.

هذه المفاهيم من المفترض أنّك رأيتها مستخدَمة في الجزء الأول؛ هنا نشرح الآليّة
تحتها لا الاستخدام.

% ─────────────────────────────────────────────────────────────
\subsection{المشكلة أولًا: لماذا لا يكفي \en{HTTP} العاديّ؟}

الـ \en{HTTP} التقليديّ يعمل بنمط «طلب ثم ردّ»: المتصفّح يفتح اتصالًا، يرسل
طلبًا، يستقبل الردّ، ثم يُغلَق الاتصال. الخادم لا يملك أي طريقة ليبدأ
الحديث من نفسه؛ هو يردّ فقط حين يُسأل. طيب، ماذا لو وصلت رسالة شات جديدة وأنا لا
أسأل؟ لن أعرف بها.

الحلّ الساذج اسمه \en{polling}: اجعل المتصفّح يسأل الخادم كل ثانية «في جديد؟».
هذا يعمل، لكنّه مكلف: آلاف الطلبات الفارغة (\en{no data}) لكل مستخدم، كل واحد منها
يحمل ترويسات \en{HTTP} كاملة ومصافحة \en{TLS}، وتأخير يصل إلى ثانية كاملة قبل أن
يرى المستخدم أي تحديث. تخيّلها مضروبة في آلاف المتّصلين.

\begin{why}
اللحظيّة الحقيقيّة تحتاج أن يدفع الخادم البيانات متى شاء — لا أن ينتظر
سؤالًا. الـ \en{polling} يقلب المعادلة: يجعل العميل يسأل باستمرار، فيدفع
ثمن التأخير والحِمل معًا. نحتاج قناةً يبقى الخادم قادرًا فيها على الكلام.
\end{why}

% ─────────────────────────────────────────────────────────────
\subsection{\en{WebSocket}: مصافحة تُرقّي الاتصال إلى قناة دائمة}

الـ \en{WebSocket} بروتوكول يحلّ هذا بالضبط. يبدأ حياته كطلب \en{HTTP} عاديّ، لكن
بترويسة خاصة \code{Upgrade: websocket}. الخادم يوافق بردّ الحالة
\code{101 Switching Protocols}، وهذه اللحظة تُسمّى «المصافحة» \en{(handshake)}.
بعدها لا يُغلَق الاتصال، بل يُرقّى إلى قناة دائمة ثنائية الاتجاه
\en{(full-duplex)}: الطرفان يرسلان في أي وقت وعبر نفس الاتصال المفتوح.

النتيجة المهمة: صار الخادم قادرًا على أن يدفع \en{(push)} رسالةً للعميل بلا
أي طلب سابق. لا مزيد من \en{polling}، ولا تأخير، واتصال واحد يخدم كل التحديثات.

\begin{center}\begin{LTR}
\begin{tikzpicture}
  % ---- polling (المشكلة) ----
  \node[actor] (cp) at (0,0) {Client};
  \node[actor] (sp) at (6,0) {Server};
  \draw[arr] ([yshift=5pt]cp.east) -- node[lbl,above]{GET /updates?} ([yshift=5pt]sp.west);
  \draw[arrd] ([yshift=-7pt]sp.west) -- node[lbl,below]{204 no data} ([yshift=-7pt]cp.east);
  \node[fbox,text width=6cm] at (3,-1.25) {polling: يتكرّر كل ثانية = حِمل + تأخير};

  % ---- websocket (الحل) ----
  \node[actor] (cw) at (0,-3.1) {Client};
  \node[actor] (sw) at (6,-3.1) {Server};
  \draw[arr] ([yshift=9pt]cw.east) -- node[lbl,above]{HTTP Upgrade (handshake)} ([yshift=9pt]sw.west);
  \draw[arrb] (cw.east) -- (sw.east);
  \draw[arr] ([yshift=-9pt]sw.west) -- node[lbl,below]{server push (بلا طلب)} ([yshift=-9pt]cw.east);
  \node[tbox,text width=6cm] at (3,-4.4) {قناة واحدة دائمة ثنائية الاتجاه (full-duplex)};
\end{tikzpicture}
\end{LTR}\end{center}

في الأعلى: الـ \en{polling} يرسل طلبًا بعد طلب فيرجع أغلبها فارغًا. في الأسفل:
مصافحة واحدة تفتح قناة دائمة، ثم يدفع الخادم متى شاء عبرها.

% ─────────────────────────────────────────────────────────────
\subsection{\en{SignalR} فوق الـ \en{WebSocket}: تجريد يخفي النقل}

الـ \en{WebSocket} قويّ لكنّه منخفض المستوى: أنت تدير أُطُر البايتات
\en{(frames)}، وإعادة الاتصال عند القطع، وماذا تفعل على المتصفّحات أو الشبكات
التي تحظر الـ \en{WebSocket}. هنا يدخل \en{SignalR}: مكتبة من \en{ASP.NET Core}
تجلس فوق الـ \en{WebSocket} وتخفي كل هذا التعقيد.

\begin{idea}[الفكرة: \en{SignalR} = تجريد النقل + إدارة الاتصالات]
\en{SignalR} يفاوض تلقائيًّا \en{(negotiate)} على أفضل نقل متاح: يستخدم الـ
\en{WebSocket} أولًا، وإن تعذّر يسقط \en{(fallback)} إلى \en{Server-Sent Events}
ثم \en{Long Polling}. كودك لا يتغيّر إطلاقًا مهما كان النقل تحته. زيادةً على ذلك،
يدير \en{SignalR} تعريف كل اتصال \en{(ConnectionId)}، وتجميع الاتصالات في
مجموعات، وإعادة الاتصال بعد القطع.
\end{idea}

القلب المفاهيميّ في \en{SignalR} اسمه الـ \en{Hub}: صنف \en{C\#} على الخادم يعمل
كنقطة نداء ثنائية الاتجاه. العميل يقدر يستدعي دوالّ هذا الصنف (مثل
\code{JoinConversation})، والخادم يقدر يدفع أحداثًا مسمّاة للعميل (مثل
\code{"notification"}). كأنّه هاتف: كلا الطرفين يرفع السماعة ويتكلّم.

\begin{keyp}[نقطة مهمة: الـ \en{Hub} ثنائيّ الاتجاه]
لا تتخيّل الـ \en{Hub} كـ \en{controller} يردّ على الطلبات فقط. هو قناة في
الاتجاهين: صعودًا العميل ينادي دوالّ الـ \en{Hub}، ونزولًا
الخادم يدفع أحداثًا عبر \code{Clients} إلى عميل أو مجموعة عملاء.
\end{keyp}

في \en{ScholarPath} ثلاثة \en{hubs} مربوطة على مسارات ثابتة: \code{/hubs/chat}
و\code{/hubs/notifications} و\code{/hubs/community}.

% ─────────────────────────────────────────────────────────────
\subsection{دورة حياة الاتصال والهُويّة من الـ \en{JWT}}

لكل \en{hub} دالّتان تُستدعيان تلقائيًّا: \code{OnConnectedAsync} لحظة اتصال
عميل جديد، و\code{OnDisconnectedAsync} لحظة انقطاعه. هذا المكان المثاليّ لضمّ
الاتصال إلى مجموعاته أو إعلان حالة الحضور.

لكن قبل أي شيء: مَن هذا المتّصل؟ كل \en{hubs} النظام ترث من صنف أساس واحد يفرض
التوثيق. هنشوف الملف الحقيقيّ.

\begin{english}
\begin{lstlisting}[language={[Sharp]C}]
[Authorize]                              // only authenticated connections
public abstract class AuthenticatedHub : Hub
{
    public Task<string> Ping() =>
        Task.FromResult($"pong from {GetType().Name}");
}

public sealed class NotificationHub : AuthenticatedHub
{
    public override async Task OnConnectedAsync()
    {
        // AddToGroup MUST be awaited (see the last section on the race)
        if (Context.UserIdentifier is not null)
        {
            await Groups.AddToGroupAsync(
                Context.ConnectionId, $"user:{Context.UserIdentifier}");
        }
        await base.OnConnectedAsync();
    }
}
\end{lstlisting}
\end{english}

الشرح سطرًا بسطر لأهم الأسطر:
\begin{itemize}
\item \code{[Authorize]} فوق \code{AuthenticatedHub}: بوّابة توثيق على مستوى
  الاتصال كلّه. أي عميل يحاول فتح قناة بلا \en{JWT} صالح يُرفَض في المصافحة نفسها،
  فلا يصل إلى أي دالّة.
\item \code{Context.UserIdentifier}: هوية المتّصل مستخرَجة تلقائيًّا من الـ
  \en{JWT} (من مطالبة المعرّف). هذه القيمة موثوقة لأن الـ \en{token} موقَّع؛
  العميل لا يستطيع تزويرها.
\item \code{Groups.AddToGroupAsync(Context.ConnectionId, "user:\{id\}")}: بمجرّد
  الاتصال، نضمّ هذا الاتصال إلى مجموعته الشخصية \code{user:\{id\}}. بعد هذا السطر
  يصبح العميل جاهزًا لاستقبال أي إشعار يُدفَع إلى تلك المجموعة.
\end{itemize}

\begin{secnote}
التوثيق على مستوى الـ \en{hub} فارق أمنيّ جوهريّ: القناة الحيّة ليست ثغرة خلفية
تتجاوز التوثيق. نفس الـ \en{JWT} الذي يحمي الـ \en{REST API} يحمي الـ
\en{WebSocket}، والهُويّة تأتي من الـ \en{token} لا من أي معطى يرسله العميل.
\end{secnote}

الرسم التالي يلخّص دورة الحياة من الاتصال إلى الجاهزية للاستقبال:

\begin{center}\begin{LTR}
\begin{tikzpicture}[node distance=0.42cm]
  \node[wbox] (s1) {Client connects: withUrl("/hubs/notifications") + JWT};
  \node[wbox,below=0.42cm of s1] (s2) {[Authorize] gate: verify JWT, set Context.UserIdentifier};
  \node[wbox,below=0.42cm of s2] (s3) {OnConnectedAsync()};
  \node[wbox,below=0.42cm of s3] (s4) {await Groups.AddToGroupAsync(ConnectionId, "user:id")};
  \node[wbox,below=0.42cm of s4,fill=spLightTeal] (s5) {جاهز لاستقبال أي push فوريًّا};
  \draw[arr] (s1)--(s2); \draw[arr] (s2)--(s3);
  \draw[arr] (s3)--(s4); \draw[arr] (s4)--(s5);
\end{tikzpicture}
\end{LTR}\end{center}

% ─────────────────────────────────────────────────────────────
\subsection{المجموعات والتوجيه \en{(Groups \& Routing)}}

المجموعة \en{(Group)} في \en{SignalR} مجرّد اسم نصّي يجمع تحته عدّة اتصالات.
الدفع يصير للاسم لا للاتصال الفرديّ، وهذا يحلّ مشكلتين: المستخدم الواحد قد يفتح
عدّة تبويبات (عدّة اتصالات، مجموعة واحدة \code{user:\{id\}})، والحدث الواحد قد
يخصّ عدّة مستخدمين (مثل محادثة أو تصنيف منتدى).

\en{ScholarPath} يستعمل ثلاث عائلات أسماء واضحة:

\begin{center}
\begin{tabular}{L{4cm} L{8cm}}
\toprule
{\color{spPrimary}المجموعة} & {\color{spPrimary}المعنى} \\
\midrule
\code{user:\{id\}} & كل اتصالات مستخدم واحد — قناته الشخصية للإشعارات \\
\code{conversation:\{id\}} & طرفا محادثة شات — رسائل ومؤشّر الكتابة \\
\code{forum-category:\{slug\}} & متابعو تصنيف في المنتدى — مواضيع جديدة \\
\code{forum-post:\{id\}} & مشاهدو موضوع واحد — الردود الجديدة \\
\bottomrule
\end{tabular}
\end{center}

أما التوجيه فيتمّ عبر خصائص \code{Clients} داخل الـ \en{hub}:
\code{Clients.Group(name)} للدفع لمجموعة كاملة، و\code{Clients.User(id)} لمستخدم
بعينه، و\code{Clients.Others} لكل المتّصلين ما عدا المُرسِل، و
\code{Clients.OthersInGroup(name)} لكل أعضاء المجموعة ما عدا المُرسِل — مثاليّ
لمؤشّر «يكتب الآن» حتى لا يرى المرسِل نفسه يكتب.

خلّينا نشوف \code{ChatHub} حقيقيًّا: كيف يضمّ اتصالًا لمحادثة، وكيف يبثّ الكتابة،
والأهمّ: بوّابة تمنع التنصّت.

\begin{english}
\begin{lstlisting}[language={[Sharp]C}]
public async Task JoinConversation(string conversationId)
{
    await EnsureParticipantAsync(conversationId);          // IDOR gate
    await Groups.AddToGroupAsync(
        Context.ConnectionId, $"conversation:{conversationId}");
}

public async Task TypingStart(string conversationId)
{
    await EnsureParticipantAsync(conversationId);
    await Clients.OthersInGroup($"conversation:{conversationId}")
        .SendAsync("TypingStart", conversationId, Context.UserIdentifier);
}

private async Task EnsureParticipantAsync(string conversationId)
{
    // the JWT caller must be one of the two participants
    // ... parse ids ...
    var isParticipant = await db.Conversations.AsNoTracking()
        .AnyAsync(c => c.Id == convId
            && (c.ParticipantOneId == userId || c.ParticipantTwoId == userId));
    if (!isParticipant)
        throw new HubException("You are not a participant.");
}
\end{lstlisting}
\end{english}

الشرح سطرًا بسطر لأهم الأسطر:
\begin{itemize}
\item \code{JoinConversation}: العميل ينادي هذه الدالّة صعودًا حين يفتح محادثة.
  لاحظ أنّها تستدعي \code{EnsureParticipantAsync} قبل الضمّ للمجموعة.
\item \code{Clients.OthersInGroup(...)}: مؤشّر الكتابة يُبثّ لكل أعضاء
  \code{conversation:\{id\}} ما عدا المُرسِل — فلا أرى نفسي «أكتب».
\item \code{EnsureParticipantAsync}: البوّابة الأمنية. تتأكّد من قاعدة البيانات
  أنّ صاحب الـ \en{JWT} أحد طرفَي المحادثة فعلًا؛ وإلا ترمي \code{HubException}.
\end{itemize}

\begin{spotlight}[ركّز هنا لإبهار اللجنة: بوّابة \en{IDOR} على القناة الحيّة]
اسم المجموعة \code{conversation:\{id\}} قابل للتخمين. لولا
\code{EnsureParticipantAsync} لاستطاع أي مستخدم موثَّق أن ينادي
\code{JoinConversation} بمعرّف محادثة غيره ويتنصّت على رسائل خاصّة تُبثّ لتلك
المجموعة. الدرس: التوجيه بالمجموعات لا يكفي وحده كأمان — لا بدّ من بوّابة
تتحقّق أنّك طرفٌ فعليّ قبل الانضمام. نفس المنطق مطبَّق في \code{CommunityHub} على
التصنيفات والمواضيع (تصنيف نشِط، وموضوع مرئيّ غير محذوف).
\end{spotlight}

% ─────────────────────────────────────────────────────────────
\subsection{أين المشتركون؟ على الخادم عبر \en{IHubContext}، وعلى العميل عبر \code{.on}}

سؤال جوهريّ: الدفع رأيناه داخل الـ \en{hub}. لكن الإشعار يُولَد عادةً في
مكان بعيد — بعد نجاح دفعة، أو قبول حجز. كيف ندفع من خارج الـ \en{hub}؟ الجواب:
\code{IHubContext<THub>}، وهو مقبض للـ \en{hub} تحقنه أي خدمة عبر حقن التبعيّات.

هذا بالضبط ما يفعله \code{NotificationDispatcher}: يحفظ صفوف الإشعار في قاعدة
البيانات، ثم يدفع النسخة الحيّة عبر \code{IHubContext} إلى مجموعة المستخدم.

\begin{english}
\begin{lstlisting}[language={[Sharp]C}]
public sealed class NotificationDispatcher(
    IHubContext<NotificationHub> hub, /* ... */) : INotificationDispatcher
{
    // ... persist a Notification row per channel first ...

    private async Task DeliverAsync(Notification row, /* ... */)
    {
        // InApp: real-time toast to the recipient's personal group
        if (!silenced)
            await hub.Clients.Group($"user:{row.RecipientUserId}")
                .SendAsync("notification", new
                {
                    id = row.Id,
                    type = row.Type.ToString(),
                    row.TitleAr, row.BodyAr, row.DeepLink,
                }, ct);
    }
}
\end{lstlisting}
\end{english}

الشرح: \code{IHubContext<NotificationHub>} محقون في الباني \en{(constructor)}؛
الخدمة لا ترث من الـ \en{hub} ولا تعرف تفاصيله، فقط تدفع عبره. السطر
\code{hub.Clients.Group("user:\{id\}").SendAsync("notification", ...)} يبثّ حدثًا
اسمه \code{"notification"} لمجموعة المستخدم، فيصل لكل تبويباته المفتوحة دفعةً
واحدة. لاحظ الحقول العربية \code{TitleAr} و\code{BodyAr}: الحمولة ثنائية اللغة.

أما «المشترك» على الطرف الآخر، فيجلس على العميل ويربط اسم الحدث بدالّة معالجة عبر
\code{connection.on("event", handler)}. هنشوف كود الواجهة الحقيقيّ (\en{TypeScript}):

\begin{english}
\begin{lstlisting}[language=Java]
const connection = new signalR.HubConnectionBuilder()
  .withUrl(apiBase + "/hubs/notifications", {
    // current token on every (re)connect — survives token refresh
    accessTokenFactory: () =>
      useAuthStore.getState().tokens?.accessToken ?? "",
    // offer all transports so it works where WebSockets are blocked
    transport:
      signalR.HttpTransportType.WebSockets |
      signalR.HttpTransportType.ServerSentEvents |
      signalR.HttpTransportType.LongPolling,
  })
  .withAutomaticReconnect()
  .build();

// subscriber: fires whenever the server pushes "notification"
connection.on("notification", (payload) => {
  toast(payload.titleAr, { description: payload.bodyAr });
});
await connection.start();
\end{lstlisting}
\end{english}

الشرح سطرًا بسطر لأهم الأسطر:
\begin{itemize}
\item \code{withUrl(apiBase + "/hubs/notifications")}: نفس مسار الـ \en{hub}
  المربوط على الخادم؛ الـ \en{apiBase} يتبع الـ \en{API} المنشور على \en{Azure}.
\item \code{accessTokenFactory}: يُقرأ الـ \en{token} الحاليّ من المخزن في كل
  اتصال أو إعادة اتصال، فلا ينكسر الاتصال بعد تدوير الـ \en{token}.
\item \code{transport: WebSockets | ServerSentEvents | LongPolling}: نعرض كل
  وسائل النقل ونُبقي خطوة التفاوض، فلو تعذّر الـ \en{WebSocket} يسقط تلقائيًّا.
\item \code{connection.on("notification", ...)}: هنا المشترك. أي دفعة يبثّها
  \code{NotificationDispatcher} بهذا الاسم تُشغّل هذه الدالّة فورًا فيظهر الـ
  \en{toast}.
\end{itemize}

الرسم التالي يجمع الصورة: حدث واحد على الخادم يتفرّع \en{(fan-out)} إلى كل عملاء
المجموعة.

\begin{center}\begin{LTR}
\begin{tikzpicture}[node distance=0.7cm and 1.25cm]
  \node[srcbox,text width=2.7cm] (evt) {NotificationDispatcher\\(server event)};
  \node[tbox,right=1.25cm of evt,text width=2.7cm] (ctx) {IHubContext\\Clients.Group(...)};
  \node[gbox,right=1.25cm of ctx,text width=2.1cm] (grp) {group\\user:id};
  \node[actor,above right=0.35cm and 1.25cm of grp] (c1) {browser tab};
  \node[actor,right=1.25cm of grp] (c2) {second tab};
  \node[actor,below right=0.35cm and 1.25cm of grp] (c3) {phone};
  \draw[arr] (evt) -- node[lbl,above]{SendAsync} (ctx);
  \draw[arr] (ctx) -- (grp);
  \draw[arr] (grp) -- (c1);
  \draw[arr] (grp) -- (c2);
  \draw[arr] (grp) -- (c3);
\end{tikzpicture}
\end{LTR}\end{center}

\begin{spotlight}[ركّز هنا لإبهار اللجنة: اتصال واحد دائم يخدم كل التحديثات]
كل مستخدم يفتح اتصالًا حيًّا واحدًا لكل \en{hub}، ويبقى مفتوحًا طوال
الجلسة. الإشعارات، والشات، ومؤشّر الكتابة، والحضور، وردود المنتدى — كلها تسري في
نفس القناة الدائمة. صفر \en{polling}، صفر تأخير، وحِمل أدنى بكثير من سؤال
الخادم كل ثانية. هذا هو الفرق بين واجهة «تحسّها حيّة» وأخرى تحتاج تحديثًا يدويًّا.
\end{spotlight}

% ─────────────────────────────────────────────────────────────
\subsection{الـ \en{hubs} الثلاثة: خريطة الجروبات والأحداث المدفوعة}

نلخّص المعمار كلّه في جدول واحد: كل \en{hub}، ومساره، ومجموعاته، وأحداثه المدفوعة
نزولًا إلى العميل.

\begin{center}
\begin{tabular}{L{2.9cm} L{3.9cm} L{5.2cm}}
\toprule
{\color{spPrimary}الـ \en{Hub}} & {\color{spPrimary}المجموعات} & {\color{spPrimary}أحداث تُدفَع للعميل} \\
\midrule
\code{NotificationHub}\newline \code{/hubs/notifications} & \code{user:\{id\}} &
\code{"notification"} (إشعار حيّ لكل تبويبات المستخدم) \\
\code{ChatHub}\newline \code{/hubs/chat} & \code{conversation:\{id\}} &
\code{"MessageReceived"}، \code{"TypingStart"}، \code{"TypingStop"}،
\code{"UserOnline"}، \code{"UserOffline"} \\
\code{CommunityHub}\newline \code{/hubs/community} & \code{forum-category:\{slug\}}،
\code{forum-post:\{id\}} & \code{"ReplyCreated"} وتحديثات المواضيع الحيّة \\
\bottomrule
\end{tabular}
\end{center}

\begin{note}
لاحظ الاتساق: الأحداث النازلة أسماء نصّية (\code{"MessageReceived"}) يشترك فيها
العميل عبر \code{.on}، والدوالّ الصاعدة (\code{JoinConversation}) يناديها العميل
عبر \code{invoke}. اسم الحدث على الطرفين يجب أن يتطابق حرفيًّا؛ خطأ حرف واحد =
مشترك لا يسمع شيئًا.
\end{note}

% ─────────────────────────────────────────────────────────────
\subsection{تفصيلة دقيقة: لماذا \code{await} على \code{AddToGroupAsync}؟}

عُد لكود \code{NotificationHub.OnConnectedAsync}: الضمّ للمجموعة مسبوق بـ
\code{await}. لماذا لا نتركه «أطلِق وانسَ» \en{(fire-and-forget)}؟ لأن هذا يفتح
باب \en{race condition} خبيث.

\begin{qa}[سؤال متوقّع: ماذا يحدث لو لم نضع \code{await}؟]
تخيّل السيناريو: العميل اتصل، وبدأ ضمّه للمجموعة \code{user:\{id\}} في الخلفية بلا
انتظار. في نفس اللحظة أطلق حدثٌ آخر إشعارًا فوريًّا لتلك المجموعة. لو سبق الإشعارُ
اكتمالَ الانضمام، يُدفَع لمجموعة لا يوجد بها هذا الاتصال بعد — فيضيع
الإشعار الحيّ. المستخدم لن يراه إلا بعد إعادة اتصال لاحقة، وربما يظهر عندها كأنّه
«مكرَّر». وضع \code{await} يضمن أن الاتصال صار عضوًا كامل العضوية قبل أن
يعود التحكّم، فلا نافذة سباق.
\end{qa}

هذه التفصيلة الصغيرة درسٌ عام: في الكود غير المتزامن، ترتيب اكتمال العمليّات ليس
تفصيلًا تجميليًّا — قد يكون الفرق بين إشعار يصل وإشعار يضيع.

% ─────────────────────────────────────────────────────────────
\begin{summary}[الخلاصة في سطور]
\begin{itemize}
\item الـ \en{HTTP} العاديّ طلب-وردّ ثم يُغلَق؛ لِعرض التحديثات نحتاج \en{polling}
  مكلفًا وبطيئًا لأن الخادم لا يبدأ الحديث.
\item الـ \en{WebSocket} يبدأ بمصافحة \en{(HTTP Upgrade)} تُرقّي الاتصال إلى قناة
  دائمة ثنائية الاتجاه، فيدفع الخادم \en{(push)} متى شاء بلا سؤال.
\item \en{SignalR} تجريد فوق الـ \en{WebSocket}: يفاوض على أفضل نقل ويسقط للأقدم
  عند الحاجة، ويدير الاتصالات والمجموعات. القلب هو الـ \en{Hub}: قناة ثنائية —
  العميل ينادي دوالّه، والخادم يدفع أحداثه.
\item التوثيق عبر \code{[Authorize]} على الـ \en{hub}، والهُويّة من الـ \en{JWT}
  عبر \code{Context.UserIdentifier} (لا يمكن تزويرها).
\item التوجيه بالمجموعات: \code{user:\{id\}}، \code{conversation:\{id\}}،
  \code{forum-category:\{slug\}}؛ والدفع عبر \code{Clients.Group} / \code{User} /
  \code{Others} / \code{OthersInGroup}.
\item المشتركون: على العميل عبر \code{connection.on("event")}، وعلى الخادم من خارج
  الـ \en{hub} عبر \code{IHubContext} (كما في \code{NotificationDispatcher} و
  \code{ChatRealtimeNotifier}).
\item بوّابة \en{IDOR} (لازم تكون طرفًا) تمنع التنصّت على قنوات الغير، و
  \code{await} على \code{AddToGroupAsync} يمنع سباقًا يفوّت إشعارًا.
\end{itemize}
\end{summary}

\clearpage
% ch23-hangfire-jobs.tex — الفصل: Hangfire والمهامّ الخلفيّة (Background Jobs)
\section{\en{Hangfire} والمهامّ الخلفيّة (\en{Background Jobs})}

في الفصل ده هنفهم حاجة بسيطة بس مهمّة جدًّا لأي نظام حقيقيّ: مش كل شغل لازم يتعمل
«وإحنا ماسكين المستخدم منتظر». في شغل تقيل أو مؤجَّل الأحسن نعمله في الخلفيّة، بعيدًا عن
الطلب اللي جاي من المتصفّح. هنعرف ليه نحتاج ده، وإيه هي أداة \en{Hangfire} اللي بنستعملها،
وإيه المهامّ الحقيقيّة اللي شغّالة في \en{ScholarPath}، وإزاي المهمّة الواحدة بتمشي خطوة
بخطوة من لحظة ما نطلبها لحد ما تخلص. القارئ يفترض إنه فاهم من قبل فكرة الطبقات
(\en{layers})، والـ\en{ports}، والـ\en{idempotency} — وهنعتمد عليهم من غير ما نعيد شرحهم
بعمق.

% ═══════════════════════════════════════════════════════════════════════
\subsection{ليه نحتاج مهامّ خلفيّة أصلًا؟}

خلّينا نتخيّل إن طالب ضغط زرّ «قدِّم على المنحة». في الخلفيّة لازم يحصل: نحفظ الطلب،
ونبعت له إيميل تأكيد، ويمكن نحدّث بعض الإحصاءات. لو عملنا كل ده جوّه الطلب الواحد،
الطالب هيفضل يبصّ في شاشة بتلفّ لحد ما الإيميل يتبعت والإحصاء يتحسب. وده بطيء وبيتعطّل
لو خدمة الإيميل زعلت.

الفكرة إن الشغل نوعين. في شغل «متزامن» (\en{synchronous}) لازم يتعمل جوّه الطلب لأن الردّ
يعتمد عليه — زي التحقّق من الصلاحيّة وحفظ الطلب نفسه. وفي شغل «خلفيّ» (\en{background})
ممكن يتأجّل لأن المستخدم مش محتاج ينتظره — زي إرسال الإيميل، أو معالجة تسجيل فيديو
الجلسة، أو تجهيز تقرير. النوع التاني بنرميه لطابور خلفيّ ونردّ على المستخدم فورًا.

\begin{idea}[الفكرة: افصل الثقيل عن الطلب]
القاعدة البسيطة: لو المستخدم مش محتاج نتيجة العمليّة عشان يكمّل، متخلّيهوش يستنّاها.
احفظ اللي لازم يتحفظ فورًا، وارمي الباقي (الإيميل، الفيديو، التقرير) للخلفيّة، وردّ عليه
بسرعة.
\end{idea}

وفي كمان نوع تالت مهمّ: شغل «مجدوَل» (\en{scheduled}) بيشتغل من نفسه على مواعيد، من غير
ما حد يضغط أي زرّ. زي تذكيرات ليليّة باقتراب مواعيد المنح، أو كنس دوريّ للجلسات اللي حصل
فيها غياب (\en{no-show}). ده شغل مالوش «طلب مستخدم» أصلًا — محتاج حاجة تصحّيه بنفسها كل
فترة.

\begin{center}\begin{LTR}
\begin{tikzpicture}[node distance=0.9cm and 1.3cm]
  \node[actor] (user) {User};
  \node[pbox, right=of user] (api) {API handler};
  \node[tbox, right=of api] (resp) {fast 202 reply};
  \node[sbox, below=1.2cm of api] (enq) {enqueue};
  \node[dbbox, right=of enq] (store) {job storage};
  \node[pbox, right=of store] (worker) {worker: email / video};

  \draw[arr] (user) -- node[lbl,above]{request} (api);
  \draw[arr] (api) -- node[lbl,above]{fast} (resp);
  \draw[arr] (resp.south) to[bend left=25] node[lbl,right]{back} (user.east);
  \draw[arrd] (api) -- node[lbl,right]{offload heavy} (enq);
  \draw[arr] (enq) -- (store);
  \draw[arr] (store) -- (worker);
\end{tikzpicture}
\end{LTR}\end{center}

الرسمة فوق بتقول القصّة كلها: الطلب بيدخل، الـ\en{API} بيعمل اللازم بسرعة ويردّ (السهم
العلويّ للمستخدم)، وفي نفس الوقت بيرمي الشغل التقيل لطابور خلفيّ (السهم المتقطّع لأسفل)،
والعامل (\en{worker}) بيلتقطه ويشتغل عليه بعيدًا. المستخدم مستنّاش الجزء التقيل خالص.

\begin{spotlight}[ركّز هنا لإبهار اللجنة: ردّ سريع للمستخدم]
أهمّ رسالة تقولها للّجنة: «إحنا بنردّ على المستخدم بسرعة لأن الشغل التقيل بيتأجّل
للخلفيّة». تقديم طلب منحة بيرجع للطالب في أجزاء من الثانية، بينما إرسال الإيميل ومعالجة
الفيديو وتجهيز التقارير بتحصل بعيدًا عن الطلب. النتيجة: واجهة سريعة لا تتعطّل لو خدمة
خارجيّة (زي بريد أو \en{Stripe}) بطُؤت لحظيًّا.
\end{spotlight}

% ═══════════════════════════════════════════════════════════════════════
\subsection{يعني إيه Hangfire؟}

\en{Hangfire} مكتبة لإدارة المهامّ الخلفيّة في \en{.NET}. بدل ما نبني طابور بإيدينا، هي
بتدّينا أربع حاجات جاهزة:

\begin{itemize}
  \item طابور مدعوم بتخزين دائم — المهامّ مش بتتخزّن في الذاكرة (اللي بتضيع لو
  السيرفر وقع)، لكن في جداول جوّه قاعدة البيانات (\en{SQL Server} عندنا). يعني لو السيرفر
  اتعمل له \en{restart}، المهامّ لسه موجودة وهتكمّل.
  \item عمّال (\en{workers}) — خيوط بتفضل تسأل الطابور: «في شغل؟» وتلتقطه وتنفّذه.
  \item جدولة — تقدر تجدول مهمّة تشتغل بعد شويّة، أو تتكرّر بجدول \en{cron}.
  \item لوحة تحكّم (\en{dashboard}) — صفحة بتوريك المهامّ الشغّالة والفاشلة
  والمعادة، ومحميّة عندنا بصلاحيّة \en{Admin} بس.
\end{itemize}

وفي تلات طرق أساسيّة نطلب بيها مهمّة:

\begin{center}
\begin{tabular}{L{3.4cm} L{11.2cm}}
\toprule
النوع & المعنى \\
\midrule
\en{enqueue} (\en{fire-and-forget}) & نفّذها دلوقتي في الخلفيّة، وانسَها — زي إرسال إيميل بعد حفظ الطلب. \\
\en{schedule} & نفّذها مرّة واحدة بعد وقت محدّد (مثلًا بعد ساعة). \\
\en{recurring} & كرّرها على جدول \en{cron} — زي كل 15 دقيقة أو كل ليلة. \\
\bottomrule
\end{tabular}
\end{center}

وأهمّ ميزة: لو المهمّة فشلت (استثناء)، \en{Hangfire} بيعيد محاولتها تلقائيًّا كذا مرّة
بفواصل متزايدة. وده بيفتح سؤال مهمّ: لو المهمّة اتنفّذت نصّها وفشلت وأُعيدت، هتتعمل مرّتين؟
عشان كده لازم كل مهمّة تكون \en{idempotent} — يعني تنفيذها مرّتين بنفس المدخلات يدّي نفس
النتيجة من غير ضرر (لا إيميل مكرّر ولا دفعة مزدوجة). هنرجع للنقطة دي بالتفصيل في آخر الفصل.

\begin{keyp}[نقطة مهمة: التخزين هو اللي بيخلّي الطابور ينجو من السقوط]
لأن \en{Hangfire} بيكتب المهامّ في قاعدة البيانات مش في الذاكرة، أي مهمّة اتحطّت في الطابور
بتفضل موجودة حتى لو السيرفر وقع أو اتعمل له \en{deploy}. ده الفرق بين طابور «هشّ» في
الذاكرة وطابور «دائم» يعتمد عليه.
\end{keyp}

الكود ده بيبيّن إزاي بنركّب \en{Hangfire} على تخزين \en{SQL Server}. لاحظ إن كل ده متعلّق
بمفتاح إعداد — لو الميزة مقفولة، مفيش تخزين ولا عمّال أصلًا.

\begin{english}
\begin{lstlisting}[language={[Sharp]C}]
builder.Services.AddHangfire(cfg =>
{
    cfg.UseSimpleAssemblyNameTypeSerializer()
       .UseRecommendedSerializerSettings();

    if (!string.IsNullOrWhiteSpace(connectionString))
    {
        cfg.UseSqlServerStorage(connectionString, new SqlServerStorageOptions
        {
            CommandBatchMaxTimeout = TimeSpan.FromMinutes(5),
            QueuePollInterval = TimeSpan.FromSeconds(30),
        });
    }
    // ...
});

builder.Services.AddHangfireServer();   // start the workers
\end{lstlisting}
\end{english}

شرح سريع لأهمّ الأسطر: \code{UseSqlServerStorage} بيقول لـ\en{Hangfire} يخزّن جداوله في
نفس قاعدة البيانات (نفس \code{connectionString}). و\code{QueuePollInterval} هو كل قدّ إيه
العامل يسأل الطابور عن شغل جديد (كل 30 ثانية). و\code{AddHangfireServer} هو اللي بيشغّل
العمّال اللي بيلتقطوا المهامّ فعلًا. لو التخزين مش متاح، بيرجع لتخزين ذاكرة مؤقّت
(\en{memory storage}) للتطوير.

% ═══════════════════════════════════════════════════════════════════════
\subsection{المهامّ الحقيقيّة في ScholarPath}

كل المهامّ المتكرّرة متحكَّم فيها بمفتاح واحد: \code{Hangfire:Enabled}. لو \code{false}
(القيمة الافتراضيّة في التطوير)، مفيش عمّال بيتشغّلوا ولا مهامّ بتتجدول خالص — ده بيخلّي
التشغيل المحلّي خفيف. في الإنتاج بنشغّله. الجدول ده بيلخّص أهمّ المهامّ الفعليّة المسجَّلة،
وجدولها، وبتعمل إيه:

\begin{center}
\begin{tabular}{L{3.7cm} L{2.7cm} L{7.6cm}}
\toprule
المهمّة & الجدول & بتعمل إيه \\
\midrule
\en{deadline-reminders} & يوميًّا 09:00 & تذكّر الطلاب باقتراب موعد المنح المحفوظة أو المتقدَّم لها. \\
\en{draft-reminders} & يوميًّا 10:00 & تنكز مَن عنده مسودّة طلب لسه ما اتبعتتش. \\
\en{scholarship-auto-close} & كل ساعة & تقفل المنح تلقائيًّا بعد ما موعدها يعدّي. \\
\en{meeting-no-show-sweep} & كل 15 دقيقة & ترصد غياب أحد طرفَي الجلسة وتفتح \en{NoShowReport}. \\
\en{booking-reminders} & كل 15 دقيقة & تبعت تذكير قبل الجلسة بـ 24 ساعة وبساعة. \\
\en{stripe-payouts} & يوميًّا 02:00 & تصرف دفعات \en{Stripe} للمستشارين على دفعة مجمّعة. \\
\en{data-export-sweep} & كل ساعة & تنفّذ طلبات تصدير بيانات \en{GDPR}. \\
\en{data-delete-sweep} & يوميًّا 03:00 & تنفّذ طلبات حذف بيانات \en{GDPR}. \\
\en{redaction-audit-sampling} & شهريًّا (اليوم 1، 02:00) & تسحب عيّنة تدقيق لتنقية بيانات \en{PII}. \\
\en{integrity-check} & يوميًّا 04:00 & فحص سلامة يكشف الصفوف اليتيمة وغير المتّسقة. \\
\bottomrule
\end{tabular}
\end{center}

\begin{note}
فيه شغل خلفيّ مش متكرّر بجدول، لكنه بيتحطّ في الطابور «عند الطلب» — أشهر مثال إعادة بناء
قاعدة المعرفة للذكاء الاصطناعيّ (\en{RebuildKnowledgeBase}). ده بيتفعّل لمّا الأدمن يضغط
الزرّ (أو مرّة عند إقلاع النظام) بدل ما يتكرّر كل ليلة، لأنه شغل تقيل مش محتاج يتعمل كل يوم.
\end{note}

الكود ده هو تسجيل المهامّ المتكرّرة نفسه. لاحظ إن كله جوّه شرط \code{Hangfire:Enabled}،
وإن كل مهمّة ليها اسم ثابت وجدول \en{cron}:

\begin{english}
\begin{lstlisting}[language={[Sharp]C}]
if (hangfireOpts.Enabled)
{
    var recurring = app.Services.GetRequiredService<IRecurringJobManager>();

    recurring.AddOrUpdate<IMeetingNoShowSweepJob>(
        "meeting-no-show-sweep", j => j.RunAsync(CancellationToken.None), "*/15 * * * *");
    recurring.AddOrUpdate<IStripePayoutJob>(
        "stripe-payouts", j => j.RunAsync(CancellationToken.None), Cron.Daily(2));   // 02:00 UTC
    recurring.AddOrUpdate<IScholarshipAutoCloseJob>(
        "scholarship-auto-close", j => j.RunAsync(CancellationToken.None), Cron.Hourly);
    recurring.AddOrUpdate<IDeadlineReminderJob>(
        "deadline-reminders", j => j.RunAsync(CancellationToken.None), Cron.Daily(9));
    // ... the rest of the jobs follow the same pattern
}
\end{lstlisting}
\end{english}

شرح: \code{AddOrUpdate} بيسجّل مهمّة متكرّرة باسم ثابت — لو موجودة بيحدّثها بدل ما يعمل
نسخة تانية. الوسيط التالت هو الجدول: \code{"*/15 * * * *"} تعبير \en{cron} معناه «كل 15
دقيقة»، و\code{Cron.Daily(2)} معناها الساعة 02:00 بتوقيت \en{UTC}. لاحظ إن كل مهمّة بتُسجَّل
عبر \en{interface} زي \code{IStripePayoutJob} — وده مهمّ هنرجع له في آخر الفصل. الرسمة دي
بتوزّع المهامّ على «خطّ زمنيّ» حسب تكرارها:

\begin{center}\begin{LTR}
\begin{tikzpicture}[node distance=0.55cm and 0.5cm]
  \node[gbox] (t1) {every 15 min};
  \node[wbox, right=of t1] (r1) {no-show sweep, booking reminders, booking completion, session expiry};
  \node[gbox, below=of t1] (t2) {hourly};
  \node[wbox, right=of t2] (r2) {scholarship auto-close, GDPR export sweep};
  \node[gbox, below=of t2] (t3) {nightly (02:00 - 10:00)};
  \node[wbox, right=of t3] (r3) {Stripe payouts, GDPR delete, integrity check, deadline and draft reminders};
  \node[gbox, below=of t3] (t4) {monthly (1st, 02:00)};
  \node[wbox, right=of t4] (r4) {PII redaction audit sampling};

  \draw[arr] (t1) -- (t2);
  \draw[arr] (t2) -- (t3);
  \draw[arr] (t3) -- (t4);
\end{tikzpicture}
\end{LTR}\end{center}

الأعمدة اليسرى (بنفسجيّ) هي الإيقاع: من الأسرع (كل 15 دقيقة) للأبطأ (شهريًّا)، والسهم النازل
بيمثّل مرور الزمن. الأشرطة الرماديّة الطويلة بتقول المهامّ اللي بتشتغل في كل إيقاع.

% ═══════════════════════════════════════════════════════════════════════
\subsection{كيف تشتغل مهمّة خطوة بخطوة}

خلّينا نتتبّع رحلة مهمّة واحدة من أوّلها لآخرها. الرحلة بسيطة: حاجة بتحطّ المهمّة في الطابور
(\en{enqueue})، التخزين بيسجّلها، عامل بيلتقطها، بينفّذها، وبعدين يا إمّا تنجح يا إمّا تفشل
وتُعاد تلقائيًّا.

\begin{center}\begin{LTR}
\begin{tikzpicture}[node distance=0.8cm and 1.05cm]
  \node[sbox] (enq) {enqueue};
  \node[dbbox, right=of enq] (q) {storage / queue};
  \node[pbox, right=of q] (w) {worker executes};
  \node[dbox, right=of w] (dec) {success?};
  \node[tbox, right=of dec] (done) {done};
  \node[fbox, below=1.1cm of dec] (retry) {retry (back-off)};

  \draw[arr] (enq) -- (q);
  \draw[arr] (q) -- (w);
  \draw[arr] (w) -- (dec);
  \draw[arr] (dec) -- node[lbl,above]{yes} (done);
  \draw[arr] (dec) -- node[lbl,right]{no} (retry);
  \draw[arrd] (retry) -| node[lbl,below]{re-queued} (q);
\end{tikzpicture}
\end{LTR}\end{center}

المهمّ في الرسمة هو الفرع النازل: لو فشلت المهمّة، مبتضيعش — بترجع للطابور (السهم المتقطّع)
عشان تتحاول تاني بعد فترة انتظار متزايدة. ودي بالظبط النقطة اللي بتخلّي الـ\en{idempotency}
شرطًا مش رفاهيّة.

الكود ده مقتطف من مهمّة تذكيرات المواعيد الحقيقيّة (\en{DeadlineReminderJob}). شوف إزاي
بتبعت التذكير — وركّز في السطر بتاع \code{idempotencyKey}:

\begin{english}
\begin{lstlisting}[language={[Sharp]C}]
foreach (var recipientId in recipients)
{
    await notifications.DispatchAsync(
        recipientId,
        NotificationType.ApplicationDeadlineApproaching,
        new NotificationParams
        {
            TitleEn = scholarship.TitleEn,
            TitleAr = scholarship.TitleAr,
            Count = daysLeft,
        },
        deepLink: $"/student/scholarships/{scholarship.Id}",
        idempotencyKey: $"deadline-reminder:{scholarship.Id:N}:{recipientId:N}:{deadlineStamp}",
        ct).ConfigureAwait(false);
}
\end{lstlisting}
\end{english}

شرح: المهمّة بتلفّ على كل مستلم وتبعت له تذكير عبر \code{DispatchAsync}. السطر الأهمّ هو
\code{idempotencyKey}: مفتاح مبنيّ من رقم المنحة، ورقم المستلم، وتاريخ الموعد
(\code{deadlineStamp}). يعني نفس التذكير لنفس الشخص لنفس الموعد ليه بصمة واحدة ثابتة. لو
المهمّة اليوميّة اشتغلت مرّتين بالغلط، أو أُعيدت بعد فشل، طبقة الإشعارات هتشوف إن البصمة دي
اتبعتت قبل كده فمش هتبعت تاني. تذكير واحد بالظبط لكل موعد — لا تكرار.

% ═══════════════════════════════════════════════════════════════════════
\subsection{Idempotency وإعادة المحاولة: لا صرف مزدوج}

دلوقتي نجمع الخيطين: إعادة المحاولة التلقائيّة نعمة، لكنها بتفتح خطر التكرار. الحلّ إن كل
مهمّة تتصمّم بحيث لو اتنفّذت مرّتين متعملش ضرر. في التذكيرات عملنا ده بمفتاح
\code{idempotencyKey}. في صرف الفلوس الموضوع أخطر بكتير — لو صرفنا لمستشار مرّتين، دي فلوس
حقيقيّة ضاعت.

عشان كده مهمّة صرف \en{Stripe} بتتحمي بطبقتين. الأولى قفل بيمنع تشغيلين متوازيين لنفس
المهمّة في نفس اللحظة. شوف إزاي القفل متحطّ على الـ\en{interface} نفسه:

\begin{english}
\begin{lstlisting}[language={[Sharp]C}]
public interface IStripePayoutJob
{
    // The lock must live HERE on the interface method — Hangfire reads
    // job filters off the registered (interface) method, not the impl.
    [DisableConcurrentExecution(timeoutInSeconds: 30 * 60)]
    Task RunAsync(CancellationToken ct);
}
\end{lstlisting}
\end{english}

شرح: \code{DisableConcurrentExecution} بيقول لـ\en{Hangfire} «مسموح بتشغيل واحد بس لهذه
المهمّة في المرّة». الطبقة التانية (في التنفيذ) هي إن كل دفعة بتتعلّم عليها علامة قبل الصرف
باستخدام رمز تزامن (\en{RowVersion}) على صفّ الدفعة — يعني حتى لو اتنين خادمين حاولوا، واحد
بس هينجح في «حجز» الدفعة والصرف. الاتنين مع بعض: لا صرف مزدوج مهما تكرّرت المهمّة.

\begin{spotlight}[ركّز هنا لإبهار اللجنة: \en{idempotency} + إعادة المحاولة = صفر تكرار]
هنا بيت القصيد اللي بيفرّق نظام هاوي عن نظام جادّ. لأن المهامّ بتتعاد تلقائيًّا عند الفشل،
صمّمنا كل مهمّة تكون \en{idempotent}: التذكيرات ببصمة \code{idempotencyKey} فمفيش إشعار
مكرّر، والدفعات بقفل \code{DisableConcurrentExecution} + رمز التزامن على الصفّ فمفيش صرف
مزدوج ولا فلوس ضايعة. الرسالة للّجنة: «إعادة المحاولة تحمينا من الفشل، والـ\en{idempotency}
تحمينا من إعادة المحاولة نفسها».
\end{spotlight}

\begin{secnote}
لوحة تحكّم \en{Hangfire} على المسار \code{/hangfire} مش مفتوحة للكل — محميّة بفلتر صلاحيّة
يسمح لـ\en{Admin} بس، وبتتفعّل بمفتاح منفصل \code{DashboardEnabled}. يعني حتى الأداة اللي
بتوريك قلب النظام الخلفيّ نفسها متقفلة خلف الصلاحيّات.
\end{secnote}

\begin{qa}[سؤال متوقّع: طب لو نفس المهمّة اليوميّة اشتغلت مرّتين؟]
مفيش مشكلة. المهامّ عندنا مصمّمة على مبدأ «آمنة عند التكرار». التذكيرات هتتصفّى بالبصمة،
والدفعات هيتصدّها القفل ورمز التزامن، والإغلاق التلقائيّ للمنح بيعتمد على الحالة الحاليّة
(\en{Open} فقط) فتشغيله تاني على منحة اتقفلت خلاص مبيعملش حاجة. ده بالظبط تعريف الـ
\en{idempotency}.
\end{qa}

\begin{why}
ليه نتكبّد كل ده بدل ما نشغّل الشغل جوّه الطلب؟ لأن الفصل ده بيدّينا تلات مكاسب مع بعض:
واجهة سريعة (المستخدم مستنّاش الثقيل)، وصمود (المهمّة في التخزين تنجو من سقوط السيرفر
وتُعاد لو فشلت)، وأتمتة (شغل بيحصل من نفسه من غير تدخّل بشريّ زي إغلاق المنح والتذكيرات).
\end{why}

% ═══════════════════════════════════════════════════════════════════════
\begin{summary}[الخلاصة في سطور]
\begin{itemize}
  \item الشغل التقيل أو المؤجَّل مايصحّش يعطّل ردّ المستخدم — بنرميه للخلفيّة ونردّ فورًا.
  \item \en{Hangfire} بيدّينا طابورًا مدعومًا بتخزين دائم في \en{SQL Server}، وعمّالًا
  يلتقطوا المهامّ، وجدولة، ولوحة تحكّم محميّة بـ\en{Admin}.
  \item تلات طرق للطلب: \en{enqueue} فوري، و\en{schedule} مؤجَّل، و\en{recurring} بجدول
  \en{cron}. وكله متحكَّم فيه بمفتاح \code{Hangfire:Enabled}.
  \item عندنا مهامّ حقيقيّة: تذكيرات المواعيد والمسودّات، الإغلاق التلقائيّ للمنح، كنس
  الـ\en{no-show}، صرف \en{Stripe}، تصدير وحذف \en{GDPR}، عيّنات تدقيق تنقية \en{PII}،
  وفحوص السلامة.
  \item دورة حياة المهمّة: \en{enqueue} $\rightarrow$ تخزين $\rightarrow$ عامل يلتقط
  $\rightarrow$ تنفيذ $\rightarrow$ نجاح أو إعادة محاولة تلقائيّة.
  \item لأن الفشل بيسبّب إعادة محاولة، كل مهمّة لازم تكون \en{idempotent}: بصمة
  \code{idempotencyKey} للإشعارات، وقفل + رمز تزامن للدفعات — فلا إشعار مكرّر ولا صرف
  مزدوج.
\end{itemize}
\end{summary}

\clearpage
% ch24-community-block
\section{المجتمع والحظر: كيف يعمل بالكامل}

مجتمع \en{ScholarPath} هو المساحة التي يتناقش فيها الطلاب: يسألون عن المنح، يشاركون
تجاربهم، ويصوّتون على أفضل الإجابات. لكن أيّ مساحة مفتوحة تحتاج قواعد تحكمها: مَن
يكتب؟ كيف نمنع المحتوى الضارّ؟ وماذا يحدث حين يريد طالبٌ ألّا يرى شخصًا آخر إطلاقًا؟

في هذا الفصل هنفهم ثلاثة أشياء متتالية. أولًا: كيانات المجتمع — الفئات والمنشورات
والردود والتصويت والإبلاغ والإشارات المرجعيّة — وكيف تترابط. ثانيًا: دورة حياة منشور
كاملة، من لحظة كتابته حتى إخفائه تلقائيًّا لو أبلغ عنه عدد كافٍ من الناس. ثالثًا: نظام
الحظر (\en{Block}) — وهو سؤالنا الأهمّ: كيف يُنشَأ الحظر، والأهمّ: أين يُنفَّذ؟
هنكتشف أنّ الحظر لا يُطبَّق في مكان واحد، بل دفاعيًّا في كل مسار: فحص عند الكتابة،
وفلتر عند القراءة.

نفترض أنّك تعرف من الفصول السابقة فكرة \en{CQRS} (كل عملية إمّا أمر أو استعلام)،
و\en{soft delete} (الحذف المنطقيّ بعلَم بدل المحو الفعليّ)، و\en{SignalR} وقنوات البثّ
الحيّ (\en{hubs} ومجموعات المشتركين)، وتنقية \en{HTML} قبل التخزين. سنشير إليها بسطر
عند الحاجة فقط.

% ====================================================================
\subsection{كيانات المجتمع: الخريطة الكاملة}

المجتمع مبنيّ على حفنة كيانات بسيطة. أذكى قرار في التصميم أنّ المنشور والردّ نفس
الكيان: كلاهما \code{ForumPost}. الفرق الوحيد حقلٌ اسمه \code{ParentPostId}: لو كان
\code{null} فهذا منشور جذر (\en{thread})، ولو كان يحمل مُعرّفًا فهذا ردّ على منشور آخر.
هذا يوفّر علينا جدولًا ومنطقًا مكرّرًا.

\begin{center}
\begin{tabular}{L{3.2cm} L{9cm}}
\toprule
الكيان & الدور \\
\midrule
\en{ForumCategory} & فئة نقاش (منح، تأشيرات، سكن...) لها \code{Slug} وعلَم \code{IsActive}. \\
\en{ForumPost} & منشور جذر أو ردّ (يفرّقهما \code{ParentPostId}). يحمل عدّادات مخزّنة و\code{ModerationStatus}. \\
\en{ForumVote} & صوت واحد لكل مستخدم على منشور: \code{Up} أو \code{Down}. \\
\en{ForumFlag} & بلاغ من مستخدم عن منشور، بسبب و\code{IsValid} يقدر الأدمن يبطله. \\
\en{ForumBookmark} & إشارة مرجعيّة تربط طالبًا بمنشور حفظه ليرجع إليه. \\
\bottomrule
\end{tabular}
\end{center}

المقتطف التالي من ملف \code{Community.cs} يبيّن قلب \code{ForumPost} — لاحظ حقلين مصيريّين:
\code{ModerationStatus} (حالة الإشراف) و\code{IsAutoHidden} (الإخفاء التلقائيّ):

\begin{english}
\begin{lstlisting}[language={[Sharp]C}]
public class ForumPost : AuditableEntity, ISoftDeletable
{
    public Guid AuthorId { get; set; }
    public Guid? CategoryId { get; set; }
    public Guid? ParentPostId { get; set; } // null = root thread; else a reply

    public PostModerationStatus ModerationStatus { get; set; }
        = PostModerationStatus.Visible;

    // Cached aggregates (kept in sync on write, cheap to read)
    public int UpvoteCount { get; set; }
    public int FlagCount { get; set; }
    public int ReplyCount { get; set; }

    public bool IsAutoHidden { get; set; } // hidden by the threshold rule
    public bool IsDeleted { get; set; }    // soft delete
    // ... navigations: Replies, Votes, Flags, Bookmarks ...
}
\end{lstlisting}
\end{english}

شرح سريع لأهمّ الأسطر:
\begin{itemize}
\item \code{ParentPostId} هو المفتاح الذكيّ: منشور وردّ في نفس الجدول، والفرق حقل واحد.
\item \code{ModerationStatus} حالة إشراف من نوع \code{PostModerationStatus}: \code{Visible}
  أو \code{Removed} (حذفه الأدمن) أو \code{PendingReview} (بانتظار مراجعة) أو \code{Hidden}.
\item \code{IsAutoHidden} علَمٌ منفصل: صحيح حين تتخطّى البلاغاتُ العتبةَ فيُخفى المنشور
  تلقائيًّا قبل أن يتدخّل أيّ إنسان.
\item العدّادات \code{UpvoteCount} و\code{FlagCount} و\code{ReplyCount} مخزّنة مسبقًا
  (\en{cached aggregates}): نحدّثها لحظة الكتابة حتى تكون القراءة رخيصة بلا عدٍّ متكرّر.
\item \code{IsDeleted} هو الـ \en{soft delete}: المنشور المحذوف يبقى صفًّا في القاعدة لكن
  يختفي من كل الاستعلامات.
\end{itemize}

الرسم التالي يجمع العلاقات: الفئة تحوي منشورات، والمنشور له ردود، وحوله التصويت والإبلاغ
والإشارة المرجعيّة.

\begin{center}
\begin{LTR}
\begin{tikzpicture}[node distance=0.6cm and 1.4cm]
  \node[dbbox] (cat) {Category};
  \node[dbbox,right=1.7cm of cat] (post) {Post\\(root)};
  \node[dbbox,right=1.7cm of post] (reply) {Reply};
  \node[sbox,text width=2.0cm,above=0.95cm of post] (vote) {Vote\\Up / Down};
  \node[sbox,text width=2.0cm,below=0.95cm of post] (flag) {Flag};
  \node[sbox,text width=2.0cm,above=0.95cm of reply] (book) {Bookmark};
  \draw[arr] (cat) -- (post) node[lbl,midway,above] {posts};
  \draw[arr] (post) -- (reply) node[lbl,midway,above] {ParentPostId};
  \draw[arrb] (post) -- (vote);
  \draw[arrb] (post) -- (flag);
  \draw[arrb] (post) -- (book);
\end{tikzpicture}
\end{LTR}
\end{center}

% ====================================================================
\subsection{دورة حياة منشور: من الكتابة إلى العرض}

المنشور يبدأ حياته بأمر \code{CreatePostCommand}. أوّل خطوة حارسة: مَن يكتب؟ فقط الطالب.
المستشارون والجهات لا ينتجون محتوى هنا. ثانيًا: كل نصٍّ يمرّ على مُنقِّي \en{HTML} قبل
التخزين، حتى لا يُحقَن سكربت خبيث. المقتطف من \code{CreatePostCommand} يبيّن الحارسَين
معًا، ثم إطلاق حدث البثّ الحيّ:

\begin{english}
\begin{lstlisting}[language={[Sharp]C}]
if (!currentUser.IsInRole("Student"))
    throw new ForbiddenAccessException("Only students can create community posts.");

var sanitizer = new Ganss.Xss.HtmlSanitizer();
var titleEn = sanitizer.Sanitize(request.TitleEn);
var bodyEn  = sanitizer.Sanitize(request.BodyEn);
// ... TitleAr / BodyAr sanitized the same way (posts are bilingual) ...

var post = new ForumPost { AuthorId = authorId, CategoryId = request.CategoryId,
    TitleEn = titleEn, BodyEn = bodyEn /* ... */ };
db.ForumPosts.Add(post);

// Raise BEFORE save so the dispatcher (runs after SaveChanges) broadcasts it.
post.RaiseDomainEvent(new ForumPostCreatedEvent(post.Id, authorId, request.CategoryId));
await db.SaveChangesAsync(ct);
\end{lstlisting}
\end{english}

الأسطر المهمّة:
\begin{itemize}
\item الفحص \code{IsInRole("Student")} يرفض أيّ دور آخر برمي \code{ForbiddenAccessException}.
\item \code{HtmlSanitizer} ينظّف العنوان والنصّ باللغتين — منشورات المجتمع ثنائيّة اللغة
  مثل المنح تمامًا.
\item \code{RaiseDomainEvent} يُطلَق قبل \code{SaveChangesAsync} عمدًا: الموزّع
  يلتقط الحدث بعد الحفظ فيبثّ الإشعار الحيّ. لو أطلقناه بعد الحفظ لضاع البثّ.
\end{itemize}

بعد أن يُعرَض المنشور، يتفاعل الطلاب معه: يردّون، ويصوّتون. التصويت أمرٌ لطيف اسمه
\code{ToggleVoteCommand}: نفس الزرّ يضيف الصوت أو يزيله أو يقلبه. لاحظ قاعدتين: لا تصوّت
على محتواك، وعدّاد المنشور يتحدّث في نفس اللحظة.

\begin{english}
\begin{lstlisting}[language={[Sharp]C}]
if (post.AuthorId == currentUser.UserId)
    throw new ConflictException("You can't vote on your own content.");

var existingVote = post.Votes.FirstOrDefault(v => v.UserId == currentUser.UserId);
if (existingVote != null && existingVote.VoteType == request.VoteType)
{
    db.ForumVotes.Remove(existingVote);            // same vote again -> toggle off
    if (request.VoteType == VoteType.Up) post.UpvoteCount--; else post.DownvoteCount--;
}
// ... else: swap direction, or add a brand-new vote ...
await db.SaveChangesAsync(ct);
\end{lstlisting}
\end{english}

الفكرة أنّ العدّاد المخزّن \code{UpvoteCount} يتغيّر مع كل عمليّة، فلا نضطرّ لعدّ الأصوات
عند كل قراءة. والقيد الفريد على (المنشور، المستخدم) في القاعدة يضمن صوتًا واحدًا لكل شخص
مهما تكرّرت الطلبات.

\begin{keyp}[نقطة مهمة: العدّادات المخزّنة]
\code{UpvoteCount} و\code{FlagCount} و\code{ReplyCount} كلّها أرقام محفوظة على صفّ المنشور،
نحدّثها لحظة الكتابة. النتيجة: صفحة المجتمع تقرأ عددًا جاهزًا بدل \code{COUNT} على آلاف
الأصوات في كل تحميل. الثمن أن نبقيها متّسقة في كل أمر — وهو ما نراه في التصويت والإبلاغ.
\end{keyp}

% ====================================================================
\subsection{الإبلاغ والإخفاء التلقائيّ: حماية بلا تدخّل بشريّ فوريّ}

هنا تكمن أذكى قاعدة في المجتمع. حين يرى طالبٌ منشورًا مسيئًا، يبلّغ عنه بأمر
\code{FlagPostCommand}. لكن بلاغًا واحدًا لا يكفي لإخفاء منشور — قد يكون كيديًّا. القاعدة:
لو أبلغ ثلاثة أشخاص مختلفين (بلاغات صالحة متمايزة) عن نفس المنشور، يُخفى تلقائيًّا
ويُحوَّل إلى \code{PendingReview} بانتظار قرار الأدمن.

المقتطف من \code{FlagPostCommand} هو قلب القاعدة:

\begin{english}
\begin{lstlisting}[language={[Sharp]C}]
// distinct valid flaggers already persisted, + our own new one
var distinctValidFlags = priorDistinctValidFlags + 1;
post.FlagCount++;

// Auto-hide rule: 3 distinct valid flags route the post to moderation.
if (distinctValidFlags >= 3 && !post.IsAutoHidden)
{
    post.IsAutoHidden = true;
    post.AutoHiddenAt = DateTimeOffset.UtcNow;
    post.ModerationStatus = PostModerationStatus.PendingReview;
    post.RaiseDomainEvent(new PostAutoHiddenEvent(post.Id, distinctValidFlags));
}
await db.SaveChangesAsync(ct);
\end{lstlisting}
\end{english}

ملاحظات دقيقة:
\begin{itemize}
\item نعدّ المُبلِّغين المتمايزين (\code{Distinct}) لا عدد البلاغات — حتى لا يخفي
  شخصٌ واحد منشورًا بتكرار البلاغ. والفحص المسبق على تكرار المُبلِّغ نفسه يمنع ذلك أصلًا.
\item عند بلوغ الثلاثة يُرفع علَم \code{IsAutoHidden} وتتحوّل الحالة إلى \code{PendingReview}،
  ويُطلَق حدث \code{PostAutoHiddenEvent} لإخطار الأدمن.
\item الشرط \code{!post.IsAutoHidden} يمنع الإخفاء وإطلاق الحدث مرّتين لو أُعيدت المحاولة.
\end{itemize}

الرسم يلخّص القرار: هل بلغت البلاغاتُ المتمايزة الصالحة العتبةَ؟

\begin{center}
\begin{LTR}
\begin{tikzpicture}[node distance=0.6cm and 1.1cm]
  \node[pbox] (flag) {FlagPost};
  \node[sbox,text width=2.2cm,right=1.1cm of flag] (count) {distinct valid flags};
  \node[dbox,right=1.1cm of count] (dec) {flags $\geq 3$ ?};
  \node[fbox,above right=0.35cm and 1.2cm of dec] (hide) {AutoHidden\\PendingReview};
  \node[tbox,below right=0.35cm and 1.2cm of dec] (stay) {stays Visible};
  \draw[arr] (flag) -- (count);
  \draw[arr] (count) -- (dec);
  \draw[arr] (dec) -- node[lbl,above] {yes} (hide);
  \draw[arr] (dec) -- node[lbl,below] {no} (stay);
\end{tikzpicture}
\end{LTR}
\end{center}

بعد الإخفاء التلقائيّ يأتي دور الأدمن. أمامه مساران: \code{DismissPostFlags} (يقرّر أنّ
المنشور سليم فيُبطل كل البلاغات ويعيده \code{Visible})، أو \code{DeletePost} (يحذفه
منطقيًّا). المقتطف من \code{DismissPostFlags} يبيّن الرجوع للحالة الطبيعيّة:

\begin{english}
\begin{lstlisting}[language={[Sharp]C}]
foreach (var flag in post.Flags) flag.IsValid = false; // invalidate every flag
post.FlagCount = 0;
post.IsAutoHidden = false;
post.AutoHiddenAt = null;
post.ModerationStatus = PostModerationStatus.Visible;
post.ModeratedByAdminId = adminId;
post.ModeratedAt = DateTimeOffset.UtcNow;
await db.SaveChangesAsync(ct);
\end{lstlisting}
\end{english}

\begin{spotlight}[ركّز هنا لإبهار اللجنة: حماية ذاتيّة بلا انتظار]
عتبة البلاغ الثلاثيّ تجعل المجتمع يحمي نفسه. المحتوى المسيء يختفي لحظة يتّفق ثلاثة
مستخدمين مستقلّين عليه — بلا انتظار أدمن يقظ على مدار الساعة. وفي الوقت نفسه، اشتراط
ثلاثة مُبلِّغين متمايزين يمنع إساءة استخدام الزرّ لإسكات رأي واحد. إنّها موازنة
هندسيّة أنيقة بين سرعة الحماية وعدالة القرار، والقرار النهائيّ يبقى للأدمن (إبقاء أو حذف).
\end{spotlight}

% ====================================================================
\subsection{البثّ الحيّ عبر CommunityHub وبوّابات الرؤية}

حين يُنشر ردّ جديد على منشور، نريد أن يظهر فورًا لكل مَن يشاهد ذلك الموضوع دون تحديث
الصفحة. هذا يتمّ عبر \code{CommunityHub} (قناة \en{SignalR}). العميل يشترك في مجموعتين:
مجموعة الفئة (\code{JoinCategory}) ليعرف بالمنشورات الجديدة، ومجموعة المنشور
(\code{JoinPost}) ليعرف بالردود الجديدة على موضوع بعينه.

لكن الاشتراك ليس مفتوحًا للجميع. هنا تعمل بوّابة الرؤية \en{SEC-02}: لا تشترك إلّا في فئة
نشطة فعلًا، ولا في منشور مرئيّ فعلًا — غير محذوف، وغير مخفيّ تلقائيًّا، وحالته
\code{Visible}. المقتطف من \code{CommunityHub} يبيّن حارس \code{JoinPost}:

\begin{english}
\begin{lstlisting}[language={[Sharp]C}]
public async Task JoinPost(string postId)
{
    if (!Guid.TryParse(postId, out var id))
        throw new HubException("Invalid post id.");

    var visible = await db.ForumPosts.AsNoTracking()
        .AnyAsync(p => p.Id == id && !p.IsDeleted && !p.IsAutoHidden
            && p.ModerationStatus == PostModerationStatus.Visible);
    if (!visible)
        throw new HubException("Post not found.");

    await Groups.AddToGroupAsync(Context.ConnectionId, $"forum-post:{postId}");
}
\end{lstlisting}
\end{english}

لماذا نفحص الرؤية قبل الاشتراك؟ لأنّ بدون هذا الحارس يقدر أيّ مستخدم مُصادَق أن ينضمّ
لمجموعة \code{forum-post} لمنشور محذوف أو مخفيّ، فيتنصّت على ردوده الحيّة. الحارس يغلق
الباب عند لحظة الانضمام.

بعد الاشتراك، حين يُنشأ ردّ، يُطلَق حدث \code{ForumReplyCreatedEvent}، ومعالِجه ينادي
\code{CommunityRealtimeNotifier} الذي يبثّ للمجموعة فقط:

\begin{english}
\begin{lstlisting}[language={[Sharp]C}]
public Task NotifyNewReplyAsync(Guid replyId, Guid parentPostId, CancellationToken ct) =>
    hub.Clients
       .Group($"forum-post:{parentPostId}")
       .SendAsync("ReplyCreated", new { replyId, parentPostId }, ct);
\end{lstlisting}
\end{english}

الرسم يوضّح الـ \en{fan-out}: ردٌّ واحد يصل لكل المشتركين في مجموعة المنشور دفعة واحدة.

\begin{center}
\begin{LTR}
\begin{tikzpicture}[node distance=0.5cm and 1.0cm]
  \node[actor] (author) {Student reply};
  \node[pbox,right=1.0cm of author] (cmd) {CreateReply};
  \node[gbox,right=1.0cm of cmd] (hub) {CommunityHub};
  \node[dbbox,text width=2.4cm,right=1.0cm of hub] (grp) {group\\forum-post:id};
  \node[sbox,text width=1.9cm,above right=0.5cm and 1.0cm of grp] (s1) {subscriber};
  \node[sbox,text width=1.9cm,right=1.1cm of grp] (s2) {subscriber};
  \node[sbox,text width=1.9cm,below right=0.5cm and 1.0cm of grp] (s3) {subscriber};
  \draw[arr] (author) -- (cmd);
  \draw[arr] (cmd) -- node[lbl,above] {ReplyCreated} (hub);
  \draw[arr] (hub) -- (grp);
  \draw[arr] (grp) -- (s1);
  \draw[arr] (grp) -- (s2);
  \draw[arr] (grp) -- (s3);
\end{tikzpicture}
\end{LTR}
\end{center}

% ====================================================================
\subsection{الحظر: كيف يُنشَأ وأين يُنفَّذ}

نصل الآن للسؤال الأهمّ. حين يريد طالبٌ ألّا يتعامل مع مستخدم آخر، يُصدر
\code{BlockUserCommand}. هذا الأمر بسيط: ينشئ صفًّا واحدًا في جدول \code{UserBlock} يربط
الحاجب بالمحجوب.

\begin{english}
\begin{lstlisting}[language={[Sharp]C}]
if (request.UserIdToBlock == currentUser.UserId)
    throw new ConflictException("You cannot block yourself.");
// ... return early if the block row already exists (idempotent) ...

var block = new UserBlock
{
    BlockerId = currentUser.UserId ?? throw new ForbiddenAccessException(),
    BlockedUserId = request.UserIdToBlock,
    Reason = request.Reason
};
db.UserBlocks.Add(block);
await db.SaveChangesAsync(ct);
\end{lstlisting}
\end{english}

لكنّ إنشاء الصفّ نصف القصّة. الصفّ مجرّد بيان بأنّ هناك حظرًا؛ السؤال الحقيقيّ:
مَن يقرأ هذا الصفّ ويطبّقه؟ الجواب المهمّ: التنفيذ موزّع دفاعيًّا على كل مسار، وينقسم
لنوعين — فحص عند الكتابة، وفلتر عند القراءة.

\begin{idea}[الفكرة: تنفيذ الحظر = كتابة + قراءة]
عند \en{الكتابة}: كل أمر يُنتج تفاعلًا (إرسال رسالة، كتابة ردّ) يفحص جدول \code{UserBlock}
قبل التنفيذ ويرفض لو كان هناك حظر. عند \en{القراءة}: كل استعلام يعرض محتوى شخص آخر
يمرّ على فلتر يُسقط محتوى المحجوبين. لو غاب أيّ طرف، تسرّب المحتوى — لذا نطبّقه في الطرفين.
\end{idea}

المسار الأوّل (كتابة) — أمر \code{SendMessageCommand} يفحص الحظر قبل حفظ الرسالة:

\begin{english}
\begin{lstlisting}[language={[Sharp]C}]
var isBlocked = await db.UserBlocks.AnyAsync(b =>
    (b.BlockerId == recipientId          && b.BlockedUserId == currentUser.UserId) ||
    (b.BlockerId == currentUser.UserId   && b.BlockedUserId == recipientId), ct);

if (isBlocked)
    throw new ConflictException("Cannot send message. User is blocked.");
\end{lstlisting}
\end{english}

لاحظ أنّ الفحص ثنائيّ الاتجاه: يرفض سواء حجبتُ أنا الطرف الآخر أو حجبني هو. نفس
المنطق في المجتمع: أمر \code{CreateReplyCommand} ينادي \code{AreBlockedAsync} فيرفض الردّ
لو كان بينهما حظر في أيّ اتجاه.

المسار الثاني (قراءة) — الفلتر المشترك \code{CommunityBlockFilter.NotBlockedWith}. هو تعبير
\code{Expression} واحد تستخدمه كل استعلامات المجتمع، حتى لا تنحرف دلالة الحظر بين مكان
وآخر:

\begin{english}
\begin{lstlisting}[language={[Sharp]C}]
public static Expression<Func<ForumPost, bool>> NotBlockedWith(
    IApplicationDbContext db, Guid viewerId) =>
    p => !db.UserBlocks.Any(b =>
        (b.BlockerId == viewerId    && b.BlockedUserId == p.AuthorId) ||
        (b.BlockerId == p.AuthorId  && b.BlockedUserId == viewerId));
\end{lstlisting}
\end{english}

هذا التعبير يعني: «أرِني المنشور فقط إن لم يكن صاحبه في علاقة حظر معي في أيّ اتّجاه».
تستدعيه \code{GetPostsQuery} على قائمة المنشورات، و\code{GetMyBookmarksQuery} يطبّق نفس
الشرط بيده فيُسقِط من إشاراتي المرجعيّة أيّ منشور صار صاحبه محجوبًا. الجدول التالي يجمع
نقاط التنفيذ:

\begin{center}
\begin{tabular}{L{3.4cm} L{2.6cm} L{6cm}}
\toprule
الموضع & النوع & السلوك \\
\midrule
\en{SendMessage} & كتابة & يرفض إرسال الرسالة. \\
\en{CreateReply} & كتابة & يرفض الردّ (\code{AreBlockedAsync}). \\
\en{GetPosts} & قراءة & يُخفي منشورات المحجوب من القائمة. \\
\en{GetMyBookmarks} & قراءة & يُسقط المحجوب من الإشارات المرجعيّة. \\
\en{GetConversations} & قراءة & يعلّم مَن حجبتُ لعرض زرّ إلغاء الحظر. \\
\bottomrule
\end{tabular}
\end{center}

الرسم يجسّد الفكرة: أمرٌ واحد ينشئ صفًّا واحدًا، ثم يُقرأ في مسارين متعامدين.

\begin{center}
\begin{LTR}
\begin{tikzpicture}[node distance=0.55cm and 1.1cm]
  \node[actor] (u) {Blocker};
  \node[pbox,right=1.1cm of u] (cmd) {BlockUser};
  \node[dbbox,text width=2.3cm,right=1.1cm of cmd] (row) {UserBlock row};
  \node[fbox,text width=2.9cm,above right=0.45cm and 1.3cm of row] (write) {write path\\SendMessage rejected};
  \node[tbox,text width=2.9cm,below right=0.45cm and 1.3cm of row] (read) {read path\\BlockFilter hides};
  \draw[arr] (u) -- (cmd);
  \draw[arr] (cmd) -- (row);
  \draw[arr] (row) -- node[lbl,above] {on write} (write);
  \draw[arr] (row) -- node[lbl,below] {on read} (read);
\end{tikzpicture}
\end{LTR}
\end{center}

\begin{secnote}
الفحص ثنائيّ الاتجاه مقصود ومهمّ. لو كان الفلتر يفحص «مَن حجبتُ أنا» فقط، لتمكّن مَن حجبني
من متابعة محتواي والردّ عليّ. توحيد الدلالة في تعبير \code{NotBlockedWith} واحد يضمن أنّ
كل استعلام يطبّق النسخة المتبادلة نفسها، فلا تنشأ ثغرة من اختلاف التطبيق بين شاشة وأخرى.
\end{secnote}

\begin{spotlight}[ركّز هنا لإبهار اللجنة: الحظر دفاعيّ في كل مسار]
لو سألتك اللجنة «أين يُنفَّذ الحظر؟» فالإجابة القويّة ليست «في مكان واحد»، بل: في كل
مسار كتابة وقراءة على حِدة. أمر واحد (\code{BlockUser}) يكتب صفًّا واحدًا، لكن هذا الصفّ
يُقرأ في \code{SendMessage} و\code{CreateReply} (كتابة) وفي \code{GetPosts} و\code{GetMyBookmarks}
و\code{GetConversations} (قراءة). هذا نمط \en{defense-in-depth}: لا نعتمد على طبقة واحدة
تحرس كل شيء، بل نغلق الباب عند كل نقطة تماسّ. النتيجة: لا مسار واحد يتسرّب منه محتوى المحجوب.
\end{spotlight}

\begin{qa}[سؤال متوقّع: لماذا لا تحرس مكانًا مركزيًّا واحدًا فقط؟]
لأنّ المسارات مختلفة الطبيعة: الكتابة تحتاج رفضًا صريحًا برسالة خطأ، والقراءة تحتاج
إسقاطًا صامتًا من النتائج. لا توجد نقطة مركزيّة واحدة يمرّ منها الاثنان. لذلك نضع
الفحص عند كل أمر كتابة، والفلتر (\code{NotBlockedWith}) عند كل استعلام قراءة — مع توحيد
الدلالة المتبادلة في مكان واحد حتى لا تتكرّر ولا تنحرف.
\end{qa}

% ====================================================================
\begin{summary}[الخلاصة في سطور]
\begin{itemize}
\item المنشور والردّ نفس الكيان \code{ForumPost}؛ يفرّقهما \code{ParentPostId}.
\item حول المنشور: تصويت (\code{ForumVote}) وبلاغ (\code{ForumFlag}) وإشارة مرجعيّة
  (\code{ForumBookmark})، وعدّادات مخزّنة تجعل القراءة رخيصة.
\item الإنشاء يحرسه دورٌ (طالب فقط) وتنقية \en{HTML}، ويُطلق حدثًا للبثّ الحيّ قبل الحفظ.
\item الإبلاغ يحمي المجتمع ذاتيًّا: ثلاثة مُبلِّغين متمايزين يُخفون المنشور تلقائيًّا
  (\code{IsAutoHidden}) ويحوّلونه \code{PendingReview}، والقرار النهائيّ للأدمن.
\item البثّ الحيّ عبر \code{CommunityHub} مع بوّابة رؤية \en{SEC-02}: لا اشتراك إلّا في
  فئة نشطة أو منشور مرئيّ فعلًا.
\item الحظر يُنشَأ بصفّ \code{UserBlock} واحد، لكنّه يُنفَّذ دفاعيًّا في كل مسار: فحص عند
  الكتابة (\code{SendMessage}، \code{CreateReply})، وفلتر عند القراءة (\code{NotBlockedWith})
  — وكلّه ثنائيّ الاتجاه فلا ثغرة.
\end{itemize}
\end{summary}

\clearpage
% ch25-events-notifications-audit
\section{الأحداث الدوميّة، الإشعارات، والتدقيق (\en{Audit}) — وليه بنعمل كده؟}

في الفصل ده هنجاوب على سؤال بسيط في ظاهره، عميق في باطنه: لمّا يحصل «شيء مهمّ» في
النظام — اتأكّد حجز، وصلت رسالة، اتقدّم طالب على منحة — إزاي بنخلّي بقيّة النظام
«يعرف» ويردّ، من غير ما نكدّس كل ردود الفعل دي جوّه المنطق الأساسيّ؟ الجواب اسمه
\en{Domain Event} (الحدث الدوميّ). وهنكمّل بعدها على أخته التوأم في المساءلة:
\en{Audit} (التدقيق) — إزاي بنسجّل «مين عمل إيه وامتى» على كل فعل حسّاس.

القارئ هنا يُفترَض إنه عارف من قبل ثلاث حاجات، فمش هنعيد شرحها بعمق: \en{MediatR}
كناقل رسائل داخليّ (شفناه في الجزء الأوّل)، والإشعارات على مستوى الاستخدام —
يعني إزاي \code{INotificationDispatcher} بيوصّل الإشعار للجرس والبريد (فصل الإشعارات)،
والمعترِض على \code{SaveChanges} اللي بيلفّ عمليّة الحفظ. اللي جديد النهارده هو
الحلقة اللي بتربط دول ببعض: من لحظة وقوع الحدث لحدّ ظهور نتيجته.

\begin{idea}[الفكرة: الحدث يعلن، والمشترِكون يردّون]
بدل ما «معالِج الحجز» يعرف كل حاجة عن الإيميل و\en{SignalR} والتحليلات، هو بس
بيعلن: «حصل حجز مؤكَّد». مين يهتمّ يشترك ويردّ. ده نفس مبدأ «افصل مَن يُعلِن
عن مَن يتصرّف»، وهو اللي بيخلّي كل قطعة صغيرة، مستقلّة، وسهلة الاختبار.
\end{idea}

% ─────────────────────────────────────────────────────────────────────────────
\subsection{يعني إيه «حدث دوميّ» (\en{Domain Event})؟}

الحدث الدوميّ هو مجرّد إعلان: «حاجة ذات معنى في مجال العمل حصلت خلاص». مش أمر
(\en{command}) بيطلب تنفيذ، ولا استعلام (\en{query}) بيطلب بيانات — إعلان بحقيقة
وقعت. في \en{ScholarPath} الأحداث دي \en{records} صغيرة ثابتة (\en{immutable})،
كلّها بترث نوعًا واحدًا اسمه \code{DomainEvent}، واللي بدوره بيحقّق واجهة
\code{INotification} بتاعة \en{MediatR} — يعني \en{MediatR} يعرف ينشره لكل مشترِك:

\begin{english}
\begin{lstlisting}[language={[Sharp]C}]
// every event inherits DomainEvent -> which is a MediatR.INotification
public abstract record DomainEvent : MediatR.INotification;

public record ChatMessageReceivedEvent(
    Guid ConversationId,
    Guid MessageId,
    Guid SenderId,
    Guid RecipientId,
    string BodyPreview) : DomainEvent;

public record BookingConfirmedEvent(
    Guid BookingId, Guid StudentId, Guid ConsultantId) : DomainEvent;

public record ApplicationSubmittedEvent(
    Guid ApplicationId, Guid StudentId, Guid? ScholarshipId) : DomainEvent;
\end{lstlisting}
\end{english}

هنشوف إيه في المقتطف ده:
\begin{itemize}
\item \code{DomainEvent} نوع مجرّد (\en{abstract record}) بيرث \code{INotification}
  — دي التذكرة اللي بتخلّي \en{MediatR} قادر ينشر أيّ حدث لكل مَن اشترك فيه.
\item كل حدث بيحمل الحقائق بس، مش السلوك: \code{ChatMessageReceivedEvent}
  فيه مُعرِّف المحادثة والرسالة والمُرسِل والمستقبِل ومعاينة نصّ — لا إيميل ولا
  دفع ولا أيّ «ردّ فعل».
\item الأسماء دايمًا بصيغة الماضي: \code{BookingConfirmed}،
  \code{ApplicationSubmitted}، \code{ChatMessageReceived} — لأنها بتوصف حاجة
  حصلت فعلًا، مش حاجة عايزين تحصل.
\end{itemize}

\begin{keyp}[نقطة مهمة: حدث بيحمل حقائق، مش سلوك]
لو لقيت نفسك بتحطّ منطق إرسال بريد جوّه تعريف الحدث، فأنت خلطت الإعلان بالتصرّف.
الحدث «رسالة في زجاجة» — بيقول ما حصل وخلاص. التصرّف مكانه المشترِك، وهنشوفه دلوقتي.
\end{keyp}

% ─────────────────────────────────────────────────────────────────────────────
\subsection{بعد الحفظ فقط: ليه ننشر الحدث بعد \code{SaveChanges} الناجح؟}

في قاعدة ذهبيّة: الحدث ما يتنشرش إلّا بعد ما \code{SaveChangesAsync} ينجح
ويثبّت التغيير في قاعدة البيانات. ليه؟ لأنّ ردّ الفعل بيبني على «حقيقة». لو نشرنا
\code{ChatMessageReceivedEvent} قبل الحفظ، وبعدين الحفظ فشل أو اترجّع
(\en{rollback})، هنكون بعتنا إشعار وبريد عن رسالة مش موجودة أصلًا. ده كذب
على المستخدم، وخطأ يصعُب تتبّعه.

عشان كده الترتيب في المعالِج ثابت: احفظ الأوّل، وبعد نجاح الحفظ ابعت الدفع الحيّ،
وبعدها انشر الحدث. المقتطف التالي من \code{SendMessageCommandHandler} بيوضّح
اللحظة دي بالضبط (مقتصّ):

\begin{english}
\begin{lstlisting}[language={[Sharp]C}]
// 1) persist the truth first: the message row + conversation pointers
await db.SaveChangesAsync(ct);

// 2) live push to whoever has the conversation open now (SignalR group)
await chatNotifier.NotifyNewMessageAsync(
    conversation.Id, message.Id, message.SenderId, message.Body, message.SentAt, ct);

// 3) THEN, and only then: publish the domain event (bell + email for the absent party)
await publisher.Publish(
    new ChatMessageReceivedEvent(
        ConversationId: conversation.Id,
        MessageId:      message.Id,
        SenderId:       senderId,
        RecipientId:    request.RecipientId,
        BodyPreview:    sanitizedBody),
    ct);
\end{lstlisting}
\end{english}

هنشوف إيه:
\begin{itemize}
\item السطر \code{publisher.Publish(new ChatMessageReceivedEvent(...))} هو النقطة
  اللي بيتفرّع منها كل شيء. \code{publisher} هو \code{IPublisher} بتاع
  \en{MediatR}؛ بمجرّد ما تناديه، \en{MediatR} بيروح يدوّر على كل مشترِك في
  \code{ChatMessageReceivedEvent} ويناديه.
\item لاحظ إنّ الحفظ (\code{SaveChangesAsync}) قبله بسطور. لو رمى استثناء، السطر
  ده أصلًا ما بيتنفّذش — فمفيش إشعار عن رسالة اترجّعت.
\end{itemize}

\begin{why}
ليه ننشر بعد الحفظ مش قبله؟ عشان «الاتّساق»: كل ردّ فعل (إشعار، بريد، تحليلات)
لازم يعكس حالة حقيقيّة ثابتة في القاعدة. النشر بعد النجاح بيضمن إنّ أيّ ردّ
فعل يشوف نفس الحقيقة اللي شافها العميل. لو عكسنا الترتيب، هنعيش في احتمال دائم إننا
نبلّغ عن حاجة اتلغت — وده أسوأ من إشعار متأخّر بثوانٍ.
\end{why}

% ─────────────────────────────────────────────────────────────────────────────
\subsection{فين المشتركون؟ معالِجات الأحداث (\en{EventHandlers})}

السؤال الطبيعيّ: طيب مين بيسمع الحدث ده؟ الإجابة في مكان واحد متوقَّع: تحت
\code{Application/*/EventHandlers/*}. كل مشترِك كلاس بيحقّق \code{INotificationHandler}
لنوع حدث معيّن، وبيحطّ ردّ فعله في دالة \code{Handle}. أمثلة حيّة في النظام:

\begin{itemize}
\item \code{ChatMessageNotificationEventHandler} — بيبعت جرس/بريد لرسالة جديدة.
\item \code{BookingConfirmedEventHandler} — بيبلّغ الطالب إنّ حجزه اتقبل.
\item \code{ApplicationSubmittedNotificationEventHandler} — بيبلّغ إنّ الطلب اتقدّم.
\end{itemize}

ده جسم المشترِك بتاع الشات، مقتصًّا للجوهر:

\begin{english}
\begin{lstlisting}[language={[Sharp]C}]
public sealed class ChatMessageNotificationEventHandler(
    IApplicationDbContext db,
    INotificationDispatcher notifications,
    IChatPresenceQuery chatPresence,
    ILogger<ChatMessageNotificationEventHandler> logger)
    : INotificationHandler<ChatMessageReceivedEvent>          // subscribes to the event
{
    public async Task Handle(ChatMessageReceivedEvent notification, CancellationToken ct)
    {
        // recipient is watching chat now -> the live push already delivered it, skip the bell
        if (chatPresence.IsOnline(notification.RecipientId))
            return;
        // ... look up sender name + build the role-based deep link ...
        await notifications.DispatchAsync(
            notification.RecipientId,
            NotificationType.ChatMessageReceived,
            new NotificationParams { /* CounterpartyName, Preview */ },
            deepLink: deepLink,
            idempotencyKey: $"chat-message:{notification.MessageId:N}",
            ct);
    }
}
\end{lstlisting}
\end{english}

هنشوف إيه سطرًا-بسطر:
\begin{itemize}
\item السطر \code{: INotificationHandler<ChatMessageReceivedEvent>} هو
  الاشتراك نفسه. مجرّد ما الكلاس بيقول «أنا معالِج لهذا الحدث»، \en{MediatR}
  بيسجّله تلقائيًّا وبيناديه لمّا الحدث ينُشَر.
\item دالة \code{Handle} بتاخد الحدث و\code{CancellationToken}، وفيها ردّ الفعل
  كامل. لاحظ إنّ المعالِج ده ما يعرفش مين نشر الحدث ولا ليه — بس بيتصرّف.
\item الحارس \code{if (chatPresence.IsOnline(...)) return;} بيمنع الإزعاج
  المزدوج: لو الطرف فاتح الشات، الرسالة وصلته حيّة، فمفيش داعي لجرس.
\item آخر خطوة: \code{notifications.DispatchAsync(...)} — وهنا المشترِك بيسلّم
  الراية لـ\code{INotificationDispatcher}، اللي شرحنا آليّته بالكامل في فصل
  الإشعارات (منع تكرار \en{$\rightarrow$} نصّ باللغتين \en{$\rightarrow$} قنوات
  \en{$\rightarrow$} حفظ \en{$\rightarrow$} تسليم). المشترِك بيقول «أبلِغ»،
  والموزّع بيتكفّل بالباقي.
\end{itemize}

\begin{note}
لاحظ التغليف بـ\code{try/catch} حوالين \code{DispatchAsync} في الكود الأصليّ:
فشل قناة إشعار ما يصحّش يكسر خطّ الإرسال. لو الإيميل وقع، الرسالة اتخزّنت
خلاص والدفع الحيّ اتبعت — فنسجّل تحذيرًا ونكمّل. ردّ الفعل «إضافيّ»، مش شرطًا
لنجاح الفعل الأساسيّ.
\end{note}

% ─────────────────────────────────────────────────────────────────────────────
\subsection{ليه نعمل «حدث دوميّ» أصلًا؟}

دلوقتي السؤال الصريح: ما دام «احجز» ممكن ببساطة ينادي «ابعت إشعار» على طول، ليه
نلفّ الموضوع بحدث ومشترِك؟ الجواب في كلمة واحدة: الفصل — فصل الاهتمامات
(\en{separation of concerns}). خلّينا نعدّد المكاسب:

\begin{itemize}
\item المعالِج الأساسيّ يفضل نظيفًا: معالِج «احجز» شغلته يعمل الحجز ويحفظه
  وخلاص. ما يعرفش حاجة عن \en{SMTP} ولا \en{SignalR} ولا مقاييس التحليلات.
\item حدث واحد يشغّل عدّة ردود: من نفس \code{BookingConfirmedEvent} ممكن
  يتفرّع إشعار للطالب، وتسجيل مقياس تحليليّ، وبريد تذكير — كلٌّ في مشترِكه.
\item إضافة ردّ جديد بلا لمس المنطق: عايز تضيف «حدّث لوحة التحكّم» بكرة؟
  اكتب مشترِكًا جديدًا لنفس الحدث. مفيش سطر واحد بيتغيّر في معالِج الحجز.
\item اختبار أسهل: تختبر معالِج الحجز إنه نشر الحدث وبس؛ وتختبر كل
  مشترِك لوحده. مفيش تشابك يخلّي اختبارًا واحدًا يلمس نصف النظام.
\end{itemize}

الرسم التالي بيوضّح فكرة \en{fan-out}: نقطة نشر واحدة بتتفرّع لعدّة مشترِكين
مستقلّين. المشترِك الأوّل (الإشعار) موجود فعلًا في النظام، والباقي أمثلة على ردود
ممكن تضيفها لاحقًا بلا ما تلمس معالِج الأمر:

\begin{center}\begin{LTR}
\begin{tikzpicture}[node distance=0.55cm and 1.25cm]
  \node[sbox,text width=2.2cm] (save) {SaveChanges OK};
  \node[gbox,right=1.1cm of save,text width=3.6cm,font=\scriptsize\bfseries] (evt)
    {Publish event\\ChatMessageReceivedEvent};
  \node[srcbox,right=1.4cm of evt,text width=3.0cm] (h1)
    {Notify Handler\\bell + email};
  \node[tbox,above=0.55cm of h1,text width=3.0cm] (h2)
    {Analytics Handler\\usage metric};
  \node[fbox,below=0.55cm of h1,text width=3.0cm] (h3)
    {Digest Handler\\daily summary};
  \draw[arr] (save) -- (evt);
  \draw[arr] (evt.east) -- (h1.west);
  \draw[arr] (evt.east) -- (h2.west);
  \draw[arr] (evt.east) -- (h3.west);
  \node[lbl,below=0.08cm of evt] {one publish};
  \node[lbl,below=0.08cm of h3] {many reactions};
\end{tikzpicture}
\end{LTR}\end{center}

هنشوف إيه في الرسم: العقدة البنفسجيّة في النصّ هي نشر الحدث بعد نجاح الحفظ.
منها بتخرج ثلاثة أسهم لثلاثة مشترِكين — الأخضر (الإشعار) حقيقيّ اليوم، والأزرق
(التحليلات) والبرتقاليّ (الملخّص اليوميّ) أمثلة تُضاف بلا لمس معالِج الأمر. لاحظ:
مصدر واحد، ردود كثيرة، بلا خيط يربط بينهم.

\begin{spotlight}[ركّز هنا لإبهار اللجنة: حدث واحد، ردود كثيرة، بلا تشابك]
لو اللجنة سألت «إزاي بتضيفوا سلوكًا جديدًا من غير ما تكسروا القديم؟» — ده الجواب
النموذجيّ. النظام مبنيّ بحيث تضيف مشترِكًا جديدًا لحدث قائم فيشتغل فورًا، من
غير ما تفتح ملفّ معالِج الأمر الأصليّ أصلًا. ده تطبيق عمليّ لمبدأ «مفتوح للامتداد،
مغلق للتعديل» (\en{Open/Closed}) — وبيتحوّل من شعار في كتاب لخاصّية تلمسها في الكود.
\end{spotlight}

% ─────────────────────────────────────────────────────────────────────────────
\subsection{ربط الدائرة: من الأمر إلى الجرس}

يلّا نوصّل القطع كلّها في مسار واحد كامل. المستخدم يبعت رسالة \en{$\rightarrow$}
معالِج الأمر يحفظ \en{$\rightarrow$} ينشر الحدث \en{$\rightarrow$} المشترِك يلتقطه
\en{$\rightarrow$} المشترِك ينادي \code{DispatchAsync} \en{$\rightarrow$} يتخزّن صفّ
إشعار ويومض الجرس. الرسم ده بيتتبّع الحلقة من أوّلها لآخرها:

\begin{center}\begin{LTR}
\begin{tikzpicture}[node distance=0.7cm and 0.7cm]
  \node[actor,text width=1.7cm] (u) {User};
  \node[pbox,right=0.7cm of u,text width=2.5cm] (cmd) {SendMessage Handler};
  \node[sbox,right=0.6cm of cmd,text width=1.9cm] (save) {SaveChanges};
  \node[gbox,right=0.6cm of save,text width=2.2cm] (pub) {Publish Event};
  \node[srcbox,below=0.85cm of pub,text width=2.6cm] (eh) {EventHandler};
  \node[pbox,left=0.8cm of eh,text width=2.4cm] (disp) {DispatchAsync};
  \node[dbbox,left=0.8cm of disp,text width=2.2cm] (row) {Bell row};
  \draw[arr] (u) -- (cmd);
  \draw[arr] (cmd) -- (save);
  \draw[arr] (save) -- (pub);
  \draw[arr] (pub) -- (eh) node[lbl,right=1pt,midway] {MediatR};
  \draw[arr] (eh) -- (disp);
  \draw[arr] (disp) -- (row);
  \draw[arr] (row.west) to[out=180,in=270] node[lbl,left] {bell blinks} (u.south);
\end{tikzpicture}
\end{LTR}\end{center}

هنشوف إيه: الصفّ العلويّ هو «الفعل» — المستخدم، المعالِج، الحفظ، ثمّ النشر. من عقدة
النشر البنفسجيّة ينزل السهم لـ\en{MediatR} اللي يسلّم الحدث للمشترِك الأخضر
(\code{EventHandler}). المشترِك ينادي \code{DispatchAsync} فيتخزّن صفّ الإشعار
(العقدة ذات الحدّ المزدوج = قاعدة بيانات)، وأخيرًا يومض جرس المستخدم. لاحظ إنّ
النصف السفليّ (الردّ) منفصل تمامًا عن النصف العلويّ (الفعل) — الرابط الوحيد بينهم
هو الحدث.

\begin{qa}[سؤال متوقّع: مش ده تعقيد زيادة؟]
ممكن حدّ يقول: «كان ممكن معالِج الرسالة ينادي الإشعار على طول في سطر، وخلصنا».
صح للحظة، لكن جرّب تضيف تاني وتالت ردّ فعل: كل ردّ هيدخل جوّه معالِج الرسالة، لحدّ
ما يبقى المعالِج «يعرف» عن الإيميل والتحليلات والبريد والدفع — وده بالظبط التشابك
اللي بنهرب منه. الحدث بيدفع تكلفة صغيرة مقدّمًا (نشر + اشتراك) عشان يوفّر تكلفة
ضخمة لاحقًا (كل تعديل بيلمس ملفًّا واحدًا فقط).
\end{qa}

% ─────────────────────────────────────────────────────────────────────────────
\subsection{التدقيق (\en{Audit}): مين عمل إيه وامتى؟}

لو الأحداث بتجاوب على «إيه اللي حصل وإزاي نردّ»، فالتدقيق بيجاوب على سؤال تانيّ
تمامًا: «مين عمل الفعل الحسّاس ده، وامتى، ومن فين؟». استرداد دفعة، تغيير دور
مستخدم، حذف مستند، تعديل إعداد ماليّ — كلّها أفعال محتاجة مساءلة كاملة. ولو
حطّينا كود التسجيل ده جوّه كل معالِج، هيتكرّر وينتشر ويُنسى. الحلّ نفس فلسفة
الأحداث: افصله في طبقة اعتراض واحدة.

الآليّة مكوّنة من قطعتين متعاونتين. الأولى سِمة (\en{attribute}) اسمها
\code{[Auditable]} بتوسم الأمر اللي عايزينه يتسجّل، والتانية سلوك في خطّ أنابيب
\en{MediatR} اسمه \code{AuditBehavior} بيلفّ كل أمر: لو الأمر موسوم، يكتب صفّ تدقيق
بعد نجاح المعالِج؛ ولو مش موسوم، يعدّي بلا عمل. ده جوهر السلوك:

\begin{english}
\begin{lstlisting}[language={[Sharp]C}]
public async Task<TResponse> Handle(
    TRequest request, RequestHandlerDelegate<TResponse> next, CancellationToken ct)
{
    var response = await next(ct);          // run the real handler first

    var attr = (AuditableAttribute?)Attribute.GetCustomAttribute(
        typeof(TRequest), typeof(AuditableAttribute));
    if (attr is null) return response;              // not marked -> no-op

    var targetId = ResolveTargetId(response, request, attr.TargetIdProperty);
    var summary  = BuildSummary(attr.SummaryTemplate, targetId, response, request);
    await audit.WriteAsync(attr.Action, attr.TargetType, targetId,
        beforeJson: null, afterJson: null, summary, ct);   // write the audit row

    return response;
}
\end{lstlisting}
\end{english}

هنشوف إيه:
\begin{itemize}
\item \code{await next(ct)} بينفّذ المعالِج الأصليّ الأوّل. التدقيق بيتكتب
  بعد نجاح الفعل — نفس منطق «سجّل الحقيقة، مش النيّة».
\item \code{Attribute.GetCustomAttribute(...)} بيفحص هل الأمر موسوم بـ
  \code{[Auditable]}. لو \code{null}، السلوك بيرجّع فورًا بلا أيّ أثر — فمفيش
  حمل زائد على الأوامر العاديّة.
\item \code{ResolveTargetId} و\code{BuildSummary} بيستخرجوا «على مين» الفعل ووصفًا
  مقروءًا (زيّ «\en{Refunded payment ...}») من ردّ المعالِج أو الأمر نفسه.
\item \code{audit.WriteAsync(...)} بيكتب الصفّ. لاحظ إنّ الكتابة كلّها ملفوفة بـ
  \code{try/catch} في الأصل: فشل التدقيق ما يصحّش يلغي فعلًا ناجحًا.
\end{itemize}

والوسم نفسه سطر واحد فوق تعريف الأمر. مثال حسّاس — استرداد دفعة:

\begin{english}
\begin{lstlisting}[language={[Sharp]C}]
[Auditable(AuditAction.PaymentRefunded, "Payment",
    TargetIdProperty = nameof(PaymentId),
    SummaryTemplate = "Refunded payment {PaymentId} - amount={AmountCents}")]
public sealed record RefundPaymentCommand(
    Guid PaymentId, long? AmountCents, string? Reason = null) : IRequest<bool>;
\end{lstlisting}
\end{english}

الوسم بيقول لـ\code{AuditBehavior} ثلاث حاجات: نوع الفعل (\code{PaymentRefunded})،
نوع الكيان المستهدَف (\code{Payment})، وأيّ خاصيّة تُؤخَذ كمُعرِّف للهدف
(\code{PaymentId})، وقالب وصف بشريّ. في النظام أكتر من ثمانين أمرًا موسومًا كده —
كل فعل حسّاس تقريبًا: تسجيل الدخول والخروج، تغيير الأدوار، الموافقة والرفض، الحذف،
والاستردادات — بلا سطر تدقيق واحد داخل أيّ معالِج منهم.

طيب، الصفّ اللي بيتكتب بيحتوي إيه بالظبط؟ ده قلب \code{AuditService.WriteAsync}:

\begin{english}
\begin{lstlisting}[language={[Sharp]C}]
var entry = new AuditLog
{
    ActorUserId   = currentUser.UserId,          // WHO   (the actor)
    Action        = action,                      // WHAT  (the verb)
    TargetType    = targetType,                  // on which entity type
    TargetId      = targetId,                     // and which exact row
    BeforeJson    = beforeJson,                   // state BEFORE
    AfterJson     = afterJson,                    // state AFTER
    IpAddress     = Trim(currentUser.IpAddress, 64),      // FROM WHERE
    CorrelationId = Trim(currentUser.CorrelationId, 128), // which request
    OccurredAt    = clock.UtcNow,                 // WHEN
};
db.AuditLogs.Add(entry);   // append-only: only ever Add, never Update or Delete
await db.SaveChangesAsync(ct);
\end{lstlisting}
\end{english}

هنشوف إيه: الصفّ بيجمع كل أبعاد المساءلة في مكان واحد — مَن (\code{ActorUserId})،
ماذا (\code{Action} + \code{TargetType} + \code{TargetId})، متى
(\code{OccurredAt})، من أين (\code{IpAddress} و\code{UserAgent})، وأيّ طلب
ربط الأحداث ببعضها (\code{CorrelationId}). وفيه كمان عمودان \code{BeforeJson} و
\code{AfterJson} لالتقاط لقطة الحالة قبل الفعل وبعده لمّا نحتاج التفصيل ده. الرسم
التالي بيلخّص مسار الاعتراض:

\begin{center}\begin{LTR}
\begin{tikzpicture}[node distance=0.7cm and 0.8cm]
  \node[actor,text width=1.8cm] (adm) {Admin};
  \node[pbox,right=0.75cm of adm,text width=2.6cm] (cmd)
    {[Auditable]\\RefundPayment};
  \node[sbox,right=0.7cm of cmd,text width=2.2cm] (hnd) {Handler runs\\refund};
  \node[gbox,right=0.7cm of hnd,text width=2.4cm] (beh) {AuditBehavior\\after next()};
  \node[dbbox,below=0.9cm of beh,text width=2.3cm] (log) {AuditLog\\append-only};
  \node[sbox,left=1.0cm of log,text width=4.0cm,align=left,font=\scriptsize]
    (fields) {who | what | when\\IpAddress | CorrelationId\\Before / After JSON};
  \draw[arr] (adm) -- (cmd);
  \draw[arr] (cmd) -- (hnd);
  \draw[arr] (hnd) -- (beh);
  \draw[arr] (beh) -- (log) node[lbl,right=1pt,midway] {write};
  \draw[arr] (log) -- (fields) node[lbl,above,midway] {columns};
\end{tikzpicture}
\end{LTR}\end{center}

هنشوف إيه في الرسم: المدير يبعت أمر استرداد موسومًا بـ\code{[Auditable]}. المعالِج
ينفّذ الاسترداد الحقيقيّ، وبعد نجاحه بيلتقط \code{AuditBehavior} الأمر ويكتب صفّ
\code{AuditLog} (الحدّ المزدوج = قاعدة بيانات) بكل الأعمدة على اليسار. لاحظ إنّ
الصفّ \en{append-only}: بنضيفه ومنلمسوش تاني.

\begin{secnote}
ليه \en{append-only} (إضافة فقط، لا تعديل ولا حذف)؟ لأنّ سجلّ التدقيق ما ينفعش
يبقى قابلًا للتزوير — وإلّا مقصده كلّه يسقط. مسار الكود بيضيف صفوفًا بس؛ مفيش أمر
واحد في النظام بيعدّل أو بيمسح صفّ \code{AuditLog}، وجانب القراءة (استعلام السجلّ)
بيقرأ بس. النتيجة: أثر ثابت لا رجعة فيه — حتى المدير اللي عمل الفعل ما يقدرش
يمسح دليله.
\end{secnote}

\begin{spotlight}[ركّز هنا لإبهار اللجنة: كل فعل حسّاس متتبَّع، وبلا رجعة]
اللجنة بتحبّ سؤال «إزاي تعرفوا مين عمل إيه؟». الجواب: أيّ أمر حسّاس بيتوسم بسطر
واحد \code{[Auditable]}، فيلتقطه اعتراض مركزيّ ويكتب صفًّا يحمل الفاعل والفعل
والوقت و\en{IP} ومُعرِّف الطلب، وقابليّة لقطة قبل/بعد — كلّه \en{append-only}. يعني:
مساءلة كاملة، من غير ما نبعثر كود تسجيل في ثمانين معالِجًا، ومن غير ما نثق في إنّ
حدًّا فاكر يكتبه. ده الفرق بين «عندنا \en{logs}» و«عندنا نظام مساءلة».
\end{spotlight}

% ─────────────────────────────────────────────────────────────────────────────
\subsection{الخلاصة}

\begin{summary}[الخلاصة في سطور]
\begin{itemize}
\item الحدث الدوميّ إعلانٌ بأنّ «شيئًا ذا معنى حصل» (\code{BookingConfirmed}،
  \code{ChatMessageReceived}، \code{ApplicationSubmitted})؛ \en{record} ثابت يرث
  \code{DomainEvent} وهو \code{INotification} في \en{MediatR}.
\item يُنشَر بعد \code{SaveChanges} الناجح فقط — عشان مانبلّغش عن تغيير
  اترجّع؛ السطر \code{publisher.Publish(new ...Event(...))} هو نقطة التفرّع.
\item المشترِكون تحت \code{Application/*/EventHandlers/*}؛ كل واحد
  \code{INotificationHandler} لحدث، بيحطّ ردّ فعله في \code{Handle} وبينادي
  \code{DispatchAsync} — بلا معرفة بمَن نشر الحدث.
\item ليه أحداث؟ فصل الاهتمامات: المعالِج الأساسيّ نظيف، حدث واحد يشغّل ردودًا
  كثيرة، إضافة ردّ جديد بلا لمس المنطق، واختبار أبسط.
\item التدقيق: سطر \code{[Auditable]} فوق الأمر + \code{AuditBehavior} يعترض
  بعد النجاح ويكتب \code{AuditLog}: مَن/ماذا/متى/\en{IP}/\en{correlation}، مع
  قابليّة قبل/بعد \en{JSON}، و\en{append-only} — لا تعديل ولا حذف.
\item الفكرة الجامعة: افصل الإعلان عن التصرّف، وافصل المساءلة عن
  المنطق — فيبقى كل شيء مستقلًّا، قابلًا للامتداد، وسهل الاختبار.
\end{itemize}
\end{summary}

\clearpage
% ch26-identity-tokens-sso
\section{الهُوية (\en{Identity})، التوكنات، و\en{SSO} بعمق}

في فصل التوثيق السابق عرفنا فكرة الرمزين على مستوى عالٍ: رمز وصول (\en{access
JWT}) قصير العمر يُرفَق مع كل طلب، ورمز تحديث (\en{refresh token}) أطول عمرًا
وظيفته استبدال الأول عند انتهائه. في هذا الفصل ننزل طبقةً أعمق ونجيب على أسئلة
عملية دقيقة: أين تُخزَّن التوكنات وبأيّ شكل؟ كيف نُبطلها في منتصف الجلسة؟ لماذا يرث
سياق قاعدة البيانات صنف \en{Identity} تحديدًا؟ وكيف يعمل الدخول عبر \en{Google}
و\en{Microsoft} خطوةً بخطوة؟ كل ما هنا مقتبس من الكود الحيّ في \code{server/src}،
بمُعرِّفات حقيقية.

\begin{idea}[الفكرة: هوية واحدة، مصدرها نحن]
\en{ScholarPath} يقبل عدّة طرق للدخول (كلمة مرور، \en{Google}، \en{Microsoft})،
لكنه في النهاية يحوّلها جميعًا إلى هوية محلية واحدة يملكها هو ويصدر لها توكناته
الخاصّة. أيًّا كان باب الدخول، بعد لحظة إثبات الهوية يصبح النظام هو المتحكّم الكامل
في الجلسة: من يدخل، وكم يبقى، ومتى يُطرَد.
\end{idea}

% ====================================================================
\subsection{ليه \code{ApplicationDbContext} يرث \code{IdentityDbContext} لا \code{DbContext}؟}

هذا سؤال يتكرّر في المناقشات. الجواب المختصر: لأن \en{ASP.NET Core Identity}
يحتاج مجموعة جداول جاهزة (المستخدمون، الأدوار، الـ\en{claims}، تسجيلات الدخول
الخارجية، التوكنات، وربط المستخدم بالدور)، والوراثة من \code{IdentityDbContext}
تُعرِّف هذه الجداول وتربط علاقاتها تلقائيًّا. هذا هو سطر إعلان الصنف الحقيقيّ:

\begin{english}
\begin{lstlisting}[language={[Sharp]C}]
public class ApplicationDbContext(
    DbContextOptions<ApplicationDbContext> options,
    IMediator? mediator = null,
    IFieldEncryptionService? encryption = null)
    : IdentityDbContext<ApplicationUser, ApplicationRole, Guid>(options),
      IApplicationDbContext
{
    public DbSet<RefreshToken> RefreshTokens => Set<RefreshToken>();
    // ... business DbSets: Scholarships, Applications, Reviews, ...
}
\end{lstlisting}
\end{english}

نقرأه سطرًا سطرًا:
\begin{itemize}
\item \code{: IdentityDbContext<ApplicationUser, ApplicationRole, Guid>} — هذه
  الوراثة تُخبر \en{EF Core} بأن يبني جداول الهوية بأنواعنا نحن: المستخدم
  \code{ApplicationUser}، والدور \code{ApplicationRole}، ومفتاح كلاهما من نوع
  \code{Guid}.
\item \code{IApplicationDbContext} — منفذ طبقة \en{Application}؛ يجعل المُعالِجات
  ترى السياق كواجهة مجرّدة لا كصنف ملموس (اتّساقًا مع \en{Clean Architecture}).
\item \code{DbSet<RefreshToken>} وبقيّة الجداول — هي جداول العمل التي نضيفها نحن
  فوق جداول الهوية.
\end{itemize}

الفرق الجوهريّ يظهر داخل \code{OnModelCreating}. صنفنا يستدعي أولًا
\code{base.OnModelCreating(builder)} فتُبنى جداول الهوية السبعة بعلاقاتها، ثم نضيف
تهيئتنا فوقها:

\begin{english}
\begin{lstlisting}[language={[Sharp]C}]
protected override void OnModelCreating(ModelBuilder builder)
{
    builder.ApplyConfigurationsFromAssembly(Assembly.GetExecutingAssembly());
    base.OnModelCreating(builder);   // <-- builds ALL Identity tables + keys
    // ... our own model tweaks after Identity is wired
}
\end{lstlisting}
\end{english}

\begin{why}
لو ورثنا \code{DbContext} العاديّ لاضطُررنا لتعريف كل جدول هوية يدويًّا: جدول
المستخدمين، وجدول الأدوار، وجدول ربط المستخدم بالدور، وجداول الـ\en{claims}
وتسجيلات الدخول الخارجية والتوكنات — مع مفاتيحها الأجنبية وفهارسها الفريدة (مثل
البريد المُطبَّع الفريد). هذا عمل طويل وسهل الخطأ. الوراثة من
\code{IdentityDbContext} تعطينا كل ذلك مضبوطًا ومختبَرًا من مايكروسوفت، ونكتفي نحن
بإضافة جداول العمل.
\end{why}

الجدول التالي يلخّص جداول الهوية التي تأتي مجّانًا مع الوراثة:

\begin{center}
\begin{tabular}{L{4.2cm} L{10.3cm}}
\toprule
{\bfseries\color{spPrimary}الجدول} & {\bfseries\color{spPrimary}ما يخزّنه} \\
\midrule
\en{AspNetUsers} & الحسابات: البريد، كلمة المرور المُجزَّأة، عدّاد المحاولات، القفل. \\
\en{AspNetRoles} & الأدوار المتاحة في المنصّة. \\
\en{AspNetUserRoles} & ربط كل مستخدم بأدواره (علاقة كثير-لكثير). \\
\en{AspNetUserClaims} & \en{claims} إضافية مرتبطة بالمستخدم. \\
\en{AspNetUserLogins} & تسجيلات الدخول الخارجية (\en{Google}/\en{Microsoft}) المربوطة بالحساب. \\
\en{AspNetUserTokens} & توكنات تُخزّنها \en{Identity} لأغراض مثل إعادة التعيين. \\
\bottomrule
\end{tabular}
\end{center}

% ====================================================================
\subsection{نموذج الأدوار والـ\en{claims}: إزاي أهندل الرولز}

\en{ScholarPath} يعرّف خمسة أدوار: \code{SuperAdmin} و\code{Admin}
و\code{Student} و\code{Consultant} و\code{ScholarshipProvider} (مزوّد المنح،
وكان اسمه سابقًا \en{Company}). المهم أن نفهم كيف تنتقل هذه الأدوار من قاعدة
البيانات إلى داخل الطلب.

\begin{keyp}[نقطة مهمة: الدور يُحمَل داخل التوكن]
لا يرجع النظام إلى قاعدة البيانات في كل طلب ليسأل «ما أدوار هذا المستخدم؟». بدلًا
من ذلك، لحظة الإصدار تُوضَع الأدوار داخل الـ\en{JWT} على هيئة \en{claims} موقّعة.
أي طلب لاحق يحمل التوكن يحمل معه أدواره، فيصبح فحص الصلاحية سريعًا وبلا استعلام.
\end{keyp}

هذا مقتطف حقيقيّ من \code{TokenService.CreateAccessToken} يبني قائمة الـ\en{claims}:

\begin{english}
\begin{lstlisting}[language={[Sharp]C}]
var claims = new List<Claim>
{
    new(JwtRegisteredClaimNames.Sub,   user.Id.ToString()),
    new(JwtRegisteredClaimNames.Email, user.Email ?? string.Empty),
    new("account_status", user.AccountStatus.ToString()),
    new(ClaimTypes.Name, user.FullName),
};
if (!string.IsNullOrEmpty(activeRole))
    claims.Add(new Claim("active_role", activeRole));   // the session's ACTIVE role
foreach (var r in roles)
    claims.Add(new Claim(ClaimTypes.Role, r));          // every role the user holds
\end{lstlisting}
\end{english}

الشرح:
\begin{itemize}
\item \code{Sub} و\code{Email} و\code{Name} — هوية المستخدم الأساسية داخل التوكن.
\item \code{active\_role} — نقطة ذكية: المستخدم قد يملك أكثر من دور (طالب ومستشار
  مثلًا)، لكن جلسته الحالية تعمل بدور واحد فعّال. نضع هذا الدور في \en{claim}
  مستقلّ.
\item \code{foreach ... ClaimTypes.Role} — نضيف كل أدوار المستخدم كـ\en{claims}
  من نوع الدور، فتُعرَف صلاحياته كاملة عند الحاجة.
\end{itemize}

الآن السؤال العمليّ: كيف تُحمى نقطة النهاية بهذا؟ نُهيّئ التحقّق ليعتبر \en{claim}
اسمه \code{active\_role} هو «claim الدور»، ثم نضع سمة \code{[Authorize(Roles=...)]}
فوق المتحكّم:

\begin{english}
\begin{lstlisting}[language={[Sharp]C}]
opts.TokenValidationParameters = new TokenValidationParameters
{
    ValidAlgorithms = [SecurityAlgorithms.RsaSha256],
    RoleClaimType   = "active_role",   // [Authorize(Roles=...)] checks the ACTIVE role
};

[Authorize(Roles = "Admin,SuperAdmin")]   // e.g. AdminController
public sealed class AdminController : ControllerBase { /* ... */ }
\end{lstlisting}
\end{english}

الشرح:
\begin{itemize}
\item \code{RoleClaimType = "active\_role"} — يجعل \code{[Authorize(Roles=...)]}
  يفحص الدور الفعّال للجلسة لا مجموعة الأدوار كلها. فمستخدم يملك دورَي طالب ومستشار،
  ويعمل الآن كطالب، لا يستطيع فتح صفحات المستشار حتى يبدّل دوره.
\item \code{[Authorize(Roles = "Admin,SuperAdmin")]} — هذه هي الحماية الفعلية على
  \code{AdminController} في الكود: فقط من دوره الفعّال \en{Admin} أو
  \en{SuperAdmin} يمرّ.
\end{itemize}

الرسم التالي يوضّح كيف «يتفرّع» التوكن الواحد إلى عدّة \en{claims}، وأيّها يستعمله
التفويض:

\begin{center}\begin{LTR}
\begin{tikzpicture}[node distance=0.55cm]
  \node[pbox, text width=2.7cm, minimum height=1.1cm] (jwt) {Access JWT\\(RS256, 15 min)};
  \node[sbox, text width=3.2cm] (c1) at (5.4,2.3)  {sub = user id};
  \node[sbox, text width=3.2cm] (c2) at (5.4,1.35) {email};
  \node[sbox, text width=3.2cm] (c3) at (5.4,0.4)  {account status};
  \node[tbox, text width=3.2cm] (c4) at (5.4,-0.55){active-role};
  \node[sbox, text width=3.2cm] (c5) at (5.4,-1.5) {role: Student};
  \node[sbox, text width=3.2cm] (c6) at (5.4,-2.45){role: Consultant};
  \foreach \c in {c1,c2,c3,c4,c5,c6} \draw[arr] (jwt.east) -- (\c.west);
  \node[fbox, text width=3.0cm, right=0.5cm of c4] (use) {used by\\Authorize(Roles)};
  \draw[arrd] (c4.east) -- (use.west);
\end{tikzpicture}
\end{LTR}\end{center}

% ====================================================================
\subsection{تخزين التوكنات: ليه الـ\en{refresh} مُجزَّأ بـ\en{SHA-256}؟}

رمز الوصول لا يُخزَّن أصلًا في القاعدة — هو موقّع ويُفحَص بمفتاح، فلا حاجة لحفظه.
أمّا رمز التحديث فيجب أن نتذكّره لنستطيع إبطاله لاحقًا. والنقطة الحسّاسة: لا نُخزّنه
كما هو، بل نُخزّن تجزئته فقط في الحقل \code{RefreshToken.TokenHash}. هذا مقتطف
حقيقيّ من \code{TokenService.IssueTokens} يوضّح لحظة الحفظ:

\begin{english}
\begin{lstlisting}[language={[Sharp]C}]
var refreshToken = GenerateSecureRefreshToken();     // 64 crypto-random bytes
var refreshHash  = HashRefreshToken(refreshToken);   // SHA-256(token) -> hex
db.RefreshTokens.Add(new RefreshToken
{
    UserId    = user.Id,
    TokenHash = refreshHash,        // the HASH is persisted, never the raw token
    ExpiresAt = refreshExpires,
    CreatedAt = now,
});
db.SaveChanges();
return new TokenPair(accessToken, refreshToken, accessExpires, refreshExpires);
\end{lstlisting}
\end{english}

الشرح:
\begin{itemize}
\item \code{GenerateSecureRefreshToken()} — يولّد 64 بايت عشوائيًّا تشفيريًّا، فالرمز
  غير قابل للتخمين إطلاقًا.
\item \code{TokenHash = refreshHash} — ما يدخل القاعدة هو تجزئة \en{SHA-256} فقط.
\item \code{return new TokenPair(..., refreshToken, ...)} — الرمز الخام يُعاد للعميل
  مرّة واحدة فقط في الاستجابة، ولا يُكتَب في القاعدة أبدًا.
\end{itemize}

عند التجديد، يأخذ النظام الرمز الخام القادم من العميل، يجزّئه بنفس الدالة، ثم يبحث
عن تجزئة مطابقة. المقارنة إذن بين تجزئتين، لا بين نصّين.

\begin{secnote}
لماذا التجزئة؟ لنفس سبب تجزئة كلمات المرور: لو تسرّبت قاعدة البيانات كاملةً، لن يجد
المهاجم إلا \en{SHA-256} لا يمكن عكسها إلى الرمز الأصلي. وبما أن النظام يقارن بتجزئة
ما يُقدَّم إليه، فالتجزئة المسروقة عديمة الفائدة — لا تصلح لانتحال جلسة. هذا يجعل
قاعدة البيانات هدفًا أقلّ جاذبية بكثير.
\end{secnote}

% ====================================================================
\subsection{الإبطال والانتهاء: إزاي أبوّظها بعد مرحلة معيّنة}

هذا لبّ الفصل. سؤال متكرّر: «طيّب لو عايز أبطّل جلسة في نصّها، أعمل إيه؟». الجواب
يختلف بين الرمزين لأن طبيعتهما مختلفة تمامًا.

رمز الوصول ينتهي ذاتيًّا. داخله \en{claim} اسمه \code{exp} يحمل لحظة الانتهاء (بعد
15 دقيقة). عند كل طلب يفحص المتحقّق هذا الوقت دون الرجوع إلى قاعدة البيانات؛ فإذا
مضى الوقت رُفِض الرمز. الثمن هو أننا لا نستطيع سحب رمز وصول قبل انتهائه — لكن نافذته
15 دقيقة فقط، فالخطر محدود.

رمز التحديث يُبطَل فورًا. لأنه محفوظ (مُجزَّأً) في القاعدة، يمكننا وضع علامة عليه في
أي لحظة. وله ثلاثة نطاقات إبطال:

\begin{center}
\begin{tabular}{L{4.3cm} L{3.2cm} L{6.8cm}}
\toprule
{\bfseries\color{spPrimary}سبب الإبطال} & {\bfseries\color{spPrimary}النطاق} & {\bfseries\color{spPrimary}الأثر} \\
\midrule
تسجيل الخروج & الرمز الحالي & يُعلَّم \code{IsRevoked} على هذا الرمز فقط؛ بقيّة الأجهزة تبقى. \\
كشف إعادة استخدام رمز مُبطَل & كل رموز المستخدم & اشتباه سرقة $\rightarrow$ إلغاء كل الجلسات فورًا. \\
تغيير كلمة المرور & كل رموز المستخدم & كل الجلسات القديمة تسقط؛ يلزم دخول جديد. \\
انتهاء المدّة & الرمز المعنيّ & \code{ExpiresAt} في الماضي $\rightarrow$ يُرفَض دون علامة. \\
التدوير عند الاستخدام & الرمز القديم & يُبطَل القديم لحظة إصدار بديله (لا رمزان صالحان معًا). \\
\bottomrule
\end{tabular}
\end{center}

الكود الذي ينفّذ الإبطال «الفرديّ» و«لكل العائلة». لاحظ منطق كشف إعادة الاستخدام:

\begin{english}
\begin{lstlisting}[language={[Sharp]C}]
// Reuse detection: a revoked token presented again means it was already
// rotated -- a replay attack or a stale client. Revoke the whole family.
if (stored.IsRevoked)
{
    await RevokeAllForUserAsync(stored.UserId, "Refresh token reuse detected", ct);
    return null;
}

// Global revoke: mark every live token for this user as revoked.
var tokens = await db.RefreshTokens
    .Where(t => t.UserId == userId && !t.IsRevoked)
    .ToListAsync(ct);
foreach (var t in tokens)
{
    t.IsRevoked = true;
    t.RevokedAt = clock.UtcNow;
    t.RevokedReason = reason;
}
\end{lstlisting}
\end{english}

الشرح:
\begin{itemize}
\item \code{if (stored.IsRevoked)} — لو وصلنا رمزٌ مُبطَل سلفًا، فهذا يعني أنه دُوِّر
  من قبل ثم عاد للظهور: إمّا هجوم إعادة تشغيل وإمّا عميل قديم. الأأمن أن نفترض السرقة.
\item \code{RevokeAllForUserAsync(...)} — الإبطال «لكل العائلة»: يقطع كل جلسات
  المستخدم دفعةً واحدة، فحتى لو كان المهاجم يملك رمزًا صالحًا آخر يسقط أيضًا.
\item \code{t.RevokedReason = reason} — نسجّل سبب الإبطال، فيبقى أثر يمكن تدقيقه.
\end{itemize}

الرسم التالي هو آلة حالات رمز التحديث: من الإصدار، مرورًا بالتدوير، وصولًا إلى
الحالتين النهائيتين (مُبطَل أو منتهٍ):

\begin{center}\begin{LTR}
\begin{tikzpicture}[node distance=1.3cm and 2.3cm]
  \node[tbox, text width=2.6cm] (iss) {issued\\(valid)};
  \node[sbox, text width=2.9cm, right=of iss] (rot) {rotated\\(old to revoked)};
  \node[fbox, text width=2.6cm, below=of iss] (exp) {expired};
  \node[fbox, text width=2.6cm, below=of rot] (rev) {revoked\\(terminal)};
  \draw[arr] (iss) -- node[midway,above,lbl] {on use} (rot);
  \draw[arrd] (rot) to[bend right=32] node[midway,below,lbl] {fresh pair issued} (iss);
  \draw[arr] (iss) -- node[midway,left,lbl] {ExpiresAt passed} (exp);
  \draw[arr] (iss) -- node[midway,right,lbl,align=left] {logout / reuse /\\pwd change} (rev);
  \draw[arrd] (rot.south) -- node[midway,right,lbl] {old} (rev.north);
\end{tikzpicture}
\end{LTR}\end{center}

\begin{spotlight}[ركّز هنا لإبهار اللجنة: الـ\en{refresh} مُجزَّأ وقابل للإبطال فورًا]
اجمع الصورتين معًا: رمز التحديث مُجزَّأ في القاعدة (تسريبها لا يفيد المهاجم)، وفي
الوقت نفسه قابل للإبطال لحظيًّا بثلاثة نطاقات — فرديًّا (خروج من جهاز)، ولكل العائلة
(اشتباه سرقة أو تغيير كلمة مرور)، وعالميًّا. هذا يجمع بين ميزتين نادرًا ما تجتمعان:
أمان تحت التسريب، وتحكّم كامل في الجلسة. حين تُسأل «إزاي تطرد جلسة مسروقة؟» تكون
الإجابة جاهزة: علامة واحدة على صفٍّ في القاعدة تُسقط الجلسة عند أول محاولة تجديد.
\end{spotlight}

% ====================================================================
\subsection{الدخول الموحّد (\en{SSO}) عبر \en{OAuth 2.0}}

الآن الباب الثاني للدخول: \en{Google} و\en{Microsoft}. المعيار المستخدَم هو
\en{OAuth 2.0} بنمط «كود التفويض» (\en{authorization-code}). الفكرة العامّة أن
المستخدم يثبت هويته عند المزوّد لا عندنا، فيعود متصفّحه إلينا حاملًا كودًا مؤقّتًا،
ونحن نبادله في الخلفية بمعلومات المستخدم.

\begin{idea}[الفكرة: الكود قناة آمنة]
لماذا لا يعيد المزوّد التوكن مباشرةً للمتصفّح؟ لأن المتصفّح بيئة غير موثوقة (قد يقرأ
المحتوى برنامج خبيث). بدل ذلك يعيد كودًا قصير العمر لا قيمة له وحده، ثم نبادله نحن
من الخادم — حيث يكمن سرّ تطبيقنا (\en{client secret}) — للحصول على التوكن. هكذا لا
يمرّ التوكن الحسّاس عبر المتصفّح.
\end{idea}

هذا الرسم هو التسلسل الكامل بين المتصفّح والمزوّد وواجهتنا (الأسهم مرقّمة بالترتيب):

\begin{center}\begin{LTR}
\begin{tikzpicture}
  \node[actor] (B) at (0,0) {Browser};
  \node[gbox]  (P) at (5,0) {Google /\\Microsoft};
  \node[pbox]  (A) at (10,0) {our API};
  \draw[arrd] (B.south) -- (0,-7.9);
  \draw[arrd] (P.south) -- (5,-7.9);
  \draw[arrd] (A.south) -- (10,-7.9);
  \draw[arr] (0,-1.1) -- (10,-1.1) node[midway,above,lbl] {(1) start Google sign-in};
  \draw[arr] (10,-2.0) -- (0,-2.0) node[midway,above,lbl] {(2) authorize URL + state};
  \draw[arr] (0,-2.9) -- (5,-2.9) node[midway,above,lbl] {(3) authorize + state};
  \draw[arr] (5,-3.8) -- (0,-3.8) node[midway,above,lbl] {(4) code + state};
  \draw[arr] (0,-4.7) -- (10,-4.7) node[midway,above,lbl] {(5) code + state};
  \draw[arr] (10,-5.6) -- (5,-5.6) node[midway,above,lbl] {(6) exchange code (state ok)};
  \draw[arr] (5,-6.5) -- (10,-6.5) node[midway,above,lbl] {(7) provider token + userinfo};
  \draw[arr] (10,-7.4) -- (0,-7.4) node[midway,above,lbl] {(8) OUR access + refresh};
\end{tikzpicture}
\end{LTR}\end{center}

نلاحظ في الخطوة (6) أن الواجهة تتحقّق من الـ\en{state} قبل مبادلة الكود. هذا الـ
\en{state} هو حارس مكافحة \en{CSRF}: قيمة نصدرها في الخطوة (2) ويجب أن تعود كما هي
في الخطوة (5). ندير هذا عبر \code{DataProtectionSsoStateStore}:

\begin{english}
\begin{lstlisting}[language={[Sharp]C}]
public string Issue()
{
    // payload = <expiry-unix-seconds>:<random-nonce>, protected (tamper-proof).
    var expiresAt = DateTimeOffset.UtcNow.Add(Ttl).ToUnixTimeSeconds();
    var payload   = $"{expiresAt}:{Guid.NewGuid():N}";
    return _protector.Protect(payload);          // signed + encrypted
}

public bool Validate(string? state)
{
    var payload = _protector.Unprotect(state);   // throws if forged/tampered
    var sep = payload.IndexOf(':');
    return sep > 0 && long.TryParse(payload.AsSpan(0, sep), out var exp)
        && DateTimeOffset.FromUnixTimeSeconds(exp) > DateTimeOffset.UtcNow;
}
\end{lstlisting}
\end{english}

الشرح:
\begin{itemize}
\item \code{Protect(payload)} — الـ\en{state} ليس رقمًا نحفظه في الخادم، بل حزمة
  موقّعة ومُشفّرة تحمل تاريخ انتهائها بنفسها. لا نحتاج تخزينًا من طرف الخادم، فيبقى
  التحقّق صالحًا حتى مع إعادة تشغيل التطبيق أو تعدّد نُسخه.
\item \code{Unprotect(state)} — يفشل (يرمي استثناءً) إن كان الـ\en{state} مزوّرًا
  أو معدَّلًا؛ فأيّ محاولة تزييف تسقط فورًا.
\item فحص الوقت — نرفض أي \en{state} انتهت مدّته (30 دقيقة)، فلا يُعاد استعماله لاحقًا.
\end{itemize}

بعد نجاح المبادلة وجلب بيانات المستخدم (\en{userinfo})، يتولّى
\code{SsoLoginCommandHandler} إيجاد المستخدم المحليّ أو إنشاءه، ثم — وهذه النقطة
الأهمّ — يصدر توكناتنا نحن:

\begin{english}
\begin{lstlisting}[language={[Sharp]C}]
// Resolve by the recorded external identity first, then by verified email.
var linkedUserId = await userAdministration
    .FindUserIdByExternalLoginAsync(info.Provider, info.ProviderUserId, ct);
var user = linkedUserId is { } lid
    ? await db.Users.FirstOrDefaultAsync(u => u.Id == lid, ct)
    : await db.Users.FirstOrDefaultAsync(u => u.NormalizedEmail == normalizedEmail, ct);

if (user is null) { /* first SSO sign-in: provision an Unassigned account */ }

// Link the external identity (idempotent), then issue OUR tokens.
await userAdministration.AddExternalLoginAsync(user.Id, info.Provider, info.ProviderUserId, ct);
var tokens = tokenService.IssueTokens(user, roles, user.ActiveRole, rememberMe: false);
\end{lstlisting}
\end{english}

الشرح:
\begin{itemize}
\item \code{FindUserIdByExternalLoginAsync} — نبحث أولًا بالهوية الخارجية المسجّلة
  (المزوّد $+$ معرّف المستخدم عنده)، لا بالبريد؛ فربط الحساب يصبح صريحًا ومستقرًّا
  حتى لو غيّر المزوّد البريد لاحقًا.
\item \code{if (user is null)} — أوّل دخول عبر \en{SSO} يُنشئ حسابًا محليًّا جديدًا
  (بحالة \code{Unassigned}) لأن البريد موثَّق من المزوّد.
\item \code{AddExternalLoginAsync} — نسجّل رابط الهوية الخارجية (عملية عديمة الأثر
  عند التكرار)، فيُحَل الدخول التالي إلى نفس الحساب.
\item \code{tokenService.IssueTokens(...)} — السطر الحاسم: بعد إثبات الهوية نصدر
  زوجنا نحن (\en{access} $+$ \en{refresh})، ولا نستعمل توكن \en{Google} بعد ذلك أبدًا.
\end{itemize}

\begin{spotlight}[ركّز هنا لإبهار اللجنة: بعد \en{SSO} نُصدر توكناتنا نحن]
توكن \en{Google} أو \en{Microsoft} يُستعمل مرّة واحدة فقط: لإثبات أن هذا البريد
يخصّ صاحبه. بمجرّد أن نتأكّد، نرميه ونصدر توكناتنا الخاصّة. لماذا هذا مهمّ؟ لأنه
يبقي التحكّم كاملًا في أيدينا: دورة حياة الجلسة، والإبطال بنطاقاته الثلاثة،
والأدوار — كلها تعمل بنفس المسار سواء دخل المستخدم بكلمة مرور أو عبر \en{Google}.
لا نعتمد على جلسة طرف ثالث ولا على انتهائها؛ نحن من يقرّر متى تنتهي الجلسة ومتى
تُطرَد.
\end{spotlight}

\begin{qa}[سؤال متوقّع: لو حساب جوجل اتسرق؟]
«لو حساب \en{Google} للمستخدم اتسرق، هل هيدخل عندنا؟» — الدخول عبر \en{Google} يثبت
فقط ملكيّة البريد لحظتها. لو فقد المستخدم حساب \en{Google} فالمشكلة عند \en{Google}.
أمّا عندنا فبمجرّد أن يغيّر كلمة مروره المحلية أو نُبطِل جلساته، تسقط كل التوكنات
التي أصدرناها فورًا — لأننا لا نعتمد على جلسة \en{Google}، بل على رموزنا القابلة
للإبطال.
\end{qa}

% ====================================================================
\begin{summary}[الخلاصة في سطور]
\begin{itemize}
\item \code{ApplicationDbContext} يرث \code{IdentityDbContext<ApplicationUser,
  ApplicationRole, Guid>}، فيحصل على جداول الهوية السبعة وعلاقاتها جاهزة عبر
  \code{base.OnModelCreating}؛ الوراثة من \code{DbContext} كانت ستتطلّب تعريفها كلّها يدويًّا.
\item الأدوار الخمسة (\en{SuperAdmin/Admin/Student/Consultant/ScholarshipProvider})
  تُحمَل داخل الـ\en{JWT} كـ\en{claims}، والتفويض عبر \code{[Authorize(Roles=...)]}
  مضبوط على \code{active\_role} — الدور الفعّال للجلسة.
\item رمز التحديث لا يُخزَّن خامًا، بل تُخزَّن تجزئته \en{SHA-256} في
  \code{RefreshToken.TokenHash}؛ تسريب القاعدة لا يعطي رمزًا صالحًا.
\item رمز الوصول ينتهي ذاتيًّا عبر \code{exp} (15 دقيقة، بلا رجوع للقاعدة)؛ ورمز
  التحديث يُدوَّر عند كل استخدام ويُبطَل فرديًّا/للعائلة/عالميًّا عند الخروج أو كشف
  إعادة الاستخدام أو تغيير كلمة المرور.
\item \en{SSO} يتبع \en{OAuth 2.0} (كود التفويض): كود $\rightarrow$ مبادلة
  $\rightarrow$ \en{userinfo} $\rightarrow$ إيجاد/إنشاء مستخدم، مع \code{state}
  موقّع (\code{DataProtectionSsoStateStore}) لمكافحة \en{CSRF}.
\item بعد \en{SSO} نصدر توكناتنا نحن ولا نعتمد على توكن المزوّد — تحكّم كامل في
  الجلسة والإبطال.
\end{itemize}
\end{summary}

\clearpage
% ch27-testing — الاختبارات: كيف نضمن أن كل شيء يعمل
\section{الاختبارات \en{(Testing)}: كيف نضمن أنّ كلّ شيء يعمل}

في الفصل ده هنفهم إزاي بنتأكّد إنّ \en{ScholarPath} شغّال صحّ، مش بس أوّل يوم، لكن بعد كلّ تعديل بنعمله. هنعرف ليه نكتب اختبارات أصلًا، وإيه الفرق بين اختبار \en{unit} واختبار \en{integration}، وإيه «هرم الاختبار». وأهمّ حاجة: هنشوف إزاي فكرة الـ \en{ports/adapters} اللي اتكلّمنا عنها قبل كده (المنفذ = واجهة \en{interface}، والمحوِّل = تنفيذ حقيقيّ أو وهميّ خلفه) هي اللي بتخلّي الاختبار ممكنًا أصلًا: نحقن بديلًا وهميًّا \en{(stub)} بدل الخدمة الحقيقيّة، فالاختبار يشتغل بلا إنترنت وبنتيجة ثابتة كلّ مرّة. وفي الآخر هنشرّح اختبارًا حقيقيًّا من مشروع الاختبارات سطرًا بسطر.

\begin{idea}[الفكرة: الاختبار شبكة أمان]
كلّ اختبار هو جملة تقول: «لو النظام اشتغل صحّ، لازم تطلع النتيجة دي». لمّا نجمّع مئات الجُمَل دي ونشغّلها في ثوانٍ بعد أيّ تعديل، بنمسك أيّ كسر جديد فورًا قبل ما يوصل للمستخدم. الاختبار مش رفاهية؛ هو اللي بيخلّينا نعدّل في الكود ونحن مطمئنّين.
\end{idea}

\subsection{ليه نختبر أصلًا؟}

النظام فيه قواعد عمل كتير وحسّاسة: حساب الاسترداد \en{(refund)}، صلاحيّات مين يقدر يعمل إيه، خصم تقييم المستشار، ساعات الهدوء للإشعارات. أيّ خطأ صغير في القواعد دي ممكن يرجّع فلوس غلط أو يكشف بيانات مستخدم لمستخدم تاني. لو اعتمدنا على «نجرّب بإيدينا» كلّ مرّة، ده بطيء وبينسى حالات، وبيتعب مع كِبَر النظام.

\begin{why}
الاختبار الآليّ بيحوّل «أعتقد إنّه شغّال» إلى «أثبتُّ إنّه شغّال». وبيدّينا جُرأة نعيد هيكلة الكود \en{(refactor)} من غير خوف: طالما الاختبارات خضراء، السلوك اتحافظ عليه. ده بالظبط اللي خلّانا نصلّح عشرات الأخطاء في المشروع من غير ما نكسر حاجة تانية.
\end{why}

\subsection{أنواع الاختبارات وهرم الاختبار}

عندنا تلات مستويات، كلّ واحد بيجاوب على سؤال مختلف:

\begin{center}
\begin{tabular}{L{3.2cm} L{9.2cm}}
\toprule
النوع & بيتأكّد من إيه؟ \\
\midrule
\en{Unit} (وحدة) & قطعة صغيرة معزولة: \en{handler} أو \en{validator} أو خدمة حساب. سريع جدًّا، بلا قاعدة بيانات حقيقيّة ولا شبكة. \\
\en{Integration} (تكامل) & كذا طبقة مع بعض عبر قاعدة بيانات وخدمات وهميّة: هل الطلب \en{HTTP} بيمرّ صحّ من الـ \en{API} للـ \en{handler} للقاعدة؟ \\
\en{E2E} (طرف لطرف) & رحلة المستخدم كاملة في المتصفّح. أبطأ وأغلى، فبنعمل منه أقلّ عدد. \\
\bottomrule
\end{tabular}
\end{center}

الشكل ده بيترجم القاعدة المعروفة: «كتير \en{unit}، أقلّ \en{integration}، قليل \en{e2e}». كلّما نزلنا لتحت، الاختبار أسرع وأرخص وأكتر عددًا؛ وكلّما طلعنا لفوق، أبطأ وأغلى فبنقلّله.

\begin{center}\begin{LTR}
\begin{tikzpicture}[node distance=3pt]
  \node[fbox, text width=3.2cm, minimum height=0.75cm] (e2e) {E2E $\approx$ few};
  \node[tbox, text width=6.2cm, minimum height=0.78cm, below=of e2e] (int) {Integration $\approx$ 35};
  \node[pbox, text width=9.4cm, minimum height=0.82cm, below=of int] (unit) {Unit $\approx$ 800};
  % سهم جانبي يبيّن الاتجاه: أبطأ/أغلى لفوق، أسرع/أرخص لتحت
  \draw[arr] ([xshift=6mm]unit.east) -- ([xshift=6mm]e2e.east)
     node[lbl, midway, right, align=left] {slower\\ costlier};
  \draw[arr] ([xshift=-6mm]e2e.west) -- ([xshift=-6mm]unit.west)
     node[lbl, midway, left, align=right] {faster\\ more};
\end{tikzpicture}
\end{LTR}\end{center}

\begin{note}
عندنا في المشروع أكثر من \textenglish{800} اختبار \en{unit} (على \textenglish{123} ملفًّا) وحوالي \textenglish{35} اختبار \en{integration}. النسبة دي هي شكل الهرم بالظبط: القاعدة عريضة والقمّة ضيّقة.
\end{note}

\subsection{إزاي الـ \en{ports} بيخلّوا الاختبار ممكنًا؟}

المشكلة: الـ \en{handler} بتاع إلغاء الحجز محتاج يكلّم \en{Stripe} علشان يرجّع الفلوس، وغرفة الجلسة محتاجة \en{Azure}. لو الاختبار كلّم الخدمات دي حقيقةً هيبقى بطيء، ومحتاج إنترنت ومفاتيح، وأسوأ حاجة: نتيجته مش ثابتة (ممكن \en{Stripe} يرجّع خطأ شبكة النهاردة). الحلّ اللي شرحناه في الجزء الأوّل: الكود بيعتمد على منفذ \en{(interface)} زي \code{IStripeService}، مش على التنفيذ الحقيقيّ. في الإنتاج نحقن المحوِّل الحقيقيّ، وفي الاختبار نحقن \en{stub} وهميّ يرجّع نتيجة جاهزة فورًا.

الكود ده من تسجيل الخدمات بيوضّح المبدأ: نفس الواجهة، تنفيذان، والاختيار بيتمّ حسب الإعداد.

\begin{english}
\begin{lstlisting}[language={[Sharp]C}]
// Stripe: the real client only when a genuine secret key is configured;
// placeholder / missing values fall through to the deterministic dev stub.
var key = config.GetValue<string>($"{StripeOptions.SectionName}:SecretKey");
if (key is not null && key.StartsWith("sk_", StringComparison.Ordinal))
    services.AddSingleton<IStripeService, StripeService>();       // prod adapter
else
    services.AddSingleton<IStripeService, StubStripeService>();   // offline stub

// Same config-driven fallback for the video-meeting provider:
if (!string.IsNullOrWhiteSpace(acsConnString))
    services.AddSingleton<IMeetingService, AzureCommunicationMeetingService>();
else
    services.AddSingleton<IMeetingService, StubMeetingService>();
\end{lstlisting}
\end{english}

شرح الأسطر المهمّة:
\begin{itemize}
  \item السطر بتاع \code{key.StartsWith(...)}: لو فيه مفتاح \en{Stripe} حقيقيّ نستخدم \code{StripeService}؛ غير كده نقع على \code{StubStripeService}.
  \item نفس النمط بالظبط للـ \code{IMeetingService}: \en{Azure} الحقيقيّ لو فيه اتّصال، وإلّا \code{StubMeetingService} الحتميّ.
  \item الكلاس اللي بيستهلك الواجهة مش بيفرق عنده مين اللي اتحقن — ده كلّ سرّ قابليّة الاختبار.
\end{itemize}

الـ \en{stub} نفسه بسيط: بيرجّع كائن نتيجة جاهز من غير ما يلمس الشبكة. مثال \code{StubStripeService} لمّا يُطلب منه استرداد بيرجّع فورًا نتيجة بحالة \code{"succeeded"}. وكمان \code{NoOpFileScanService} بيقول عن أيّ ملفّ إنّه «نظيف» علشان تختبر رفع الملفّات بلا خادم فحص فيروسات. الرسم ده بيلخّص الحكاية:

\begin{center}\begin{LTR}
\begin{tikzpicture}[node distance=10mm and 16mm]
  \node[actor] (sut) {System under test};
  \node[gbox, right=of sut] (port) {IStripeService (port)};
  \node[srcbox, above right=6mm and 18mm of port, text width=2.7cm] (stub) {StubStripeService};
  \node[sbox, below right=6mm and 18mm of port, text width=2.7cm] (real) {StripeService};
  \node[tbox, right=of stub, text width=2.6cm] (mem) {in-memory offline};
  \node[fbox, right=of real, text width=2.6cm] (cloud) {Stripe cloud (network)};
  \draw[arr] (sut) -- (port);
  \draw[arr] (port) -- (stub) node[lbl, midway, above, sloped] {test injects};
  \draw[arr] (port) -- (real) node[lbl, midway, below, sloped] {prod injects};
  \draw[arr] (stub) -- (mem);
  \draw[arr] (real) -- (cloud);
\end{tikzpicture}
\end{LTR}\end{center}

\begin{spotlight}[ركّز هنا لإبهار اللجنة: الاختبارات حتميّة وبلا إنترنت]
لأنّ كلّ تكامل خارجيّ (\en{Stripe}، \en{Azure}، فحص الفيروسات، البريد) مخفيّ خلف منفذ، بنحقن مكانه \en{stub} في الاختبار. النتيجة: مجموعة الاختبارات كلّها بتشتغل \en{offline}، في ثوانٍ، وبتدّي نفس النتيجة كلّ مرّة — مفيش «فشل عشوائيّ» بسبب الشبكة. ده هو العائد العمليّ المباشر لمعماريّة \en{ports/adapters}: مش بس كود نظيف، ده اختبار ممكن أصلًا.
\end{spotlight}

\subsection{ماذا اختبرنا فعلًا؟}

مش بنختبر كلّ سطر بالتساوي؛ بنركّز على القواعد اللي لو غلطت تعمل ضررًا. أهمّ العائلات اللي غطّيناها:

\begin{center}
\begin{tabular}{L{4cm} L{8.4cm}}
\toprule
المجال & أمثلة اختبارات حقيقيّة \\
\midrule
حساب الاسترداد & \code{RefundCalculatorService}: استرداد \textenglish{100}\% قبل \textenglish{24} ساعة، \textenglish{50}\% بعدها، ورفض من مش طرف في الحجز. \\
درجات الترشيح/الأهليّة & \code{LocalAiService}: تحويل نتائج المعايير إلى حكم \en{Eligible / PartiallyEligible / NotEligible}. \\
المُتحقِّقات \en{(validators)} & رفض التطبيق بحالة غلط، ومنع الرفض بلا سبب، ورفض نوع إشعار غير معرّف. \\
بوّابات الملكيّة/الصلاحيّة & مستخدم لا يرى تفضيلات إشعارات مستخدم تاني، والطلب بلا تسجيل دخول يُرفض. \\
منطق ساعات الهدوء \en{(DND)} & المستخدم الصامت يحتفظ بإشعار داخل التطبيق لكن بلا بريد. \\
\bottomrule
\end{tabular}
\end{center}

\begin{keyp}[نقطة مهمة: نختبر السلوك لا التفاصيل]
لاحظ إنّ الاختبارات بتتكلّم بلغة العمل: «قبل \textenglish{24} ساعة» و«استرداد كامل» و«ممنوع بلا سبب». مش بنختبر أسماء المتغيّرات الداخليّة؛ بنختبر القاعدة اللي المستخدم بيهتمّ بيها. كده لو غيّرنا التنفيذ الداخليّ والقاعدة اتحافظ عليها، الاختبار يفضل أخضر.
\end{keyp}

\subsection{تشريح اختبار حقيقيّ: نمط \en{Arrange / Act / Assert}}

كلّ اختبار عندنا بيمشي على تلات خطوات: \en{Arrange} (نجهّز المدخلات)، \en{Act} (نستدعي السلوك مرّة واحدة)، \en{Assert} (نتأكّد من النتيجة). الرسم ده بيوضّح التدفّق التلاتيّ ده:

\begin{center}\begin{LTR}
\begin{tikzpicture}[node distance=14mm]
  \node[sbox, text width=3cm] (arr) {Arrange\\ set up inputs};
  \node[pbox, text width=3cm, right=of arr] (act) {Act\\ one call};
  \node[gbox, text width=3cm, right=of act] (asr) {Assert\\ check result};
  \node[wbox, below=8mm of act] (det) {deterministic $\cdot$ offline $\cdot$ same result every run};
  \draw[arr] (arr) -- (act);
  \draw[arr] (act) -- (asr);
  \draw[arrd] (act) -- (det);
\end{tikzpicture}
\end{LTR}\end{center}

ده اختبار حقيقيّ من \code{RefundCalculatorServiceTests} بيتحقّق من الحالة الحدّيّة: إلغاء حجز مؤكّد قبل \textenglish{24} ساعة بالظبط لازم يرجّع نصّ المبلغ.

\begin{english}
\begin{lstlisting}[language={[Sharp]C}]
private readonly RefundCalculatorService _sut = new();   // fresh service, no deps

[Fact]
public void Calculate_ConfirmedBookingExactly24HoursBefore_ReturnsHalfRefund()
{
    // Arrange - a fixed "now" so the test never depends on the wall clock.
    var nowUtc = new DateTimeOffset(2026, 4, 25, 10, 0, 0, TimeSpan.Zero);

    // Act - one call to the behaviour under test.
    var result = _sut.Calculate(
        bookingStatus: BookingStatus.Confirmed,
        cancelledByUserId: StudentId, studentId: StudentId, consultantId: ConsultantId,
        scheduledStartAt: nowUtc.AddHours(24),      // exactly on the 24h boundary
        amountCents: 10_000, nowUtc: nowUtc);

    // Assert - the business rule, stated as expectations.
    Assert.Equal(50, result.RefundPercentage);
    Assert.Equal(5_000, result.RefundAmountCents);
    Assert.Equal(CancellationReason.StudentCancelledLessThan24HoursBefore,
                 result.CancellationReason);
}
\end{lstlisting}
\end{english}

نقرأه سطرًا بسطر:
\begin{itemize}
  \item السطر الأوّل: الخدمة دي حساب صافٍ بلا اعتماديّات؛ منفيش قاعدة بيانات ولا شبكة، فالاختبار سريع وبسيط.
  \item \code{nowUtc} ثابت: ده أهمّ سطر. لأنّنا بنمرّر «الآن» كوسيط بدل ما الخدمة تقرأ ساعة النظام، الاختبار بقى حتميًّا — مش هيفشل بكرة الساعة \textenglish{3} فجرًا.
  \item خطوة \en{Act} استدعاء واحد لـ \code{Calculate} — الاختبار بيركّز على سلوك واحد.
  \item خطوات \en{Assert}: بنتحقّق من تلات حاجات مترابطة — النسبة \textenglish{50}\%، والمبلغ \textenglish{5000} سنت، وسبب الإلغاء الصحيح. لو أيّ واحدة غلطت، الاختبار يحمرّ ويقول لنا بالظبط فين الكسر.
\end{itemize}

\begin{qa}[سؤال متوقّع: ليه تمرّر «الآن» كوسيط بدل قراءة ساعة النظام؟]
لأنّ أيّ اختبار بيقرأ الساعة الحقيقيّة بيبقى «هشًّا» \en{(flaky)}: ينجح دلوقتي ويفشل عند حدّ زمنيّ. بحقن الوقت من برّه، بنثبّت الزمن فنقدر نختبر الحالات الحدّيّة (بالظبط \textenglish{24} ساعة) بثقة. نفس الفلسفة بتاعة حقن الـ \en{stub}: نتحكّم في كلّ مدخل عشان النتيجة تبقى ثابتة.
\end{qa}

مثال تاني على قاعدة عمل حسّاسة: ساعات الهدوء. لمّا المستخدم يفعّل «عدم الإزعاج»، لازم البريد يتوقّف لكن الإشعار داخل التطبيق يفضل محفوظ (مش نخسر حاجة). الاختبار ده من \code{NotificationDispatcherTests} بيثبّت القاعدة دي:

\begin{english}
\begin{lstlisting}[language={[Sharp]C}]
[Fact]
public async Task Muted_user_keeps_inapp_row_but_gets_no_email()
{
    // Arrange - a user who switched Do-Not-Disturb on.
    var userId = Guid.NewGuid();
    db.UserProfiles.Add(new UserProfile { UserId = userId, NotificationsMuted = true });
    await db.SaveChangesAsync();

    // Act - something tries to notify them.
    await Sut(db).DispatchAsync(userId, NotificationType.ResourceApproved,
        NotificationParams.Empty, null, null, default);

    // Assert - the in-app row survives; the email channel is suppressed.
    var rows = await db.Notifications.Where(n => n.RecipientUserId == userId).ToListAsync();
    rows.Should().ContainSingle().Which.Channel.Should().Be(NotificationChannel.InApp);
}
\end{lstlisting}
\end{english}

هنا بنستخدم قاعدة \en{InMemory} من \en{EF Core} (قاعدة بيانات في الذاكرة) بدل \en{SQL Server} حقيقيّ، فالاختبار سريع ومعزول. والتوكيدات مكتوبة بأسلوب \code{FluentAssertions} القريب من اللغة الطبيعيّة: «لازم يكون فيه صفّ واحد، وقناته \en{InApp}».

\subsection{اختبارات التكامل: نفس الفكرة عبر الطبقات}

اختبار الوحدة بيختبر قطعة معزولة. اختبار التكامل بيختبر الرحلة الكاملة: طلب \en{HTTP} حقيقيّ بيدخل الـ \en{API}، يعدّي على المصادقة، الـ \en{handler}، وقاعدة بيانات حقيقيّة. بنستخدم \code{WebApplicationFactory} علشان نشغّل التطبيق كلّه في الذاكرة، و\en{Testcontainers} علشان نطلّع \en{SQL Server} حقيقيّ مؤقّت في حاوية \en{Docker}. وبرضه بنحقن بدائل اختبار للأجزاء اللي مش عايزين نعتمد فيها على العالم الخارجيّ:

\begin{english}
\begin{lstlisting}[language={[Sharp]C}]
builder.ConfigureServices(services =>
{
    // Point EF Core at a throwaway SQL Server container, not production.
    services.AddDbContext<ApplicationDbContext>(o =>
        o.UseSqlServer(_sqlContainer.GetConnectionString()));

    // Swap the real identity for a deterministic, controllable test user.
    services.RemoveAll<ICurrentUserService>();
    services.AddSingleton<TestCurrentUserService>();
    services.AddSingleton<ICurrentUserService>(
        sp => sp.GetRequiredService<TestCurrentUserService>());

    // ... plus a fake auth scheme so requests authenticate fully offline.
    services.AddAuthentication(/* TestAuthHandler.SchemeName */)
        .AddScheme<AuthenticationSchemeOptions, TestAuthHandler>(
            TestAuthHandler.SchemeName, _ => { });
});
\end{lstlisting}
\end{english}

النقاط المهمّة:
\begin{itemize}
  \item توجيه \en{EF Core} لقاعدة الحاوية المؤقّتة اللي بتتولد وتتمسح مع الاختبار — قاعدة حقيقيّة لكن معزولة ونظيفة.
  \item \code{RemoveAll} ثمّ \code{TestCurrentUserService}: بنشيل المستخدم الحقيقيّ ونحطّ مستخدم اختبار نتحكّم في دوره وهويّته — نفس مبدأ حقن الـ \en{stub}، بس للهويّة.
  \item \code{TestAuthHandler}: مخطّط مصادقة وهميّ يخلّي الطلب «مسجّل دخول» بلا \en{JWT} حقيقيّ، فنختبر البوّابات الأمنيّة \en{offline}.
\end{itemize}

كده اختبار التكامل بيتأكّد إنّ القطع مركّبة صحّ مع بعض (التوجيه، المصادقة، الصلاحيّات، حفظ القاعدة)، والوقت نفسه بيفضل \en{offline} وحتميّ لأنّ \en{Stripe} و\en{Azure} لسّه خلف الـ \en{stubs}.

\begin{spotlight}[ركّز هنا لإبهار اللجنة: كلّ قاعدة مال/صلاحيّة مغطّاة]
أيّ قاعدة تلمس الفلوس أو الصلاحيّات لها اختبار صريح: نسبة الاسترداد قبل/بعد \textenglish{24} ساعة، خصم \textenglish{20}\% من تقييم المستشار، حَجْب الطالب \textenglish{3} أيّام، منع رفض التطبيق بلا سبب، ومنع مستخدم من رؤية بيانات غيره. مش شرط تكون كثافة الاختبار متساوية في كلّ مكان — الكثافة بتتركّز حيث الخطر. ده اللي بيدّي اللجنة ثقة إنّ الأجزاء الحرجة محميّة بشبكة أمان، مش بحُسن النيّة.
\end{spotlight}

\begin{summary}[الخلاصة في سطور]
\begin{itemize}
  \item الاختبار بيحوّل «أعتقد إنّه شغّال» إلى «أثبتُّ إنّه شغّال»، وبيدّينا جُرأة نعدّل بلا خوف.
  \item هرم الاختبار: كتير \en{unit} في القاعدة، أقلّ \en{integration} في الوسط، قليل \en{e2e} في القمّة — عندنا أكثر من \textenglish{800} اختبار \en{unit}.
  \item الـ \en{ports/adapters} هي اللي بتخلّي الاختبار ممكنًا: نحقن \en{stub} بدل \en{Stripe}/\en{Azure}، فالاختبارات تشتغل \en{offline} وبنتيجة ثابتة.
  \item بنركّز التغطية على القواعد الحرجة: الاسترداد، الأهليّة، المُتحقِّقات، بوّابات الصلاحيّة، ساعات الهدوء.
  \item كلّ اختبار = \en{Arrange / Act / Assert}؛ وحقن «الآن» كوسيط بيخلّي الحالات الزمنيّة الحدّيّة حتميّة.
  \item اختبار التكامل بيركّب الطبقات مع \code{WebApplicationFactory} و\en{Testcontainers}، وبيفضل \en{offline} لأنّ الخدمات الخارجيّة خلف \en{stubs}.
\end{itemize}
\end{summary}

\clearpage
% ch28-student-journey
\section{يوم في حياة النظام: رحلة طالبة كاملة (وكيف يردّ الآخرون)}

كل فصل سابق شرح جزءًا: الطبقات، و\en{CQRS}، والأمان، والأحداث، والدفع،
والذكاء الاصطناعيّ، والزمن الحقيقيّ. هذا الفصل مختلف: لن نضيف مفهومًا جديدًا،
بل نربط كل ما سبق في قصّة واحدة متّصلة. سنتابع طالبة اسمها «سارة» من
لحظة إنشاء حسابها حتى تسأل مساعد الذكاء الاصطناعيّ سؤالًا، وعند كل خطوة نسأل
ثلاثة أسئلة: أيّ \en{command} أو \en{query} يعمل؟ أيّ \en{handler} ينفّذه؟ وماذا
يفعل الدور الآخر (أدمن، أو استشاريّ، أو جهة مانحة) ردًّا على فِعل سارة؟

\begin{idea}[الفكرة: النظام سلسلة أفعال وردود]
كل ما يحدث في \en{ScholarPath} هو أمرٌ يغيّر الحالة (\en{command})، أو استعلامٌ
يقرأها (\en{query}). لكن الأمر لا يعيش وحده: كثيرٌ منه يُطلق حدثًا دوميًّا
(\en{domain event}) بعد نجاحه، وهذا الحدث يوقظ دورًا آخر في الطرف المقابل. رحلة
سارة إذًا ليست خطًّا مستقيمًا، بل حوارٌ متبادَل بين أربعة أدوار.
\end{idea}

\subsection{خريطة اليوم: الرحلة في لمحة}

قبل التفاصيل، هذه خريطة الرحلة كلها. إحدى عشرة محطّة، ملوّنة حسب المرحلة: التسجيل،
ثم الاكتشاف والذكاء الاصطناعيّ، ثم التقديم، ثم الحجز والدفع، ثم الجلسة، ثم المجتمع.
اقرأيها من أعلى لأسفل — كل صندوق يذكر الأمر الذي يعمل والدور الذي يردّ.

\begin{center}
\begin{LTR}
\begin{tikzpicture}[node distance=0.26cm]
  \node[wbox,fill=spLightGreen,draw=spWhy!50] (a1)
    {\small 1 — سارة تُنشئ حسابها وتؤكّد بريدها \en{(ASP.NET Identity)}};
  \node[wbox,below=0.26cm of a1,fill=spLightGreen,draw=spWhy!50] (a2)
    {\small 2 — تُكمل ملفها والـ \en{onboarding}: المستوى والمجالات والدول المفضّلة};
  \node[wbox,below=0.26cm of a2,fill=spLightPurple,draw=spQA!45] (a3)
    {\small 3 — الذكاء الاصطناعيّ يقترح منحًا ويفحص الأهليّة \en{(LocalAiService)}};
  \node[wbox,below=0.26cm of a3,fill=spLightBlue,draw=spPrimary!45] (a4)
    {\small 4 — تكتشف منحة وتقدّم طلبًا \en{(SubmitApplicationCommand)}};
  \node[wbox,below=0.26cm of a4,fill=spLightBlue,draw=spPrimary!45] (a5)
    {\small 5 — الجهة المانحة تراجع وتردّ (قبول أو رفض بسبب) والإشعار يصل للجرس};
  \node[wbox,below=0.26cm of a5,fill=spLightOrange,draw=spSec!45] (a6)
    {\small 6 — تحجز جلسة استشاريّ وتدفع بـ \en{Stripe}: الحجز يُجمِّد المبلغ};
  \node[wbox,below=0.26cm of a6,fill=spLightOrange,draw=spSec!45] (a7)
    {\small 7 — الاستشاريّ يقبل \en{(AcceptBookingCommand)}: القبض يأخذ المبلغ};
  \node[wbox,below=0.26cm of a7,fill=spLightTeal,draw=spAccent!45] (a8)
    {\small 8 — جلسة فيديو \en{ACS} داخل المتصفّح مع تسجيل تلقائيّ};
  \node[wbox,below=0.26cm of a8,fill=spLightTeal,draw=spAccent!45] (a9)
    {\small 9 — سارة تقيّم الجلسة \en{(SubmitConsultantRatingCommand)}};
  \node[wbox,below=0.26cm of a9,fill=spLightYellow,draw=spYellow!70!black] (a10)
    {\small 10 — تنشر سؤالًا في المجتمع ويصل الردّ لحظيًّا عبر \en{SignalR}};
  \node[wbox,below=0.26cm of a10,fill=spLightYellow,draw=spYellow!70!black] (a11)
    {\small 11 — تسأل مساعد الذكاء الاصطناعيّ فيردّ بإجابة مع مصادرها \en{(RAG)}};

  \foreach \f/\t in {a1/a2,a2/a3,a3/a4,a4/a5,a5/a6,a6/a7,a7/a8,a8/a9,a9/a10,a10/a11}
    \draw[arr] (\f) -- (\t);
\end{tikzpicture}
\end{LTR}
\end{center}

الأخضر مرحلة التسجيل، والبنفسجيّ مرحلة الذكاء الاصطناعيّ، والأزرق مرحلة التقديم،
والبرتقاليّ مرحلة الحجز والدفع، والتركوازيّ مرحلة الجلسة، والأصفر مرحلة المجتمع.
لاحظي أنّ خطوات مثل 5 و7 و10 يقوم بها دورٌ غير سارة — هذه هي الردود.

\subsection{المحطّة الأولى: التسجيل وتأكيد البريد}

تبدأ سارة بنموذج تسجيل بسيط: اسم، وبريد، وكلمة مرور. الواجهة ترسل طلبًا إلى
\code{AuthController}، فيُنشئ \en{ASP.NET Identity} صفَّ مستخدم جديدًا بكلمة مرور
مجزّأة (\en{hashed}) لا مخزّنة كنصّ. النظام «مبوَّب» (\en{gated}): لا شيء يعمل قبل
تأكيد البريد. لذا يُرسَل بريد فيه رابط رمزيّ (\en{token})، وحين تضغط عليه سارة
يتأكّد الحساب ويصبح دورها الافتراضيّ \en{Student}.

\begin{why}
لماذا نُبوِّب النظام خلف تأكيد البريد؟ لأنّ حسابات وهميّة تُلوّث التوصيات، وتُغرق
الجهات المانحة بطلبات مزيّفة، وتستهلك ميزانيّة الذكاء الاصطناعيّ. تأكيد البريد
بوّابةٌ رخيصة تُبقي مجتمع المنصّة حقيقيًّا منذ اللحظة الأولى.
\end{why}

\begin{note}
هنا لا يوجد «دور آخر يردّ» بعد — سارة وحدها. لكنّ كل فِعل حسّاس من الآن فصاعدًا
سيُسجَّل في سجلّ التدقيق (\en{audit log}) عبر سِمة \code{[Auditable]} كما رأينا سابقًا.
\end{note}

\subsection{المحطّة الثانية: استكمال الملف والـ onboarding}

بعد الدخول تمرّ سارة بخطوات \en{onboarding} تجمع بيانات ملفها: مستواها الأكاديميّ
(\en{AcademicLevel})، ومجالات دراستها المفضّلة (\en{PreferredFields})، والدول التي
تريد الدراسة فيها (\en{PreferredCountries}). قد تبدو خطوة شكليّة، لكنّها في الحقيقة
وقود الذكاء الاصطناعيّ: بدونها لا يملك المحرّك ما يطابق به المنح.

\begin{keyp}[نقطة مهمة: الملف هو مُدخَل الخوارزميّة]
كل حقل تملؤه سارة الآن سيُقرَأ حرفيًّا في المحطّة التالية: \code{AcademicLevel} يعطي
٣٠ نقطة عند تطابق المستوى، و\code{PreferredFields} حتى ٤٠ نقطة، و\code{PreferredCountries}
حتى ٢٥ نقطة. الملف الفارغ يعني توصيات بصفر — لذا نحثّ الطالب على إكماله.
\end{keyp}

\subsection{المحطّة الثالثة: الذكاء الاصطناعيّ يقترح ويفحص الأهليّة}

الآن يتدخّل الذكاء الاصطناعيّ لأوّل مرّة. حين تفتح سارة لوحتها، ترسل الواجهة
\code{GenerateRecommendationsCommand}. يستقبله \code{GenerateRecommendationsCommandHandler}
الذي يمرّ أوّلًا ببوّابة التكلفة (\code{AiCostGate}) لئلّا يتجاوز حدّ اليوم، ثمّ
ينادي منفذ \code{IAiService}. في بيئتنا الحاليّة يُحقَن المحوّل \code{LocalAiService}:
محرّكٌ حسابيّ بحت بلا نداء لـ \en{GPT}، يعطي نتائج قابلة للتفسير والاختبار.

هذه هي قلب الخوارزميّة — حلقة تحسب نقاط كل منحة مفتوحة مقابل ملف سارة:

\begin{english}
\begin{lstlisting}[language={[Sharp]C}]
foreach (var c in candidates)
{
    var score = 0;
    // level match = 30
    if (userLevel.HasValue && userLevel.Value == c.TargetLevel) score += 30;
    // field overlap = up to 40, country overlap = up to 25
    score += OverlapScore(preferredFields, fosFields, maxPoints: 40);
    score += OverlapScore(NormalizeCountries(preferredCountries),
                          NormalizeCountries(countries), maxPoints: 25);
    // funding bonus = 5 for large awards
    if (c.FundingAmountUsd >= 10_000) score += 5;

    scored.Add((c.Id, Math.Min(score, 100), c.TitleEn, c.TitleAr));
}
var top = scored.OrderByDescending(x => x.Score).Take(topN).ToList();
\end{lstlisting}
\end{english}

هنشوف إيه في الأسطر دي، واحدة واحدة:
\begin{itemize}
\item \code{if (userLevel == c.TargetLevel) score += 30}: تطابق المستوى الأكاديميّ
يعطي ٣٠ نقطة ثابتة — أقوى إشارة مفردة.
\item \code{OverlapScore(preferredFields, ...)}: التقاطع بين مجالات سارة ومجالات
المنحة يعطي حتى ٤٠ نقطة، موزَّعة بحسب عدد المجالات المشتركة.
\item \code{NormalizeCountries(...)}: نوحّد أسماء الدول («مصر» و\en{Egypt} و\en{EG})
إلى مفتاح واحد قبل المقارنة، حتى لا يفشل التطابق بسبب اختلاف الكتابة.
\item \code{Math.Min(score, 100)}: نسقُف النتيجة عند ١٠٠ لتبقى نسبةً مفهومة.
\item \code{OrderByDescending(...).Take(topN)}: نرتّب تنازليًّا ونأخذ الأعلى فقط.
\end{itemize}

\begin{idea}[الفكرة: ذكاء حسابيّ يمكن الدفاع عنه]
لأنّ الحساب صريح (٣٠ + ٤٠ + ٢٥ + ٥)، تستطيع سارة أن تفهم لماذا ظهرت لها منحة بعينها،
وتستطيع اللجنة أن تتحقّق من النتيجة بالورقة والقلم. لا صندوق أسود، ولا هلوسة، ولا
تكلفة استدعاء نموذج خارجيّ. النتيجة تُخزَّن في صفّ \code{AiInteraction} لتتبّع التكلفة
والاستخدام.
\end{idea}

بجانب الترشيح، تضغط سارة على «افحص الأهليّة» في صفحة منحة، فيعمل
\code{CheckEligibilityCommand} ويقارن \code{LocalAiService} ملفها بثلاثة معايير:
المستوى، والدولة، والمجال. ثمّ تجمع الدالة \code{DeriveVerdict} هذه المعايير في حكم
واحد بقاعدة بسيطة: أيّ «لا» صريحة تعني \en{NotEligible}، وأيّ «جزئيّ» أو «غير معروف»
يعني \en{PartiallyEligible}، وكلّها «نعم» يعني \en{Eligible}.

\begin{center}
\begin{LTR}
\begin{tikzpicture}[node distance=0.7cm and 1.0cm]
  \node[sbox,text width=3.0cm] (in) {\small profile + listing criteria};
  \node[dbox,right=1.0cm of in,text width=3.0cm] (d) {\small any criterion = ``no''?};
  \node[fbox,right=1.1cm of d] (no) {\small NotEligible};
  \node[dbox,below=0.7cm of d,text width=3.0cm] (d2) {\small any ``partial'' / ``unknown''?};
  \node[tbox,right=1.1cm of d2] (part) {\small PartiallyEligible};
  \node[srcbox,below=0.7cm of d2,text width=2.6cm] (yes) {\small Eligible};

  \draw[arr] (in) -- (d);
  \draw[arr] (d) -- node[lbl,above] {yes} (no);
  \draw[arr] (d) -- node[lbl,right] {no} (d2);
  \draw[arr] (d2) -- node[lbl,above] {yes} (part);
  \draw[arr] (d2) -- node[lbl,right] {no} (yes);
\end{tikzpicture}
\end{LTR}
\end{center}

\begin{keyp}[نقطة مهمة: الحكم مُحافِظ عمدًا]
لاحظي أنّ «غير المعروف» يُعامَل مثل «الجزئيّ» لا مثل «النعم». لو كان ملف سارة ناقصًا،
لن يعِدَها النظام بأهليّة كاملة زورًا، بل يقول «جزئيّ» ويطلب منها إكمال الملف. هذا
اختيار تصميميّ محافظ: نُفضّل تحفُّظًا صادقًا على تفاؤل كاذب.
\end{keyp}

\subsection{المحطّة الرابعة: اكتشاف منحة وتقديم طلب}

تختار سارة منحة وتبدأ طلبًا. الطلب يعيش أوّلًا كمسوَّدة (\en{Draft})، تملأ فيها
النموذج وترفع المستندات. حين تضغط «قدِّم»، ترسل الواجهة \code{SubmitApplicationCommand}
فينفّذه \code{SubmitApplicationCommandHandler}. لا يكتفي المعالج بتغيير الحالة، بل
يحرس عدّة قواعد: يتحقّق من ملكيّة الطلب، ويمنع التقديم بعد انتهاء الموعد النهائيّ،
ويتأكّد أنّ كل مستند مرفق يعود فعلًا لملف رفعته سارة عبر خزنة المستندات المفحوصة.

\begin{english}
\begin{lstlisting}[language={[Sharp]C}]
// state machine: Draft -> Pending only if the transition is legal
ApplicationStateMachine.EnsureTransition(oldStatus, ApplicationStatus.Pending);
// deadline guard, even for a Draft that sat past the deadline
if (application.Scholarship is { } sch && sch.Deadline < DateTimeOffset.UtcNow)
    throw new ConflictException("This scholarship's deadline has passed.");

application.Status = ApplicationStatus.Pending;
application.SubmittedAt = DateTimeOffset.UtcNow;
application.RaiseDomainEvent(new ApplicationSubmittedEvent(
    application.Id, application.StudentId, application.ScholarshipId));
await db.SaveChangesAsync(ct);
\end{lstlisting}
\end{english}

الفكرة في الأسطر دي:
\begin{itemize}
\item \code{EnsureTransition}: آلة حالة صغيرة ترفض أيّ انتقال غير قانونيّ (مثلًا من
مقبول إلى مسوَّدة) — القاعدة في مكان واحد لا مبعثرة.
\item حارس الموعد النهائيّ: يمنع تقديم مسوَّدة قديمة بعد فوات الأوان.
\item \code{RaiseDomainEvent(new ApplicationSubmittedEvent(...))}: يسجّل الحدث على
الكيان، ثمّ يُنشَر بعد \code{SaveChangesAsync} فيوقظ الجهة المانحة.
\end{itemize}

\subsection{المحطّة الخامسة: الجهة المانحة تراجع وتردّ (أوّل ردّ فعل)}

الآن يدخل دورٌ آخر للمرّة الأولى. حين قدّمت سارة طلبها، وصل إشعار للجهة المانحة
صاحبة المنحة. تفتح الجهة الطلب وتقرّر: قبول، أو وضعه تحت المراجعة، أو رفض بسبب.
القرار يمرّ عبر \code{ReviewApplicationCommand} الذي ينفّذه المعالج التالي —
لاحظي حارس الملكيّة والإشعار الذي يُرسَل لسارة في نهايته:

\begin{english}
\begin{lstlisting}[language={[Sharp]C}]
// only the scholarship's owning provider may review its applications
if (application.Scholarship.OwnerScholarshipProviderId != currentUser.UserId)
    throw new ForbiddenAccessException();

ApplicationStateMachine.EnsureTransition(oldStatus, request.Status);
application.Status = request.Status;
application.DecisionAt = DateTimeOffset.UtcNow;
application.DecisionReason = request.DecisionReason;   // e.g. reject reason

await db.SaveChangesAsync(ct);
await notifications.DispatchAsync(
    application.StudentId, NotificationType.ApplicationStatusChanged,
    new NotificationParams { StatusText = request.Status.ToString(),
                             Reason = request.DecisionReason }, null, null, ct);
\end{lstlisting}
\end{english}

هنشوف الردّ ماشي إزّاي:
\begin{itemize}
\item \code{OwnerScholarshipProviderId != currentUser.UserId}: حارسٌ يمنع أيّ جهة من
لمس طلبات منحة لا تملكها — عزلٌ صارم بين الجهات.
\item \code{DecisionReason}: عند الرفض تُلزَم الجهة بسببٍ نصّيّ، فتراه سارة صراحةً بدل
رفضٍ صامت — قاعدة عمل حقيقيّة في النظام.
\item \code{notifications.DispatchAsync(application.StudentId, ...)}: يُوجَّه الإشعار
إلى سارة تحديدًا، فيظهر في جرسها لحظيًّا عبر \en{SignalR} مع نصّ الحالة والسبب.
\end{itemize}

\begin{spotlight}[ركّز هنا لإبهار اللجنة: كل خطوة لها ردّ فعل من دورٍ آخر]
هذه أوّل نقطة تُظهر فيها الرحلةُ جوهرَها: فِعلُ سارة (تقديم الطلب) أطلق حدثًا،
والحدثُ أيقظ الجهة المانحة، وردُّ الجهة أطلق إشعارًا رجع إلى سارة. أربعة أدوار
(طالب، وجهة، وأدمن، واستشاريّ) لا تتحدّث مباشرةً، بل تتحاور عبر أحداث وإشعارات
مصمَّمة. أظهر للجنة هذا «التعاون غير المتزامن»: تكاملٌ بلا اقتران متين.
\end{spotlight}

\subsection{المحطّة السادسة: حجز استشاريّ ودفع Stripe (الحجز يُجمّد المبلغ)}

تريد سارة مشورةً بشريّة، فتحجز جلسة مع استشاريّ. عند الحجز يعمل
\code{RequestBookingCommand}. هنا تظهر لمسة الدفع الذكيّة: النظام لا يسحب المال فورًا،
بل يُنشئ نيّة دفع \en{Stripe} بأسلوب القبض اليدويّ
(\code{captureMethod: "manual"}). هذا يعني أنّ بطاقة سارة تُصرَّح (\en{authorized})
والمبلغ يُجمَّد، لكنّه لا يُؤخَذ بعد. يبقى الحجز في حالة \en{Requested} وينتظر ردّ
الاستشاريّ، ويصل إشعار للاستشاريّ بأنّ طلبًا ينتظره.

\begin{why}
لماذا نجمّد المبلغ بدل سحبه؟ لأنّ الاستشاريّ قد يرفض أو لا يردّ. لو سحبنا المال
فورًا لاضطررنا لاستردادٍ لاحقٍ معقّد. التجميد (\en{authorize}) يحجز المبلغ دون أخذه،
فإن قبل الاستشاريّ نقبض، وإن رفض نُفرِج عن التجميد بلا حركة ماليّة فعليّة. المال
«في الضمان» (\en{escrow}) حتى تُحسَم الجلسة.
\end{why}

\subsection{المحطّة السابعة: الاستشاريّ يقبل — نتبّع الأمر من أوّله لآخره}

هذه أهمّ محطّة تقنيًّا، فلنتتبّع أمرًا واحدًا كاملًا من الواجهة إلى الإشعار. حين يضغط
الاستشاريّ «اقبل»، يعمل \code{AcceptBookingCommand}. أوّلًا تعريف الأمر نفسه — سطرٌ
واحد يحمل مُعرِّف الحجز، ومزيّنٌ بسِمة التدقيق:

\begin{english}
\begin{lstlisting}[language={[Sharp]C}]
[Auditable(AuditAction.BookingAccepted, "ConsultantBooking",
    TargetIdProperty = nameof(BookingId),
    SummaryTemplate = "Booking accepted: {TargetId}")]
public sealed record AcceptBookingCommand(Guid BookingId) : IRequest;
\end{lstlisting}
\end{english}

سِمة \code{[Auditable]} تجعل \en{MediatR} يكتب صفَّ تدقيق تلقائيًّا قبل تنفيذ المعالج،
بلا سطر واحد داخل المنطق. الآن قلب المعالج \code{AcceptBookingCommandHandler}:

\begin{english}
\begin{lstlisting}[language={[Sharp]C}]
// 1. authorization: only the owning consultant, only a Requested booking
if (!_currentUser.IsInRole("Consultant")) throw new UnauthorizedAccessException(...);
if (booking.ConsultantId != consultantId) throw new UnauthorizedAccessException(...);
if (booking.Status != BookingStatus.Requested) throw new BookingDomainException(...);

// 2. capture the money that was only held at request time
var captureResult = await _stripeService.CapturePaymentIntentAsync(
    booking.StripePaymentIntentId!, amountCents, idempotencyKey, ct);
payment.Status = PaymentStatus.Captured;   // + platform / payee split

// 3. flip state and provision the ACS video room
booking.Status = BookingStatus.Confirmed;
var room = await _meetingService.CreateRoomAsync(booking.Id, ct);
booking.MeetingRoomId = room.RoomId;

// 4. persist, THEN publish the domain event
await _context.SaveChangesAsync(cancellationToken);
await _publisher.Publish(
    new BookingConfirmedEvent(booking.Id, booking.StudentId, booking.ConsultantId), ct);
\end{lstlisting}
\end{english}

هنشوف المعالج ماشي إزّاي، مرحلة مرحلة:
\begin{itemize}
\item التصريح: ثلاث حراسات تتأكّد أنّ المُنادِي استشاريّ، وأنّه صاحب الحجز، وأنّ
الحجز ما زال \en{Requested} — لا يُقبَل حجزٌ مرّتين.
\item \code{CapturePaymentIntentAsync}: هنا يتحوّل التجميد إلى قبضٍ فعليّ. المبلغ
الذي جُمِّد في المحطّة السابقة يُؤخَذ الآن، وتُحدَّث حصّة المنصّة وحصّة الاستشاريّ.
\item \code{CreateRoomAsync}: تُهيَّأ غرفة \en{ACS} للجلسة ويُحفَظ \code{MeetingRoomId}
على الحجز — نفس الغرفة الكسولة التي شرحناها سابقًا.
\item \code{SaveChangesAsync} ثمّ \code{Publish}: أوّلًا نثبّت الحقائق في قاعدة البيانات،
وبعد نجاح الحفظ فقط نُطلق \code{BookingConfirmedEvent}. الترتيب مقصود: لا نُشعِر أحدًا
بشيء لم يُحفَظ.
\end{itemize}

هذا هو تدفّق الأمر الواحد مرسومًا — من الواجهة حتى الجرس:

\begin{center}
\begin{LTR}
\begin{tikzpicture}[node distance=0.5cm and 0.55cm]
  \node[actor,text width=2.2cm] (fe) {React\\frontend};
  \node[pbox,right=0.55cm of fe,text width=2.7cm] (ctrl) {BookingsController\\.Accept};
  \node[gbox,right=0.55cm of ctrl,text width=2.4cm] (med) {IMediator\\.Send};
  \node[pbox,right=0.55cm of med,text width=2.7cm] (hdlr) {AcceptBooking\\Handler};

  \node[tbox,below=0.75cm of hdlr,text width=2.7cm] (stripe) {Stripe capture\\+ ACS room};
  \node[dbbox,left=0.6cm of stripe,text width=2.2cm] (db) {SQL Server\\SaveChanges};
  \node[tbox,left=0.55cm of db,text width=2.6cm] (evt) {Publish\\BookingConfirmedEvent};
  \node[sbox,left=0.55cm of evt,text width=2.5cm] (notif) {NotificationHub\\bell -> student};

  \draw[arr] (fe) -- node[lbl,above]{POST} (ctrl);
  \draw[arr] (ctrl) -- node[lbl,above]{1 line} (med);
  \draw[arr] (med) -- (hdlr);
  \draw[arr] (hdlr) -- (stripe);
  \draw[arr] (stripe) -- (db);
  \draw[arr] (db) -- node[lbl,above]{then} (evt);
  \draw[arr] (evt) -- (notif);
  \draw[arrd] (notif.south) to[out=270,in=270] node[lbl,below]{push} (fe.south);
\end{tikzpicture}
\end{LTR}
\end{center}

من يلتقط \code{BookingConfirmedEvent}؟ معالجُ حدثٍ منفصل تمامًا،
\code{BookingConfirmedEventHandler}، لا يعرف شيئًا عن \en{Stripe} ولا عن الغرفة —
مهمّته الوحيدة إشعار سارة أنّ حجزها تأكّد:

\begin{english}
\begin{lstlisting}[language={[Sharp]C}]
public async Task Handle(BookingConfirmedEvent notification, CancellationToken ct)
{
    var view = await BookingNotificationContext.LoadAsync(db, notification.BookingId, ct);
    if (view is null) { /* log + return */ return; }

    await notifications.DispatchAsync(
        notification.StudentId,                       // recipient = the student
        NotificationType.BookingConfirmed,
        new NotificationParams { CounterpartyName = view.ConsultantName,
            StartAtText = BookingNotificationContext.FormatStartAt(view.ScheduledStartAt) },
        deepLink: $"/student/bookings/{notification.BookingId}",
        idempotencyKey: $"booking-confirmed:{notification.BookingId:N}", ct);
}
\end{lstlisting}
\end{english}

النقاط المهمّة في ردّ الفعل ده:
\begin{itemize}
\item \code{DispatchAsync(notification.StudentId, ...)}: الإشعار يُوجَّه لسارة تحديدًا،
فيصل جرسها لحظيًّا عبر \en{SignalR}.
\item \code{deepLink}: رابطٌ عميق يفتح صفحة الحجز مباشرة عند الضغط — لا بحث يدويّ.
\item \code{idempotencyKey}: مفتاحٌ فريد يمنع تكرار نفس الإشعار لو أُعيد إطلاق الحدث.
\item فصلٌ نظيف: معالج القبول لا يعرف كيف يُشعِر أحدًا، ومعالج الإشعار لا يعرف شيئًا
عن الدفع. كلٌّ منهما مسؤوليّة واحدة، والحدث هو الجسر بينهما.
\end{itemize}

\begin{spotlight}[ركّز هنا لإبهار اللجنة: رحلة واحدة تلمس كل أجزاء النظام]
تتبّعنا أمرًا واحدًا لمس كلَّ طبقة وخدمة تقريبًا: الواجهة، والمتحكّم، و\en{MediatR}،
والمعالج، و\en{Stripe} (دفع)، و\en{ACS} (فيديو)، وقاعدة البيانات، والحدث الدوميّ،
و\en{SignalR} (إشعار)، وسجلّ التدقيق (\en{[Auditable]}). هذا هو التماسك المعماريّ:
كل قطعة في مكانها، ومسارٌ واحد يشبك بينها جميعًا بلا فوضى. أظهر للجنة هذه اللقطة —
إنّها ملخّص النظام كلّه في أمرٍ واحد.
\end{spotlight}

\subsection{المحطّتان الثامنة والتاسعة: الجلسة والتقييم}

في موعد الجلسة يضغط الطرفان «انضمّ» ضمن نافذة \en{$\pm$15} دقيقة. يعمل
\code{RecordMeetingJoinCommand} فيُصدِر رمز وصول \en{ACS} قصير العمر لكل مشارك،
وتجري المكالمة بين المتصفّحَين مباشرةً بينما يبقى خادمنا مجرّد منسّق. عند الاتصال
يبدأ التسجيل تلقائيًّا ويُخزَّن في \en{Blob Storage}. النظام يسجّل أيضًا وقت انضمام
كل طرف، وهذه الإشارة يقرؤها لاحقًا مسحُ الغياب الآليّ (\en{no-show}) عبر \en{Hangfire}.

بعد الجلسة تقيّم سارة الاستشاريّ عبر \code{SubmitConsultantRatingCommand}: نجوم من
واحد إلى خمسة مع تعليق اختياريّ. التقييم يُغذّي متوسّط الاستشاريّ الظاهر لكل الطلاب
مستقبلًا، وبذلك يردّ فعلُ سارة الصغير على تجربة طلابٍ لم تلتقِهم بعد.

\subsection{المحطّة العاشرة: منشور في المجتمع وردّ حيّ}

تطرح سارة سؤالًا في المجتمع، فيُنشَأ منشور ثنائيّ اللغة. طالبٌ آخر متّصلٌ الآن يرى
الردّ يظهر لحظيًّا دون تحديث الصفحة، لأنّ \en{CommunityHub} يدفع التحديث عبر
\en{WebSocket}. هنا يظهر دورٌ آخر من الأدوار: لو خالف منشورٌ القواعد، يستطيع الأدمن
إخفاءه أو الردّ على بلاغ — نفس نمط الحدث والإشعار الذي رأيناه، لكن في سياق الإشراف.

\begin{note}
لاحظي أنّ الفكرة نفسها تتكرّر في كل محطّة: أمرٌ يغيّر الحالة، ثمّ حدثٌ يُنشَر، ثمّ
دفعٌ لحظيّ عبر \en{SignalR} لمن يعنيه الأمر. تعلّمُ نمطٍ واحد يفتح لك النظام كلّه.
\end{note}

\subsection{المحطّة الحادية عشرة: سؤال لمساعد الذكاء الاصطناعيّ بمصادر}

تختم سارة يومها بسؤال مساعد الشات: «ما الأوراق المطلوبة عادةً؟». يعمل أمر الشات،
وينادي \code{LocalAiService.AskAsync}. هنا يتجلّى \en{RAG} (التوليد المُعزَّز
بالاسترجاع): أوّلًا يسترجع المنفذ \code{IKnowledgeRetriever} أقرب الوثائق من قاعدة
المعرفة، ثمّ يُبقي فقط ما تجاوزت درجتُه العتبة \code{RagMinScore}، ثمّ يبني الإجابة
من أعلى وثيقة مطابقة — ويُرفِق مصادرها.

\begin{english}
\begin{lstlisting}[language={[Sharp]C}]
var docs = await retriever.RetrieveAsync(msg, opts.Value.RagTopK, ct);
var relevant = docs.Where(d => d.Score >= opts.Value.RagMinScore).ToList();

string reply = relevant.Count > 0
    ? $"{lead}\n\n{(arabic ? relevant[0].ContentAr : relevant[0].ContentEn).Trim()}"
    : RouteChat(msg);            // keyword fallback when nothing scores high enough

return new AiChatResponse(Message: reply, /* tokens, cost */
    Sources: RagSupport.ToSources(relevant, arabic));   // cited sources
\end{lstlisting}
\end{english}

المهمّ في الأسطر دي:
\begin{itemize}
\item \code{RetrieveAsync}: يبحث في قاعدة المعرفة عن أقرب \code{RagTopK} وثيقة للسؤال.
\item \code{d.Score >= RagMinScore}: عتبةُ ثقة — لو لم تتجاوز أيّ وثيقة العتبة، لا نخترع
إجابة، بل نرجع لموجِّه كلماتٍ مفتاحيّة (\code{RouteChat}) يعطي نصيحة عامّة آمنة.
\item \code{Sources: RagSupport.ToSources(...)}: تُرفَق المصادر مع الردّ، فترى سارة
من أين جاءت الإجابة — لا صندوق أسود.
\end{itemize}

\begin{keyp}[نقطة مهمة: نفس التصميم يعمل بلا مفتاح خارجيّ]
هذا المسار «المحلّيّ» يشغّل \en{RAG} كاملًا بلا استدعاء \en{GPT} إطلاقًا. لو ضُبِط
مزوّد \en{Azure OpenAI}، يُحقَن محوّلٌ آخر يولّد إجابةً حرّة فوق نفس الوثائق المسترجَعة.
تبديل المزوّد قرارُ إعداد لا قرار كود — نفس نمط المنافذ في كل النظام.
\end{keyp}

\subsection{من فعل ماذا: خريطة الأدوار في الرحلة}

لنجمع الرحلة في جدول واحد يبيّن الأمر عند كل خطوة والدور الذي يردّ:

\begin{center}
\begin{tabular}{L{3.4cm} L{4.3cm} L{4.3cm}}
\toprule
المرحلة & الأمر / المعالج & الدور الذي يردّ \\
\midrule
التسجيل & \en{Identity} + تأكيد البريد & لا أحد (بوّابة دخول) \\
التوصية والأهليّة & \en{GenerateRecommendations} & \en{LocalAiService} (حسابيّ) \\
تقديم الطلب & \en{SubmitApplicationCommand} & الجهة المانحة تُشعَر \\
مراجعة الطلب & \en{ReviewApplicationCommand} & سارة تُشعَر بالنتيجة والسبب \\
حجز الجلسة & \en{RequestBookingCommand} & الاستشاريّ يُشعَر (المبلغ مُجمَّد) \\
قبول الحجز & \en{AcceptBookingCommand} & سارة تُشعَر (القبض يتمّ) \\
المجتمع والإشراف & منشور / بلاغ & طالب آخر أو أدمن يردّ \\
\bottomrule
\end{tabular}
\end{center}

\begin{qa}[سؤال متوقّع: هل الأدوار مقترنة اقترانًا متينًا؟]
لا. لاحظ أنّ أيّ معالج لا ينادي دورًا آخر مباشرةً. معالج القبول يُطلق حدثًا فقط،
ولا يعرف أنّ هناك معالج إشعار يستمع. لو حذفنا الإشعار غدًا، لن يتغيّر سطرٌ في منطق
القبول. هذا اقتران رخو (\en{loose coupling}) عبر الأحداث: الأدوار تتعاون دون أن
يعرف بعضُها تفاصيل بعض.
\end{qa}

\begin{summary}[الخلاصة في سطور]
\begin{itemize}
\item تابعنا سارة عبر إحدى عشرة محطّة، فلمسنا كل جزء من النظام في قصّة واحدة متّصلة.
\item كل خطوة أمرٌ (\en{command}) يغيّر الحالة، وكثيرٌ منها يُطلق حدثًا دوميًّا بعد
\code{SaveChangesAsync}، والحدثُ يوقظ دورًا آخر عبر إشعار لحظيّ.
\item تتبّعنا \code{AcceptBookingCommand} كاملًا: واجهة، فمتحكّم، فمعالج، فقبضُ \en{Stripe}،
فغرفة \en{ACS}، فحفظٌ، فحدثٌ، فإشعارٌ للجرس — رحلة أمرٍ واحدة تلخّص المعمار كلّه.
\item الذكاء الاصطناعيّ ظهر ثلاث مرّات: توصيةٌ حسابيّة قابلة للتفسير (٣٠+٤٠+٢٥+٥)،
وفحصُ أهليّة بحكمٍ محافظ، وشاتٌ يردّ بإجابةٍ مع مصادرها عبر \en{RAG}.
\item الأدوار الأربعة (طالب، وجهة، وأدمن، واستشاريّ) تتحاور عبر أحداث وإشعارات لا عبر
نداءات مباشرة — تكاملٌ مصمَّم باقترانٍ رخو، وهو سرّ تماسك النظام.
\end{itemize}
\end{summary}

\clearpage
% ch10-defense-demo — دليل المناقشة والعرض الحيّ
\section{دليل المناقشة والعرض: نتكلم عن إيه ونعرض إزاي}

الفصول السابقة شرحت \en{ScholarPath} من الداخل: الطبقات، و\en{CQRS}، والتوثيق،
والذكاء الاصطناعيّ، والفيديو، والمدفوعات، والأمان. هذا الفصل مختلف تمامًا: هو دليل
عمليّ يجهّزك لِلحظة الحقيقة — الوقوف أمام لجنة الدكاترة. لن نضيف هنا معلومة جديدة
عن النظام؛ بدل ذلك نجيب عن سؤالين لكل ميزة: «نتكلم فيها عن إيه؟» و«نعرضها حيّة
إزاي؟» — أي ماذا تفتح على الشاشة، وأين تضغط، وما الذي على اللجنة أن تراه بأمّ عينها.

لكل محور في هذا الفصل أربعة عناصر ثابتة: جملة تفتتح بها الكلام، نقاط القوّة التي
تتكلّم فيها، خطوات عرضٍ حيّ واقعية، وسؤال متوقّع من اللجنة مع إجابة قصيرة واثقة.
اقرأه مرّة قبل المناقشة، ثمّ استعمله كخريطة أثناءها.

% ====================================================================
\subsection{كيف تدير المناقشة؟}

مبدأ واحد يحكم الجلسة كلها، احفظه جيدًا: «لا تدّعِ شيئًا لا تستطيع إثباته على
الشاشة». اللجنة سمعت مئات العروض التي تصف ميزات لا وجود لها؛ ما يميّزك أنّ كل جملة
تقولها لها دليل — إمّا سطر كود تفتحه، أو زرّ تضغطه فيحدث الشيء أمامهم.

\begin{keyp}[نقطة مهمة: كل ادعاء له دليل]
رتّب كلامك على قاعدة «أقول ثمّ أُرِي». لا تقل «النظام يدعم الإشعارات اللحظيّة» ثمّ
تنتقل؛ قُلها ثمّ افتح نافذتين وابعث رسالة فيظهر الإشعار فورًا. الادعاء المصحوب
بعرضٍ حيّ يساوي عشرة ادعاءات بلا دليل، ويقطع الطريق على سؤال «طيّب وإثبات كده فين؟».
\end{keyp}

معظم عروضنا الحيّة تعتمد خدمات سحابية (\en{Azure OpenAI}, \en{ACS}, \en{Stripe}).
والسحابة قد تخذلك في أسوأ لحظة: النت يقطع، أو الخدمة تتأخّر. لهذا صُمِّم النظام منذ
البداية بمنطق «المنافذ والمحوّلات» (\en{ports/adapters})، وله لكل خدمة سحابية بديلٌ
محلّيّ (\en{Stub}) يشتغل \en{offline}.

\begin{note}[ملاحظة: خطة بديلة إن سقط النت]
لو انقطع الإنترنت أثناء ديمو سحابيّ، لا تتلعثم. اذكر بثقة أنّ النظام يُشغَّل بمنافذ
مجرّدة، وأنّ كل مزوّد خارجيّ (\en{Stripe}, \en{ACS}, \en{Azure OpenAI}) خلفه محوّل
\en{Stub} محلّيّ يعيد نتائج جاهزة بلا إنترنت — فتكمل عرض بقية الشاشات دون توقّف. هذا
نفسه نقطة قوّة معماريّة، لا اعتذار.
\end{note}

% ====================================================================
\subsection{خريطة عرضٍ في عشرين دقيقة}

قبل الدخول في التفاصيل، هذه خطّة زمنيّة تحكم إيقاع العرض كله. الفكرة أن توزّع
العشرين دقيقة على المحاور بحيث يبقى نصفها تقريبًا للعرض الحيّ لا للكلام. الجدول
يبيّن المدى الزمنيّ، والمحور، وما تفتحه على الشاشة في كل مرحلة:

\begin{center}
\begin{tabular}{C{2.2cm} L{5.2cm} L{7.0cm}}
\toprule
{\color{spPrimary}المدى الزمنيّ} & {\color{spPrimary}المحور} & {\color{spPrimary}تفتح إيه على الشاشة} \\
\midrule
0:00–3:00 & المقدمة والمشكلة & شرائح العرض: مشكلة تشتّت المنح والإرشاد، ثمّ الحلّ المقترَح \\
3:00–6:00 & المعماريّة & شجرة مجلّدات المشروع في المحرّر (\en{Domain/Application/Infrastructure/API}) \\
6:00–8:00 & التوثيق والتوكنات & المتصفّح مع \en{DevTools} على تبويبَي \en{Network} و\en{Application} \\
8:00–12:00 & الذكاء الاصطناعيّ (\en{RAG}) & صفحة إدارة قاعدة المعرفة $+$ نافذة الشات \\
12:00–15:00 & الفيديو (\en{ACS}) & نافذتان: حساب طالب وحساب مستشار في نفس الحجز \\
15:00–17:00 & الإشعارات والشات & نافذتان جنبًا لجنب لمستخدمَين مختلفَين \\
17:00–19:00 & المدفوعات (\en{Stripe}) & لوحة \en{Stripe} (وضع الاختبار) $+$ سجلّ الحجز \\
19:00–20:00 & الأمان & جدول \code{AuditLog} $+$ عمود مشفّر في القاعدة \\
\bottomrule
\end{tabular}
\end{center}

\begin{note}
راقب الساعة بطرف عينك. أكثر ما يفسد المناقشات هو الاستغراق في المقدمة حتى لا يبقى
وقت للعرض الحيّ — وهو بالضبط ما يريد الدكاترة رؤيته. لو ضاق الوقت، اختصر الكلام لا
الديمو.
\end{note}

% ====================================================================
\subsection{المحور الأول: المعماريّة (Clean Architecture و CQRS)}

افتتح بجملة: «قبل أي ميزة، نظامنا مبنيّ على معماريّة نظيفة تجعل كل جزء يعرف مكانه —
وهذا ما سمح لخمسة مطوّرين أن يعملوا متوازين بأقلّ تعارض ممكن».

نقاط القوّة التي تتكلّم فيها:
\begin{itemize}
\item الاعتماد يتّجه للداخل دائمًا: \en{Domain} في القلب لا يعرف شيئًا عن قاعدة
  البيانات ولا عن \en{ASP.NET}، والبنية التحتيّة هي التي تعتمد على القلب لا العكس.
\item الـ\en{controller} رفيع جدًّا: لا يحمل منطق عمل، فقط يمرّر الطلب.
\item كل حالة استخدام لها \en{handler} مستقلّ يملكه مطوّر واحد؛ فالتعديل في ميزة
  لا يلمس ملفّات ميزة أخرى، ممّا قلّل تعارضات الدمج (\en{merge conflicts}) بين
  الخمسة إلى أدنى حدّ.
\end{itemize}

\begin{idea}[اعرض كده: رِقّة الـ controller ثم الـ handler]
\begin{itemize}
\item افتح شجرة المجلّدات في المحرّر، ووسّع \code{server/src} لتُظهِر الطبقات
  الأربع جنبًا لجنب: \en{Domain}، \en{Application}، \en{Infrastructure}، \en{API}.
\item افتح أيّ \en{controller} وأشِر إلى أنّ الدالة سطر أو ثلاثة أسطر فقط، ثمّ افتح
  الـ\en{handler} المقابل لتُظهِر أن منطق العمل كلّه هناك، معزولًا وقابلًا للاختبار.
\end{itemize}
\end{idea}

المقتطف التالي يكفي لتوضيح رقّة الـ\en{controller}: إنّه يحوّل الطلب فقط إلى
\en{MediatR} ولا يقرّر شيئًا:

\begin{english}
\begin{lstlisting}[language={[Sharp]C}]
[HttpPost]
public async Task<IActionResult> Create(CreateBookingCommand cmd)
    => Ok(await sender.Send(cmd));   // the controller only forwards the request
\end{lstlisting}
\end{english}

\begin{qa}[سؤال متوقّع: CQRS مش overkill لمشروع تخرّج؟]
لا، وهذا ما يثبته حجم المشروع نفسه. عندنا نحو 211 متطلّبًا وخمسة مطوّرين؛ ومع هذا
الحجم كانت الحدود الواضحة (أمر لكل كتابة، استعلام لكل قراءة، \en{handler} مملوك
لمطوّر واحد) هي ما منع تعارضات الدمج وسمح بالعمل المتوازي. أضِف إلى ذلك أنّ نمط
الـ\en{pipeline} في \en{MediatR} أعطانا التحقّق والتسجيل والتدقيق «مجانًا» في مكان
واحد لكل الأوامر، بدل تكرارها في كل \en{controller}.
\end{qa}

% ====================================================================
\subsection{المحور الثاني: التوثيق والتوكنات (access مقابل refresh)}

هذا أكثر محور يُسأل عنه، فأتقِنه جيدًا. افتتح بجملة: «عند الدخول لا نصدر رمزًا
واحدًا بل رمزين، لكلٍّ دورٌ ونافذةُ خطرٍ مختلفة تمامًا».

نقاط القوّة:
\begin{itemize}
\item \en{Access JWT}: رمز حامل قصير العمر جدًّا — 15 دقيقة فقط — يُرفَق مع كل طلب،
  ويُفحَص محلّيًّا دون الرجوع للقاعدة (لذلك أبقيناه قصيرًا).
\item \en{Refresh token}: رمز عشوائيّ أطول عمرًا، وظيفته الوحيدة استبدال رمز الوصول
  المنتهي؛ يُدوَّر عند كل استخدام، ولا يُخزَّن خامًا بل مُجزَّأً (\en{SHA-256}) في
  الحقل \code{RefreshToken.TokenHash}.
\end{itemize}

\begin{idea}[اعرض كده: دورة حياة التوكن في DevTools]
\begin{itemize}
\item افتح \en{DevTools} على تبويب \en{Network}، سجّل الدخول، وأظهِر أن كل طلب يحمل
  ترويسة \code{Authorization: Bearer ...} فيها الـ\en{access JWT}.
\item انسخ الرمز إلى موقع فكّ \en{JWT} (أو أظهِر حمولته)، وأشِر إلى حقل انتهاء
  الصلاحيّة القصير.
\item انتظر انتهاء الصلاحيّة (أو أظهِر طلبًا يعيد \en{401})؛ عندها ينادي العميلُ
  \code{/api/auth/refresh} فيعود بزوجٍ جديد وتستمرّ الجلسة بلا إعادة إدخال كلمة المرور.
\item على تبويب قاعدة البيانات أظهِر أنّ عمود \code{RefreshToken.TokenHash} يحوي
  تجزئة (سلسلة \en{hex}) لا الرمز الخام.
\end{itemize}
\end{idea}

\begin{qa}[سؤال متوقّع: لو اتسرّب توكن يحصل إيه؟]
نفرّق بين الرمزين. لو تسرّب رمز الوصول فنافذته 15 دقيقة فقط ثمّ يصير عديم الفائدة.
ولو تسرّبت قاعدة البيانات فما فيها إلا تجزئات لا يمكن عكسها، فالتجزئة المسروقة لا
تُقبَل عند التحقّق. وفوق ذلك، رمز التحديث يُدوَّر عند كل استخدام مع كشف إعادة
الاستخدام: إن ظهر رمزٌ مُبطَل مرّة أخرى، نُلغي كامل عائلة رموز المستخدم فورًا كإجراء
ضدّ السرقة.
\end{qa}

\begin{qa}[سؤال متوقّع: ليه RS256 مش HS256؟]
لأنّ \en{RS256} يفصل التوقيع عن التحقّق: مفتاح خاصّ يوقّع، ومفتاح عامّ يتحقّق فقط.
الخدماتُ التي تتحقّق من الرمز تحتاج المفتاح العامّ وحده — ويمكن مشاركته بلا خطر لأنّه
لا يستطيع التوقيع. أمّا \en{HS256} فسرّه الواحد يوقّع ويتحقّق معًا، فمن يقرأه في أيّ
خدمة يستطيع تزوير رموز. \en{RS256} أيضًا يجعل تدوير المفاتيح عمليّة تشغيليّة سهلة.
\end{qa}

\begin{secnote}
الرسالة الأمنيّة هنا في جملة: نفصل «ما يُرسَل للعميل» عن «ما يُخزَّن عندنا». العميل
يأخذ الرمز الخام، والقاعدة تحفظ تجزئته فقط. هذا المبدأ نفسه نطبّقه على رمز إعادة
تعيين كلمة المرور، فتسريب القاعدة لا يمنح المهاجم أيّ رمزٍ قابلٍ للاستعمال.
\end{secnote}

% ====================================================================
\subsection{المحور الثالث: الذكاء الاصطناعيّ (RAG)}

افتتح بجملة: «مساعدنا الذكيّ لا يخترع الإجابات — إنّه يقرأ من قاعدة معرفة مبنيّة من
بيانات منصّتنا نفسها، ويستشهد بمصادره».

نقاط القوّة:
\begin{itemize}
\item خطّ \en{RAG} حقيقيّ: نبني قاعدة معرفة من بيانات المنصّة، نبحث فيها عن السياق
  الأقرب للسؤال، ثمّ نطلب من النموذج إجابةً مستندةً إلى ذلك السياق فقط.
\item الردّ يعرض «المصادر» أسفله، فيستطيع المستخدم التحقّق بنفسه.
\item هناك سقف تكلفة يوميّ يحمي الميزانيّة، كما أنّ ميزتَي التوصيات والأهليّة
  حسابيّتان محلّيًّا بلا أيّ نداء لِـ\en{GPT}.
\end{itemize}

\begin{idea}[اعرض كده: إعادة البناء ثم سؤال بمصادر]
\begin{itemize}
\item افتح صفحة إدارة قاعدة المعرفة، واضغط زرّ \en{Rebuild}، وأظهِر الأعداد العائدة:
  كم مستندًا \en{upserted} و\en{reembedded} و\en{removed} — فهذا يثبت أنّ البناء
  حيّ لا وهميّ.
\item انتقل إلى الشات واسأل سؤالًا عن منحة بعينها، وأظهِر الإجابة ثمّ كتلة «المصادر»
  أسفلها التي تشير إلى المستندات التي استندت إليها.
\item اذكر شفهيًّا سقفَ التكلفة اليوميّ، وأنّ التوصيات محسوبة محلّيًّا بلا \en{GPT}.
\end{itemize}
\end{idea}

\begin{qa}[سؤال متوقّع: المساعد بيهلوس (يخترع معلومات)؟]
احتمالُ الاختلاق مكبوح بطبقتين. أولًا، الـ\en{system prompt} يمنع النموذج صراحةً من
اختراع ما ليس في السياق، ويُلزمه أن يستند إلى كتلة سياقٍ مرقّمة تُمرَّر إليه. ثانيًا،
لو لم يوجد أيّ مستند تجاوز عتبة التشابه (نحو 0.6) نُعيد ردًّا عامًّا («لا تتوفّر لديّ
معلومة كافية») بدل التخمين. فالمساعد مقيَّد بالمصادر، لا يتكلّم من فراغ.
\end{qa}

\begin{qa}[سؤال متوقّع: ليه مش قاعدة بيانات متجهات (vector DB)؟]
لأنّ الحجم لا يبرّرها. عندنا نحو 692 مستندًا فقط، ومسحُها في الذاكرة لحساب التشابه
يستغرق أقلّ من 10 أجزاء من الألف من الثانية. إدخالُ قاعدة متجهات مستقلّة هنا يزيد
التعقيد والتكلفة التشغيليّة بلا مكسبٍ محسوس. لو كبرت القاعدة إلى مئات الآلاف من
المستندات، عندها يصبح القرار مختلفًا — والمعماريّة تسمح بالتبديل خلف منفذٍ واحد.
\end{qa}

% ====================================================================
\subsection{المحور الرابع: الجلسات المرئيّة (ACS)}

افتتح بجملة: «الاستشارة نفسها تحدث داخل المنصّة — مكالمة فيديو مسجَّلة بين الطالب
والمستشار، بلا الحاجة لأيّ أداة خارجيّة».

نقاط القوّة:
\begin{itemize}
\item مكالمة الفيديو تجري بين المتصفّحَين وخدمة \en{Azure Communication Services}؛
  خادمُنا لا يمرّر بايتات الصوت والصورة، بل ينسّق فقط.
\item لكل مشارك رمز وصول قصير العمر خاصّ به، والجلسة تُسجَّل تلقائيًّا.
\item نتتبّع الحضور بدقّة (لحظة دخول كلٍّ من الطالب والمستشار)، فنكتشف الغياب
  (\en{no-show}) آليًّا لا يدويًّا.
\end{itemize}

\begin{idea}[اعرض كده: جلسة حيّة من نافذتين]
\begin{itemize}
\item افتح حجزًا مؤكَّدًا، وادخل من نافذتين منفصلتين: واحدة بحساب الطالب وأخرى بحساب
  المستشار.
\item اضغط «انضمّ» في كلٍّ منهما، وأظهِر أنّ الاتصال قام، وأنّ كلّ طرفٍ يرى كاميرا
  الآخر ويسمعه.
\item أظهِر بدء التسجيل عند اتصال المكالمة، واذكر أنّ أوقات الدخول تُسجَّل وأنّ
  مهمّةً خلفيّة تعلّم أيّ حجزٍ لم يحضره طرفُه بـ\en{no-show}.
\end{itemize}
\end{idea}

\begin{qa}[سؤال متوقّع: ليه توكن لكل مشارك مش توكن واحد للغرفة؟]
لأنّ رمز الفيديو هو بيانات اعتماد: من يحمله يدخل المكالمة. الرمز المشترَك لا يمكن
إبطاله لطرفٍ دون الآخر ولا يميّز مَن دخل. أمّا رمزٌ قصير العمر لكل مشارك فيربط الدخول
بهويّة، ويجعل تتبّع الحضور دقيقًا، ويقلّص أثر أيّ تسريب إلى مستخدمٍ واحد ولدقائق
معدودة — فيصير الرمز المسروق عديم الفائدة بعد انتهاء الجلسة.
\end{qa}

% ====================================================================
\subsection{المحور الخامس: الإشعارات والشات (أقوى ديمو حيّ)}

هذا هو العرض الذي يُبهر اللجنة أكثر من غيره لأنّه لحظيّ ومرئيّ. افتتح بجملة: «كل شيء
في المنصّة يحدث في الزمن الحقيقيّ — الرسالة تصل، والجرس يرنّ، بلا تحديث صفحة».

\begin{idea}[اعرض كده: نافذتان جنبًا لجنب]
\begin{itemize}
\item افتح نافذتين لمستخدمَين مختلفَين جنبًا لجنب على الشاشة (يفضَّل نافذة عاديّة
  وأخرى خفيّة \en{incognito} لتفصل الجلستين).
\item ابعث رسالة شات من الأولى، وأظهِر ظهورها فورًا عند الثانية مع رنين الجرس
  (الإشعار) — دون أن تلمس زرّ التحديث.
\item فعّل «عدم الإزعاج» (\en{Do Not Disturb}) عند المستخدم الثاني، ثمّ ابعث إشعارًا
  آخر: أظهِر أنّه «يُخزَّن» في قائمة الإشعارات لكنّه «لم يقاطع» المستخدم — فالميزة
  تحترم تفضيله.
\end{itemize}
\end{idea}

\begin{qa}[سؤال متوقّع (أمان): إزاي تمنع مستخدمًا يتنصّت على محادثة مش بتاعته؟]
عبر بوّابة تمنع الوصول المباشر لمورد غير مملوك (\en{IDOR}) داخل \code{ChatHub}: قبل
أن ينضمّ أيّ عميلٍ إلى مجموعة محادثة، يتحقّق الـ\en{hub} أنّ صاحب الـ\en{JWT} هو
فعلًا أحد طرفَي تلك المحادثة، وإلّا يرفض الانضمام. أي أنّ الاشتراك نفسه محميّ، لا
مجرّد الدفع؛ فلا يستطيع مستخدمٌ مصادَقٌ الاشتراك في محادثة الآخرين والتنصّت عليها.
\end{qa}

\begin{secnote}
أضِف نقطةً تُظهِر عمق التفكير الأمنيّ: نصّ أيّ رسالة يُنقَّى من وسوم \en{HTML}
الخبيثة قبل تخزينه، لا بعد عرضه فقط. هذا دفاعٌ في العمق، لأنّ النصّ نفسه يسافر إلى
ثلاث وجهات — واجهة \en{React}، و\en{SignalR}، والبريد الإلكترونيّ — والتنقية عند
المصدر تحمي الوجهات الثلاث دفعةً واحدة.
\end{secnote}

% ====================================================================
\subsection{المحور السادس: المدفوعات (Stripe)}

افتتح بجملة: «فلوس الطالب تُحجَز ولا تُؤخَذ — لا تُخصَم فعليًّا إلا حين يقبل المستشار
الجلسة». هذا نموذج شبيه بالوسيط الأمين (\en{escrow}) مبنيّ على \en{Stripe}.

\begin{idea}[اعرض كده: تفويض ثم قبض]
\begin{itemize}
\item احجز جلسة كطالب، ثمّ افتح لوحة \en{Stripe} (وضع الاختبار) أو سجلّ الحجز في
  النظام.
\item أظهِر أنّ الـ\en{PaymentIntent} حالته \en{authorized} (تفويض): المبلغ محجوز
  على البطاقة لا مخصوم بعد.
\item بدّل إلى حساب المستشار واقبل الجلسة، ثمّ عُد لِلّوحة وأظهِر أنّ الحالة صارت
  \en{captured} — الآن فقط تحوّلت الفلوس فعليًّا.
\end{itemize}
\end{idea}

\begin{qa}[سؤال متوقّع: فين الفلوس بين الحجز والجلسة؟]
محجوزة على بطاقة الطالب عبر التفويض اليدويّ للقبض (\en{manual capture})، لا في
حسابنا ولا مخصومة منه. مصيرها يتحدّد بالحدث التالي: تُقبَض عند قبول المستشار،
تُفَكّ عند الرفض أو انتهاء المهلة، وتُسترَدّ (كاملة أو جزئيًّا حسب السياسة) عند
الغياب (\en{no-show}). فالطالب لا يخسر قرشًا مقابل خدمةٍ لم تُؤكَّد.
\end{qa}

\begin{secnote}
اذكر ثلاث ضمانات تُظهِر نضج التعامل مع المال: أوّلًا، كل عمليّة تحمل مفتاح
لا-تكرار (\en{idempotency}) فلا يُخصَم المبلغ مرّتين لو تكرّر الطلب. ثانيًا، نتحقّق من
توقيع كل \en{webhook} قادم من \en{Stripe} قبل تصديقه. ثالثًا، بيانات البطاقة لا تلمس
خوادمنا إطلاقًا — تُدخَل مباشرةً في عناصر \en{Stripe} الآمنة، فلا نخزّن رقم بطاقة قطّ.
\end{secnote}

% ====================================================================
\subsection{المحور السابع: الأمان (سريع ومؤثّر)}

الوقت هنا دقيقة واحدة، فاجعلها مركّزة على ثلاثة أدلّة ملموسة تُثبِت أنّ الأمان ليس
شعارًا. افتتح بجملة: «الأمان عندنا ليس ميزة مضافة، بل طبقات مدمجة يمكنك رؤيتها في
القاعدة نفسها».

\begin{idea}[اعرض كده: ثلاثة أدلّة أمنيّة]
\begin{itemize}
\item افتح جدول \code{AuditLog} وأظهِر صفًّا يسجّل عمليّة حسّاسة بقيمتها «قبل» و«بعد»
  ومَن نفّذها ومتى — فكل تغيير مهمّ متروكٌ له أثر.
\item أظهِر في القاعدة عمودًا مشفَّرًا: القيمة تظهر نصًّا مشفّرًا (\en{ciphertext})
  لا نصًّا مقروءًا، فحتى من يفتح القاعدة لا يقرأ البيانات الحسّاسة.
\item ارفع ملفًّا مصابًا (عيّنة اختبار) وأظهِر أنّ فحص \en{ClamAV} يرفضه — والسلوك
  «الفشل المغلق» (\en{fail-closed}): لو تعذّر الفحص لأيّ سبب، يُرفَض الملفّ لا يُقبَل.
\end{itemize}
\end{idea}

\begin{keyp}[نقطة مهمة: الأمان مرئيّ لا موصوف]
كل ادعاء أمنيّ في هذا المحور له أثرٌ ملموس على الشاشة — صفٌّ في سِجلّ، أو نصٌّ
مشفّر، أو ملفٌّ مرفوض. هذا أقوى بكثير من قول «النظام آمن»، ويُنهي المناقشة على انطباع
احترافيّ.
\end{keyp}

% ====================================================================
\subsection{نصائح ختاميّة للعرض}

قبل الجلسة بقليل، جهّز مسرح العرض حتى لا تضيّع دقائق ثمينة في الترتيب أمام اللجنة:

\begin{itemize}
\item سجّل الدخول مسبقًا بالحسابات التجريبيّة المزروعة في النظام قبل بدء المناقشة؛
  لا تخترع كلمات مرور على الهواء — استخدم الحسابات التجريبيّة المزروعة الجاهزة.
\item جهّز التبويبات مقدّمًا: تبويب لكل محور (المحرّر، صفحة قاعدة المعرفة، الشات،
  الحجز، لوحة \en{Stripe})، مرتّبة بترتيب خريطة العشرين دقيقة.
\item للديمو الحيّ ثنائيّ المستخدم (الفيديو، والشات)، رتّب النافذتين جنبًا لجنب على
  الشاشة قبل البدء حتى تظهر اللحظيّة بوضوح.
\item حضّر خطة بديلة لو انقطع النت: تذكّر أنّ لكل خدمة سحابية محوّل \en{Stub} محلّيّ،
  فتُكمِل عرض بقيّة الشاشات \en{offline} دون ارتباك.
\item لا تُطِل في المقدمة على حساب العرض الحيّ؛ اللجنة تريد أن ترى النظام يعمل.
\end{itemize}

\begin{note}
تدرّب على العرض كاملًا مرّة واحدة على الأقلّ بالساعة في يدك، من تسجيل الدخول حتى
الأمان. التدريب يكشف التبويب الناقص، والحساب غير المزروع، والخطوة التي تأخذ وقتًا
أطول ممّا تظنّ — قبل أن تكشفها اللجنة.
\end{note}

% ====================================================================
\begin{summary}[الخلاصة في سطور]
\begin{itemize}
\item أدِر المناقشة بقاعدة «أقول ثمّ أُرِي»: كل ادعاء يُدعَم بكودٍ أو عرضٍ حيّ، ولكل
  خدمة سحابية بديلٌ محلّيّ يعمل \en{offline} لو سقط النت.
\item وزّع العشرين دقيقة بخريطة واضحة، واترك نصفها تقريبًا للعرض الحيّ لا للكلام.
\item ثلاث رسائل يجب أن تبقى عالقة في ذهن اللجنة بعد خروجك:
  \begin{enumerate}
    \item تطويرٌ مدفوعٌ بالمواصفات (نحو 211 متطلّبًا) وكلّ شيءٍ مُدقَّق عبر
      \code{AuditLog} و\en{pipeline} موحَّد.
    \item ذكاءٌ اصطناعيّ إنتاجيّ بتكلفة حدّية قريبة من الصفر (توصيات محليّة، بحثٌ في
      الذاكرة) مع سقوطٍ آمن يمنع الاختلاق.
    \item بنية فيديو (\en{ACS}) وزمنٍ حقيقيّ (\en{SignalR}) متكاملة داخل المنصّة، مع
      أمانٍ مرئيّ لا موصوف (توكنات قصيرة، تشفير، فحص ملفّات، بوّابات \en{IDOR}).
  \end{enumerate}
\end{itemize}
\end{summary}

\clearpage
% ch29-code-map (agent writing)
\section{خريطة الكود للمناقشة}

	extit{قيد الإنشاء.}

\clearpage
% ch30-appendix-diagrams — ملحق الديجرامات الرسمية (embedded من official-diagrams.pdf)
\section{ملحق: الديجرامات الرسميّة للمشروع}

الديجرامات دي هي المخطّطات الرسميّة المعتمَدة للمشروع، مُدمَجة هنا بجودة \en{vector}
(تقدر تكبّرها في العارض من غير ما تتكسّر). كلّ ديجرام في صفحة عريضة \en{(landscape)}
لوحده. الترتيب: مخطّطات العلاقات \en{(ERD)}، ثمّ المُحسَّنة \en{(EERD)}، ثمّ الربط
العلائقيّ، ثمّ الأصناف \en{(Class)}، ثمّ المكوّنات والنشر \en{(Component)}. في فصل
«خريطة الكود للمناقشة» بنقول لك أي ديجرام تعرضه مع أي موضوع.

% مولّد صفحة ديجرام: #1 رقم الصفحة في official-diagrams.pdf ، #2 العنوان العربي
\newcommand{\diag}[2]{%
  \begin{landscape}%
  \begin{center}%
  {\small\bfseries\color{spPrimary} #2}\par\vspace{3pt}%
  \includegraphics[page=#1,width=\linewidth,height=0.86\textheight,keepaspectratio]{official-diagrams.pdf}%
  \end{center}%
  \end{landscape}}

\subsection{مخطّطات العلاقات (\en{ERD})}
\diag{1}{\en{ERD} --- كلّ الموديولات (نظرة شاملة)}
\diag{2}{\en{ERD} --- موديول 1: الهوية والوصول والملف الشخصيّ}
\diag{3}{\en{ERD} --- موديول 2: المنح والطلبات والمستندات}
\diag{4}{\en{ERD} --- موديول 3: حجز الاستشاريين والتقييمات}
\diag{5}{\en{ERD} --- موديول 4: المجتمع والدردشة}
\diag{6}{\en{ERD} --- موديول 5: مركز الموارد والإشعارات}
\diag{7}{\en{ERD} --- موديول 6: الذكاء الاصطناعيّ وقاعدة المعرفة والمنصّة}

\subsection{المخطّطات المُحسَّنة (\en{EERD})}
\diag{8}{\en{EERD} --- كلّ الموديولات (نظرة شاملة)}
\diag{9}{\en{EERD} --- تخصيص المستخدم (\en{USER Specialization}): طالب / استشاريّ / جهة مانحة / أدمن}

\subsection{الربط العلائقيّ (\en{Relational Mapping})}
\diag{16}{الربط العلائقيّ --- كلّ الموديولات (تحويل الكيانات إلى جداول)}

\subsection{مخطّطات الأصناف (\en{Class Diagrams})}
\diag{24}{مخطّط الأصناف --- موديول 1: الهوية والملف الشخصيّ}
\diag{25}{مخطّط الأصناف --- موديول 2: المنح والطلبات والمستندات}
\diag{26}{مخطّط الأصناف --- موديول 3: الحجز والاستشارات والتقييمات}
\diag{27}{مخطّط الأصناف --- موديول 4: المجتمع والدردشة}
\diag{28}{مخطّط الأصناف --- موديول 5: مركز الموارد والإشعارات}
\diag{29}{مخطّط الأصناف --- موديول 6: الذكاء الاصطناعيّ والمعرفة والمنصّة}

\subsection{مخطّطات المكوّنات والنشر (\en{Component \& Deployment})}
\diag{30}{مخطّط المكوّنات والنشر الرئيسيّ (\en{free-mode})}
\diag{31}{مخطّط المكوّنات --- موديول 1: الهوية والوصول والملف الشخصيّ}
\diag{32}{مخطّط المكوّنات --- موديول 2: المنح والطلبات والمستندات}
\diag{33}{مخطّط المكوّنات --- موديول 3: حجز الاستشاريين والتقييمات}
\diag{34}{مخطّط المكوّنات --- موديول 4: المجتمع والدردشة}
\diag{35}{مخطّط المكوّنات --- موديول 5: مركز الموارد والإشعارات}
\diag{36}{مخطّط المكوّنات --- موديول 6: الذكاء الاصطناعيّ وقاعدة المعرفة والمنصّة}


\vspace{10pt}\hrule height 0.4pt\vspace{4pt}
{\small\itshape\color{spInk}
كلّ الآليّات المشروحة في هذا الكتاب مُنفَّذة فعليًّا في \code{server/src/...}
و\code{client/src/...}، ومُتحقَّقٌ منها مقابل الكود المصدريّ الحيّ.}

\end{document}
